\appendix

\section{Proof Details}

\subsection{Parameter Closure and Totalization}\label{app:step_001}


\begin{proposition}[Exact zero-dimensional law]\label{prop:step-001-zero}
Under Assumption~\ref{assump:finite-littlestone}, if \(d=0\), then \(C\)
and \(\bar C\) are singletons. With \(N=0\), the law that releases their
unique quotient concept is a Markov kernel into \((H_C,\mathcal H_C)\), is
\((0,0)\)-DP for replace-one raw adjacency, and has population binary error
zero for every probability measure \(D\) and every \(c\in C\).


\end{proposition}

\begin{proof}
Suppose two concepts \(c,c'\in C\) were distinct. There would be
an \(x\in X\) with \(c(x)\ne c'(x)\). A depth-one Littlestone tree whose
root is labeled by \(x\) would then have both binary branches realized by
members of \(C\), which would imply \(\operatorname{LD}(C)\ge1\). Thus
\(d=0\) forces every member of the nonempty set \(C\) to be the same
function; hence \(C=\{c_0\}\) and \(\bar C=\{\bar c_0\}\).

The input space \(Z_X^0\) consists of the single empty data set. Define
\(A_0(\varnothing,E)=\mathbf1\{\bar c_0\in E\}\) for
\(E\in\mathcal H_C\). This is the Dirac probability measure at
\(\bar c_0\), and every coordinate of a kernel on a singleton input is
measurable. The only replace-one comparison at sample size zero compares the
empty input with itself, so the two output laws are identical and the kernel
is \((0,0)\)-DP, hence also \((\varepsilon,\delta)\)-DP for every allowed
positive \((\varepsilon,\delta)\). Finally,
\[
\operatorname{Dec}_C(\bar c_0)(x)
=\bar c_0(\kappa(x))=c_0(x)
\quad\text{for every }x\in X,
\]
so the error is identically zero and the event that it exceeds \(\alpha\)
has probability zero, uniformly in \(D\). No positive-branch parameter,
source mechanism, failure event, or small-\(\delta\) schedule is activated.
\(\square\)


\end{proof}

\begin{lemma}[Quotient restrictions inherit the raw Littlestone bound]\label{lem:step-001-quotient-ld}
Under Assumption~\ref{assump:finite-littlestone}, if \(d\ge1\), then every
hypothesis class \(\mathcal H\subseteq\bar C\), including the empty class,
satisfies
\[
\operatorname{LD}(\mathcal H)\le d.
\tag{QuotientLDBound}
\]
In particular, every source restriction \(H_i^r\subseteq\bar C\) formed by
the positive-dimensional pointwise quotient procedure satisfies the class-
dimension premise in Lyu~\citep{lyu2025} Definition 4.2.


\end{lemma}

\begin{proof}
The empty class shatters no positive-depth Littlestone tree, so
under either standard empty-class dimension convention its dimension is
strictly less than one and hence at most \(d\). Now let
\(\mathcal H\ne\varnothing\), let \(s\ge1\), and suppose a
finite depth-\(s\) Littlestone tree \(T_Q\), whose nodes are labeled in
\(Q_C\), is shattered by \(\mathcal H\). Index a node at depth \(j-1\) by
the preceding edge string \(b_{<j}\in\{0,1\}^{j-1}\), and write its label
as \(q_{b_{<j}}\).

For each of the finitely many node occurrences, choose an arbitrary
representative
\[
x_{b_{<j}}\in\kappa^{-1}(\{q_{b_{<j}}\}).
\]
Every such fiber is nonempty because \(\kappa\) is the quotient map. Replace
the label of each node occurrence by its chosen representative, leaving the
tree and all edge labels unchanged; call the resulting \(X\)-labeled tree
\(T_X\). This uses only finitely many existential choices and does not
define a selector or a measurable object.

Fix an arbitrary path \(b=(b_1,\ldots,b_s)\in\{0,1\}^s\). Since \(T_Q\)
is shattered, there is \(\bar h_b\in\mathcal H\) such that
\[
\bar h_b(q_{b_{<j}})=b_j\qquad(1\le j\le s).
\tag{6.1}
\]
Because \(\mathcal H\subseteq\bar C\), choose \(c_b\in C\) with
\(\bar c_b=\bar h_b\). The quotient factorization in the definition of
\(\bar c_b\) gives, at every node on the path,
\[
c_b(x_{b_{<j}})
=\bar c_b(\kappa(x_{b_{<j}}))
=\bar h_b(q_{b_{<j}})
=b_j.
\tag{6.2}
\]
Thus \(c_b\) realizes path \(b\) in \(T_X\). Since \(b\) was arbitrary,
\(C\) shatters \(T_X\), so \(s\le\operatorname{LD}(C)=d\).

Repeated quotient labels and representative collisions cause no loss in
this pullback. If two node occurrences have the same quotient label, their
chosen representatives may coincide or may be different points of the same
fiber; by the definition of \(\equiv_C\), every \(c\in C\) is constant on
that entire fiber, so (6.2) is unchanged. Moreover, if the same quotient
label occurs twice on one path, a conflicting pair of edge labels could not
be realized by any function on \(Q_C\), and hence could not occur on a path
of the shattered tree \(T_Q\). Representatives of distinct quotient labels
cannot collide because a point has a unique \(\kappa\)-image. Therefore the
argument permits arbitrary fibers, repeated labels, and all possible
representative coincidences. Since every finite shattered depth \(s\) is at
most \(d\), (QuotientLDBound) follows. \(\square\)


\end{proof}

\begin{lemma}[Legal positive-branch dictionary and source calibration]\label{lem:step-001-calibration}
Under Assumptions~\ref{assump:finite-littlestone}
and~\ref{assump:approximate-dp-regime} and
Lemma~\ref{lem:step-001-quotient-ld}, if \(d\ge1\), then
\(1\le v\le d\). For every integer \(t\ge2\), all of the following
quantities are finite and lie in their displayed domains:
\[
\gamma:=\alpha/16,
\qquad
\beta_{\rm tr}=\beta_{\rm AT}=\beta_{\rm SS}=\beta_{\rm gen}:=\beta/4,
\qquad
\delta_{\rm AT}=\delta_{\rm SS}:=\delta/2,
\tag{2.1}
\]
\[
g_\delta:=\log(4/\delta),
\qquad
\eta:=\frac{\varepsilon}
{4c_{\rm AT}(\sqrt{g_\delta}+g_\delta)},
\qquad
\varepsilon_{\rm SS}:=\varepsilon/8,
\tag{2.2}
\]
\[
a(t):=v+\log(4t/\beta),
\qquad
Q(t):=e+\frac{etd^2a(t)}{\alpha v},
\tag{2.3}
\]
\[
m(t):=\left\lceil
C_{\rm blk}\frac{d^2}{\alpha}a(t)\log Q(t)
\right\rceil,
\qquad n(t):=tm(t),
\qquad p_r(t):=2^rn(t)d\quad(0\le r\le d),
\tag{2.4}
\]
\[
L(t):=p_d(t)^d2^{d^2},
\qquad
B(t):=\left\lceil
\frac{10\log(L(t)/\delta_{\rm SS})}{\varepsilon_{\rm SS}}
\right\rceil,
\tag{2.5}
\]
\[
\tau_{\rm AT}:=\eta^{-1}
\log((d+1)/\beta_{\rm AT}),
\qquad
\tau_{\rm SS}(t):=\varepsilon_{\rm SS}^{-1}
\log((tL(t)+1)/\beta_{\rm SS}).
\tag{2.6}
\]
For every \(0\le r\le d\) and \(1\le i\le t\), if
\(H_i^r\subseteq\bar C\) is nonempty, then
\(p_r(t),d\in\mathbb N\) and
\(\operatorname{LD}(H_i^r)\le d\). Consequently, if
\(\mathcal L_i^r\) is exactly an ordering without repetition of all Lyu~\citep{lyu2025}
\((p_r(t),d)\)-essential hypotheses of \(H_i^r\), then
\[
|\mathcal L_i^r|
\le p_r(t)^d2^{d^2}
\le p_d(t)^d2^{d^2}
=L(t).
\tag{ListCap}
\]
If \(H_i^r=\varnothing\), define \(\mathcal L_i^r:=\varnothing\), and
(ListCap) still holds. Conditional on each stage input list being one of
these exact valid lists or a totalized empty-restriction list, the
exact Sparse Sample threshold, the
\((\varepsilon/4,\delta/2)\) Sparse Sample allocation, and the
\((\varepsilon/4,\delta/2)\) conditional AboveThreshold allocation all hold
with the current quotient objects and source conventions. The deterministic
confidence calibrations satisfy
\[
(d+1)e^{-\eta\tau_{\rm AT}}=\beta_{\rm AT},
\qquad
(tL(t)+1)e^{-\varepsilon_{\rm SS}\tau_{\rm SS}(t)}=\beta_{\rm SS}.
\tag{2.7}
\]


\end{lemma}

\begin{proof}
A VC-shattered set of size \(s\) yields a complete depth-\(s\)
Littlestone tree by querying those \(s\) points in one fixed order; every
binary root-to-leaf labeling is realized because the set is VC-shattered.
Thus \(v\le d\). Conversely, \(d\ge1\) supplies a shattered depth-one tree,
so two concepts give opposite labels at its root and that singleton is
VC-shattered. Hence \(v\ge1\).

The parameter ranges give \(g_\delta>\log4>1\), \(\eta>0\),
\(\varepsilon_{\rm SS}>0\), \(a(t)>0\), and \(Q(t)>e\). Therefore
\(m(t),n(t),p_r(t),L(t)\) are positive integers, \(B(t)\) is a nonnegative
integer, and both tolerances in (2.6) are finite and positive. Equation
(QuotientLDBound) gives
\(\operatorname{LD}(H_i^r)\le d\) for every nonempty
\(H_i^r\subseteq\bar C\). Together with the positive integrality of
\(p_r(t)\) and \(d\), this discharges all hypotheses of Lyu~\citep{lyu2025} Definition 4.2
and Corollary 4.1. Thus an exact essential list obeys the first inequality
in (ListCap); the second follows from \(r\le d\), hence
\(p_r(t)\le p_d(t)\). A totalized empty-restriction list has cardinality zero
and obeys the same bound without invoking the source corollary.

For an arbitrary domain \(\mathcal U\), finite lists
\(\mathcal L_1,\ldots,\mathcal L_t\subseteq\mathcal U\),
\(\varepsilon_s>0\), and \(B\ge0\), Lyu~\citep{lyu2025} Algorithm~1
assigns scores
\[
s(u)=|\{i:u\in\mathcal L_i\}|,
\qquad s(\perp)=B,
\]
and samples from \((\bigcup_i\mathcal L_i)\cup\{\perp\}\) according to
\[
\Pr[w=u]
=\frac{e^{\varepsilon_s s(u)}}
{e^{\varepsilon_s B}+
  \sum_{z\in\cup_i\mathcal L_i}e^{\varepsilon_s s(z)}}.
\tag{SparseLaw}
\]
The exponent has no factor \(1/2\), and \(\perp\) is a distinct support
point. Its privacy lemma states that, if every list has size at most
\(L\) and
\[
B\ge \frac{10\log(L/\delta_s)}{\varepsilon_s},
\tag{SparsePrivacyThreshold}
\]
then this law is \((2\varepsilon_s,\delta_s)\)-DP under addition, removal,
or replacement of one entire list coordinate.

The definition of \(B(t)\) is at least this exact threshold for every exact
or totalized-empty list, including when its unrounded value is an integer, so
the list-level privacy cost is
\((2\varepsilon_{\rm SS},\delta_{\rm SS})
=(\varepsilon/4,\delta/2)\).

For sensitivity-one adaptive queries, the AboveThreshold process
with independent \(\operatorname{Lap}(1/\eta)\) query noises and at most
\(K\) reported crossings has privacy cost
\[
\left(
\eta\,O\!\left(\sqrt{K\log(1/\delta_a)}
     +\log(1/\delta_a)\right),\delta_a
\right).
\tag{ATPrivacy}
\]
Fixing \(c_{\rm AT}\ge1\) to dominate the implicit numerical constant for
\(K=1\), substituting \(\delta_a=\delta_{\rm AT}=\delta/2\), and using
\(g_\delta\ge\log(2/\delta)=\log(1/\delta_{\rm AT})\) with the value of
\(\eta\) in (2.2) gives the stated conditional
\((\varepsilon/4,\delta/2)\) allocation. This conclusion is conditional on
the sensitivity-one interface proved later and makes no raw-input privacy
claim. Finally, (2.7) follows by substituting (2.6); no probabilistic event
is assumed in this algebra. \(\square\)


\end{proof}

\begin{lemma}[Ceiling-aware candidate envelope]\label{lem:step-001-envelope}
Under Assumptions~\ref{assump:finite-littlestone}
and~\ref{assump:approximate-dp-regime}, suppose \(d\ge1\), and use the
dictionary of Lemma~\ref{lem:step-001-calibration}. Define
\[
\ell:=\log\frac{64}{\delta\beta},
\qquad
R_T:=\frac{d^2\ell\Lambda^2}{\varepsilon},
\qquad
A_{\log}:=80+\log(1+C_{\rm blk}),
\tag{3.1}
\]
and, for a universal scalar \(C\ge1\),
\[
u_C:=1+\log C,
\qquad t_C:=\lceil CR_T\rceil.
\tag{3.2}
\]
Then
\[
\begin{aligned}
\log t_C&\le A_{\log}u_C\Lambda,&
\log Q(t_C)&\le A_{\log}u_C\Lambda,\\
\log m(t_C)&\le A_{\log}u_C\Lambda,&
\log n(t_C)&\le A_{\log}u_C\Lambda,\\
\log L(t_C)&\le A_{\log}u_Cd^2\Lambda.
\end{aligned}
\tag{3.3}
\]
If
\[
A_{\rm def}:=256(A_{\log}+c_{\rm AT}+1),
\tag{3.4}
\]
then the entire defect opposing the teacher score obeys
\[
\tau_{\rm AT}+B(t_C)+\tau_{\rm SS}(t_C)
\le A_{\rm def}u_CR_T.
\tag{3.5}
\]


\end{lemma}

\begin{proof}
First, \(\ell>\log256>5\), \(\Lambda\ge1\), and
\(R_T>5\). The definition of \(\Lambda\) gives
\[
\log d,\ \log v,\ \log(1/\alpha),\ \log(1/\beta),\
\log(1/\varepsilon)\le\Lambda.
\tag{3.6}
\]
It also gives
\(\log(e/\delta)\le e^\Lambda\). Consequently
\[
\ell
=\log64+\log(1/\delta)+\log(1/\beta)
\le 7e^\Lambda,
\qquad
\log\ell\le3\Lambda.
\tag{3.7}
\]
Here the last inequality uses \(\Lambda\ge1\). Therefore
\[
\log R_T
=2\log d+\log\ell+2\log\Lambda+\log(1/\varepsilon)
\le8\Lambda.
\tag{3.8}
\]

Because \(CR_T\ge1\), the ceiling in (3.2) gives
\(CR_T\le t_C\le2CR_T\), not an asymptotic replacement. Hence
\[
\log t_C\le\log2+\log C+\log R_T\le9u_C\Lambda.
\tag{3.9}
\]
Using \(1\le v\le d\), (3.6), and (3.9),
\[
\log(4t_C/\beta)\le12u_C\Lambda,
\qquad
\log a(t_C)\le6u_C\Lambda.
\tag{3.10}
\]
For the second inequality, use
\(\log(x+y)\le\log2+\max\{\log x,\log y\}\), followed by
\(\log u_C\le u_C\) and \(\log\Lambda\le\Lambda\).

Writing \(x=t_Cd^2a(t_C)/(\alpha v)>0\),
\(Q(t_C)=e(1+x)\) and
\(\log(1+x)\le\log2+\max\{0,\log x\}\). Equations (3.6), (3.9),
and (3.10) give
\[
1\le\log Q(t_C)\le20u_C\Lambda.
\tag{3.11}
\]
The unrounded block expression
\(d^2\alpha^{-1}a(t_C)\log Q(t_C)\) is greater than one. Therefore the
ceiling is retained through
\[
m(t_C)
\le(1+C_{\rm blk})d^2\alpha^{-1}a(t_C)\log Q(t_C).
\]
Taking logarithms and using
\(\log\log Q(t_C)\le\log Q(t_C)\) yields
\[
\log m(t_C)
\le(29+\log(1+C_{\rm blk}))u_C\Lambda.
\tag{3.12}
\]
Combining (3.9) and (3.12) bounds \(\log n(t_C)\) by
\((38+\log(1+C_{\rm blk}))u_C\Lambda\). Finally, the exact list formula,
with no ceiling or power suppressed, is
\[
\begin{aligned}
\log L(t_C)
&=d\log(2^dn(t_C)d)+d^2\log2\\
&=2d^2\log2+d\log n(t_C)+d\log d\\
&\le(41+\log(1+C_{\rm blk}))u_Cd^2\Lambda.
\end{aligned}
\tag{3.13}
\]
The definition of \(A_{\log}\) makes every inequality in (3.3) follow
from (3.9), (3.11), (3.12), and (3.13).

It remains to bound the three teacher defects. Since
\(\varepsilon_{\rm SS}=\varepsilon/8\) and
\(\delta_{\rm SS}=\delta/2\), the retained ceiling gives
\[
\begin{aligned}
B(t_C)
&\le1+\frac{80}{\varepsilon}
   \left(\log L(t_C)+\log(2/\delta)\right)\\
&\le161A_{\log}u_CR_T,
\end{aligned}
\tag{3.14}
\]
where \(\log(2/\delta)\le\ell\) and each displayed summand, including the
ceiling remainder \(1\), is bounded by its corresponding multiple of
\(R_T=d^2\ell\Lambda^2/\varepsilon\).

Because \(t_CL(t_C)\ge2\),
\[
\begin{aligned}
\tau_{\rm SS}(t_C)
&\le\frac8\varepsilon
 \left(\log(8/\beta)+\log t_C+\log L(t_C)\right)\\
&\le24A_{\log}u_CR_T.
\end{aligned}
\tag{3.15}
\]
For AboveThreshold, \(g_\delta>1\), so
\(\sqrt{g_\delta}+g_\delta\le2g_\delta\le2\ell\), while
\[
\log(4(d+1)/\beta)
\le\log4+\log(d+1)+\log(1/\beta)\le4\Lambda.
\]
Thus
\[
\tau_{\rm AT}
\le\frac{32c_{\rm AT}\ell\Lambda}{\varepsilon}
\le32c_{\rm AT}R_T.
\tag{3.16}
\]
Summing (3.14)-(3.16) and using (3.4) proves (3.5). Notice that the
candidate tuple was fully defined before any feasibility statement was
used; this is the noncircular part of the argument. \(\square\)


\end{proof}

\begin{proposition}[Finite public teacher witness and least feasible count]\label{prop:step-001-teacher}
Under Assumptions~\ref{assump:finite-littlestone}
and~\ref{assump:approximate-dp-regime}, suppose \(d\ge1\). With
\(A_{\log}\) and \(A_{\rm def}\) from
Lemma~\ref{lem:step-001-envelope}, fix the concrete finite universal
constant
\[
C_{\rm teach}:=2^{12}A_{\rm def}^2
\tag{4.1}
\]
and define
\[
\bar k:=\left\lceil
C_{\rm teach}
\frac{d^2\log(64/(\delta\beta))\Lambda^2}{\varepsilon}
\right\rceil.
\tag{4.2}
\]
Then \(\bar k\ge2\) and
\[
\frac{\bar k}{2}-\tau_{\rm AT}
\ge B(\bar k)+\tau_{\rm SS}(\bar k)+2.
\tag{4.3}
\]
Consequently the exact least feasible teacher count
\[
k:=\min\left\{t\in\mathbb Z:t\ge2,\quad
\frac t2-\tau_{\rm AT}
\ge B(t)+\tau_{\rm SS}(t)+2\right\}
\tag{4.4}
\]
exists and satisfies
\[
2\le k\le\bar k
\le2C_{\rm teach}
\frac{d^2\log(64/(\delta\beta))\Lambda^2}{\varepsilon}.
\tag{4.5}
\]
After this deterministic minimization, define exactly
\[
m:=m(k),\qquad n_0:=km(k),\qquad N:=n_0,
\qquad p_r:=2^rn_0d\ (0\le r\le d),
\tag{4.6}
\]
\[
L:=p_d^d2^{d^2}=L(k),\qquad B:=B(k),
\qquad \tau_{\rm SS}:=\tau_{\rm SS}(k).
\tag{4.7}
\]
These are finite public quantities, all integer ceilings from (2.4), (2.5),
and (4.2) are retained, and
\[
p_0=n_0d\ge\max\{n_0,d+1\}.
\tag{4.8}
\]


\end{proposition}

\begin{proof}
Put \(C=C_{\rm teach}\) in
Lemma~\ref{lem:step-001-envelope}. Since \(A_{\rm def}\ge1\),
\[
1+\log C_{\rm teach}
=1+12\log2+2\log A_{\rm def}
\le15A_{\rm def}.
\]
Therefore
\[
A_{\rm def}(1+\log C_{\rm teach})
\le15A_{\rm def}^2
\le\frac{C_{\rm teach}}4.
\tag{4.9}
\]
Equations (3.5) and (4.9) imply
\[
\tau_{\rm AT}+B(\bar k)+\tau_{\rm SS}(\bar k)
\le\frac{C_{\rm teach}R_T}{4}.
\tag{4.10}
\]
Also \(C_{\rm teach}R_T/4>2\), while
\(\bar k/2\ge C_{\rm teach}R_T/2\). Thus
\[
\frac{\bar k}{2}
\ge\tau_{\rm AT}+B(\bar k)+\tau_{\rm SS}(\bar k)+2,
\]
which is (4.3). The feasible set in (4.4) is a nonempty subset of the
integers bounded below by two, so its least element exists and is no larger
than \(\bar k\). Since \(C_{\rm teach}R_T\ge1\), the ceiling in (4.2) gives
\(\bar k\le C_{\rm teach}R_T+1\le2C_{\rm teach}R_T\), proving (4.5).

All objects in (4.6)-(4.7) are now evaluated at a fixed public integer
\(k\); there is no remaining equation in which \(k\) is defined using a
quantity that itself awaits the choice of \(k\). Finally, \(n_0\ge k\ge2\)
and \(d\ge1\), so \(p_0=n_0d\ge n_0\). Also
\(n_0d\ge2d\ge d+1\), proving (4.8). \(\square\)


\end{proof}

\begin{proposition}[Pre-sampling totalization of the quotient procedure]\label{prop:step-001-totalization}
Under Assumptions~\ref{assump:finite-littlestone}
and~\ref{assump:approximate-dp-regime},
Lemma~\ref{lem:step-001-calibration}, and
Proposition~\ref{prop:step-001-teacher}, the positive-dimensional
pointwise quotient procedure can be specified before the master sample,
partition, source events, or mechanism coins are drawn so that it returns an
element of \(H_C\) on every labeled quotient input and every internal path.
The completion assigns the empty list to every empty restriction; preserves
every exact valid list and successful path; replaces every invalid,
nonfinite, non-\(H_C\), or oversized purported list by the empty list; gives
every remaining routine-failure state a defined transcript continuation; and
sends every terminal mechanism-failure, no-success, or residual fallback
path to the setting-defined \(\bar c_0\).


\end{proposition}

\begin{proof}
Fix the tuple (2.1)-(2.6), (4.4), and (4.6)-(4.7), the default
\(\bar c_0\), and deterministic tie-breaking conventions once and for all.
For every nonempty restriction on which the certified routine supplies valid
objects, fix one optimal decomposition and an ordering without repetition of
its exact essential set. By (ListCap), every such exact list is finite, lies
in \(H_C\), and has cardinality at most \(L\); hence the cap and finiteness
checks below never change it, whether that exact list is empty or nonempty.
Every empty restriction is assigned the empty list directly. Identical local
states are assigned identical choices. This is a pointwise definition only.
Its measurable-kernel realization on the countable quotient input is
established independently after the pointwise law has been fixed;
measurability is not assumed here.

For a totalized list tuple at stage \(r\), define its occurrence score by
\[
q_r:=
\begin{cases}
\displaystyle\max_{\bar h\in\cup_{i=1}^k\mathcal L_i^r}
|\{i:\bar h\in\mathcal L_i^r\}|,
&\bigcup_{i=1}^k\mathcal L_i^r\ne\varnothing,\\[4pt]
0,&\bigcup_{i=1}^k\mathcal L_i^r=\varnothing.
\end{cases}
\tag{5.1}
\]
Feed \(q_0,\ldots,q_d\), in order, to one stopped source
AboveThreshold process with threshold \(k/2\), noise scale \(1/\eta\), and
the fixed first-success rule. At the first legal Above stage \(r_*\), invoke
the exact law (SparseLaw) on
\((\mathcal L_1^{r_*},\ldots,\mathcal L_k^{r_*})\) with parameters
\((\varepsilon_{\rm SS},B)\). If it returns an actual union item, return
that item; otherwise use the terminal convention below.

The pointwise procedure is totalized as follows. An empty stage restriction
receives the empty essential list. If a source decomposition or list routine
is undefined, returns no valid object, or fails a declared check, the local
list is replaced by the empty list. A purported list that is invalid,
nonfinite, outside \(H_C\), or larger than \(L\) is treated the same way;
valid lists are unchanged. If all lists at a stage are empty, its occurrence
score is \(0\) and the fixed transcript continues. If AboveThreshold reports
no legal stage, the transcript is invalid, or stage processing is exhausted,
the output is \(\bar c_0\). When a legal stage is selected and Sparse Sample
returns an actual union item, that item in \(H_C\) is returned. A failure
symbol, an out-of-support value, an invalid normalizer, or any remaining
mechanism failure returns \(\bar c_0\).

For finite capped lists, (SparseLaw) in fact always has a positive finite
normalizer because the \(\perp\) weight is positive; its listed failure
completion is included to make the pointwise program total independently of
that later path verification. On an exact valid path, every nonempty
restriction's exact list satisfies (ListCap), every empty restriction already
has the prescribed empty list, AboveThreshold names a legal stage, and Sparse
Sample returns an actual union item. Thus neither the finite/cap sanitizer nor
any fallback changes that path. Conversely, any invalid, nonfinite,
non-\(H_C\), or oversized purported list is totalized to the empty list, so
it cannot leave the pointwise procedure undefined.

The partition rule is also fixed before data: after the complete quotient
master sample is obtained, draw one data-independent uniform partition into
the \(k\) indexed blocks of size \(m\), and reuse that same partition at all
\(d+1\) stages. The table itself, all tie breaking, all thresholds, and the
fallback output are fixed before this draw. Thus no good event, realizability
condition, mechanism-success event, privacy conclusion, PAC conclusion, or
event-dependent choice is used to define the law.

Proposition~\ref{prop:step-001-zero} proves the complete \(d=0\) branch:
the exact sample size is zero, the output is the unique quotient concept,
the kernel is deterministic, privacy is \((0,0)\), and the uniform PAC
failure probability is zero.

On the \(d\ge1\) branch,
Lemma~\ref{lem:step-001-quotient-ld} pulls every finite quotient
Littlestone tree back to \(X\) and proves
\(\operatorname{LD}(H_i^r)\le d\) for every current restriction, including
all repeated-label and arbitrary-fiber cases. With that exact class premise
available, Lemma~\ref{lem:step-001-calibration} proves \(1\le v\le d\),
defines every confidence, privacy, block, list, failure, and mechanism
tolerance for each candidate integer \(t\ge2\), and checks the exact Lyu~\citep{lyu2025}
essential-list cap and Sparse Sample formulas in the current quotient
conventions. Lemma~\ref{lem:step-001-envelope} retains both ceilings and
proves a finite logarithmic envelope for the complete coupled tuple before
assuming a feasible teacher. Proposition~\ref{prop:step-001-teacher} then
gives the concrete finite public witness (4.2), proves its margin, takes the
least feasible integer only after nonemptiness is known, fixes
\((k,m,n_0,N,p_r,L,B)\), and exports the explicit bound (4.5), whose leading
teacher dependence is \(d^2/\varepsilon\). Finally, the totalization
construction above fixes every empty,
routine-failure, mechanism-failure, no-success, and fallback outcome before
any random event while preserving all exact valid paths.

This establishes the complete parameter dictionary, finite teacher witness,
least-feasible count, and all-input totalization used by the later measurable,
privacy, and utility results.



\end{proof}

\subsection{Evaluation Quotient and Exact Transfers}\label{app:step_002}


The propositions below give the measurable quotient, product-law, and risk
interfaces.

\begin{proposition}[Measurable evaluation factorization]\label{prop:step-002-factorization}
Under Assumption~\ref{assump:countable-evaluation-quotient}, for the nonempty
class \(C\) in the formalized basic setup, the relation
\[
x\equiv_Cx'\quad\Longleftrightarrow\quad
c(x)=c(x')\quad\hbox{for every }c\in C
\]
is an equivalence relation. Its quotient \(Q_C\) is nonempty and finite or
countably infinite, the canonical surjection
\(\kappa(x)=[x]_{\equiv_C}\) is measurable into \((Q_C,2^{Q_C})\), and
\[
\Phi:C\longrightarrow\bar C,\qquad \Phi(c)=\bar c,
\quad \bar c(\kappa(x))=c(x),
\]
is a bijection. Moreover, every \(c\in C\) and every decoded hypothesis
\(\operatorname{Dec}_C(\bar h)=\bar h\circ\kappa\),
\(\bar h\in H_C\), is \(\Sigma/2^{\{0,1\}}\)-measurable.


\end{proposition}

\begin{proof}
Reflexivity, symmetry, and transitivity follow coordinatewise from
equality of the evaluation vector \((c(x))_{c\in C}\), so the relation is an
equivalence relation. Since \(X\ne\varnothing\), its quotient is nonempty,
and its finite-or-countable cardinality is exactly the first clause of
Assumption~\ref{assump:countable-evaluation-quotient}. The quotient map is
surjective by construction.

For every \(A\subseteq Q_C\),
\[
\kappa^{-1}(A)=\bigcup_{q\in A}\kappa^{-1}(\{q\}).
\tag{1}
\]
The union is finite or countable, and each cell belongs to \(\Sigma\) by the
assumption. Hence \(\kappa\) is measurable. This is a direct cell argument;
it makes no choice of a representative of any cell.

If \(\kappa(x)=\kappa(x')\), then \(c(x)=c(x')\) for every \(c\in C\), so
\(\bar c(\kappa(x)):=c(x)\) is well-defined. Surjectivity of \(\Phi\) is the
definition of \(\bar C\). If \(\Phi(c)=\Phi(c')\), then for every \(x\in X\),
\[
c(x)=\bar c(\kappa(x))=\bar c'(\kappa(x))=c'(x),
\]
and therefore \(c=c'\) as functions. Thus \(\Phi\) is injective and hence
bijective.

Every map from the discrete space \(Q_C\) to \(\{0,1\}\) is measurable. More
explicitly, for \(b\in\{0,1\}\),
\[
c^{-1}(\{b\})
=\bigcup_{q:\bar c(q)=b}\kappa^{-1}(\{q\}),
\qquad
(\bar h\circ\kappa)^{-1}(\{b\})
=\bigcup_{q:\bar h(q)=b}\kappa^{-1}(\{q\}).
\tag{2}
\]
Both are countable unions of measurable cells. Thus the originally given
concepts and every possibly improper decoded output are measurable. The
argument uses the values of a quotient function on cells, never a selector
from the cells. \(\square\)


\end{proof}

\begin{proposition}[Countable quotient and output Borel structure]\label{prop:step-002-borel}
Under Assumption~\ref{assump:countable-evaluation-quotient} and
Proposition~\ref{prop:step-002-factorization}, \(Q_C\),
\(Z_Q=Q_C\times\{0,1\}\), and \(Z_Q^N\) for every
\(N\in\mathbb N_0\) are countable discrete standard-Borel spaces (with
\(Z_Q^0\) the one-point empty-tuple space). The output measurable space
\[
(H_C,\mathcal H_C)
=\left(\{0,1\}^{Q_C},\bigotimes_{q\in Q_C}2^{\{0,1\}}\right)
\]
is standard Borel, and \(\mathcal H_C\) is generated by the finite-coordinate
cylinders. In particular, each coordinate evaluation
\(e_q(\bar h)=\bar h(q)\) is measurable.


\end{proposition}

\begin{proof}
A finite or countably infinite discrete space is Polish under the
discrete metric: the metric is complete, and the whole countable space is a
countable dense set. Hence \(Q_C\), \(Z_Q\), and every finite power \(Z_Q^N\)
are standard Borel. Their Borel sigma-fields are their full power sets. For
\(N=0\), the only point is the empty tuple.

If \(Q_C\) is finite, then \(H_C\) is finite discrete and the conclusion is
immediate. If \(Q_C\) is countably infinite, fix an enumeration
\(Q_C=\{q_1,q_2,\ldots\}\) and put
\[
\rho(\bar h,\bar g)
:=\sum_{j=1}^{\infty}2^{-j}
  |\bar h(q_j)-\bar g(q_j)|.
\tag{3}
\]
A \(\rho\)-Cauchy sequence is eventually constant in each coordinate; the
coordinatewise limit belongs to \(H_C\), and the tail of the series in (3)
shows convergence to it. Thus \(\rho\) is complete. Functions that vanish
outside a finite initial segment form a countable dense subset. The topology
of \(\rho\) is exactly the product topology: fixing finitely many coordinates
gives a clopen cylinder, and sufficiently long initial-coordinate agreement
makes the tail in (3) arbitrarily small. Consequently this topology is
Polish and its Borel sigma-field is precisely the sigma-field generated by
finite-coordinate cylinders, namely \(\mathcal H_C\). Each \(e_q\) is a
coordinate map and therefore measurable. The enumeration is only an
enumeration of the already formed quotient; no raw representative is chosen.
\(\square\)


\end{proof}

\begin{lemma}[VC dimension is invariant under the evaluation quotient]\label{lem:step-002-vc}
Under Assumptions~\ref{assump:finite-littlestone} and
\ref{assump:countable-evaluation-quotient} and
Proposition~\ref{prop:step-002-factorization}, a finite subset of \(X\) is
shattered by \(C\) only if its quotient image has the same cardinality and is
shattered by \(\bar C\); conversely, every finite subset of \(Q_C\) shattered
by \(\bar C\) has a same-cardinality raw lift shattered by \(C\). Therefore
\[
\operatorname{VC}(\bar C)=\operatorname{VC}(C)=v.
\tag{4}
\]


\end{lemma}

\begin{proof}
Let \(S=\{x_1,\ldots,x_s\}\subseteq X\) be shattered by \(C\). If
two distinct points \(x_i,x_j\) had \(\kappa(x_i)=\kappa(x_j)\), shattering
would provide a concept taking value \(0\) at one and \(1\) at the other,
contrary to the definition of \(\equiv_C\). Hence \(\kappa\) is injective on
\(S\), so \(|\kappa(S)|=s\). Given any labeling
\(\ell:\kappa(S)\to\{0,1\}\), shattering of \(S\) supplies a \(c\in C\)
with \(c(x)=\ell(\kappa(x))\) on \(S\). Its image \(\bar c=\Phi(c)\) realizes
\(\ell\) on \(\kappa(S)\). Thus \(\kappa(S)\) is shattered by \(\bar C\).

Conversely, let \(A=\{q_1,\ldots,q_s\}\subseteq Q_C\) be shattered by
\(\bar C\). For these finitely many cells, choose one
\(x_i\in\kappa^{-1}(\{q_i\})\). The points are distinct because the cells are
disjoint. For any labeling of \(\{x_1,\ldots,x_s\}\), use the same labels on
\(A\); quotient shattering supplies \(\bar c\in\bar C\), and the unique
\(c=\Phi^{-1}(\bar c)\) realizes the raw labels. Thus the raw lift is
shattered. Only finitely many representatives are chosen for this one
combinatorial witness; no representative map on \(Q_C\), measurable or
otherwise, is constructed. The attainable finite shattering cardinalities
are the same in both classes, proving (4). \(\square\)


\end{proof}

\begin{lemma}[Littlestone dimension is invariant under the evaluation quotient]\label{lem:step-002-ld}
Under Assumptions~\ref{assump:finite-littlestone} and
\ref{assump:countable-evaluation-quotient} and
Proposition~\ref{prop:step-002-factorization}, a finite complete binary tree
is shattered by \(C\) if and only if a tree of the same depth obtained by
quotienting or finitely lifting its node labels is shattered by \(\bar C\).
Consequently,
\[
\operatorname{LD}(\bar C)=\operatorname{LD}(C)=d.
\tag{5}
\]
If \(d=0\), then \(C\) and \(\bar C\) are singletons and \(v=0\).


\end{lemma}

\begin{proof}
A complete binary \(X\)-labeled tree of depth \(t\) is a family
\((x_\sigma)_{\sigma\in\{0,1\}^{<t}}\). It is shattered by \(C\) when, for
each \(b=(b_1,\ldots,b_t)\in\{0,1\}^t\), some \(c_b\in C\) satisfies
\[
c_b(x_{b_{<r}})=b_r,
\qquad r=1,\ldots,t,
\tag{6}
\]
where \(b_{<r}=(b_1,\ldots,b_{r-1})\).

Suppose the raw tree is shattered and set
\(q_\sigma:=\kappa(x_\sigma)\). For each path \(b\), the quotient concept
\(\Phi(c_b)\) satisfies
\[
\Phi(c_b)(q_{b_{<r}})=c_b(x_{b_{<r}})=b_r,
\]
so the quotient-labeled tree is shattered. If one's tree convention requires
distinct instances along a path, that condition is automatic. Indeed, if a
quotient label repeated at two nodes on a root-to-node path, extend the path
at the later node using the bit opposite to the bit taken at the earlier
node. Equation (6) would then require one concept to assign two values to the
same quotient cell, a contradiction. Repetitions across incomparable
branches are harmless and are allowed by the usual tree definition.

Conversely, let \((q_\sigma)_{\sigma\in\{0,1\}^{<t}}\) be a quotient tree
shattered by \(\bar C\). It has only \(2^t-1\) nodes. Choose, separately for
these finitely many nodes, an \(x_\sigma\in\kappa^{-1}(\{q_\sigma\})\). For
each path, choose the witnessing \(\bar c_b\in\bar C\) and let
\(c_b=\Phi^{-1}(\bar c_b)\). Then
\[
c_b(x_{b_{<r}})=\bar c_b(q_{b_{<r}})=b_r,
\]
so the lifted raw tree is shattered at the same depth. Again, this finite
witness lift is not a global selector. The same finite depths are attainable
on both sides, proving (5).

Finally suppose \(d=0\). If two members \(c_0,c_1\in C\) were distinct, they
would disagree at some \(x\). Because the labels are binary, the one-node
tree labeled by \(x\) would have both branch labels realized and hence would
be shattered, contradicting \(d=0\). Nonemptiness therefore makes \(C\) a
singleton, and the bijection \(\Phi\) makes \(\bar C\) a singleton. A
singleton binary class shatters no one-point set, so Lemma~\ref{lem:step-002-vc}
gives \(v=0\). Notice that \(Q_C\) itself need not be a singleton: the unique
concept may take both binary values, in which case it has two evaluation
cells. \(\square\)


\end{proof}

\begin{proposition}[Measurable record transport and replacement adjacency]\label{prop:step-002-record-map}
Under Assumption~\ref{assump:countable-evaluation-quotient} and
Propositions~\ref{prop:step-002-factorization}--\ref{prop:step-002-borel},
the one-record map
\[
t:Z_X\to Z_Q,\qquad t(x,y)=(\kappa(x),y),
\tag{7}
\]
and every data map \(T_N=t^N\), \(N\in\mathbb N_0\), are measurable. If
\(s,s'\in Z_X^N\) differ in at most one coordinate, then their images are
equal or differ in exactly one coordinate. This conclusion holds on every
raw labeled input, including arbitrary nonrealizable labels and same-cell
replacements.


\end{proposition}

\begin{proof}
For each atom \((q,b)\in Z_Q\),
\[
t^{-1}(\{(q,b)\})=\kappa^{-1}(\{q\})\times\{b\}
\in\Sigma\otimes2^{\{0,1\}}.
\tag{8}
\]
The codomain \(Z_Q\) is countable discrete, so the preimage of an arbitrary
subset is a countable union of the sets in (8). Thus \(t\) is measurable.
For \(N\ge1\), the preimage of an atom of \(Z_Q^N\) is the measurable
rectangle
\[
T_N^{-1}\!\left(
 \left\{((q_r,b_r))_{r=1}^N\right\}
\right)
=\prod_{r=1}^N
  \bigl(\kappa^{-1}(\{q_r\})\times\{b_r\}\bigr).
\tag{9}
\]
Since \(Z_Q^N\) is countable discrete, another countable-union argument
proves measurability of \(T_N\). For \(N=0\), \(T_0\) is the unique map
between one-point empty-tuple spaces and is measurable.

If \(s=s'\), their images are equal. Otherwise, if \(s\) and \(s'\) differ
only at coordinate \(j\), all image coordinates
other than \(j\) agree. At coordinate \(j\), replacing \((x,y)\) by
\((x',y')\) has the following exhaustive behavior:

\noindent \ if \(\kappa(x)=\kappa(x')\) and \(y=y'\), the quotient record is unchanged,
  so the full image datasets are equal;
\noindent \ if \(\kappa(x)=\kappa(x')\) and \(y\ne y'\), only the quotient label at
  coordinate \(j\) changes;
\noindent \ if \(\kappa(x)\ne\kappa(x')\), only coordinate \(j\) changes, regardless of
  the two labels.

Thus the images are equal or strict replace-one neighbors and, under the
setting's "differ in at most one coordinate" convention, are always
replace-one neighbors. No realizability condition was used. \(\square\)


\end{proof}

\begin{proposition}[Exact realizable iid pushforward]\label{prop:step-002-iid-pushforward}
Under Assumptions~\ref{assump:countable-evaluation-quotient} and
\ref{assump:realizable-iid}, Proposition~\ref{prop:step-002-factorization},
and Proposition~\ref{prop:step-002-record-map}, for every probability measure
\(D\) on \((X,\Sigma)\), every \(c\in C\), and every
\(N\in\mathbb N_0\),
\[
t_\#P_{D,c}=P_{\bar D,\bar c},
\qquad
(T_N)_\#P_{D,c}^N=P_{\bar D,\bar c}^N,
\quad \bar D=\kappa_\#D,
\tag{10}
\]
where \(\bar c=\Phi(c)\).


\end{proposition}

\begin{proof}
Proposition~\ref{prop:step-002-factorization} shows that \(c\) is
measurable. Hence \(x\mapsto(x,c(x))\) is measurable: the preimage of a
generating rectangle \(A\times\{b\}\) is
\(A\cap c^{-1}(\{b\})\). Thus \(P_{D,c}\) is well-defined. Likewise,
\(q\mapsto(q,\bar c(q))\) is measurable on the discrete quotient, so
\(P_{\bar D,\bar c}\) is well-defined. For each atom
\((q,b)\in Z_Q\),
\[
\begin{aligned}
(t_\#P_{D,c})(\{(q,b)\})
&=D\{x:\kappa(x)=q,\ c(x)=b\}\\
&=\mathbf 1\{b=\bar c(q)\}\,D(\kappa^{-1}(\{q\}))\\
&=\mathbf 1\{b=\bar c(q)\}\,\bar D(\{q\})
=P_{\bar D,\bar c}(\{(q,b)\}).
\end{aligned}
\tag{11}
\]
The second equality uses that \(c\) is exactly constant with value
\(\bar c(q)\) on the whole fiber \(\kappa^{-1}(\{q\})\). Equality on atoms
implies equality on all subsets of the countable discrete space \(Z_Q\).

For \(N\ge1\), evaluate both sides of the second identity in (10) on an atom
\(((q_r,b_r))_{r=1}^N\). By (9), the product-measure formula, and (11), its
probability under the left side is
\[
\prod_{r=1}^N
P_{D,c}\bigl(t^{-1}(\{(q_r,b_r)\})\bigr)
=\prod_{r=1}^N P_{\bar D,\bar c}(\{(q_r,b_r)\}),
\tag{12}
\]
which is its probability under \(P_{\bar D,\bar c}^N\). Since \(Z_Q^N\) is
countable discrete, atomwise equality proves equality of the measures.
Equation (12) also shows directly that deterministic quotienting preserves
iid sampling even when several raw draws lie in the same fiber or produce
repeated quotient records. For \(N=0\), both sides are the point mass at the
empty tuple. \(\square\)


\end{proof}

\begin{proposition}[Exact risk transfer for every quotient output]\label{prop:step-002-risk}
Under Assumption~\ref{assump:countable-evaluation-quotient} and
Propositions~\ref{prop:step-002-factorization}--\ref{prop:step-002-borel},
for every probability measure \(D\) on \((X,\Sigma)\), every \(c\in C\),
and every \(\bar h\in H_C\), with no requirement that
\(\bar h\in\bar C\), the disagreement sets are measurable and
\[
\operatorname{err}_D(\operatorname{Dec}_C(\bar h),c)
=\operatorname{err}_{\bar D}(\bar h,\bar c),
\qquad \bar D=\kappa_\#D,\quad \bar c=\Phi(c).
\tag{13}
\]
For fixed \(D,c\), the common value in (13), viewed as a function of
\(\bar h\in(H_C,\mathcal H_C)\), is measurable.


\end{proposition}

\begin{proof}
Define the quotient disagreement set
\[
A_{\bar h,\bar c}:=\{q\in Q_C:\bar h(q)\ne\bar c(q)\}.
\tag{14}
\]
It is measurable because \(Q_C\) is discrete. Proposition~\ref{prop:step-002-factorization}
gives \(c=\bar c\circ\kappa\) and a measurable decoder, and therefore
\[
\{x:\operatorname{Dec}_C(\bar h)(x)\ne c(x)\}
=\kappa^{-1}(A_{\bar h,\bar c}).
\tag{15}
\]
Taking \(D\)-measure and using the definition of pushforward gives
\[
\begin{aligned}
\operatorname{err}_D(\operatorname{Dec}_C(\bar h),c)
&=D(\kappa^{-1}(A_{\bar h,\bar c}))\\
&=(\kappa_\#D)(A_{\bar h,\bar c})
=\bar D\{q:\bar h(q)\ne\bar c(q)\}\\
&=\operatorname{err}_{\bar D}(\bar h,\bar c),
\end{aligned}
\]
which proves (13) with no fiber-size, injectivity, propriety, or support
condition.

For measurability in \(\bar h\), the finite-\(Q_C\) case is a finite sum of
coordinate-cylinder indicators. In the countably infinite case, use the
enumeration from Proposition~\ref{prop:step-002-borel}. Then the right side
of (13) is the pointwise increasing limit
\[
\lim_{J\to\infty}
\sum_{j=1}^J \bar D(\{q_j\})
\mathbf 1\{\bar h(q_j)\ne\bar c(q_j)\}.
\tag{16}
\]
Each partial sum is \(\mathcal H_C\)-measurable because it depends on finitely
many coordinate evaluations. Hence the limit is measurable. Equation (16)
uses quotient atoms and cylinders only; it does not select a raw point from
any fiber. Also, for every fixed \(x\in X\), the output-coordinate map
\(\bar h\mapsto\operatorname{Dec}_C(\bar h)(x)=e_{\kappa(x)}(\bar h)\) is
\(\mathcal H_C\)-measurable by Proposition~\ref{prop:step-002-borel}.
Proposition~\ref{prop:step-002-factorization} constructs the canonical
evaluation quotient and the bijection \(\Phi:C\leftrightarrow\bar C\), while
also deriving measurability of all raw concepts and every decoded quotient
output. Proposition~\ref{prop:step-002-borel} establishes the exact
countable-discrete input spaces and standard-Borel product output required by
the setting. Lemma~\ref{lem:step-002-vc} and
Lemma~\ref{lem:step-002-ld} transfer, in both directions, every finite VC
witness and every finite Littlestone-tree witness, proving
\(\operatorname{VC}(\bar C)=v\) and \(\operatorname{LD}(\bar C)=d\), with
the singleton-class conclusion at \(d=0\).

Proposition~\ref{prop:step-002-record-map} proves measurability of the record
map and every \(T_N\) by atom/cylinder preimages, and proves exact
equal-or-replace-one transport for same-cell changes, label changes, and all
other arbitrary raw replacements. Proposition~\ref{prop:step-002-iid-pushforward}
then proves the realizable iid product-law identity for every finite sample
size, including zero. Finally, the disagreement-set calculation above identifies
the raw disagreement set as the preimage of the quotient disagreement set
and proves the exact population-risk equality for every \(\bar h\in H_C\),
including improper outputs and arbitrary fibers. These propositions give the
common quotient/raw data, adjacency, sampling, and risk interfaces.



\end{proof}

\subsection{Countable-Domain Kernel Construction}\label{app:step_003}


No paper result is used in this argument. The preceding kernel construction and
the standard measurable-space facts supply the interfaces used below.

 \begin{lemma}[Countable-atom promotion of pointwise coordinates and laws]\label{lem:step-003-countable-promotion}
Under Assumption~\ref{assump:countable-evaluation-quotient} and Proposition~\ref{prop:step-002-borel}, fix any
\(N\in\mathbb N_0\) and put \(S_N:=Z_Q^N\). Then:

1. for every measurable space \((Y,\mathcal Y)\), every function
   \(f:S_N\to Y\) is \(2^{S_N}/\mathcal Y\)-measurable;
2. for every family \((\mu_{\bar s})_{\bar s\in S_N}\) of probability
   measures on \((Y,\mathcal Y)\),
   \[
   K(\bar s,B):=\mu_{\bar s}(B),
   \qquad \bar s\in S_N,\quad B\in\mathcal Y,
   \tag{1}
   \]
   is a Markov kernel from \(S_N\) to \(Y\); and
3. the same conclusions hold after adjoining any finite or countable
   discrete collection of partition, stage, or history indices to \(S_N\).

No compatibility, enumeration, or selector condition is required of the
pointwise family in (1).


\end{lemma}

\begin{proof}
Proposition~\ref{prop:step-002-borel} gives
\(\sigma(S_N)=2^{S_N}\), including when \(N=0\), where \(S_0\) is a
singleton. For \(B\in\mathcal Y\), the preimage \(f^{-1}(B)\) is simply a
subset of \(S_N\), hence belongs to \(2^{S_N}\). This proves the first
claim.

For the second claim, \(B\mapsto K(\bar s,B)\) is the probability measure
\(\mu_{\bar s}\) for each fixed \(\bar s\). For fixed \(B\in\mathcal Y\),
\(\bar s\mapsto\mu_{\bar s}(B)\) is an arbitrary function from \(S_N\) to
\([0,1]\), so the first claim makes it measurable. These are exactly the two
Markov-kernel axioms. A finite or countable product of \(S_N\) with a
discrete index set is again countable discrete, so the same proof gives the
third claim. At no point is a choice made from a quotient cell or from an
uncountable family of source classes. \(\square\)


\end{proof}

\begin{proposition}[Measurable coding of restrictions, finite lists, and
transcripts]\label{prop:step-003-coding}
Under Assumption~\ref{assump:countable-evaluation-quotient}, 
Propositions~\ref{prop:step-001-totalization} and
\ref{prop:step-002-borel}, and
Lemma~\ref{lem:step-003-countable-promotion}, fix
\(N\in\mathbb N_0\). Every restriction, decomposition, finite-list,
validity, event-status, finite-stage transcript, and terminal-output
coordinate of the totalized quotient procedure has a measurable interface.
More precisely:

1. every realized restriction or decomposition may be coded in its
   countable realized range with the discrete sigma-field;
2. every sanitized finite list is a measurable coordinate with values in
   the standard-Borel space
   \[
   \mathsf{List}(H_C):=\bigsqcup_{\ell=0}^{\infty}H_C^\ell,
   \tag{2}
   \]
   and its empty-list, length, item, and membership coordinates are
   measurable;
3. the finite transcript may be coded in a standard-Borel space built from
   finite/countable discrete flags, finitely many copies of (2), finitely
   many real or uniform randomization coordinates, and
   \[
   H_C^\dagger:=H_C\sqcup\{\perp,\dagger\},
   \tag{3}
   \]
   where \(\perp\) is the Sparse Sample failure symbol and \(\dagger\) is an
   optional invalid-state token; and
4. every family of pointwise transcript laws indexed by the quotient input
   is a Markov kernel, and every measurable transcript coordinate has a
   Markov-kernel marginal.

The construction includes empty input, empty restrictions and lists,
invalid source paths, no-success exhaustion, and terminal fallback. It does
not put a sigma-field on the full hyperspace \(2^{H_C}\) and does not choose
a representative from any raw quotient cell.


\end{proposition}

\begin{proof}
Let \(\mathfrak P_N\) denote the finite set of data-independent
indexed partitions used by the procedure at its declared sample size; when
no partition is used, including the \(N=0\) branch, let
\(\mathfrak P_N\) be a singleton. Let \(\mathfrak A_N\) collect the finitely
many stage indices and finite-alphabet structural-history flags that can
precede a list or restriction coordinate. The structural state space
\[
\mathsf U_N:=S_N\times\mathfrak P_N\times\mathfrak A_N
\tag{4}
\]
is countable discrete. Duplicate records, arbitrary binary labels, and
nonrealizable samples remain ordinary atoms of \(S_N\); they do not alter
this conclusion.

Fix one restriction coordinate \(R:\mathsf U_N\to2^{H_C}\) generated by the
pointwise procedure. Its realized range
\[
\mathscr R_R:=\{R(u):u\in\mathsf U_N\}
\tag{5}
\]
is countable. Equip only this token set with its full power-set sigma-field.
Then \(R:\mathsf U_N\to\mathscr R_R\) is measurable by
Lemma~\ref{lem:step-003-countable-promotion}. The same construction codes
every realized decomposition or leaf-class object. This is the precise
measurable interface required by the algorithm: emptiness, the fixed chosen
decomposition, the sanitized list, and later finite coordinates are
functions of the token and the countable state. We neither assert nor need
that an arbitrary member of \(2^{H_C}\), or the possibly non-Borel class
\(\bar C\), is itself a Borel subset of \(H_C\).

By the standard-Borel closure facts, (2) is standard Borel. The component
\(\ell=0\) is the unique empty list. Length is the measurable disjoint-union
tag, and the \(j\)-th item map on components of length at least \(j\) is a
coordinate projection. Sending lists of length below \(j\) to the isolated
token \(\dagger\) makes this a total measurable
\(\mathsf{List}(H_C)\to H_C^\dagger\) coordinate. The list-membership
relation
\[
\mathsf M:=\{(\lambda,\bar h)\in
\mathsf{List}(H_C)\times H_C:\bar h\text{ occurs in }\lambda\}
\tag{6}
\]
is Borel: on the component \(H_C^\ell\times H_C\), it is
\[
\bigcup_{j=1}^{\ell}
\{((\bar h_1,\ldots,\bar h_\ell),\bar h):\bar h_j=\bar h\},
\tag{7}
\]
a finite union of inverse images of the Borel diagonal of \(H_C\). Taking
the countable union over \(\ell\) proves (6). Proposition~\ref{prop:step-001-totalization} ensures that every list reaching
the measured interface is either a genuine finite \(H_C\)-list or the empty
list. Thus each list coordinate
\(\lambda:\mathsf U_N\to\mathsf{List}(H_C)\) is measurable by
Lemma~\ref{lem:step-003-countable-promotion}, regardless of how the fixed
pointwise source choice varies among atoms.

Invalidity, oversize, nonfiniteness, non-\(H_C\) output, source failure,
empty union, and fallback are predicates on the countable state (4). Hence
each corresponding flag set is measurable before sanitization. After
sanitization, only the empty list, valid finite \(H_C\)-lists, actual
\(H_C\)-items, and the isolated tokens in (3) enter the transcript. This
avoids any need to impose a measurable structure on an untrusted source
object outside \(H_C\).

There are finitely many stages and, on the positive branch, finitely many
teacher/list coordinates. Consequently a complete transcript fits in a
finite product or finite disjoint union of spaces of the following types:
\[
\mathfrak P_N,\quad \mathfrak A_N,\quad
\mathsf{List}(H_C),\quad \mathbb R,\quad [0,1],\quad H_C^\dagger.
\tag{8}
\]
Call one such container \(\mathsf{Tr}_N\). It is standard Borel. Variable
stopping times are represented by a finite disjoint-union tag, not by an
infinite transcript.

For a fixed input atom \(\bar s\), the  totalized procedure gives a
pointwise probability law \(\nu_{\bar s}\) on \(\mathsf{Tr}_N\). Concretely,
the partition has a finite law, AboveThreshold uses finitely many real
noise coordinates, and a Sparse Sample call has a finite categorical law
because its sanitized input lists are finite and its failure-symbol weight
is positive. Concatenate the already fixed ordered lists and retain each
candidate at its first occurrence; this gives a deterministic finite support
order on every countable structural state and hence introduces no global
enumeration of \(H_C\). Comparisons with fixed thresholds are Borel subsets
of the real-noise coordinates; the categorical draw in that finite order can
be realized by consecutive intervals in one \([0,1]\) coordinate. Invalid
or exhausted paths use the constant terminal convention. Thus every
fixed-input transcript is a Borel random element and \(\nu_{\bar s}\) is a
Borel probability measure.

Lemma~\ref{lem:step-003-countable-promotion} now makes
\[
\Gamma_N(\bar s,B):=\nu_{\bar s}(B),
\qquad B\in\mathcal B(\mathsf{Tr}_N),
\tag{9}
\]
a Markov kernel. If \((Y,\mathcal Y)\) is a measurable space and
\(\pi:\mathsf{Tr}_N\to Y\) is any measurable list, flag, event-status,
stage, or terminal-output coordinate, then, for \(A\in\mathcal Y\),
\[
\Gamma_N^\pi(\bar s,A)
:=\Gamma_N(\bar s,\pi^{-1}(A))
\tag{10}
\]
is a kernel: it is a probability measure in \(A\) for fixed \(\bar s\), and
is measurable in \(\bar s\) for fixed \(A\) by (9). This proves all four
claims. For \(N=0\), (4), (8), and (9) reduce to singleton structural and
transcript spaces, so no positive-branch coordinate is silently activated.
\(\square\)


\end{proof}

\begin{proposition}[Measurable list, event, prediction, and error interfaces]\label{prop:step-003-events}
Under Assumption~\ref{assump:countable-evaluation-quotient}, 
Propositions~\ref{prop:step-002-borel} and
\ref{prop:step-002-risk}, and
Proposition~\ref{prop:step-003-coding}, all of the following interfaces are
measurable for every \(N\in\mathbb N_0\):

1. quotient evaluation and joint decoded evaluation;
2. empirical-error and fixed-target population-error events, including
   arbitrary labels, duplicate records, and improper outputs;
3. dynamic finite-list membership, empty-list, list-validity,
   transcript-status, invalid-path, and fallback events; and
4. on the declared block horizon \(N=n_0=km\), the exact source event
   \(E_{\mathrm{good}}(\bar S,\mathcal P)\) and every one of its finite
   stage/block sections.

More generally, any family of Borel output-event sections indexed by a
countable quotient state has a measurable joint graph.


\end{proposition}

\begin{proof}
The quotient evaluation map is jointly measurable. Indeed,
\[
\{(\bar h,q)\in H_C\times Q_C:\bar h(q)=1\}
=\bigcup_{q\in Q_C}
e_q^{-1}(\{1\})\times\{q\},
\tag{11}
\]
a countable union of measurable rectangles. Using the measurable quotient
cells, the joint decoded evaluation is also measurable because
\[
\{(\bar h,x)\in H_C\times X:\bar h(\kappa(x))=1\}
=\bigcup_{q\in Q_C}
e_q^{-1}(\{1\})\times\kappa^{-1}(\{q\}).
\tag{12}
\]
This proves the output-coordinate interface for every
\(\bar h\in H_C\); no membership in \(\bar C\) is required.

For \(\bar s=((q_r,y_r))_{r=1}^N\), define the total empirical-error map
\[
\widehat R_N(\bar s,\bar h):=
\begin{cases}
0,&N=0,\\[2pt]
\displaystyle\frac1N\sum_{r=1}^N
\mathbf1\{\bar h(q_r)\ne y_r\},&N\ge1.
\end{cases}
\tag{13}
\]
For each fixed \(\bar s\), (13) is a finite sum of measurable evaluation
indicators. Since \(S_N\) is countable, for every Borel
\(B\subseteq[0,1]\),
\[
\{(\bar s,\bar h):\widehat R_N(\bar s,\bar h)\in B\}
=\bigcup_{\bar s\in S_N}
\{\bar s\}\times
\{\bar h:\widehat R_N(\bar s,\bar h)\in B\}
\tag{14}
\]
is a countable union of measurable rectangles. Thus (13) is jointly
measurable. Repeated \(q_r\)'s and arbitrary \(y_r\)'s merely repeat or
change finitely many summands and do not affect (14). For fixed
\((\bar D,\bar c)\), Proposition~\ref{prop:step-002-risk} makes
\(\bar h\mapsto\operatorname{err}_{\bar D}(\bar h,\bar c)\) measurable,
so every threshold event for population error is in \(\mathcal H_C\), also
for improper \(\bar h\).

If \(\lambda:\mathsf U_N\to\mathsf{List}(H_C)\) is any generated list
coordinate, then
\[
\{(u,\bar h):\bar h\text{ occurs in }\lambda(u)\}
=\{(u,\bar h):(\lambda(u),\bar h)\in\mathsf M\}
\tag{15}
\]
is measurable by Proposition~\ref{prop:step-003-coding}. The empty-list and
length/cap events are inverse images under measurable list coordinates.
List validity, source validity, Above/Below status, first legal stage,
no-success exhaustion, invalid transcript, Sparse Sample failure, and
fallback are measurable discrete flags in the same proposition. In
particular, the terminal output map sends every fallback flag to the
constant \(\bar c_0\), whose singleton is Borel in the standard-Borel space
\(H_C\).

At the actual block horizon, \(S_N\times\mathfrak P_N\) is countable
discrete. The exact event
\[
\mathsf E_{\mathrm{good}}
:=\{(\bar s,\mathcal P):
E_{\mathrm{good}}(\bar s,\mathcal P)\text{ holds}\}
\tag{16}
\]
is therefore measurable as a subset of that space, as are its finitely many
stage, block, high-error, low-error, and complement-error sections. This
argument does not take a supremum over a nonmeasurable function class and
does not select a representative trace: the exact pointwise predicate in
the setting is evaluated on each countable input/partition atom.

Finally, let \((B_u)_{u\in\mathsf U}\) be any family with
\(B_u\in\mathcal H_C\), where \(\mathsf U\) is countable discrete. Then
\[
\{(u,\bar h):\bar h\in B_u\}
=\bigcup_{u\in\mathsf U}\{u\}\times B_u
\tag{17}
\]
is measurable. Formula (17) covers every remaining pointwise Borel output
event and proves the general assertion. \(\square\)


\end{proof}

\begin{proposition}[Quotient kernel for every totalized pointwise output law]\label{prop:step-003-quotient-kernel}
Under Assumption~\ref{assump:countable-evaluation-quotient}, 
Propositions~\ref{prop:step-001-totalization} and
\ref{prop:step-001-zero},
Lemma~\ref{lem:step-003-countable-promotion}, and
Proposition~\ref{prop:step-003-coding}, fix
\(N\in\mathbb N_0\). Let a totalized pointwise quotient law assign to every
\(\bar s\in Z_Q^N\) a Borel probability measure
\(\mu_{\bar s}\) on the exact output space \((H_C,\mathcal H_C)\). Then
\[
K_N(\bar s,E):=\mu_{\bar s}(E),
\qquad E\in\mathcal H_C,
\tag{18}
\]
is a Markov kernel
\[
K_N:(Z_Q^N,2^{Z_Q^N})\leadsto(H_C,\mathcal H_C).
\tag{19}
\]
This applies to every deterministic or randomized totalization, without a
support-cardinality or propriety restriction. In particular it applies to
the exact totalized quotient-first Lyu~\citep{lyu2025} law, its totalized old-Lyu~\citep{lyu2025} analogue,
and, when defined, the totalized finite-class pointwise law. If \(d=0\),
\(N=0\), (18) is exactly the Dirac law at the unique \(\bar c_0\).
At the public sample size of any setting-defined arm, its notation \(K_C\)
means the corresponding kernel \(K_N\) in (18).


\end{proposition}

\begin{proof}
A pointwise output law is, by definition, a probability measure
\(\mu_{\bar s}\) on \((H_C,\mathcal H_C)\) for each fixed input. For the
procedure exported by Proposition~\ref{prop:step-001-totalization}, this law exists directly as
the terminal-output marginal of the pointwise transcript law constructed in
Proposition~\ref{prop:step-003-coding}. All exact valid paths output
their actual \(H_C\)-item, while invalid, empty, exhausted, or failure paths
output the fixed \(\bar c_0\). Hence no mass is assigned outside \(H_C\),
and the output marginal has total mass one. More explicitly, on the
positive branch its support for fixed \(\bar s\) is contained in
\[
\{\bar c_0\}\ \cup\
\bigcup_{\mathcal P\in\mathfrak P_N}
\bigcup_{r=0}^{d}\bigcup_{i=1}^{k}
\mathcal L_i^r(\bar s,\mathcal P),
\tag{18a}
\]
which is finite by  totalization; continuous mechanism coins change
only the masses on this finite Borel set. A separately supplied totalized
pointwise law needs no transcript representation or finite-support
property; its Borel probability measures \(\mu_{\bar s}\) are already the
complete pointwise input.

Applying Lemma~\ref{lem:step-003-countable-promotion} with \(Y=H_C\) proves
(19): for fixed \(\bar s\), \(E\mapsto K_N(\bar s,E)\) is the probability
measure \(\mu_{\bar s}\), and for fixed \(E\in\mathcal H_C\),
\(\bar s\mapsto K_N(\bar s,E)\) is measurable on the countable discrete
input. This remains true if the pointwise supports vary arbitrarily with
\(\bar s\), are not uniformly finite, or contain improper hypotheses. Thus
the result does not obtain measurability by choosing a global list
enumeration or a selector.

When \(d=0\), the proposition named above,
Proposition~\ref{prop:step-001-zero}, gives
\(\mu_{\varnothing}=\delta_{\bar c_0}\), and \(Z_Q^0\) is a singleton, so
(18) includes the exact empty-input branch. When \(d\ge1\), all quotient
datasets, including duplicate records and arbitrary/nonrealizable labels,
are atoms of \(Z_Q^N\) on which the  procedure is total. Therefore
none is removed from the kernel domain. \(\square\)


\end{proof}

\begin{proposition}[Measurable raw pullback of the quotient kernel]\label{prop:step-003-raw-pullback}
Under Assumption~\ref{assump:countable-evaluation-quotient}, Proposition~\ref{prop:step-002-record-map}, and
Proposition~\ref{prop:step-003-quotient-kernel}, for every
\(N\in\mathbb N_0\) and every quotient kernel \(K_N\) produced there, define
\[
A_N(s,E):=K_N(T_N(s),E),
\qquad s\in Z_X^N,\quad E\in\mathcal H_C.
\tag{20}
\]
Then
\[
A_N:(Z_X^N,\mathcal Z_X^{\otimes N})
\leadsto(H_C,\mathcal H_C)
\tag{21}
\]
is a Markov kernel. Its output is the same quotient hypothesis in \(H_C\);
decoding occurs only after release and no propriety condition is introduced.


\end{proposition}

\begin{proof}
For fixed \(s\in Z_X^N\), \(E\mapsto A_N(s,E)\) is the
probability measure \(K_N(T_N(s),\cdot)\). For fixed
\(E\in\mathcal H_C\), the kernel property of \(K_N\) makes
\(\bar s\mapsto K_N(\bar s,E)\) measurable on \(Z_Q^N\), while Proposition~\ref{prop:step-002-record-map} makes \(T_N\) measurable on the
raw product space. Their composition in (20) is therefore
\(\mathcal Z_X^{\otimes N}\)-measurable. This proves both kernel axioms.

For \(N=0\), both input spaces are singletons and \(T_0\) is the unique
empty-tuple map, so the same argument gives the exact Dirac pullback. For
\(N\ge1\), raw duplicates, same-cell records, and arbitrary labels are all
in the domain of the measurable \(T_N\); no realizability condition is used.
Equation (20) is an exact pullback into \(H_C\), not a kernel on a trace,
list, representative, or raw-selector surrogate. The  neighbor
property of \(T_N\) and all pointwise privacy inequalities are deliberately
unused here.

Proposition~\ref{prop:step-001-totalization} supplies the complete
positive-dimensional pointwise quotient procedure and ensures that every
valid, empty, invalid, failure, exhausted, and fallback path terminates in
\(H_C\). Proposition~\ref{prop:step-001-zero} separately supplies
the exact \(d=0,N=0\) Dirac law at the unique \(\bar c_0\). Proposition~\ref{prop:step-002-borel} supplies the exact countable-discrete
input and standard-Borel output spaces.

Lemma~\ref{lem:step-003-countable-promotion} proves the central promotion
principle: on every \(Z_Q^N\), arbitrary pointwise coordinates are
measurable and arbitrary pointwise Borel probability laws are kernels.
Proposition~\ref{prop:step-003-coding} applies that principle to the actual
source state space, coding realized restrictions and decompositions only in
their countable ranges, coding finite lists in
\(\mathsf{List}(H_C)\), and coding the finite random transcript in a
standard-Borel space. This proves list/transcript measurability without a
hyperspace sigma-field or a selector.
Proposition~\ref{prop:step-003-events} then proves the exact list-membership,
source-event, transcript-status, fallback, prediction, empirical-error, and
population-error interfaces, including \(E_{\mathrm{good}}\) and
potentially improper outputs.

Proposition~\ref{prop:step-003-quotient-kernel} turns every totalized
pointwise \(H_C\)-law into the required quotient Markov kernel \(K_N\), with
the \(d=0,N=0\) Dirac law supplied through Proposition~\ref{prop:step-001-zero} and every arbitrary-label input
included. Finally, Proposition~\ref{prop:step-002-record-map} makes
\(T_N\) measurable, so composition with the quotient kernel proves that
\(A_N=K_N\circ T_N\) is a raw-input learner kernel for every
\(N\in\mathbb N_0\). Together, these propositions establish exactly the
stated conclusion.
The occurrence-mark, privacy, utility, and rate properties are established
in the subsequent named results.



\end{proof}

\subsection{Marked Kernel and Released Projection}\label{app:step_004}

Proposition~\ref{prop:step-003-quotient-kernel} realizes the exact released
law as the terminal-output marginal of the complete transcript kernel: for
every quotient input \(\bar s\in Z_Q^N\) and every
\(E\in\mathcal H_C\),
\[
 K_C(\bar s,E)
 =\int_{\mathsf{Tr}_N}
   \mathbf1_E(\bar H(\tau))\,\Gamma_N(\bar s,d\tau).
\tag{1}
\]
This identity uses the same transcript coordinate that is marked below; no
new coupling or output law is introduced.


 \begin{lemma}[Measurable finite producer-block occurrence set]\label{lem:step-004-occurrence}
Under  Propositions~\ref{prop:step-003-coding},
\ref{prop:step-003-events}, and
\ref{prop:step-003-quotient-kernel}, fix the declared branch \(d\geq1\),
its finite teacher count \(k\geq2\), and its sample size \(N=n_0\). There
are measurable transcript coordinates
\[
\lambda_{i,r}:\mathsf{Tr}_N\longrightarrow
\mathsf{List}(H_C),
\qquad i\in[k],\quad r\in\{0,\ldots,d\},
\tag{2}
\]
the terminal output \(\bar H:\mathsf{Tr}_N\to H_C\), and a measurable
actual-output event \(\mathsf{Act}\subseteq\mathsf{Tr}_N\). Define the
all-stage producer list by concatenation,
\[
\lambda_i(\tau)
:=\lambda_{i,0}(\tau)\mathbin{\Vert}\cdots
  \mathbin{\Vert}\lambda_{i,d}(\tau),
\tag{3}
\]
and define
\[
a_i(\tau):=
\mathbf1\{\bar H(\tau)\text{ occurs in }\lambda_i(\tau)\},
\quad
I(\tau):=\{i\in[k]:a_i(\tau)=1\},
\quad
R(\tau):=\sum_{i=1}^k a_i(\tau).
\tag{4}
\]
Then every \(a_i\) and \(R\) is measurable, \(I(\tau)\) is finite for
every transcript, and
\[
\tau\in\mathsf{Act}\quad\Longrightarrow\quad
1\leq R(\tau)=|I(\tau)|\leq k.
\tag{5}
\]
The set \(I(\tau)\) is exactly the setting-defined set of producer-block
coordinates whose all-stage candidate union contains the actual output.


\end{lemma}

\begin{proof}
Proposition~\ref{prop:step-003-coding} supplies every
coordinate in (2), including a sanitized empty list when a stage/block list
is absent or invalid. Concatenation in (3) is Borel. To see this directly,
restrict
\(\mathsf{List}(H_C)^{d+1}\) to a component with length vector
\((\ell_0,\ldots,\ell_d)\). On that component concatenation is merely the
coordinate map
\[
H_C^{\ell_0}\times\cdots\times H_C^{\ell_d}
\longrightarrow H_C^{\ell_0+\cdots+\ell_d}
\]
that preserves the displayed order. It is measurable. Taking the countable
disjoint union over all finite length vectors proves measurability on the
whole list space. This derivation keeps every list position, including
repeated positions; it makes no first-occurrence selection.

Proposition~\ref{prop:step-003-events} makes dynamic membership of
the measurable terminal output in a measurable finite-list coordinate a
Borel event. Therefore each \(a_i\) is a measurable
\(\{0,1\}\)-valued function. The finite sum \(R\) is measurable with range
\(\{0,1,\ldots,k\}\), and (4) makes \(I(\tau)\) a subset of the finite set
\([k]\).

It remains to prove nonemptiness on actual paths rather than assume it.
By the exact transcript construction exported by Proposition~\ref{prop:step-003-coding}, a transcript in
\(\mathsf{Act}\) records a selected stage \(r_*(\tau)\) and an actual
Sparse Sample categorical outcome. That categorical outcome is one of the
items in the sanitized tuple
\((\lambda_{1,r_*}(\tau),\ldots,
\lambda_{k,r_*}(\tau))\). Hence some \(i\in[k]\) satisfies
\(\bar H(\tau)\in\lambda_{i,r_*}(\tau)\). The selected-stage list is a
sublist of the concatenation (3), so \(a_i(\tau)=1\), proving the lower
bound in (5). The upper bound follows from \(I(\tau)\subseteq[k]\).

This reasoning is pathwise and retains adaptive selection: the realized
\(r_*\), lists, and output are read from the same transcript. If the output
occurs in several producer blocks, every corresponding block index remains
in \(I(\tau)\). Repetitions across stages or positions inside one block are
also retained in (3), but they intentionally map to that one block mark
\(i\), because the setting's mark records producer blocks rather than list
positions. If all lists are empty, (4) gives \(R=0\), and the contrapositive
of (5) shows that the transcript is nonactual. A failure or fallback output
may happen to equal \(\bar c_0\) appearing in a list; its status remains
nonactual and is handled by mark \(0\) below. \(\square\)


\end{proof}

\begin{lemma}[Uniform pathwise occurrence-mark kernel]\label{lem:step-004-mark-kernel}
Under Lemma~\ref{lem:step-004-occurrence}, put
\(\mathsf M_k:=\{0,1,\ldots,k\}\) with its power-set sigma-field and define,
with the convention that the second line is zero off \(\mathsf{Act}\),
\[
w_0(\tau):=\mathbf1_{\mathsf{Act}^c}(\tau),
\qquad
w_i(\tau):=
\mathbf1_{\mathsf{Act}}(\tau)
\frac{a_i(\tau)}{R(\tau)},\qquad i\in[k].
\tag{6}
\]
Then
\[
W(\tau,B):=\sum_{i\in B}w_i(\tau),
\qquad B\subseteq\mathsf M_k,
\tag{7}
\]
is a Markov kernel
\(W:\mathsf{Tr}_N\leadsto\mathsf M_k\). On an actual path it is uniform
on \(I(\tau)\), assigns no mass to \(0\), and retains every distinct
producer-block occurrence. On every fallback or otherwise nonactual path it
is \(\delta_0\), even when the terminal fallback value also occurs in a
candidate list.


\end{lemma}

\begin{proof}
Lemma~\ref{lem:step-004-occurrence} makes
\(\mathsf{Act}\), every \(a_i\), and \(R\) measurable, and (5) guarantees
that division in (6) is only performed when \(R\in\{1,\ldots,k\}\).
Equivalently, define the displayed quotient to be zero at \(R=0\); this is
a measurable finite-valued function. Thus every \(w_i\) is measurable.

If \(\tau\notin\mathsf{Act}\), then \(w_0=1\) and all positive weights are
zero. If \(\tau\in\mathsf{Act}\), then \(w_0=0\), and
\[
\sum_{i=1}^k w_i(\tau)
=\frac{\sum_{i=1}^k a_i(\tau)}{R(\tau)}=1.
\tag{8}
\]
Hence the weights are nonnegative and sum to one on every path. For fixed
\(\tau\), (7) is consequently a probability measure on the finite space
\(\mathsf M_k\); for fixed \(B\), it is a finite sum of measurable
functions. These are the two kernel axioms.

When \(I(\tau)=\{i\}\), (6) assigns mass one to \(i\). When
\(I(\tau)=[k]\), it assigns mass \(1/k\) to every producer. In particular,
at the boundary \(k=2\), the only actual-path denominators are \(1\) and
\(2\). If there is no occurrence, Lemma~\ref{lem:step-004-occurrence}
rules out an actual path, so (7) assigns mark \(0\). Empty-list, invalid,
no-success, and Sparse Sample failure paths are all nonactual and obey the
same rule. No representative producer is selected; the kernel splits mass
over all producer coordinates. \(\square\)


\end{proof}

\begin{proposition}[Marked lift of the exact quotient output law]\label{prop:step-004-lift}
Under Lemma~\ref{lem:step-003-countable-promotion}, Proposition~\ref{prop:step-003-coding}, and
Lemma~\ref{lem:step-004-mark-kernel}, for every
\(\bar s\in Z_Q^N\) and every
\(B\in\mathcal H_C\otimes2^{\mathsf M_k}\), let
\[
B_i:=\{\bar h\in H_C:(\bar h,i)\in B\},
\qquad i\in\mathsf M_k,
\tag{9}
\]
and define
\[
\widetilde K_C(\bar s,B)
:=\sum_{i=0}^k
\int_{\mathsf{Tr}_N}
\mathbf1_{B_i}(\bar H(\tau))w_i(\tau)
\,\Gamma_N(\bar s,d\tau).
\tag{10}
\]
Then
\[
\widetilde K_C:
(Z_Q^N,2^{Z_Q^N})
\leadsto
(H_C\times\mathsf M_k,
 \mathcal H_C\otimes2^{\mathsf M_k})
\tag{11}
\]
is a Markov kernel. It is the pathwise lift of the exact output law, not a
new released mechanism.


\end{proposition}

\begin{proof}
Since \(\mathsf M_k\) is finite and discrete, every
\(B\in\mathcal H_C\otimes2^{\mathsf M_k}\) decomposes uniquely as
\[
B=\bigcup_{i=0}^k B_i\times\{i\},
\tag{12}
\]
with each \(B_i\in\mathcal H_C\). The terminal map \(\bar H\) and every
\(w_i\) are measurable, so each integrand in (10) is bounded and
measurable.

Fix \(\bar s\). Nonnegativity is immediate. Equation (8) gives
\[
\widetilde K_C(\bar s,H_C\times\mathsf M_k)
=\int\sum_{i=0}^k w_i(\tau)\,
  \Gamma_N(\bar s,d\tau)=1.
\tag{13}
\]
If \((B^{(n)})_{n\geq1}\) are pairwise disjoint marked events, their
sections \((B_i^{(n)})_{n\geq1}\) are pairwise disjoint for every fixed
\(i\). Applying monotone convergence inside each of the finitely many
nonnegative integrals gives
\[
\widetilde K_C\!\left(\bar s,\bigcup_{n\geq1}B^{(n)}\right)
=\sum_{n\geq1}\widetilde K_C(\bar s,B^{(n)}).
\tag{14}
\]
Thus \(B\mapsto\widetilde K_C(\bar s,B)\) is a probability measure.

For fixed \(B\), (10) assigns one real number in \([0,1]\) to every
\(\bar s\in Z_Q^N\). Lemma~\ref{lem:step-003-countable-promotion} makes this coordinate
measurable on the countable-discrete quotient input. Together with the
fixed-input probability property, this proves both Markov-kernel axioms in
(11). In particular, for every \(E\in\mathcal H_C\),
\[
\widetilde K_C(\bar s,E\times\{i\})
=\int\mathbf1_E(\bar H(\tau))w_i(\tau)
\,\Gamma_N(\bar s,d\tau).
\tag{15}
\]
Formula (10) first marks each full internal path and only then integrates it
out. Therefore two paths with the same terminal \(\bar h\) but different
partitions, selected stages, repeated occurrences, or occurrence sets keep
their own weights; no conditional law given only \(\bar h\) and no global
selector is introduced. \(\square\)


\end{proof}

\begin{proposition}[Exact released projection and zero privacy residual]\label{prop:step-004-projection}
Under Proposition~\ref{prop:step-003-quotient-kernel} and
Proposition~\ref{prop:step-004-lift}, let
\(\pi_H:H_C\times\mathsf M_k\to H_C\) be the first-coordinate projection.
Then, for every quotient input \(\bar s\) and every
\(E\in\mathcal H_C\),
\[
(\pi_{H\#}\widetilde K_C)(\bar s,E)
=\sum_{i=0}^k
  \widetilde K_C(\bar s,E\times\{i\})
=K_C(\bar s,E).
\tag{16}
\]
Consequently the analysis-only mark changes neither the released output law
nor any privacy parameter of that released law. No privacy assertion is made
for the unreleased pair \((\bar H,J)\).


\end{proposition}

\begin{proof}
Sum (15) over the finite mark space and apply (8):
\[
\begin{aligned}
\sum_{i=0}^k
\widetilde K_C(\bar s,E\times\{i\})
&=\int
\mathbf1_E(\bar H(\tau))
\left(\sum_{i=0}^k w_i(\tau)\right)
\Gamma_N(\bar s,d\tau)\\
&=\int\mathbf1_E(\bar H(\tau))
\Gamma_N(\bar s,d\tau)\\
&=K_C(\bar s,E),
\end{aligned}
\tag{17}
\]
where the last equality is the exact terminal-marginal identity (1). This
proves (16), with no approximation, exceptional event, or residual mass.

For completeness, for any two quotient inputs \(\bar s,\bar s'\), any
\(\varepsilon,\delta\geq0\), and any released event \(E\), (16) gives the
identity
\[
\begin{aligned}
&(\pi_{H\#}\widetilde K_C)(\bar s,E)
-e^\varepsilon(\pi_{H\#}\widetilde K_C)(\bar s',E)-\delta\\
&\hspace{35mm}=
K_C(\bar s,E)-e^\varepsilon K_C(\bar s',E)-\delta.
\end{aligned}
\tag{18}
\]
Thus every future privacy inequality for the actually released kernel is
literally unchanged: the mark consumes zero privacy slack. Equation (18)
does not claim that privacy has already been proved, and it does not claim
privacy if \(J\) were released.

The same zero-residual statement holds on the raw interface by direct
substitution: for every raw input \(s\in Z_X^N\),
\[
\sum_{i=0}^k
\widetilde K_C(T_N(s),E\times\{i\})
=K_C(T_N(s),E).
\tag{19}
\]
Thus a later DP comparison on raw neighboring inputs sees exactly the same
two released-event probabilities before and after this analysis-only lift;
no marked raw kernel or additional privacy theorem is being asserted here.

The status split is essential. On an actual path, \(w_0=0\), including when
there is exactly one occurrence, all \(k\) occurrences, or \(k=2\). On a
nonactual path, \(w_0=1\), including when the fallback value \(\bar c_0\)
coincides with an item in one or more lists. Hence mark \(0\) is used only
for fallback/nonactual paths, and no fallback mass leaks into a positive
producer mark. 
Proposition~\ref{prop:step-003-coding} supplies the exact
standard-Borel transcript kernel, every stage/block list, the adaptive
selected-stage status, the actual/fallback split, and the terminal output.
Proposition~\ref{prop:step-003-events} supplies measurable dynamic
membership and status events. Lemma~\ref{lem:step-004-occurrence} combines
those  interfaces into the exact pathwise occurrence set from the
setting and proves, rather than assumes, that every actual output has a
nonempty finite set of producer coordinates.

Lemma~\ref{lem:step-004-mark-kernel} gives the unique required support rule:
uniform mass on all positive occurrence coordinates for an actual path and
unit mass at \(0\) for every nonactual path. It explicitly covers repeated
occurrences, adaptively selected stages, empty lists, one/all/no occurrence,
failure/fallback, and \(k=2\), without choosing one producer.
Proposition~\ref{prop:step-004-lift} integrates those pathwise weights
against the  transcript kernel and verifies countable additivity,
unit mass, and input-event measurability on the exact codomain
\(H_C\times\{0,\ldots,k\}\).

Finally, Proposition~\ref{prop:step-003-quotient-kernel} identifies
the unmarked terminal marginal with the exact released \(K_C\), and
summing the marked weights above proves the finite-sum identity
\[
\sum_{i=0}^k
\widetilde K_C(\bar s,E\times\{i\})=K_C(\bar s,E)
\]
for every input and every measurable released event. The same proposition
shows that the unreleased mark introduces exactly zero output-law and
privacy residual. Thus the mark is an analysis variable only; it adds no
selector, algorithmic branch, or released output coordinate.

The setting defines \(k\) and the occurrence analysis only on \(d\geq1\).
On the \(d=0,N=0\) branch, Proposition~\ref{prop:step-001-zero} gives the exact Dirac
release at \(\bar c_0\); the canonical inactive convention
\(\mathsf M_0=\{0\}\) and
\(\widetilde K_C(\varnothing,E\times\{0\})=K_C(\varnothing,E)\) satisfies
the same projection identity. It activates no positive-dimensional list,
stage, or teacher coordinate and preserves the no-data baseline.



\end{proof}

\subsection{VC Trace Counting}\label{app:step_005}


\begin{lemma}[Null and positive quotient-VC branches]\label{lem:step-005-branches}
Under Assumptions~\ref{assump:finite-littlestone} and
\ref{assump:countable-evaluation-quotient}, Proposition~\ref{prop:step-002-factorization},
Lemma~\ref{lem:step-002-vc}, and Lemma~\ref{lem:step-002-ld},
for every indexed labeled quotient sample
\(\bar S=((q_r,y_r))_{r=1}^{n_0}\), define
\[
\mathcal E_{\bar C}(\bar S)
=\{(\mathbf1\{\bar c(q_r)\ne y_r\})_{r=1}^{n_0}:\bar c\in\bar C\},
\]
and define the indexed-sample growth function by
\[
\Pi_{\bar C}(n)
:=\sup_{(q_1,\ldots,q_n)\in Q_C^n}
\left|\{(\bar c(q_r))_{r=1}^n:\bar c\in\bar C\}\right|,
\qquad n\in\mathbb N_0,
\]
with the unique empty-vector convention when \(n=0\). Then
\[
0\le v\le d.
\]
Moreover, the following are equivalent:
\[
v=0,\qquad d=0,\qquad |\bar C|=1.
\tag{1}
\]
In this case, for every \(n_0\in\mathbb N_0\) and every indexed labeled
sample \(\bar S\in Z_Q^{n_0}\),
\[
|\mathcal E_{\bar C}(\bar S)|=1.
\tag{2}
\]
If \(d\ge1\), then \(1\le v\le d\).


\end{lemma}

\begin{proof}
If \(A=\{q_1,\ldots,q_t\}\subseteq Q_C\) is shattered by
\(\bar C\), label every node at level \(r\) of a complete binary tree of
depth \(t\) by \(q_r\).  Every root-to-leaf bit string is a labeling of
\(A\), so shattering of \(A\) supplies a quotient concept realizing that
path.  Thus this tree is shattered by \(\bar C\).  Consequently every
finite VC witness of size \(t\) gives a Littlestone witness of depth \(t\),
and the  identities imply
\[
v=\operatorname{VC}(\bar C)
\le \operatorname{LD}(\bar C)=d<\infty.
\tag{3}
\]

The class \(\bar C\) is nonempty by the  factorization proposition.
If it contained distinct \(\bar c_0,\bar c_1\), they would disagree at some
\(q\in Q_C\).  Because the labels are binary, these two functions would
realize both labels on \(\{q\}\), so \(\{q\}\) would be shattered.  Hence
\(v=0\) forces \(|\bar C|=1\).  Conversely, a singleton binary class cannot
shatter a point or a depth-one tree, so it has \(v=d=0\).  The 
Littlestone quotient lemma supplies the remaining implication
\(d=0\Rightarrow|\bar C|=1\), proving (1).

When \(\bar C\) is a singleton, the definition of
\(\mathcal E_{\bar C}(\bar S)\) contains exactly the one vector generated by
its unique concept.  This remains true for \(n_0=0\), when that vector is
the empty vector, and proves (2).  Finally, (1) and (3) show that \(d\ge1\)
implies \(1\le v\le d\). \(\square\)


\end{proof}

\begin{proposition}[Indexed xor identity and duplicate-support collapse]\label{prop:step-005-xor}
Under Proposition~\ref{prop:step-002-factorization}, fix any
\(n_0\in\mathbb N_0\) and any indexed labeled quotient sample
\[
\bar S=((q_r,y_r))_{r=1}^{n_0}.
\]
Define its prediction-trace family and distinct quotient support by
\[
\mathcal P_{\bar C}(\bar S)
:=\{(\bar c(q_r))_{r=1}^{n_0}:\bar c\in\bar C\},
\qquad
A_{\bar S}:=\{q_r:1\le r\le n_0\}.
\tag{4}
\]
Then coordinatewise xor with the fixed label vector is a bijection
\[
\chi_{\mathbf y}:\mathcal P_{\bar C}(\bar S)
\longrightarrow\mathcal E_{\bar C}(\bar S),
\qquad
\chi_{\mathbf y}(z)_r=z_r\oplus y_r,
\tag{5}
\]
and replication along the indexed sample is a bijection
\[
\rho_{\bar S}:\bar C|_{A_{\bar S}}
\longrightarrow\mathcal P_{\bar C}(\bar S),
\qquad
\rho_{\bar S}(f)=(f(q_r))_{r=1}^{n_0}.
\tag{6}
\]
Consequently,
\[
|\mathcal E_{\bar C}(\bar S)|
=|\mathcal P_{\bar C}(\bar S)|
=\left|\bar C|_{A_{\bar S}}\right|.
\tag{7}
\]
These identities hold with arbitrary repeated quotient points, exact
duplicate labeled records, or different labels attached to repeated points.


\end{proposition}

\begin{proof}
For bits \(z_r,y_r\),
\(z_r\oplus y_r=\mathbf 1\{z_r\ne y_r\}\).  Therefore (5) sends the
prediction vector of each \(\bar c\) to exactly its error vector.  On the
whole cube \(\{0,1\}^{n_0}\), xor with \(\mathbf y=(y_r)_{r=1}^{n_0}\) is
an involution:
\[
\chi_{\mathbf y}(\chi_{\mathbf y}(z))=z.
\]
Its restriction in (5) is thus bijective onto the error-trace family.

The map in (6) is surjective by the definition of the prediction-trace
family.  If \(f,g\in\bar C|_{A_{\bar S}}\) are distinct, then they differ at
some \(q\in A_{\bar S}\).  By definition of the support, \(q=q_r\) for at
least one index \(r\), and the corresponding replicated coordinates differ.
Thus \(\rho_{\bar S}\) is injective, proving (7).

Repeated points only repeat coordinates under \(\rho_{\bar S}\) and cannot
create another trace.  The label vector is fixed before xor is applied, so
even if \(q_r=q_t\) but \(y_r\ne y_t\), the global involution argument still
preserves cardinality exactly.  When \(n_0=0\), the support is empty and all
three sets in (7) contain the unique empty function or empty vector.
\(\square\)


\end{proof}

\begin{lemma}[Sauer--Shelah~\citep{sauer1972} counting on a finite quotient restriction]\label{lem:step-005-sauer}
Under Lemma~\ref{lem:step-005-branches}, Lemma~\ref{lem:step-002-vc}, and
Proposition~\ref{prop:step-005-xor}, let \(\bar S\) be any fixed indexed
labeled quotient sample of length \(n_0\), and put
\(s:=|A_{\bar S}|\le n_0\).  Then
\[
|\mathcal E_{\bar C}(\bar S)|
\le \sum_{j=0}^{\min\{v,s\}}\binom{s}{j}
\le \sum_{j=0}^{\min\{v,n_0\}}\binom{n_0}{j}.
\tag{8}
\]
Equivalently, for the setting's growth function (defined on indexed
samples),
\[
|\mathcal E_{\bar C}(\bar S)|
\le \Pi_{\bar C}(n_0)
\le \sum_{j=0}^{\min\{v,n_0\}}\binom{n_0}{j}.
\tag{9}
\]
The conclusion does not require \(\bar C\) to be finite.


\end{lemma}

\begin{proof}
We first prove the finite-cube form of Sauer--Shelah~\citep{sauer1972}.  For an
integer \(t\ge0\), a class \(\mathcal F\subseteq\{0,1\}^B\) on a
\(t\)-point set \(B\), and an integer \(r\ge0\), suppose
\(\operatorname{VC}(\mathcal F)\le r\).  We claim
\[
|\mathcal F|\le S(t,r):=\sum_{j=0}^{r}\binom{t}{j},
\tag{10}
\]
where \(\binom{t}{j}=0\) for \(j>t\), and set \(S(t,-1):=0\).

The proof is by induction on \(t\).  At \(t=0\), the Boolean cube contains
only the empty function, so (10) holds.  For \(t\ge1\), fix \(x\in B\),
write \(B'=B\setminus\{x\}\), and define
\[
\mathcal F_b:=\{f|_{B'}:f\in\mathcal F,\ f(x)=b\},\quad b\in\{0,1\},
\]
\[
\mathcal U:=\mathcal F_0\cup\mathcal F_1,
\qquad
\mathcal I:=\mathcal F_0\cap\mathcal F_1.
\tag{11}
\]
Partitioning \(\mathcal F\) by its value at \(x\), followed by the finite
set identity \(|F_0|+|F_1|=|F_0\cup F_1|+|F_0\cap F_1|\), gives
\[
|\mathcal F|=|\mathcal U|+|\mathcal I|.
\tag{12}
\]
Any set shattered by \(\mathcal U\) is shattered by \(\mathcal F\), so
\(\operatorname{VC}(\mathcal U)\le r\).  If a set \(D\subseteq B'\) is
shattered by \(\mathcal I\), then for every labeling of \(D\), membership
in both \(\mathcal F_0\) and \(\mathcal F_1\) supplies both extensions at
\(x\).  Hence \(D\cup\{x\}\) is shattered by \(\mathcal F\), and
\(\operatorname{VC}(\mathcal I)\le r-1\).  For \(r=0\), this says
\(\mathcal I=\varnothing\): otherwise \(\{x\}\) would be shattered.
The induction hypothesis and Pascal's identity now yield
\[
|\mathcal F|
\le S(t-1,r)+S(t-1,r-1)
=S(t,r),
\tag{13}
\]
which proves (10).

Apply (10) to the restriction family
\[
\mathcal F_{\bar S}:=\bar C|_{A_{\bar S}}
\subseteq\{0,1\}^{A_{\bar S}}.
\]
This family is finite because the ambient cube has \(2^s\) elements, even
when \(\bar C\) is infinite.  If a subset of \(A_{\bar S}\) is shattered
by \(\mathcal F_{\bar S}\), it is shattered by \(\bar C\); the 
VC identity therefore gives
\(\operatorname{VC}(\mathcal F_{\bar S})\le v\).  Equations (7) and (10)
give the first inequality in (8).  Since \(s\le n_0\), for every
\(0\le j\le s\) one has \(\binom{s}{j}\le\binom{n_0}{j}\); adding the
remaining nonnegative terms proves the second inequality.

For any indexed \(n_0\)-sample, Proposition~\ref{prop:step-005-xor} shows
that xor preserves the prediction-trace count.  Taking the supremum over
the quotient-point sequences gives (9).  No concept representative is
selected for a trace, and no cardinality of \(\bar C\) enters the proof.
\(\square\)


\end{proof}

\begin{lemma}[Positive-dimensional Sauer-sum compression]\label{lem:step-005-binomial}
Under Lemma~\ref{lem:step-005-sauer}, if \(1\le v\le n_0\), then every fixed
indexed labeled quotient sample \(\bar S\) of length \(n_0\) satisfies
\[
|\mathcal E_{\bar C}(\bar S)|
\le \sum_{j=0}^{v}\binom{n_0}{j}
\le \left(\frac{e n_0}{v}\right)^v.
\tag{14}
\]
This statement includes \(v=1\), remains unchanged when \(v=d\), and uses a
nonstrict premise \(v\le n_0\), so it includes \(n_0=v\).


\end{lemma}

\begin{proof}
Put \(x:=v/n_0\in(0,1]\).  Since \(x^j\ge x^v\) for
\(0\le j\le v\), the binomial theorem gives
\[
(1+x)^{n_0}
=\sum_{j=0}^{n_0}\binom{n_0}{j}x^j
\ge x^v\sum_{j=0}^{v}\binom{n_0}{j}.
\tag{15}
\]
Using \(1+x\le e^x\) and \(n_0x=v\),
\[
\sum_{j=0}^{v}\binom{n_0}{j}
\le x^{-v}(1+x)^{n_0}
\le \left(\frac{n_0}{v}\right)^v e^{n_0x}
=\left(\frac{e n_0}{v}\right)^v.
\tag{16}
\]
Together with (8), this proves (14).

At \(v=1\), (14) reads
\(|\mathcal E_{\bar C}(\bar S)|\le1+n_0\le e n_0\) for \(n_0\ge1\).
At \(n_0=v\), the exact Sauer sum is
\(\sum_{j=0}^v\binom{v}{j}=2^v\le e^v\), which is precisely the right side
of (14).  If \(v=d\), substitution in (14) gives
\((e n_0/d)^d\) with no extra factor or exponent; no class-cardinality term
appears in any case. \(\square\)


\end{proof}

\begin{proposition}[The positive branch lies in the Sauer range]\label{prop:step-005-positive-range}
Under Assumptions~\ref{assump:finite-littlestone} and
\ref{assump:approximate-dp-regime}, Lemma~\ref{lem:step-001-calibration}, Proposition~\ref{prop:step-001-teacher}, and
Lemma~\ref{lem:step-005-branches}, suppose \(d\ge1\) and fix the setting's
stipulated sufficiently large universal block constant so that, in
particular, \(C_{\mathrm{blk}}\ge1\).  Then the preceding calibration fixes
\[
a(k)=v+\log(4k/\beta),\qquad
Q(k)=e+\frac{e k d^2a(k)}{\alpha v},
\]
\[
m=m(k)=\left\lceil C_{\mathrm{blk}}
\frac{d^2}{\alpha}a(k)\log Q(k)\right\rceil,
\qquad n_0=km(k),
\]
with \(k\ge2\), and this realized tuple satisfies \(1\le v\le n_0\).


\end{proposition}

\begin{proof}
Lemma~\ref{lem:step-001-calibration} gives \(1\le v\le d\) and
the displayed candidate dictionary for every integer \(t\ge2\).
Proposition~\ref{prop:step-001-teacher} fixes the least feasible
\(k\ge2\) and the realized identities \(m=m(k)\) and \(n_0=km(k)\).
These  conclusions agree with the positive branch isolated by
Lemma~\ref{lem:step-005-branches}.  Since \(0<\beta<1/4\) and
\(0<\alpha<1/4\), substitution of \(t=k\) in the  dictionary gives
\[
a(k)>v,\qquad Q(k)>e,\qquad
\log Q(k)>1,\qquad \alpha^{-1}>4.
\]
Since \(d\ge1\) and \(C_{\mathrm{blk}}\ge1\),
\[
m(k)
\ge C_{\mathrm{blk}}\frac{d^2}{\alpha}a(k)\log Q(k)
>4v.
\tag{17}
\]
Thus \(n_0=km(k)\ge m(k)\ge v\).  Requiring
\(C_{\mathrm{blk}}\ge1\) does not add a new parameter assumption: it is one
of the finite lower bounds included in the setting's choice of a
"sufficiently large universal" block constant. 
The zero-dimensional case is separate.  By
Lemma~\ref{lem:step-005-branches}, \(v=0\) is exactly the singleton,
\(d=0\) case, and
\(|\mathcal E_{\bar C}(\bar S)|=1\), including for the empty sample.  The
setting's no-data learner therefore bypasses the positive-dimensional trace
union.  The undefined expression \((e n_0/v)^v\) is never invoked at
\(v=0\).

On the positive branch, Lemma~\ref{lem:step-001-calibration} supplies the
exact legal dictionary, Proposition~\ref{prop:step-001-teacher} fixes the
realized \(k,m,n_0\), and the positive-range calculation above derives
\(1\le v\le n_0\) from those  outputs.  Proposition~\ref{prop:step-005-xor}
then gives the exact chain
\[
|\mathcal E_{\bar C}(\bar S)|
=|\mathcal P_{\bar C}(\bar S)|
=\left|\bar C|_{A_{\bar S}}\right|.
\tag{18}
\]
Thus neither xor nor repeated quotient records cost a multiplicative factor.
Lemma~\ref{lem:step-005-sauer} bounds this finite restriction family, and
Lemma~\ref{lem:step-005-binomial} compresses the resulting sum.  Hence
\[
|\mathcal E_{\bar C}(\bar S)|
\le \Pi_{\bar C}(n_0)
\le \sum_{j=0}^{v}\binom{n_0}{j}
\le \left(\frac{e n_0}{v}\right)^v,
\tag{19}
\]
which is the exact VC-trace bound.

The proof of Lemma~\ref{lem:step-005-sauer} operates only on the finite set
of restrictions inside \(\{0,1\}^{A_{\bar S}}\).  It therefore applies when
\(\bar C\) is infinite and does not select or count concepts.  The boundary
calculations in Lemma~\ref{lem:step-005-binomial} explicitly retain
\(v=1\), \(v=d\), and \(n_0=v\), while
Proposition~\ref{prop:step-005-xor} covers every duplicate-record pattern.
All downstream source lists may continue to contain actual quotient
functions; (19) counts only the exact error vectors used in the finite
trace union.



\end{proof}

\subsection{Fixed-Trace Concentration Event}\label{app:step_006}


The source event uses its internal accuracy parameter at thresholds one
third and one half of that parameter.  Here the parameter is
\(\gamma=\alpha/16\), so the thresholds and relative tolerance below are
the exact source thresholds in the present notation.

For every fixed positive-branch quotient master sample \(\bar S\),
Lemma~\ref{lem:step-005-binomial} and Proposition
~\ref{prop:step-005-positive-range} give the exact finite trace interface
\[
 |\mathcal E_{\bar C}(\bar S)|
 \leq\Pi_{\bar C}(n_0)
 \leq(en_0/v)^v,
 \qquad 1\leq v\leq n_0.
\tag{1}
\]
The only external tail used in the high-error branch is Lyu's
without-replacement proposition~\citep{lyu2025}: for fixed bits
\(x_1,\ldots,x_N\in\{0,1\}\), their mean
\(p=N^{-1}\sum_r x_r\), a uniformly random subset
\(A\subseteq[N]\) of size \(t\), and \(0<\zeta<1\),
\[
 \Pr\!\left[
 \left|t^{-1}\sum_{r\in A}x_r-p\right|>\zeta p
 \right]
 \leq2\exp(-\zeta^2tp/3).
\tag{2}
\]
Here \(N=n_0\), \(t=m\), \(x_r=z_r\),
\(p=\mu(z)\), and \(\zeta=1/(5d)\).  The marginal block law proved next
discharges uniform sampling; \(n_0=km\) and \(k\ge2\) give \(t\le N\),
and \(d\ge1\) gives \(0<\zeta<1\).  Equation (2) is not used for the
near-zero branch.



\begin{proposition}[Marginal law of a uniform labeled block partition]\label{prop:step-006-partition}
Under Assumption~\ref{assump:approximate-dp-regime} and the formalized
setting's positive-branch construction, let \(d\geq1\), \(k\geq2\), let
\(m\) be the exact integer produced by the displayed ceiling, and let
\(n_0=km\).  Fix any indexed master sample
\(\bar S=((q_r,y_r))_{r=1}^{n_0}\).  If
\(\mathcal P=(B_1,\ldots,B_k)\) is uniform over all labeled partitions of
\([n_0]\) with \(|B_i|=m\), then, conditional on \(\bar S\), for every
\(i\in[k]\) and every \(A\subseteq[n_0]\) of size \(m\),

\[
 \Pr_{\mathcal P}[B_i=A\mid\bar S]=\binom{n_0}{m}^{-1}.
\tag{3}
\]

This is a marginal statement only.  The blocks are pairwise disjoint and
jointly cover \([n_0]\), and hence are generally dependent.  In particular,
when \(k=2\), \(B_2=[n_0]\setminus B_1\).


\end{proposition}

\begin{proof}
The partition is drawn independently of the record values, so
conditioning on \(\bar S\) does not change its law.  The number of labeled
equal-size partitions is

\[
 \frac{n_0!}{(m!)^k}.
\]

After fixing \(B_i=A\), the remaining \(n_0-m\) indices can be assigned to
the remaining labeled blocks in

\[
 \frac{(n_0-m)!}{(m!)^{k-1}}
\]

ways.  Their ratio is

\[
 \frac{(n_0-m)!}{(m!)^{k-1}}
 \frac{(m!)^k}{n_0!}
 =\frac{m!(n_0-m)!}{n_0!}
 =\binom{n_0}{m}^{-1},
\]

which proves (3).  Disjointness and coverage are part of the definition of
a partition and give the dependence statements, including the complement
identity for \(k=2\).  No later argument replaces (3) by a joint product
law. \(\square\)


\end{proof}

\begin{lemma}[High-error relative fixed-trace tail]\label{lem:step-006-high-tail}
Under Assumption~\ref{assump:approximate-dp-regime}, the positive-branch
definitions \(d\geq1\), \(\gamma=\alpha/16\), and \(a_d=1/(5d)\),
Proposition~\ref{prop:step-006-partition}, and Lyu~\citep{lyu2025} v1
  Proposition 2, instantiated in (2), fix an arbitrary trace
\(z=(z_r)_{r=1}^{n_0}\in\{0,1\}^{n_0}\) and define

\[
 \mu(z):=\frac1{n_0}\sum_{r=1}^{n_0}z_r,
 \qquad
 \widehat\mu_i(z):=\frac1m\sum_{r\in B_i}z_r.
\tag{4}
\]

If \(\mu(z)>\gamma/3\), then for every \(i\in[k]\),

\[
\begin{aligned}
 &\Pr_{\mathcal P}
 \left[|\widehat\mu_i(z)-\mu(z)|>a_d\mu(z)\mid\bar S\right]\\
 &\qquad\leq
 2\exp\!\left(-\frac{m\mu(z)}{75d^2}\right)
 \leq
 2\exp\!\left(-\frac{m\alpha}{3600d^2}\right).
\end{aligned}
\tag{5}
\]

The good relative interval is closed, so equality at either relative
endpoint is not a failure event in (5).


\end{lemma}

\begin{proof}
Conditional on \(\bar S\), the vector \(z\) is fixed, and
Proposition~\ref{prop:step-006-partition} makes \(B_i\) a uniformly random
size-\(m\) subset of the \(n_0\) indexed coordinates.  In (2), instantiate

\[
 N=n_0,\qquad t=m,\qquad p=\mu(z),\qquad \zeta=a_d=\frac1{5d}.
\]

The conditions \(t\leq N\) and \(0<\zeta<1\) hold because
\(n_0=km\), \(k\geq2\), and \(d\geq1\).  Thus (2) gives

\[
 2\exp\!\left(-\frac{a_d^2m\mu(z)}3\right)
 =2\exp\!\left(-\frac{m\mu(z)}{75d^2}\right).
\]

The branch premise and \(\gamma=\alpha/16\) give
\(\mu(z)>\gamma/3=\alpha/48\).  Hence

\[
 \frac{m\mu(z)}{75d^2}
 >\frac{m\alpha}{75\cdot48\,d^2}
 =\frac{m\alpha}{3600d^2},
\]

which proves the second inequality in (5).  Lyu~\citep{lyu2025}'s proposition controls the
strict event \texttt{greater than the relative tolerance}; its complement is
therefore exactly the closed interval used in the event definition.
\(\square\)


\end{proof}

\begin{lemma}[Finite-population exponential-moment domination]\label{lem:step-006-mgf}
Under Proposition~\ref{prop:step-006-partition}, fix an arbitrary binary
trace \(z\), define \(\mu(z)\) and \(\widehat\mu_i(z)\) by (4), and fix a
block \(i\in[k]\).  Then for every \(\lambda\geq0\), conditional on the
fixed master sample,

\[
 \mathbb E_{\mathcal P}
 \left[e^{\lambda m\widehat\mu_i(z)}\mid\bar S\right]
 \leq \left(1-\mu(z)+\mu(z)e^\lambda\right)^m.
\tag{6}
\]

No independence between selected coordinates or between different blocks
is asserted.


\end{lemma}

\begin{proof}
Write \(N=n_0\), \(a_r=e^{\lambda z_r}\geq0\), and let
\(e_m(a_1,\ldots,a_N)\) be the degree-\(m\) elementary symmetric sum.
By the exact marginal law (3),

\[
\begin{aligned}
 \mathbb E_{\mathcal P}
 \left[e^{\lambda\sum_{r\in B_i}z_r}\mid\bar S\right]
 &=\binom Nm^{-1}
   \sum_{A\subseteq[N]:|A|=m}\prod_{r\in A}a_r\\
 &=\binom Nm^{-1}e_m(a_1,\ldots,a_N).
\end{aligned}
\tag{7}
\]

We prove the symmetric-mean inequality needed to bound (7).  For two
coordinates \(a,b\) and the vector \(\mathbf c\) of all remaining
coordinates,

\[
 e_m(a,b,\mathbf c)
 =e_m(\mathbf c)+(a+b)e_{m-1}(\mathbf c)
  +ab\,e_{m-2}(\mathbf c),
\tag{8}
\]

with the usual conventions \(e_0=1\) and \(e_j=0\) outside the available
degrees.  Replacing \(a,b\) by their average preserves their sum and changes
their product by

\[
 \left(\frac{a+b}{2}\right)^2-ab=\frac{(a-b)^2}{4}\geq0.
\]

Since every elementary symmetric sum of the nonnegative coordinates
\(\mathbf c\) is nonnegative, (8) shows that pairwise averaging cannot
decrease \(e_m\).  To make the limiting step explicit, repeatedly replace a
current maximum coordinate \(a\) and a current minimum coordinate \(b\) by
their average.  The coordinate sum is preserved, while the sum of squares
decreases by

\[
 a^2+b^2-2\left(\frac{a+b}{2}\right)^2
 =\frac{(a-b)^2}{2}.
\]

The nonnegative sum of squares cannot decrease by a fixed positive amount
infinitely often.  Hence the maximum-minus-minimum range tends to zero, and,
because the sum stays fixed, the vector converges to its constant mean
vector.  Continuity of \(e_m\) now gives

\[
 e_m(a_1,\ldots,a_N)
 \leq\binom Nm\left(\frac1N\sum_{r=1}^N a_r\right)^m.
\tag{9}
\]

For the current binary population,

\[
 \frac1N\sum_{r=1}^N a_r
 =\frac1N\sum_{r=1}^N e^{\lambda z_r}
 =1-\mu(z)+\mu(z)e^\lambda.
\]

Substitution in (7) proves (6).  This proof uses only the single-block
marginal law. \(\square\)


\end{proof}

\begin{lemma}[Low-error one-sided fixed-trace tail]\label{lem:step-006-low-tail}
Under Assumption~\ref{assump:approximate-dp-regime}, the positive-branch
definitions \(d\geq1\) and \(\gamma=\alpha/16\), and
Lemma~\ref{lem:step-006-mgf}, fix an arbitrary trace \(z\) with
\(0\leq\mu(z)\leq\gamma/3\).  Then every \(i\in[k]\) satisfies

\[
\begin{aligned}
 \Pr_{\mathcal P}
 [\widehat\mu_i(z)>\gamma/2\mid\bar S]
 &\leq \exp(-m\gamma/30)\\
 &=\exp(-m\alpha/480)\\
 &\leq\exp\!\left(-\frac{m\alpha}{3600d^2}\right).
\end{aligned}
\tag{10}
\]

If \(\mu(z)=0\), the failure event is empty.  The conclusion includes the
branch boundary \(\mu(z)=\gamma/3\).


\end{lemma}

\begin{proof}
Put \(p=\mu(z)\) and \(q=\gamma/2\).  Because
\(0<\alpha<1/4\), one has \(0<q<1\).  If \(p=0\), every coordinate of the
binary trace is zero.  Therefore \(\widehat\mu_i(z)=0\) for every partition,
and the claimed failure probability is zero.

Now suppose \(0<p\leq\gamma/3<q\).  For every \(\lambda>0\), exponential
Markov and Lemma~\ref{lem:step-006-mgf} give

\[
\begin{aligned}
 \Pr_{\mathcal P}[\widehat\mu_i(z)>q\mid\bar S]
 &\leq e^{-\lambda mq}
 \mathbb E_{\mathcal P}
 [e^{\lambda m\widehat\mu_i(z)}\mid\bar S]\\
 &\leq
 \exp\{m[-\lambda q+\log(1-p+pe^\lambda)]\}.
\end{aligned}
\tag{11}
\]

The minimizing value is

\[
 \lambda_*:=\log\frac{q(1-p)}{p(1-q)}>0.
\]

Direct substitution in (11) yields

\[
 \Pr_{\mathcal P}[\widehat\mu_i(z)>q\mid\bar S]
 \leq e^{-mD(q\Vert p)},
\tag{12}
\]

where

\[
 D(q\Vert p):=q\log\frac qp+(1-q)\log\frac{1-q}{1-p}
\tag{13}
\]

is the binary relative entropy.  For fixed \(q\) and \(0<p<q\),

\[
 \frac{\partial}{\partial p}D(q\Vert p)
 =\frac{p-q}{p(1-p)}<0.
\tag{14}
\]

Consequently (13) is minimized over
\(0<p\leq p_0:=\gamma/3\) at \(p=p_0\).  At this endpoint,
\(q=(3/2)p_0\).  For any \(0<p_0<q<1\), set
\(u=(q-p_0)/(1-p_0)\).  The elementary inequality

\[
 \log(1-u)\geq-\frac{u}{1-u},\qquad 0\leq u<1,
\tag{15}
\]

follows because the difference between the two sides is zero at \(u=0\)
and has derivative \(u/(1-u)^2\geq0\).  Since
\(1-q=(1-p_0)(1-u)\), (15) gives

\[
 (1-q)\log\frac{1-q}{1-p_0}\geq-(q-p_0).
\tag{16}
\]

It follows that

\[
\begin{aligned}
 D(q\Vert p_0)
 &\geq q\log(q/p_0)-(q-p_0)\\
 &=p_0\left(\frac32\log\frac32-\frac12\right).
\end{aligned}
\tag{17}
\]

For \(x\geq0\),

\[
 \log(1+x)\geq\frac{2x}{2+x},
\tag{18}
\]

because the difference is zero at zero and its derivative is
\(x^2/((1+x)(2+x)^2)\geq0\).  Taking \(x=1/2\) in (18) gives
\(\log(3/2)\geq2/5\).  Thus (17) implies

\[
 D(q\Vert p_0)
 \geq p_0\left(\frac32\frac25-\frac12\right)
 =\frac{p_0}{10}
 =\frac\gamma{30}.
\tag{19}
\]

Equations (12), (14), and (19) prove the first line of (10), including
\(p=p_0=\gamma/3\).  Finally,
\(\gamma/30=\alpha/480\), and \(d\geq1\) gives

\[
 \frac1{480}\geq\frac1{3600d^2}.
\]

This proves the remaining lines of (10). \(\square\)


\end{proof}

\begin{proposition}[Exact simultaneous source trace event]\label{prop:step-006-good-event}
Under Assumption~\ref{assump:approximate-dp-regime}, the setting-defined
positive branch, Lemma~\ref{lem:step-005-binomial},
Proposition~\ref{prop:step-005-positive-range},
Lemma~\ref{lem:step-006-high-tail}, and
Lemma~\ref{lem:step-006-low-tail}, fix an arbitrary labeled quotient master
sample \(\bar S\) and let \(\mathcal P=(B_1,\ldots,B_k)\) be the uniform
labeled partition.  For every
\(z\in\mathcal E_{\bar C}(\bar S)\), use the means in (4), and define

\[
\begin{aligned}
 E_{\mathrm{good}}(\bar S,\mathcal P)
 :=\bigcap_{z\in\mathcal E_{\bar C}(\bar S)}
   \bigcap_{i=1}^k
 \Big(&\{\mu(z)>\gamma/3\}
       \cap\{|\widehat\mu_i(z)-\mu(z)|\leq a_d\mu(z)\}\\
 &\quad\cup
       \{\mu(z)\leq\gamma/3\}
       \cap\{0\leq\widehat\mu_i(z)\leq\gamma/2\}\Big).
\end{aligned}
\tag{20}
\]

Then, with \(c_{\mathrm{tr}}:=1/3600\),

\[
\begin{aligned}
 \Pr_{\mathcal P}
 [E_{\mathrm{good}}(\bar S,\mathcal P)^c\mid\bar S]
 &\leq4k|\mathcal E_{\bar C}(\bar S)|
   \exp\!\left(-c_{\mathrm{tr}}\frac{m\alpha}{d^2}\right)\\
 &\leq4k\Pi_{\bar C}(n_0)
   \exp\!\left(-c_{\mathrm{tr}}\frac{m\alpha}{d^2}\right)\\
 &\leq4k\left(\frac{en_0}{v}\right)^v
   \exp\!\left(-c_{\mathrm{tr}}\frac{m\alpha}{d^2}\right).
\end{aligned}
\tag{21}
\]

The event (20) is exactly Lyu~\citep{lyu2025}'s source event after mapping the source's
internal accuracy to \(\gamma\): for every \(\bar c\in\bar C\), its error
trace gives its master and block empirical errors, and concepts with the
same trace impose the same inequalities.


\end{proposition}

\begin{proof}
For every trace-block pair, define the two branch-qualified
failure events

\[
\begin{aligned}
 F^{\mathrm H}_{z,i}
 &:=\{\mu(z)>\gamma/3\}
   \cap\{|\widehat\mu_i(z)-\mu(z)|>a_d\mu(z)\},\\
 F^{\mathrm L}_{z,i}
 &:=\{\mu(z)\leq\gamma/3\}
   \cap\{\widehat\mu_i(z)>\gamma/2\}.
\end{aligned}
\tag{22}
\]

The lower endpoint in the low branch needs no failure event because every
block mean of a binary trace is nonnegative.  The predicates in (22) split
all possible master means and assign the equality
\(\mu(z)=\gamma/3\) to the low branch.  Hence

\[
 E_{\mathrm{good}}(\bar S,\mathcal P)^c
 \subseteq
 \bigcup_{z\in\mathcal E_{\bar C}(\bar S)}
 \bigcup_{i=1}^k
 \left(F^{\mathrm H}_{z,i}\cup F^{\mathrm L}_{z,i}\right).
\tag{23}
\]

Lemma~\ref{lem:step-006-high-tail} gives

\[
 \Pr(F^{\mathrm H}_{z,i}\mid\bar S)
 \leq2\exp\!\left(-\frac{m\alpha}{3600d^2}\right),
\tag{24}
\]

where the probability is zero if its branch predicate is false.  Likewise,
Lemma~\ref{lem:step-006-low-tail} gives

\[
 \Pr(F^{\mathrm L}_{z,i}\mid\bar S)
 \leq\exp\!\left(-\frac{m\alpha}{3600d^2}\right)
 \leq2\exp\!\left(-\frac{m\alpha}{3600d^2}\right).
\tag{25}
\]

Apply the union bound to (23), retain the four-term envelope per
trace-block pair from (24)--(25), and use the trace estimate supplied by
Lemma~\ref{lem:step-005-binomial} and
Proposition~\ref{prop:step-005-positive-range} in (1).
This proves all three lines of (21).  Only the one-block marginal law (3)
entered (24) and (25).  The proof neither factors a joint probability nor
multiplies block success probabilities, so the dependence of the single
partition's blocks is irrelevant.

For exact object matching, if
\(z_r=\mathbf1\{\bar c(q_r)\ne y_r\}\), then \(\mu(z)\) is precisely
\(\operatorname{err}_{\bar S}(\bar c)\), and
\(\widehat\mu_i(z)\) is precisely
\(\operatorname{err}_{\bar S_i}(\bar c)\).  Thus (20) is neither a trace
representative event nor a changed empirical metric. 
The positive branch fixes the actual integer \(m\) through its ceiling and
sets \(n_0=km\).  Proposition~\ref{prop:step-006-partition} therefore makes
every block, conditional on an arbitrary fixed master sample, a uniform
size-\(m\) sample without replacement.  This remains true at \(k=2\), even
though then the second block is completely determined by the first.

For a fixed exact error trace above \(\gamma/3\),
Lemma~\ref{lem:step-006-high-tail} instantiates the Lyu~\citep{lyu2025} proposition
with the unchanged relative tolerance \(1/(5d)\) and gives exponent
\(m\alpha/(3600d^2)\).  For a trace at or below \(\gamma/3\),
Lemma~\ref{lem:step-006-mgf} derives the without-replacement exponential
moment directly, and Lemma~\ref{lem:step-006-low-tail} turns it into the
one-sided threshold \(\gamma/2\).  That lemma treats mean zero
deterministically and includes the equality boundary \(\mu=\gamma/3\).

Definition (20) and the union-bound calculation (23)--(25) give the exact two-clause
source event simultaneously over the finite trace family produced above and
all \(k\) blocks.  Its finite union retains exactly
the required \(4k\) envelope and proves

\[
 \Pr_{\mathcal P}
 [E_{\mathrm{good}}(\bar S,\mathcal P)^c\mid\bar S]
 \leq4k\Pi_{\bar C}(n_0)
 \exp\!\left(-\frac{m\alpha}{3600d^2}\right).
\tag{26}
\]

No factor \(d+1\) is present: the single event (20) is defined once from
master and block errors and is reused at every fixed stage. No
independence among blocks is used.  This is the conditional trace interface
used by the fixed-point calculation below.



\end{proof}

\subsection{Trace Fixed Point and Confidence Budget}\label{app:step_007}


The deterministic positive-branch dictionary furnished by
Lemma~\ref{lem:step-001-calibration} and
Proposition~\ref{prop:step-001-teacher} is
\[
 \beta_{\rm tr}:=\beta/4,
 \qquad
 a(t):=v+\log(4t/\beta),
 \qquad
 Q(t):=e+\frac{etd^2a(t)}{\alpha v},
\tag{1}
\]
\[
 m(t):=\left\lceil
 C_{\rm blk}\frac{d^2}{\alpha}a(t)\log Q(t)
 \right\rceil,
 \qquad n(t):=tm(t).
\tag{2}
\]
For the least feasible public teacher count \(k\ge2\), set
\[
 a:=a(k),\qquad Q:=Q(k)=Q_{\rm blk},\qquad
 m:=m(k),\qquad n_0:=n(k)=km.
\tag{3}
\]
No sample, partition, trace, event, stage, or mechanism output enters
(1)--(3).  Fix the universal block constant once and for all so that
\[
 C_{\rm blk}>0,
 \qquad
 \frac{C_{\rm blk}}{3600}
 \ge 4+\log(1+C_{\rm blk}).
\tag{4}
\]
For example, \(C_{\rm blk}=144000\) is admissible: the left side is
\(40\), whereas
\(1+144000<2^{18}<e^{18}\) gives
\(4+\log(1+144000)<22\).  This is a universal specialization of the
setting's sufficiently large block constant, not a parameter-dependent
condition.

Proposition~\ref{prop:step-006-good-event} supplies, for every fixed master
sample, the exact conditional estimate
\[
 \Pr_{\mathcal P}
 [E_{\rm good}(\bar S,\mathcal P)^c\mid\bar S]
 \le 4k\left(\frac{en_0}{v}\right)^v
 \exp\!\left(-\frac{m\alpha}{3600d^2}\right).
\tag{5}
\]
This is the only probabilistic input needed for the fixed-point calculation
below; all remaining steps are algebraic.



\begin{lemma}[Exact ceiling-aware trace fixed point]\label{lem:step-007-fixed-point}
Under Assumption~\ref{assump:approximate-dp-regime}, Lemma~\ref{lem:step-001-calibration}, and Proposition~\ref{prop:step-001-teacher}, suppose \(d\ge1\), and use the
exact positive-branch dictionary (1)-(3). Define the appendix-local
universal constant

\[
 C_{\rm fp}:=2+\log(1+C_{\rm blk}).
\tag{6}
\]

Then

\[
 \frac{en_0}{v}\le(1+C_{\rm blk})Q\log Q,
 \qquad
 \log\frac{en_0}{v}\le C_{\rm fp}\log Q,
\tag{7}
\]

and the exact block size also obeys

\[
 m\le(1+C_{\rm blk})
 \frac{d^2}{\alpha}a\log Q.
\tag{8}
\]


\end{lemma}

\begin{proof}
Lemma~\ref{lem:step-001-calibration} gives
\(1\le v\le d\), while the parameter regime gives
\(0<\alpha<1/4\). Moreover \(a=v+\log(4k/\beta)>v\ge1\), and
\(Q>e\), hence \(\log Q>1\). Put

\[
 x:=C_{\rm blk}\frac{d^2}{\alpha}a\log Q.
\tag{9}
\]

The ceiling is represented exactly as

\[
 m=\lceil x\rceil=x+\theta,
 \qquad 0\le\theta<1.
\tag{10}
\]

This includes the integer case \(\theta=0\); no ceiling remainder is
dropped. Since

\[
 \frac{d^2}{\alpha}a\log Q>4,
\tag{11}
\]

we have \(1\le d^2a\log Q/\alpha\). Therefore (10) gives

\[
 m\le
 C_{\rm blk}\frac{d^2}{\alpha}a\log Q+1
 \le(1+C_{\rm blk})\frac{d^2}{\alpha}a\log Q,
\]

which proves (8), including the additive ceiling contribution.

For the fixed point itself, substitute \(n_0=k(x+\theta)\) and separate the
ceiling term:

\[
\begin{aligned}
 \frac{en_0}{v}
 &=C_{\rm blk}\log Q
   \left(\frac{ekd^2a}{\alpha v}\right)
   +\frac{ek\theta}{v}\\
 &\le C_{\rm blk}Q\log Q+Q.
\end{aligned}
\tag{12}
\]

Indeed, the first term uses
\(ekd^2a/(\alpha v)\le Q\). For the second,
\(d^2a/\alpha\ge1\) implies

\[
 \frac{ek\theta}{v}\le\frac{ek}{v}
 \le\frac{ekd^2a}{\alpha v}\le Q.
\tag{13}
\]

Because \(\log Q>1\), the last \(Q\) in (12) is at most \(Q\log Q\),
which proves the first inequality in (7). Taking logarithms and using
\(\log\log Q\le\log Q\) gives

\[
\begin{aligned}
 \log\frac{en_0}{v}
 &\le\log(1+C_{\rm blk})+\log Q+\log\log Q\\
 &\le\bigl(2+\log(1+C_{\rm blk})\bigr)\log Q
 =C_{\rm fp}\log Q.
\end{aligned}
\tag{14}
\]

The last step also uses \(\log Q>1\) to bound the fixed positive constant
\(\log(1+C_{\rm blk})\) by
\(\log(1+C_{\rm blk})\log Q\). Equations (12)-(14) close the fixed point
using the same \(n_0=km\) that appears in \(Q,m\), and the trace count.
\(\square\)


\end{proof}

\begin{lemma}[Public-witness sample envelope]\label{lem:step-007-sample-envelope}
Under Assumption~\ref{assump:approximate-dp-regime}, Lemma~\ref{lem:step-001-envelope}, Proposition~\ref{prop:step-001-teacher}, and
Lemma~\ref{lem:step-007-fixed-point}, suppose \(d\ge1\). Define the 
appendix quantities

\[
 \ell:=\log\frac{64}{\delta\beta},
 \qquad
 R_T:=\frac{d^2\ell\Lambda^2}{\varepsilon},
 \qquad
 A_{\log}:=80+\log(1+C_{\rm blk}),
\tag{15}
\]

and the appendix-local universal constant

\[
 H:=A_{\log}(1+\log C_{\rm teach}).
\tag{16}
\]

Then the exact sample size obeys

\[
 n_0
 \le(1+C_{\rm blk})\bar k\frac{d^2}{\alpha}
 a(\bar k)\log Q(\bar k),
\tag{17}
\]

and, more explicitly,

\[
 n_0
 \le
 2(1+C_{\rm blk})C_{\rm teach}H
 \frac{d^4\ell\Lambda^3}{\varepsilon\alpha}
 \bigl[v+(H+3)\Lambda\bigr].
\tag{18}
\]

In particular, with the explicit universal constant

\[
 K_{\rm fp}:=
 2(1+C_{\rm blk})C_{\rm teach}H(H+3),
\tag{19}
\]

one has the fully parameter-exposed intermediate envelope

\[
 n_0\le
 K_{\rm fp}
 \frac{d^4\log(64/(\delta\beta))\Lambda^3}
      {\varepsilon\alpha}
 (v+\Lambda).
\tag{20}
\]

No hidden constant in (18)-(20) depends on
\(d,v,\alpha,\beta,\varepsilon,\delta\).


\end{lemma}

\begin{proof}
For integers \(2\le s\le t\), the definition (1) gives
\(a(s)\le a(t)\). Because all factors are positive,
\(sa(s)\le ta(t)\), and hence \(Q(s)\le Q(t)\) and
\(\log Q(s)\le\log Q(t)\). Thus the right side of (8), multiplied by
\(k\), and the  inequality \(k\le\bar k\) give

\[
\begin{aligned}
 n_0=km(k)
 &\le(1+C_{\rm blk})k\frac{d^2}{\alpha}
      a(k)\log Q(k)\\
 &\le(1+C_{\rm blk})\bar k\frac{d^2}{\alpha}
      a(\bar k)\log Q(\bar k),
\end{aligned}
\tag{21}
\]

which is (17). This comparison uses monotonicity only after the least
feasible \(k\) and the independently verified public witness \(\bar k\)
have both been produced; it does not assume that the teacher feasible set
is monotone.

The  witness is exactly
\(\bar k=\lceil C_{\rm teach}R_T\rceil\). Lemma~\ref{lem:step-001-envelope}, instantiated with
\(C=C_{\rm teach}\), and Proposition~\ref{prop:step-001-teacher} give

\[
 \bar k\le2C_{\rm teach}R_T,
 \qquad
 \log\bar k\le H\Lambda,
 \qquad
 \log Q(\bar k)\le H\Lambda.
\tag{22}
\]

Since \(\Lambda\ge1\), \(\log(1/\beta)\le\Lambda\), and
\(\log4<2\),

\[
\begin{aligned}
 a(\bar k)
 &=v+\log4+\log\bar k+\log(1/\beta)\\
 &\le v+(H+3)\Lambda.
\end{aligned}
\tag{23}
\]

Insert (22)-(23) and
\(R_T=d^2\ell\Lambda^2/\varepsilon\) into (17) to obtain (18).
Finally, \(H+3\ge1\) implies

\[
 v+(H+3)\Lambda\le(H+3)(v+\Lambda),
\]

so (19) turns (18) into (20). Every use of a ceiling is upstream in the
 bound \(\bar k\le2C_{\rm teach}R_T\) or explicitly retained in
(8); no additive integer term has been hidden. \(\square\)


\end{proof}

\begin{proposition}[Exact trace-exponent domination and conditional charge]\label{prop:step-007-conditional-charge}
Under Assumption~\ref{assump:approximate-dp-regime}, Proposition~\ref{prop:step-006-good-event},
Lemma~\ref{lem:step-007-fixed-point}, and the universal block calibration
(4), suppose \(d\ge1\). Then

\[
 \boxed{
 \frac{m\alpha}{3600d^2}
 \ge
 v\log\frac{en_0}{v}
 +\log\frac{4k}{\beta_{\rm tr}}
 }
\tag{24}
\]

and, for every fixed indexed quotient master sample \(\bar S\),

\[
 \Pr_{\mathcal P}
 [E_{\rm good}(\bar S,\mathcal P)^c\mid\bar S]
 \le\beta_{\rm tr}=\beta/4.
\tag{25}
\]


\end{proposition}

\begin{proof}
The lower side of the exact ceiling (2) is

\[
 m\ge C_{\rm blk}\frac{d^2}{\alpha}a\log Q,
\]

so

\[
 \frac{m\alpha}{3600d^2}
 \ge\frac{C_{\rm blk}}{3600}a\log Q.
\tag{26}
\]

The two positive terms that must be paid are controlled separately. First,
Lemma~\ref{lem:step-007-fixed-point} and \(v\le a\) give

\[
 v\log\frac{en_0}{v}
 \le C_{\rm fp}v\log Q
 \le C_{\rm fp}a\log Q.
\tag{27}
\]

Second, put \(b:=\log(4k/\beta)>0\), so \(a=v+b\). Since
\(\beta_{\rm tr}=\beta/4\), \(\log4<2\le2v\), and
\(\log Q>1\),

\[
\begin{aligned}
 \log\frac{4k}{\beta_{\rm tr}}
 &=\log\frac{16k}{\beta}=b+\log4\\
 &\le b+2v\le2(v+b)=2a\le2a\log Q.
\end{aligned}
\tag{28}
\]

Combining (27)-(28), and recalling
\(C_{\rm fp}=2+\log(1+C_{\rm blk})\), yields the exact opposing-term
bound

\[
\begin{aligned}
 v\log\frac{en_0}{v}
 +\log\frac{4k}{\beta_{\rm tr}}
 &\le(C_{\rm fp}+2)a\log Q\\
 &=\bigl(4+\log(1+C_{\rm blk})\bigr)a\log Q\\
 &\le\frac{C_{\rm blk}}{3600}a\log Q
 \le\frac{m\alpha}{3600d^2},
\end{aligned}
\tag{29}
\]

where the penultimate inequality is precisely (4), and the last is (26).
This proves (24) with every multiplicity and confidence term visible.

Now apply the  conditional finite-union estimate (5). Taking the
logarithm of its deterministic upper bound and using (24),

\[
\begin{aligned}
 &\log(4k)+v\log\frac{en_0}{v}
   -\frac{m\alpha}{3600d^2}\\
 &\qquad\le
 \log(4k)-\log\frac{4k}{\beta_{\rm tr}}
 =\log\beta_{\rm tr}.
\end{aligned}
\tag{30}
\]

Exponentiating (30) proves (25). The finite trace/block union is already
the one in Proposition~\ref{prop:step-006-good-event}; no new
factor, stage union, or independence assertion is inserted. \(\square\)


\end{proof}

\begin{proposition}[Tower conversion for the trace ledger]\label{prop:step-007-tower}
Under Assumption~\ref{assump:approximate-dp-regime}, Proposition~\ref{prop:step-006-good-event}, and
Proposition~\ref{prop:step-007-conditional-charge}, suppose \(d\ge1\).
Let \(\nu\) be any probability law on the setting-defined discrete quotient
master-sample space and, after drawing the sample, draw the data-independent
uniform labeled partition \(\mathcal P\) prescribed by the setting. Then

\[
 \Pr_{\bar S\sim\nu,\,\mathcal P}
 [E_{\rm good}(\bar S,\mathcal P)^c]
 \le\beta_{\rm tr}.
\tag{31}
\]


\end{proposition}

\begin{proof}
The partition space is finite. For each fixed sample, Proposition~\ref{prop:step-006-good-event} forms the bad section by a finite
union over its finite exact trace set and the \(k\) labeled blocks. Thus its
conditional section probability is the function

\[
 r(\bar s):=
 \Pr_{\mathcal P}[E_{\rm good}(\bar s,\mathcal P)^c
 \mid\bar S=\bar s].
\]

On the setting's quotient sample interface this section is measurable; in
particular, the quotient sample space is discrete in the declared branch,
and the finite-partition sum has measurable sections. Proposition
~\ref{prop:step-007-conditional-charge} proves the pointwise inequality
\(0\le r(\bar s)\le\beta_{\rm tr}\) for every \(\bar s\), without a
distributional or realizability hypothesis. Therefore the tower identity
is simply

\[
\begin{aligned}
 \Pr[E_{\rm good}^c]
 &=\int r(\bar s)\,\nu(d\bar s)\\
 &\le\int\beta_{\rm tr}\,\nu(d\bar s)
 =\beta_{\rm tr}.
\end{aligned}
\tag{32}
\]

This is a mixture of one-partition failure probabilities. It neither
conditions on a successful generated event nor multiplies block success
probabilities. The same event is reused across the fixed source stages, so
there is no additional stage union in (32). \(\square\)


\end{proof}

\begin{proposition}[Boundary integrity and null-branch bypass]\label{prop:step-007-boundaries}
Under Assumption~\ref{assump:approximate-dp-regime}, Proposition~\ref{prop:step-001-zero},
Proposition~\ref{prop:step-006-good-event},
Lemma~\ref{lem:step-007-fixed-point},
Lemma~\ref{lem:step-007-sample-envelope}, and
Proposition~\ref{prop:step-007-conditional-charge}, the following
piecewise conclusions hold: Item 7 applies when \(d=0\), and Items 1-6
apply when \(d\ge1\).

1. At \(v=1\), (7), (18), (24), and the charge (25) remain valid with no
   division by zero and with trace logarithm \(\log(en_0)\).
2. At \(v=d\ge1\), (7), (18), and (24) remain valid, the trace term is
   \(d\log(en_0/d)\), and the factor \(d^4(v+\Lambda)\) in (20) becomes
   \(d^4(d+\Lambda)\), whose polynomial \(v\)-term is \(d^5\). This is only
   the fixed-point baseline specialization, not a final sample-complexity
   theorem.
3. At \(k=2\), all scalar inequalities remain valid. No proof line requires
   block independence;  (5) already treats the two complementary
   blocks through their marginal laws and a finite union.
4. At \(d=v=1\), every denominator in (1)-(3), (7), and (24) is positive,
   \(Q=e+eka/\alpha>e\), and the same constants apply.
5. The ceiling is retained in both useful directions:
   \(m\ge C_{\rm blk}d^2a\log Q/\alpha\) in (26), while (10)-(13) pay its
   entire additive remainder in the logarithmic fixed point and (8), (17)
   pay it in the sample envelope.
6. For every \(0<\alpha<1/4\), the proof is uniform and uses no positive
   lower cutoff on \(\alpha\). Thus it remains valid along any sequence
   \(\alpha\downarrow0\); the explicit \(1/\alpha\), \(Q\), and logarithmic
   growth remain visible. The excluded value \(\alpha=0\) is not claimed.
7. If \(d=0\), Proposition~\ref{prop:step-001-zero} gives the exact
   \(N=0\) singleton law. The positive-branch fact \(v\ge1\), the partition,
   \(m,k,n_0,Q\), and \(E_{\rm good}\) are not constructed, so no expression
   containing \(v^{-1}\) or \(d^{-2}\) is evaluated and the trace ledger
   contributes zero failure.


\end{proposition}

\begin{proof}
For Item 1, the  positive-branch range is
\(1\le v\le d\), and substitution \(v=1\) in (1), (7), (18), and (24)
leaves every quantity finite. For Item 2, substitution \(v=d\) turns
\(d^4(v+\Lambda)\) in (20) into \(d^4(d+\Lambda)\), while the exact trace
power in (5) becomes \(d\); no inequality changes direction. Item 3 follows
because \(k\ge2\) is used only through positive logarithms and 
(5), whose proof did not factor the block law. Item 4 is the common
specialization of Items 1-3 and the strict parameter ranges.

Item 5 is exactly the pair of ceiling inequalities (10), (12)-(13), and
(26). For Item 6, every use of \(\alpha\) is through its strict positive
range, (11), or an explicitly displayed \(1/\alpha\) term; the constants in
(4), (6), (16), and (19) are independent of \(\alpha\). Item 7 is precisely
the branch conclusion of Proposition~\ref{prop:step-001-zero}; it is logically prior to every
positive-branch definition. 
Lemma~\ref{lem:step-001-calibration} and Proposition~\ref{prop:step-001-teacher} first produce the exact public
dictionary (1)-(3), including the least feasible \(k\), its independent
witness \(\bar k\), the ceiled block size \(m\), and \(n_0=km\), before a
sample or event exists. Lemma~\ref{lem:step-007-fixed-point} then separates
the ceiling as \(m=x+\theta\) and proves

\[
 \log(en_0/v)
 \le\bigl(2+\log(1+C_{\rm blk})\bigr)\log Q.
\tag{33}
\]

Thus (33) closes the fixed-point relation noncircularly at the exact public
quantity \(Q=Q(k)\) and the actual sample size. Lemma
~\ref{lem:step-007-sample-envelope} next uses \(k\le\bar k\) and the
 ceiling-aware witness bounds to prove both the exact auxiliary
envelope (17) and the fully exposed intermediate envelope (20).

Proposition~\ref{prop:step-006-good-event} supplies the conditional
finite-union bound (5) with the actual \(m\). Proposition
~\ref{prop:step-007-conditional-charge} proves the exact bridge inequality

\[
 \frac{m\alpha}{3600d^2}
 \ge v\log(en_0/v)+\log(4k/\beta_{\rm tr}),
\]

so that (5) is at most \(\beta_{\rm tr}\) for every fixed master sample.
Proposition~\ref{prop:step-007-tower} integrates this pointwise bound over
the master sample and proves the same unconditional trace charge. Finally,
the boundary calculations above verify \(v=1\), \(v=d\),
\(k=2\), \(d=v=1\), the complete ceiling remainder, arbitrary positive
\(\alpha\) tending to zero, and the exact \(d=0\) bypass.

Consequently, the fixed point, trace confidence charge, and explicit
intermediate sample envelope all hold with the stated parameters.



\end{proof}

\subsection{Source-Current Restrictions and Essential Lists}\label{app:step_008}

On the positive branch \(d\ge1\), put
\[
 \gamma:=\alpha/16,\qquad
 a_d:=\frac{1}{5d},\qquad
 \rho:=1-\frac{1}{2d}.
\tag{1}
\]
For an indexed quotient master sample \(\bar S\), its block
\(\bar S_i\), and \(\bar h\in\bar C\), write
\[
 e(\bar h):=\operatorname{err}_{\bar S}(\bar h),
 \qquad
 e_i(\bar h):=\operatorname{err}_{\bar S_i}(\bar h).
\tag{2}
\]
The endpoint-complete current restrictions and structural parameters are
\[
 H_i^r:=\{\bar h\in\bar C:e_i(\bar h)\leq
                  \rho^{r+1}\gamma\},
 \qquad
 p_r:=2^r n_0d,
 \qquad 0\leq r\leq d.
\tag{3}
\]
On the event produced by Proposition~\ref{prop:step-006-good-event},
simultaneously for every \(\bar h\in\bar C\) and \(i\in[k]\),
\[
\begin{cases}
 (1-a_d)e(\bar h)\leq e_i(\bar h)
       \leq(1+a_d)e(\bar h),
       &e(\bar h)>\gamma/3,\\
 0\leq e_i(\bar h)\leq\gamma/2,
       &e(\bar h)\leq\gamma/3.
\end{cases}
\tag{4}
\]
Equality at \(e(\bar h)=\gamma/3\) belongs to the second branch.

We also record the exact structural interfaces used below from
Lyu~\citep{lyu2025}.  A nonempty binary class \(\mathcal A\) on \(Q_C\)
is \(q\)-irreducible when, for every \(q\)-tuple
\(x_1,\ldots,x_q\in Q_C\),
\[
 \operatorname{LD}\!\left(
 \mathcal A|_{(x_1,\operatorname{SOA}_{\mathcal A}(x_1)),\ldots,
                    (x_q,\operatorname{SOA}_{\mathcal A}(x_q))}
 \right)
 =\operatorname{LD}(\mathcal A).
\tag{5}
\]
For positive integers \(p,d\) and nonempty
\(\mathcal A\subseteq\{0,1\}^{Q_C}\) with
\(\operatorname{LD}(\mathcal A)\le d\), a valid
\((p,d)\)-decomposition is a finite binary restriction tree whose node
classes \(\mathcal A_u\) and leaf classes \(\mathcal A_\ell\) obey
\[
 \operatorname{depth}(u)
 \le p\bigl(2^{d-\operatorname{LD}(\mathcal A_u)+1}-1\bigr),
\tag{6}
\]
\[
 \operatorname{depth}(\ell)
 \le p\bigl(2^{d-\operatorname{LD}(\mathcal A_\ell)}-1\bigr),
 \qquad
 \mathcal A_\ell\text{ is }
 p2^{d-\operatorname{LD}(\mathcal A_\ell)}
 \text{-irreducible}.
\tag{7}
\]
Its degree is the largest leaf Littlestone dimension, and
\[
 \operatorname{DDim}_{p,d}(\mathcal A)
 :=\min\{\operatorname{degree}(T):T\text{ is valid}\}.
\tag{8}
\]
Such a valid tree exists, the minimum is attained, and every valid tree has
at most
\[
 p^d2^{d^2}
\tag{9}
\]
leaves.  If nonempty \(\mathcal G\subseteq\mathcal H\) have Littlestone
dimension at most \(d\), then arbitrary optimal \((2p,d)\)- and
\((p,d)\)-decompositions satisfy
\[
 \operatorname{DDim}_{2p,d}(\mathcal G)
 \le \operatorname{DDim}_{p,d}(\mathcal H).
\tag{10}
\]
Only when
\[
 \operatorname{DDim}_{2p,d}(\mathcal G)
 =\operatorname{DDim}_{p,d}(\mathcal H)=t
\tag{11}
\]
does every dimension-\(t\) leaf \(\mathcal G_v\) of the first tree have a
dimension-\(t\) leaf \(\mathcal H_u\) of the second tree with
\[
 \operatorname{SOA}_{\mathcal H_u}
 =\operatorname{SOA}_{\mathcal G_v}
 \quad\text{as functions }Q_C\to\{0,1\}.
\tag{12}
\]
Finally, with \(t=\operatorname{DDim}_{p,d}(\mathcal A)\), an actual
function \(f:Q_C\to\{0,1\}\) is \((p,d)\)-essential to \(\mathcal A\)
when every optimal decomposition has a leaf satisfying
\[
 \operatorname{LD}(\mathcal A_\ell)=t,
 \qquad
 \operatorname{SOA}_{\mathcal A_\ell}=f
 \quad\text{pointwise on }Q_C.
\tag{13}
\]
The corresponding corollary gives at most \(p^d2^{d^2}\) essential
functions, inheritance from a subclass with the same \((p,d)\)-DDim,
existence when the \((2p,d)\)- and \((p,d)\)-DDim values agree, and, at
DDim zero, equality of the essential set with the finite nonempty class.

\begin{lemma}[Endpoint-complete source/current stage and dimension map]\label{lem:step-008-stage-map}
Under Assumption~\ref{assump:finite-littlestone} and the setting's
positive branch \(d\ge1\), use the endpoint-complete current restrictions
and structural parameters in (1)--(3).  Resolve the source endpoint by
defining, for every \(i\in[k]\),

\[
 H_{i,\mathrm{src}}^s
 :=\{\bar h\in\bar C:e_i(\bar h)\le\rho^s\gamma\},
 \qquad
 p_{s,\mathrm{src}}:=2^s n_0d,
 \qquad 1\le s\le d+1.
\tag{14}
\]

Then, for every \(0\le r\le d\),

\[
 H_i^r=H_{i,\mathrm{src}}^{r+1},
 \qquad
 p_r=2^rn_0d=\frac{p_{r+1,\mathrm{src}}}{2},
\tag{15}
\]

and, for every \(0\le r<d\),

\[
 p_{r+1}=2p_r.
\tag{16}
\]

Every nonempty \(H_i^r\) is a binary class on \(Q_C\) with
\(\operatorname{LD}(H_i^r)\le d\), and every \(p_r\) is a positive integer.
Moreover, if an actual \(\bar c^\star\in\bar C\) has
\(e_i(\bar c^\star)=0\) for every \(i\), then
\(\bar c^\star\in H_i^r\) for all \(i,r\); hence all current restrictions
are nonempty under that explicit condition.


\end{lemma}

\begin{proof}
The construction uses the range \(j\le d+1\), and in
the proof prints the restriction display for \(j\in[d]\), but its
algorithm immediately operates in \(d+1\) stages and consumes
\(H_i^{d+1}\). Thus the source's operational endpoint convention is the
same formula at \(s=d+1\). Equation (14) makes that necessary endpoint
explicit; it does not introduce an additional stage or a new threshold.

Substituting \(s=r+1\) in (14) and comparing with (3) proves both identities
in (15). Equation (16) is the exact computation
\[
 p_{r+1}=2^{r+1}n_0d=2(2^rn_0d)=2p_r.
\]
In particular, the current \((p_r,d)\)-essential list on
\(H_i^r=H_{i,\mathrm{src}}^{r+1}\) is not mislabeled as the source
algorithm's literal stage-\(r+1\) list: its parameter is exactly one half
of \(p_{r+1,\mathrm{src}}\). The source structural statements are uniform
in the positive integer \(p\), and (16) is the exact \((2p,p)\) pairing
used below.

It remains to discharge the dimension premise on the quotient without
assuming class finiteness. Suppose \(\bar C\) Littlestone-shatters a finite
binary tree of depth \(s\) whose nodes are labeled by elements of \(Q_C\).
For each of the finitely many node labels \(q\), choose an
\(x_q\in X\) with \(\kappa(x_q)=q\). Replacing every node label \(q\) by
\(x_q\) gives a tree shattered by \(C\), because each witnessing
\(\bar c\) is induced by the corresponding \(c\) and
\(\bar c(q)=c(x_q)\). Hence \(s\le\operatorname{LD}(C)=d\), so
\(\operatorname{LD}(\bar C)\le d\). Monotonicity under subclasses gives
\(\operatorname{LD}(H_i^r)\le d\) whenever the outer restriction is
nonempty. Finally, \(n_0,d\in\mathbb N\) on this branch, so \(p_r\in
\mathbb N\). If \(e_i(\bar c^\star)=0\), positivity of
\(\rho^{r+1}\gamma\) makes the last membership assertion immediate. This
last statement is a proved conditional implication, not a realizability
assumption. \(\square\)


\end{proof}

\begin{proposition}[Exact eventwise interleaving inclusion]\label{prop:step-008-inclusion}
Under Assumption~\ref{assump:finite-littlestone}, Proposition~\ref{prop:step-006-good-event}, and
Lemma~\ref{lem:step-008-stage-map}, suppose \(d\ge1\) and the exact derived
event \(E_{\mathrm{good}}\) occurs. Then, for every
\(r\in\{0,\ldots,d-1\}\) and every \(i_*\in[k]\),

\[
 H_{i_*}^{r+1}\subseteq\bigcap_{i=1}^k H_i^r.
\tag{17}
\]

This conclusion holds whether or not either side is empty.


\end{proposition}

\begin{proof}
Fix \(r<d\), \(i_*\), and
\(\bar h\in H_{i_*}^{r+1}\). By (3),

\[
 e_{i_*}(\bar h)\le\rho^{r+2}\gamma.
\tag{18}
\]

First suppose \(e(\bar h)>\gamma/3\). The high-master-error clause of the
exact  event gives, for every \(i\),

\[
 e_i(\bar h)\le(1+a_d)e(\bar h),
 \qquad
 e_{i_*}(\bar h)\ge(1-a_d)e(\bar h).
\tag{19}
\]

Since \(1-a_d>0\), (18)-(19) imply

\[
 e_i(\bar h)
 \le \frac{1+a_d}{1-a_d}e_{i_*}(\bar h)
 \le \frac{1+a_d}{1-a_d}\rho^{r+2}\gamma.
\tag{20}
\]

The scalar comparison that turns (20) into the required inclusion is

\[
 \frac{1+a_d}{1-a_d}\rho\le1.
\tag{21}
\]

Indeed,

\[
\begin{aligned}
 (1-a_d)-\rho(1+a_d)
 &=\frac{1}{2d}-\frac{1}{5d}
      \left(2-\frac{1}{2d}\right)\\
 &=\frac{d+1}{10d^2}>0.
\end{aligned}
\tag{22}
\]

Both denominators in (21) are positive, so (22) proves (21), and (20)
becomes \(e_i(\bar h)\le\rho^{r+1}\gamma\).

Now suppose \(e(\bar h)\le\gamma/3\). The low-master-error clause of the
same  event gives \(e_i(\bar h)\le\gamma/2\) for every \(i\).
For \(x\in[0,1]\) and an integer \(s\ge1\), the elementary inequality
\((1-x)^s\ge1-sx\) follows by induction: after the base case, multiplication
of the induction hypothesis by \(1-x\ge0\) gives
\[
 (1-x)^{s+1}\ge(1-sx)(1-x)
 =1-(s+1)x+sx^2\ge1-(s+1)x.
\]
Taking \(x=1/(2d)\) and \(s=d\) gives

\[
 \rho^d=\left(1-\frac{1}{2d}\right)^d
 \ge1-\frac{d}{2d}=\frac12.
\tag{23}
\]

Because \(r+1\le d\) and \(0<\rho<1\),
\(\rho^{r+1}\ge\rho^d\ge1/2\). Therefore

\[
 e_i(\bar h)\le\gamma/2\le\rho^{r+1}\gamma.
\tag{24}
\]

In both branches \(\bar h\in H_i^r\) for every \(i\), proving (17).
If \(H_{i_*}^{r+1}\) is empty, (17) is vacuous, so the proof never assumes
nonemptiness or a generated success event. \(\square\)


\end{proof}

\begin{proposition}[Finite optimal decompositions for current restrictions]\label{prop:step-008-decompositions}
Under Assumption~\ref{assump:finite-littlestone},
Lemma~\ref{lem:step-008-stage-map}, the current-notation versions
of Lyu~\citep{lyu2025} Definition 4.2, its valid-decomposition existence claim, and
Lyu~\citep{lyu2025} Lemma 4.1, let \(\mathcal A\) be any nonempty current restriction
\(H_i^r\), and let \(p\in\mathbb N\). Then:

1. at least one finite valid \((p,d)\)-decomposition of \(\mathcal A\)
   exists;
2. an optimal \((p,d)\)-decomposition exists, although it need not be
   unique;
3. every valid or optimal tree has at most \(p^d2^{d^2}\) leaves; and
4. if a nonempty leaf \(\mathcal A_\ell\) has
   \(t=\operatorname{LD}(\mathcal A_\ell)\), then it is
   \(p2^{d-t}\)-irreducible.

None of these conclusions assumes that \(Q_C,\bar C\), or
\(\mathcal A\) is finite.


\end{proposition}

\begin{proof}
Lemma~\ref{lem:step-008-stage-map} gives
\(\operatorname{LD}(\mathcal A)\le d\). Together with \(p,d\in\mathbb N\)
and nonemptiness, this discharges every premise of the source existence
claim, which supplies a valid tree. In its greedy construction, only
nonempty reducible leaves are expanded. Every such node has Littlestone
dimension at least zero. The node-depth clause of Lyu~\citep{lyu2025}
Definition~4.2 therefore bounds its depth by
\[
 B:=p(2^{d+1}-1).
\]
An empty child is never expanded. Hence the constructed binary tree has
depth at most \(B+1\), and therefore at most
\(2^{B+2}-1\) nodes. In particular, the produced valid tree is finite; no
cardinality of \(\mathcal A\) enters this conclusion.

Let
\[
 \mathscr D
 :=\{\operatorname{degree}(T):T\text{ is a valid }(p,d)
                         \text{-decomposition of }\mathcal A\}.
\]
The existence claim makes \(\mathscr D\) nonempty. Its values are integers
between \(0\) and \(d\): the tree covers the nonempty root class, so some
leaf is nonempty, and every nonempty leaf is a subclass of
\(\mathcal A\). Hence the finite set of possible integer degrees has a
least attained value \(t\), and a tree whose degree equals \(t\) is
optimal. This proves optimal attainment without compactness, enumeration,
or a cardinality assumption on \(\mathcal A\). It also leaves arbitrary
nonuniqueness intact.

Lyu~\citep{lyu2025} Lemma~4.1 gives at most \(p^d2^{d^2}\) leaves for every
valid tree. Finally, the valid-leaf clause of Definition~4.2 says directly
that a leaf of dimension \(t\) is \(p2^{d-t}\)-irreducible.
\(\square\)


\end{proof}

\begin{proposition}[Exact factor-two transition and actual SOA inheritance]\label{prop:step-008-transition}
Under Assumption~\ref{assump:finite-littlestone},
Lemma~\ref{lem:step-008-stage-map},
Propositions~\ref{prop:step-008-inclusion} and
\ref{prop:step-008-decompositions}, and Lyu~\citep{lyu2025} Lemma 4.3, suppose
\(E_{\mathrm{good}}\) occurs. Fix \(0\le r<d\) and \(i_*,i\in[k]\).
If \(H_{i_*}^{r+1}\ne\varnothing\), then \(H_i^r\ne\varnothing\), and the
only Lemma 4.3 instantiation made here is

\[
 \mathcal G=H_{i_*}^{r+1},\qquad
 \mathcal H=H_i^r,\qquad
 (2p,p)=(p_{r+1},p_r).
\tag{25}
\]

It yields

\[
 \operatorname{DDim}_{p_{r+1},d}(H_{i_*}^{r+1})
 \le
 \operatorname{DDim}_{p_r,d}(H_i^r).
\tag{26}
\]

If the two sides in (26) are both equal to \(t\), then, for every
arbitrarily chosen optimal \((p_{r+1},d)\)-decomposition of
\(H_{i_*}^{r+1}\), every dimension-\(t\) leaf \(\mathcal G_v\), and every
arbitrarily chosen optimal \((p_r,d)\)-decomposition of \(H_i^r\), there
is a dimension-\(t\) leaf \(\mathcal H_u\) satisfying

\[
 \operatorname{SOA}_{\mathcal H_u}
 =\operatorname{SOA}_{\mathcal G_v}
 \quad\text{pointwise on }Q_C.
\tag{27}
\]

Consequently, under that same equal-DDim premise only,

\[
 \operatorname{Ess}_{p_{r+1},d}(H_{i_*}^{r+1})
 \subseteq
 \operatorname{Ess}_{p_r,d}(H_i^r),
\tag{28}
\]

where \(\operatorname{Ess}_{p,d}\) denotes the actual essential-function
set of Definition 4.3.


\end{proposition}

\begin{proof}
Proposition~\ref{prop:step-008-inclusion} gives
\(H_{i_*}^{r+1}\subseteq H_i^r\); thus nonemptiness of the former proves
nonemptiness of the latter. Proposition
~\ref{prop:step-008-decompositions} supplies both arbitrary optimal trees,
and Lemma~\ref{lem:step-008-stage-map} supplies their
\(\operatorname{LD}\le d\) premises. Equation (16) gives
\[
 p_{r+1}=2p_r,
\]
so (25) is literally the source pair \((2p,p)\), with no same-scale
substitution. Lyu~\citep{lyu2025} Lemma 4.3 now gives (26) in the displayed direction.
Its second conclusion is conditional on equality, and under precisely that
premise it gives (27) as equality of actual SOA functions, not merely
equality on \(\bar S\) or a block.

For (28), let \(f\) be essential to \(H_{i_*}^{r+1}\), and fix an
arbitrary optimal \((p_r,d)\)-decomposition of \(H_i^r\). Choose any
optimal \((p_{r+1},d)\)-decomposition of \(H_{i_*}^{r+1}\). By essentiality,
some dimension-\(t\) leaf \(\mathcal G_v\) has
\(\operatorname{SOA}_{\mathcal G_v}=f\). Equation (27) produces a
dimension-\(t\) leaf in the arbitrary \(H_i^r\) tree with the same actual
function \(f\). Since the \(H_i^r\) tree was arbitrary, Definition 4.3 says
that \(f\) is essential to \(H_i^r\). This proves (28).

No conclusion stronger than (26) is drawn when the DDim values are unequal.
In particular, this proposition does not turn (26) into a strict decrease
and does not prove the later finite-potential descent. \(\square\)


\end{proof}

\begin{proposition}[Finite actual essential lists and the full source
corollary]\label{prop:step-008-essential-lists}
Under Assumption~\ref{assump:finite-littlestone},
Proposition~\ref{prop:step-008-decompositions}, Lyu~\citep{lyu2025} Definition 4.3, and
Lyu~\citep{lyu2025} Corollary 4.1, define, for \(0\le r\le d\),

\[
\mathcal L_i^r:=
 \begin{cases}
 \operatorname{Ess}_{p_r,d}(H_i^r),&H_i^r\ne\varnothing,\\
 \varnothing,&H_i^r=\varnothing.
 \end{cases}
\tag{29}
\]

The setting's pre-fixed deterministic ordering of this finite set is the
algorithmic list ordering. The mathematical set and every function identity
below are independent of that ordering.

For a nonempty \(H_i^r\), this is a finite set of actual functions
\(Q_C\to\{0,1\}\), and

\[
 |\mathcal L_i^r|\le p_r^d2^{d^2}.
\tag{30}
\]

More generally, for all nonempty quotient classes to which the source
corollary applies, its four conclusions remain exactly:

1. the bound (30);
2. same-\((p,d)\), same-DDim inheritance from a subclass to a superclass;
3. nonemptiness of the \((p,d)\)-essential set when
   \(\operatorname{DDim}_{2p,d}=\operatorname{DDim}_{p,d}\); and
4. if \(\operatorname{DDim}_{p,d}(\mathcal A)=0\), then
   \(\mathcal A\) is finite and
   \(\operatorname{Ess}_{p,d}(\mathcal A)=\mathcal A\).

The last clause is never applied to an empty class, and (29) does not
identify an empty restriction with DDim zero.


\end{proposition}

\begin{proof}
For nonempty \(H_i^r\), Proposition
~\ref{prop:step-008-decompositions} supplies an optimal decomposition.
Definition 4.3 makes every essential item an actual SOA function of a
maximal-dimensional leaf in every optimal decomposition. Fix just one such
tree. Every essential function must occur among its maximal-leaf SOA
functions. The tree has at most \(p_r^d2^{d^2}\) leaves, so this one finite
set contains every essential function and proves (30). Duplicate SOA
functions at different leaves only reduce the set cardinality.

This argument also explains why no cardinality assumption on \(H_i^r\) or
\(\bar C\) is present: a possibly infinite intersection over optimal trees
is contained in the finite SOA set of any one fixed optimal tree. The
remaining three conclusions are Corollary 4.1 of Lyu~\citep{lyu2025} after
replacing its domain by \(Q_C\). All its hypotheses are discharged by
nonemptiness, Proposition~\ref{prop:step-008-decompositions}, and the
explicit DDim premises in the respective clauses.

If \(H_i^r=\varnothing\), (29) is instead the setting's predetermined
arbitrary-input totalization. No source SOA or decomposition is evaluated.
If \(H_i^r\ne\varnothing\), its essential set may nevertheless be empty
unless a source existence clause applies; the algorithmic empty-list path
is totalized in that case as well. Thus neither class nonemptiness nor list
nonemptiness is smuggled into the proof. \(\square\)


\end{proof}

\begin{proposition}[Authoritative stage-list envelope]\label{prop:step-008-list-envelope}
Under Assumption~\ref{assump:finite-littlestone} and
Proposition~\ref{prop:step-008-essential-lists}, use the 
positive-branch dictionary

\[
 p_d=p_d(k)=2^dn_0d,\qquad
 L(k):=p_d(k)^d2^{d^2}.
\tag{31}
\]

Then, for every \(i\in[k]\) and \(0\le r\le d\),

\[
 |\mathcal L_i^r|
 \le p_r^d2^{d^2}
 \le p_d^d2^{d^2}
 =L(k).
\tag{32}
\]

The logarithm is exactly

\[
 \log L(k)
 =d\log n_0+d\log d+2d^2\log2.
\tag{33}
\]

For the actual all-stage set
\(\mathcal G_i:=\bigcup_{r=0}^d\mathcal L_i^r\),

\[
 |\mathcal G_i|\le(d+1)L(k).
\tag{34}
\]


\end{proposition}

\begin{proof}
Equation (30) proves the first inequality in (32), including the
totalized-empty case. For \(0\le r\le d\),
\[
 p_r=2^rn_0d\le2^dn_0d=p_d.
\]
Since \(d\ge1\), raising both positive sides to the \(d\)-th power and
multiplying by \(2^{d^2}\) proves the second inequality and the equality
with the authoritative definition (31). There is no additional factor,
power, asymptotic comparison, or absorbed logarithm.

Expanding (31) gives
\[
\begin{aligned}
 \log L(k)
 &=d\log(2^dn_0d)+d^2\log2\\
 &=d\log n_0+d\log d+2d^2\log2,
\end{aligned}
\]
which is (33). Finally, subadditivity of cardinality for a union of the
exactly \(d+1\) stage sets and (32) gives (34). Duplicates across stages
or within different leaf witnesses can only decrease the union size.
\(\square\)


\end{proof}

\begin{proposition}[Exact essential-leaf identity and irreducibility export]\label{prop:step-008-leaf-scale}
Under Assumption~\ref{assump:finite-littlestone} and
Propositions~\ref{prop:step-008-decompositions} and
\ref{prop:step-008-essential-lists}, fix \(i,r\) and
\(\bar f\in\mathcal L_i^r\). Then \(H_i^r\ne\varnothing\), and in every
optimal \((p_r,d)\)-decomposition of \(H_i^r\) there is a leaf
\(\mathcal A_\ell\) such that, with

\[
 t:=\operatorname{DDim}_{p_r,d}(H_i^r)
   =\operatorname{LD}(\mathcal A_\ell),
\tag{35}
\]

one has the exact actual-function identity

\[
 \bar f=\operatorname{SOA}_{\mathcal A_\ell}
 \quad\text{on all of }Q_C,
\tag{36}
\]

and \(\mathcal A_\ell\) is
\(p_r2^{d-t}\)-irreducible. Its scale satisfies

\[
 p_r2^{d-t}\ge p_0=n_0d\ge\max\{n_0,d+1\}.
\tag{37}
\]

Consequently, the same leaf is both \(n_0\)-irreducible and
\((d+1)\)-irreducible.


\end{proposition}

\begin{proof}
Membership in the first case of (29) and Definition 4.3 give (35)-(36) in
every optimal decomposition. The SOA is a function in \(H_C\); it need not
belong to \(\bar C\), so the identity preserves the source's potentially
improper actual output rather than silently asserting properness.
The valid-leaf clause of Lyu~\citep{lyu2025} Definition~4.2 gives
\(p_r2^{d-t}\)-irreducibility.

Because \(0\le t\le d\), \(2^{d-t}\ge1\). Because \(r\ge0\),
\(p_r=2^rp_0\ge p_0\). Hence the first inequality in (37) follows.
On the positive branch \(d\ge1\), so \(n_0d\ge n_0\). Also
\(n_0=km\ge2\), whence
\[
 n_0d\ge2d\ge d+1.
\]
This proves all of (37).

For completeness, \(K\)-irreducibility implies \(q\)-irreducibility for
every \(1\le q\le K\). Given a \(q\)-tuple, extend it to length \(K\) by
repeating any point of the nonempty domain \(Q_C\). If restriction by the
first \(q\) SOA-labeled points lowered Littlestone dimension, further
restrictions could not increase it, contradicting \(K\)-irreducibility.
Apply this observation with
\(K=p_r2^{d-t}\) and \(q=n_0,d+1\).

The two exported scales have distinct later roles. The \(n_0\) scale is
the input reserved for the later full-master-sample empirical
contradiction, and the \(d+1\) scale places the exact SOA function in the
later fixed irreducible-SOA family. This proposition supplies those
interfaces only; it does not prove either later empirical or fixed-family
conclusion. \(\square\)


\end{proof}

\begin{proposition}[Boundary, totalization, and object integrity]\label{prop:step-008-boundaries}
Under Assumption~\ref{assump:finite-littlestone},
Lemma~\ref{lem:step-008-stage-map}, and
Propositions~\ref{prop:step-008-inclusion},
\ref{prop:step-008-decompositions}, \ref{prop:step-008-transition},
\ref{prop:step-008-essential-lists},
\ref{prop:step-008-list-envelope}, and
\ref{prop:step-008-leaf-scale},
the endpoint, empty-restriction/list, DDim-zero, scalar/endpoint/list/scale
parts of \(d=1\), infinite countable-quotient class, \(v=d\),
duplicate-trace, nonunique-decomposition, and
actual-function-versus-trace conclusions below hold unconditionally. If
\(E_{\mathrm{good}}\) occurs, the first-transition inclusion in Item 1 and
the \(d=1\) transition inclusion in Item 5 also hold.

1. The first transition index is \(r=0\), and the exact source pair is
   \((p_1,p_0)=(2n_0d,n_0d)\) unconditionally. If
   \(E_{\mathrm{good}}\) occurs, then
   \(H_{i_*}^{1}\subseteq\cap_iH_i^0\).
2. The last current restriction is
   \(H_i^d=H_{i,\mathrm{src}}^{d+1}\). It is used as the endpoint stage;
   no undefined \(H_i^{d+1}\) or transition beyond \(r=d-1\) is claimed.
3. If an arbitrary-input restriction is empty, its list is empty and no
   DDim, SOA, decomposition, or source lemma is evaluated for it. If a
   zero-error actual concept exists, Lemma~\ref{lem:step-008-stage-map}
   proves all restrictions nonempty rather than assuming that fact.
4. DDim zero is considered only for a nonempty class and gives exactly the
   finite-class/list identity in Corollary 4.1; it is not a synonym for an
   empty restriction.
5. At \(d=1\), the map has current stages \(r=0,1\), the transition index
   \(r=0\) corresponds to the source endpoint \(s=2\), \(\rho=1/2\), and
   every scalar inequality, endpoint, leaf scale, and list bound above
   remains valid unconditionally. If \(E_{\mathrm{good}}\) occurs, the only
   transition inclusion is \(H_{i_*}^{1}\subseteq\cap_iH_i^0\).
6. The proof permits finite or countably infinite \(Q_C\) and arbitrary
   possibly infinite \(\bar C\). Its finite lists come from finite trees,
   not from enumerating or bounding \(|\bar C|\).
7. The specialization \(v=d\) changes none of the structural statements:
   \(v\) does not appear in (14)-(37), except indirectly in already fixed
   public parameters.
8. Duplicate empirical traces do not identify actual concepts or SOA
   functions. They only deduplicate the  finite union defining
   \(E_{\mathrm{good}}\).
9. Optimal decompositions need not be unique. Essentiality quantifies over
   every optimal tree, while Lemma 4.3 applies to arbitrary chosen optimal
   trees and preserves the exact actual function.


\end{proposition}

\begin{proof}
For Item 1, (16) at \(r=0\) gives the exact source pair unconditionally,
while (17) at \(r=0\) gives the stated inclusion on
\(E_{\mathrm{good}}\). Item 2 is (15) at \(r=d\) together with the range
\(r<d\) in Propositions~\ref{prop:step-008-inclusion} and
\ref{prop:step-008-transition}. Item 3 is the piecewise definition (29)
and the proved zero-error implication in Lemma
~\ref{lem:step-008-stage-map}; no realized event or mechanism success is
used. Item 4 is Proposition~\ref{prop:step-008-essential-lists}.

For Item 5, when \(d=1\), (22) remains strict, (23) is equality,
\(p_0=n_0\), \(p_1=2n_0\), and
\(p_r2^{1-t}\ge n_0\ge2=d+1\), all unconditionally. Equations (15)-(16)
identify the two current stages and the unique source-endpoint pair, while
(32)-(34) preserve the list bounds. On \(E_{\mathrm{good}}\), (17) at
\(r=0\) gives the only transition inclusion. Thus the smallest positive
dimension loses no endpoint, scalar, list, or scale interface, and its
cross-block inclusion retains the event condition. Items 6 and 9 follow
from the cardinality-free source
statements and the finite-tree arguments in Propositions
~\ref{prop:step-008-decompositions} and
\ref{prop:step-008-essential-lists}. Item 7 is immediate from inspecting
the exact displayed interfaces. Item 8 follows because (4) holds for each
actual function sharing a trace, while the SOA identities in
Propositions~\ref{prop:step-008-transition} and
\ref{prop:step-008-leaf-scale} are pointwise equalities on all of \(Q_C\).

These checks establish no strict DDim drop, common list item, mechanism
event, selected output, privacy inequality, empirical error, or population
error. 
Lemma~\ref{lem:step-008-stage-map} resolves the source's printed endpoint
omission using its explicit \(d+1\)-stage algorithm and proves, for every
\(0\le r\le d\),

\[
 H_i^r=H_{i,\mathrm{src}}^{r+1},
 \qquad
 p_r=2^rn_0d=p_{r+1,\mathrm{src}}/2.
\]

It also proves the source-required \(\operatorname{LD}\le d\) condition
directly on the quotient. Proposition~\ref{prop:step-008-inclusion} then
uses both exact branches of the  generated event. The high branch
is closed by the displayed inequality
\[
 \frac{1+a_d}{1-a_d}\rho\le1,
\]
and the low branch by
\(\gamma/2\le\rho^d\gamma\le\rho^{r+1}\gamma\). Thus it proves exactly
\[
 H_{i_*}^{r+1}\subseteq\bigcap_iH_i^r
 \quad(0\le r<d).
\]

Proposition~\ref{prop:step-008-decompositions} discharges Definition 4.2,
the valid-decomposition existence claim, optimal attainment, Lemma 4.1,
and leaf irreducibility pointwise for every nonempty current restriction.
Proposition~\ref{prop:step-008-transition} then invokes Lemma 4.3 only with
\[
 (\mathcal G,\mathcal H,2p,p)
 =(H_{i_*}^{r+1},H_i^r,p_{r+1},p_r).
\]
It preserves the direction of the DDim inequality and uses the leaf SOA
identity only under exact DDim equality. The resulting equality is of
actual functions on \(Q_C\), and under the same premise it yields the
correct cross-stage essential-list inheritance.

Proposition~\ref{prop:step-008-essential-lists} discharges Definition 4.3
and all four clauses of Corollary 4.1, defines the empty outer-restriction
path without assigning it a DDim, and proves finite actual lists without
assuming \(\bar C\) finite. Proposition
~\ref{prop:step-008-list-envelope} applies the authoritative dictionary
without hiding a power:
\[
 |\mathcal L_i^r|\le p_r^d2^{d^2}
 \le p_d^d2^{d^2}=L(k),
 \qquad
 |\mathcal G_i|\le(d+1)L(k),
\]
with the exact logarithm (33).

Finally, Proposition~\ref{prop:step-008-leaf-scale} gives every essential
actual function an exact maximal-leaf SOA witness and exports
\[
 p_r2^{d-t}\ge p_0=n_0d\ge\max\{n_0,d+1\}.
\]
The boundary calculations above cover the endpoint,
empty, zero-dimensional, infinite-class, \(v=d\), duplicate-trace,
nonuniqueness, actual-function, and scalar/list/scale parts of the
smallest-\(d\) boundary unconditionally; on \(E_{\mathrm{good}}\), it also
verifies the first-transition and \(d=1\) transition inclusions.

Thus the cross-block inclusion, exact lists and SOA witnesses, finite list
multiplicity, factor-two DDim comparison, and both irreducibility scales
hold simultaneously.



\end{proof}

\subsection{DDim Descent and Common Stage}\label{app:step_009}

Use the restrictions, scales, and actual essential-function lists
\[
 H_i^r
 =\{\bar h\in\bar C:
       \operatorname{err}_{\bar S_i}(\bar h)
       \leq\rho^{r+1}\gamma\},
 \qquad
 p_r=2^rn_0d,
 \qquad
 \mathcal L_i^r=\operatorname{Ess}_{p_r,d}(H_i^r),
\tag{1}
\]
for \(i\in[k]\) and \(r\in\{0,\ldots,d\}\), with the empty-restriction
list convention of Proposition~\ref{prop:step-008-essential-lists}.  The
structural conclusion proved in this subsection is that, on
\(E_{\mathrm{good}}\), some \(r_*\in\{0,\ldots,d\}\) and actual
\(\bar f_*\in H_C\) satisfy
\[
 \bar f_*\in\bigcap_{i=1}^k\mathcal L_i^{r_*}.
\tag{2}
\]
Proposition~\ref{prop:step-008-inclusion} gives, on the same event,
\[
 H_{i_*}^{r+1}\subseteq\bigcap_{i=1}^kH_i^r
 \quad(0\leq r<d,\ i_*\in[k]),
\tag{3}
\]
and Proposition~\ref{prop:step-008-transition} uses only the exact pairing
\[
 (\mathcal G,\mathcal H,2p,p)
 =(H_{i_*}^{r+1},H_i^r,p_{r+1},p_r).
\tag{4}
\]
In particular, Lemma~\ref{lem:step-008-stage-map} states
\[
 H_i^r
 =\{\bar h\in\bar C:
       \operatorname{err}_{\bar S_i}(\bar h)
       \leq\rho^{r+1}\gamma\},
 \qquad
 p_r=2^rn_0d,
 \qquad 0\le r\le d,
\tag{5}
\]
and every nonempty restriction has an attained integer value
\[
 D_{i,r}:=\operatorname{DDim}_{p_r,d}(H_i^r)
 \in\{0,1,\ldots,d\}.
\tag{6}
\]
For reference, the factor-two comparison from
Lyu~\citep{lyu2025} is
\[
 \operatorname{DDim}_{2p,d}(\mathcal G)
 \leq\operatorname{DDim}_{p,d}(\mathcal H),
\tag{7}
\]
and equality at a common value \(t\) gives maximum-leaf SOA identity
\[
 \operatorname{SOA}_{\mathcal H_u}
 =\operatorname{SOA}_{\mathcal G_v}
 \quad\text{as functions on }Q_C.
\tag{8}
\]
The current-notation substitution is exactly
\[
 \mathcal G=H_{i_*}^{r+1},\qquad
 \mathcal H=H_i^r,\qquad
 (2p,p)=(p_{r+1},p_r).
\tag{9}
\]
For a nonempty restriction, Definition~4.3 in
Lyu~\citep{lyu2025} is represented by
\[
 \mathcal L_i^r=\operatorname{Ess}_{p_r,d}(H_i^r).
\tag{10}
\]
Its Corollary~4.1 preserves essential functions under equal DDim and says
that at DDim zero the essential set is the entire finite nonempty class.
These are the two support interfaces used in the equality and zero-potential
arguments below.

\begin{lemma}[Legal realizable stage state, potential, and actual-function
score]\label{lem:step-009-legal-state}
Under Assumptions~\ref{assump:finite-littlestone} and
\ref{assump:realizable-iid} and Lemma~\ref{lem:step-008-stage-map} and
Propositions~\ref{prop:step-008-decompositions} and
\ref{prop:step-008-essential-lists}, suppose \(d\ge1\), fix arbitrary
\(D\) and \(c\in C\), and let the raw master sample be generated as

\[
 S=((x_j,c(x_j)))_{j=1}^{n_0},
 \qquad x_1,\ldots,x_{n_0}\stackrel{\mathrm{iid}}{\sim}D.
 \tag{11}
\]

Let \(\bar S=T_{n_0}(S)\), let \(\bar c\in\bar C\) be induced by \(c\),
and use the fixed blocks of \(\bar S\). Then

\[
 \operatorname{err}_{\bar S_i}(\bar c)=0,
 \qquad
 \bar c\in H_i^r
 \quad(i\in[k],\ 0\le r\le d).
 \tag{12}
\]

In particular, every \(H_i^r\) is nonempty. It is therefore legal to
define, only on this realizable path,

\[
 D_{i,r}:=\operatorname{DDim}_{p_r,d}(H_i^r),
 \qquad
 M_r:=\max_{i\in[k]}D_{i,r}.
 \tag{13}
\]

These maxima are attained and

\[
 D_{i,r},M_r\in\{0,1,\ldots,d\}.
 \tag{14}
\]

For every actual quotient function \(\bar f\in H_C\), define its
stage-\(r\) teacher-occurrence score and the maximum score by

\[
 s_r(\bar f):=
 \bigl|\{i\in[k]:\bar f\in\mathcal L_i^r\}\bigr|,
 \qquad
 q_r:=\max_{\bar f\in H_C}s_r(\bar f).
 \tag{15}
\]

The maximum in (15) is attained, \(q_r\in\{0,1,\ldots,k\}\), and

\[
 q_r=k
 \quad\Longleftrightarrow\quad
 \text{there is an actual }\bar f\in H_C
 \text{ with }\bar f\in\bigcap_{i=1}^k\mathcal L_i^r.
 \tag{16}
\]


\end{lemma}

\begin{proof}
For every generated record in (11), put \(q_j=\kappa(x_j)\). By the
definition of the induced quotient concept,

\[
 c(x_j)=\bar c(\kappa(x_j))=\bar c(q_j).
\]

Thus \(\bar c\) makes no error on any record and hence no error on any
block, proving the first assertion in (12). The thresholds
\(\rho^{r+1}\gamma\) in Lemma~\ref{lem:step-008-stage-map} are positive on
the positive branch, so the second assertion in (12) follows. This derives nonemptiness from the
realizable labels; it does not treat nonemptiness as a primitive or
generated-event assumption. Independence in (11) is not used in this
pointwise implication.

Proposition~\ref{prop:step-008-decompositions} now applies to every nonempty
restriction. Its integer-valued degree conclusion makes (13) well-defined on
the finite nonempty index set \([k]\), its maximum is attained, and (14)
follows.

Each list supplied by Proposition~\ref{prop:step-008-essential-lists}
consists of actual elements of \(H_C\). For any
\(\bar f\in H_C\), (15) counts block indices rather than leaf duplicates,
so \(s_r(\bar f)\in\{0,\ldots,k\}\). The space \(H_C\) is nonempty. Hence
the nonempty set of attained integer scores is a subset of the finite set
\(\{0,\ldots,k\}\) and has a largest attained value. This remains true if
every list is empty: then every score is zero. Finally, a score equals
\(k\) exactly when the same actual function is present in every one of the
\(k\) set-valued lists, proving (16). \(\square\)


\end{proof}

\begin{proposition}[Equality of consecutive maxima forces one common actual
SOA]\label{prop:step-009-equality-common}
Under Assumptions~\ref{assump:finite-littlestone} and
\ref{assump:realizable-iid}, 
Propositions~\ref{prop:step-008-inclusion},
\ref{prop:step-008-decompositions}, and
\ref{prop:step-008-transition}, Lyu~\citep{lyu2025} Lemma 4.3, and
Lemma~\ref{lem:step-009-legal-state}, suppose \(E_{\mathrm{good}}\) occurs,
\(0\le r<d\), and

\[
 M_{r+1}=M_r=:t.
 \tag{17}
\]

Then there is one actual function \(\bar f\in H_C\) that is
\((p_r,d)\)-essential to every \(H_i^r\). Consequently,

\[
 \bar f\in\bigcap_{i=1}^k\mathcal L_i^r
 \qquad\text{and}\qquad q_r=k.
 \tag{18}
\]


\end{proposition}

\begin{proof}
Because \([k]\) is finite and every next-stage DDim is defined by
Lemma~\ref{lem:step-009-legal-state}, choose

\[
 i_*\in\operatorname*{arg\,max}_{j\in[k]}D_{j,r+1},
 \qquad
 D_{i_*,r+1}=M_{r+1}=t.
 \tag{19}
\]

The same lemma gives \(H_{i_*}^{r+1}\ne\varnothing\), and the 
eventwise inclusion gives

\[
 H_{i_*}^{r+1}\subseteq H_i^r
 \quad\text{for every }i\in[k].
 \tag{20}
\]

For every current block \(i\), apply
Proposition~\ref{prop:step-008-transition} to
\((H_{i_*}^{r+1},H_i^r,p_{r+1},p_r)\). Its factor-two inequality, (17), and
(19), together with the definition of \(M_r\), give the two-sided comparison

\[
 t=M_{r+1}=D_{i_*,r+1}
 \le D_{i,r}
 \le M_r=t
 \quad\text{for every }i\in[k].
 \tag{21}
\]

Thus every quantity in (21) equals \(t\). Notice that the factor-two
comparison was made against every current block, not only a current
maximizer.

Now fix one arbitrary optimal \((p_{r+1},d)\)-decomposition
\(T_*\) of \(H_{i_*}^{r+1}\), and in that tree fix one arbitrary leaf
\(v_*\) whose Littlestone dimension equals its degree \(t\). Such a leaf
exists because the preceding proposition supplies a finite optimal tree
whose degree is the maximum leaf dimension. Define the actual pointwise
function

\[
 \bar f:=\operatorname{SOA}_{(H_{i_*}^{r+1})_{v_*}}
 \in H_C.
 \tag{22}
\]

The function in (22) may be improper, meaning \(\bar f\notin\bar C\); no
properness assertion is needed.

Fix a block \(i\in[k]\), and then fix an arbitrary optimal
\((p_r,d)\)-decomposition \(T_i\) of \(H_i^r\). Pair this arbitrary current
tree with the fixed next tree \(T_*\). The source Lemma 4.3 explicitly
allows arbitrarily chosen optimal trees. Its equality premise holds by
(21), so its second conclusion, applied to the fixed arbitrary maximum
leaf \(v_*\), produces a dimension-\(t\) leaf \(u_i\) of \(T_i\) with

\[
 \operatorname{SOA}_{(H_i^r)_{u_i}}=\bar f
 \quad\text{pointwise on all of }Q_C.
 \tag{23}
\]

Because \(T_i\) was arbitrary, Definition 4.3 says that the one function
\(\bar f\) in (22) is \((p_r,d)\)-essential to \(H_i^r\). The next tree
\(T_*\), the next maximum leaf \(v_*\), and hence \(\bar f\) were held
fixed while \(i\) and \(T_i\) varied. Therefore the same actual function
is essential to every current block. By
Proposition~\ref{prop:step-008-essential-lists}, it lies in every
\(\mathcal L_i^r\), and (16) gives \(q_r=k\).

This proof does not assert that \(\bar f\) is essential to the next-stage
class: one maximum leaf of one fixed next optimal tree is sufficient. That
distinction avoids assuming the extra same-class, two-scale equality
required by Corollary 4.1 Item 3. \(\square\)


\end{proof}

\begin{proposition}[Below-common-score stages cause one-unit DDim descent]\label{prop:step-009-unit-drop}
Under Assumptions~\ref{assump:finite-littlestone} and
\ref{assump:realizable-iid}, 
Propositions~\ref{prop:step-008-inclusion} and
\ref{prop:step-008-transition}, Lemma~\ref{lem:step-009-legal-state}, and
Proposition~\ref{prop:step-009-equality-common}, suppose
\(E_{\mathrm{good}}\) occurs, \(0\le r<d\), and

\[
 q_r<k.
 \tag{24}
\]

Then

\[
 M_{r+1}\le M_r-1.
 \tag{25}
\]


\end{proposition}

\begin{proof}
Choose a next-stage maximizer \(i_*\) exactly as in (19). For every current
block \(i\), realizable nonemptiness,
Proposition~\ref{prop:step-008-inclusion}, and the exact pair
\((p_{r+1},p_r)=(2p_r,p_r)\) discharge
Proposition~\ref{prop:step-008-transition}. Hence

\[
 M_{r+1}=D_{i_*,r+1}
 \le D_{i,r}
 \le M_r
 \quad\text{for every }i\in[k].
 \tag{26}
\]

In particular, \(M_{r+1}\le M_r\). If equality held, Proposition
~\ref{prop:step-009-equality-common} would give \(q_r=k\), contradicting
(24). Therefore

\[
 M_{r+1}<M_r.
 \tag{27}
\]

By (14), both sides are integers. Thus strict inequality (27), rather than
mere nonincrease, is exactly the signed one-unit bound (25). There is no
additive defect, stochastic forcing, or approximation term in this
recurrence. \(\square\)


\end{proof}

\begin{proposition}[DDim-zero support contains the realizable target in every
list]\label{prop:step-009-zero-support}
Under Assumptions~\ref{assump:finite-littlestone} and
\ref{assump:realizable-iid}, Proposition~\ref{prop:step-008-essential-lists}, Lyu~\citep{lyu2025}
Corollary 4.1 Item 4, and Lemma~\ref{lem:step-009-legal-state}, fix any
stage \(r\in\{0,\ldots,d\}\). If

\[
 M_r=0,
 \tag{28}
\]

then

\[
 \bar c\in\bigcap_{i=1}^k\mathcal L_i^r,
 \qquad q_r=k.
 \tag{29}
\]


\end{proposition}

\begin{proof}
Every \(D_{i,r}\) is a nonnegative integer and is at most \(M_r\). Thus
(28) gives

\[
 \operatorname{DDim}_{p_r,d}(H_i^r)=D_{i,r}=0
 \quad\text{for every }i\in[k].
 \tag{30}
\]

Lemma~\ref{lem:step-009-legal-state} proved that every \(H_i^r\) is
nonempty and that \(\bar c\in H_i^r\). Therefore the nonempty-class premise
of Corollary 4.1 Item 4 is discharged separately for every block. That item
and the  list definition give

\[
 \mathcal L_i^r
 =\operatorname{Ess}_{p_r,d}(H_i^r)
 =H_i^r
 \quad\text{for every }i\in[k].
 \tag{31}
\]

Combining target membership with (31) proves (29), and (16) gives the
score statement. In particular, the common item is supplied by the
realizable target and Item 4; it is not inferred from potential
nonnegativity alone. \(\square\)


\end{proof}

\begin{proposition}[Finite DDim budget produces a score-\(k\) stage by stage
\(d\)]\label{prop:step-009-termination}
Under Assumptions~\ref{assump:finite-littlestone} and
\ref{assump:realizable-iid}, Lemma~\ref{lem:step-009-legal-state}, and
Propositions~\ref{prop:step-009-unit-drop} and
\ref{prop:step-009-zero-support}, suppose \(d\ge1\) and
\(E_{\mathrm{good}}\) occurs. Then there is a stage

\[
 r_*\in\{0,1,\ldots,M_0\}\subseteq\{0,1,\ldots,d\}
 \tag{32}
\]

and an actual quotient function \(\bar f_*\in H_C\) such that

\[
 q_{r_*}=k,
 \qquad
 \bar f_*\in\bigcap_{i=1}^k\mathcal L_i^{r_*}.
 \tag{33}
\]

The function \(\bar f_*\) is an actual member of the stage lists; it may
be an improper SOA function unless the zero-potential branch selects the
target \(\bar c\).


\end{proposition}

\begin{proof}
By (14), let

\[
 \tau:=M_0\in\{0,1,\ldots,d\}.
 \tag{34}
\]

If \(\tau=0\), Proposition~\ref{prop:step-009-zero-support} applied at
stage (0) gives \(q_0=k\), so take \(r_*=0\).

Now suppose \(\tau\ge1\). If some \(r<\tau\) already satisfies \(q_r=k\),
take the least such \(r\). Otherwise,

\[
 q_r<k\quad\text{for every }r=0,1,\ldots,\tau-1.
 \tag{35}
\]

Because \(\tau\le d\), every index in (35) is a legal transition index:
\(r\le\tau-1\le d-1\). Proposition
~\ref{prop:step-009-unit-drop} therefore applies successively and gives,
by induction,

\[
 M_j\le M_0-j=\tau-j
 \quad(j=0,1,\ldots,\tau).
 \tag{36}
\]

At \(j=\tau\), (36) gives \(M_\tau\le0\). Nonnegativity in (14) gives
\(M_\tau=0\), so Proposition~\ref{prop:step-009-zero-support} yields

\[
 q_\tau=k,
 \qquad
 \bar c\in\bigcap_i\mathcal L_i^\tau.
 \tag{37}
\]

Take \(r_*=\tau\) and \(\bar f_*=\bar c\). Together with the earlier-success
case and (16), this proves (32)-(33).

The boundary traces are exact:

1. If \(M_0=0\), no transition is taken and stage (0) succeeds by (31).
2. If \(d=1\), then \(M_0\in\{0,1\}\). The first case succeeds at stage
   (0); in the second case, either stage (0) already has score \(k\),
   or the sole legal transition \(r=0\) gives \(M_1=0\) and stage (1)
   succeeds.
3. If \(M_0=d\) and every earlier stage has score below \(k\), (36) uses
   exactly the \(d\) transitions \(r=0,\ldots,d-1\) and reaches
   \(M_d=0\). Corollary 4.1 Item 4 is then applied at the last defined
   stage \(d\). No \(H_i^{d+1}\), \(M_{d+1}\), or transition out of stage
   \(d\) is introduced.
4. On nonrealizable arbitrary inputs, some restrictions may be empty; the
   totalization assigns them empty lists without evaluating DDim or
   SOA. This proposition makes no common-stage utility claim on those
   paths. On the realizable path, (12) proves nonemptiness at all \(d+1\)
   stages before the recurrence begins.
5. In the equality case used inside the descent proof, nonunique optimal
   trees are handled by fixing an arbitrary next tree/maximum leaf and
   quantifying over every arbitrary current tree. The resulting common
   object is the same actual function on \(Q_C\), even when it is
   improper.

Finally, the induction is deterministic conditional on the one
simultaneous event \(E_{\mathrm{good}}\). It draws no new data, introduces
no union over stages, invokes no mechanism report or noise variable, and
multiplies no sample or confidence bound by \(d+1\). 
Lemma~\ref{lem:step-009-legal-state} first derives the missing entry fact:
for every sample actually generated from \(c\), the quotient target
\(\bar c\) has zero error on every teacher block and therefore belongs to
every \(H_i^r\). This makes all quantities

\[
 D_{i,r}=\operatorname{DDim}_{p_r,d}(H_i^r),
 \qquad M_r=\max_iD_{i,r}
\]

legal on their nonempty domains and integer-valued in
\(\{0,\ldots,d\}\). The same lemma defines the exact actual-function
occurrence query \(q_r\) and proves that \(q_r=k\) is equivalent to one
actual function lying in all \(k\) lists.

For each transition \(r<d\), choose a next-stage maximizer \(i_*\).
The  interleaving inclusion and the exact factor-two source pair
give, against every current block,

\[
 M_{r+1}=D_{i_*,r+1}\le D_{i,r}\le M_r.
\]

Proposition~\ref{prop:step-009-equality-common} proves that equality of the
outer terms forces a common current-stage actual SOA. It does so with the
full source quantifiers: one arbitrary optimal next tree and one arbitrary
maximum leaf define a single function, and Lemma 4.3 maps that same
function into every arbitrary optimal current tree of every block. Hence
it is essential to every current class and has score \(k\). Therefore,
when \(q_r<k\), Proposition~\ref{prop:step-009-unit-drop} gives the strict
integer recurrence

\[
 M_{r+1}\le M_r-1.
\]

Proposition~\ref{prop:step-009-zero-support} separately supplies the
terminal mechanism source: at \(M_r=0\), Corollary 4.1 Item 4 makes each
list equal its nonempty restriction, and the already proved common target
\(\bar c\) belongs to all lists. Thus zero potential implies score \(k\)
for a concrete actual function.

Finally, the finite-potential argument above starts at the integer
\(M_0\le d\), charges one unit for every score-below-\(k\) stage, and after
at most \(M_0\) transitions reaches either an earlier common stage or a
zero-potential stage. It produces \(r_*\le M_0\le d\) and
\(\bar f_*\in\cap_i\mathcal L_i^{r_*}\). Its boundary trace covers
\(M_0=0\), \(d=1\), and first success at stage \(d\), while preserving the
empty arbitrary-input totalization and making no utility claim on
that path.

Therefore, on realizable \(E_{\mathrm{good}}\) paths with the stated list
interface, finite DDim descent produces a common actual quotient function
of score \(k\) by stage \(d\).



\end{proof}

\subsection{Mechanism Accuracy and Actual Output}\label{app:step_010}


\begin{proposition}[Calibrated quotient mechanism interfaces]\label{prop:step-010-interfaces}
Under Assumption~\ref{assump:approximate-dp-regime},
Lemma~\ref{lem:step-001-calibration},
Propositions~\ref{prop:step-001-teacher} and
\ref{prop:step-001-totalization},
Lemma~\ref{lem:step-009-legal-state},
Proposition~\ref{prop:step-009-termination},
Lemma~\ref{lem:step-004-occurrence}, and
Propositions~\ref{prop:step-004-lift} and
\ref{prop:step-004-projection}, the positive-dimensional quotient
construction has the following two mechanism interfaces.

The fixed allocations and thresholds are
\[
 \beta_{\rm AT}=\beta_{\rm SS}=\beta/4,
 \qquad \varepsilon_{\rm SS}=\varepsilon/8,
 \qquad
 \tau_{\rm AT}=\eta^{-1}\log((d+1)/\beta_{\rm AT}),
\]
\[
 \tau_{\rm SS}
 =\varepsilon_{\rm SS}^{-1}
   \log((kL+1)/\beta_{\rm SS}).
\tag{1}
\]
The least feasible teacher count satisfies the exact margin
\[
 \frac{k}{2}-\tau_{\rm AT}
 \geq B+\tau_{\rm SS}+2.
\tag{2}
\]
On every realizable path, the current lists determine attained occurrence
scores
\[
 s_r(\bar h):=|\{i\in[k]:\bar h\in\mathcal L_i^r\}|,
 \qquad
 q_r:=\max_{\bar h\in H_C}s_r(\bar h)
 \in\{0,\ldots,k\}.
\tag{3}
\]
On \(E_{\mathrm{good}}\), Proposition
~\ref{prop:step-009-termination} produces a stage
\(r^\circ\in\{0,\ldots,d\}\) and an actual
\(\bar f^\circ\in H_C\) with
\[
 \bar f^\circ\in\bigcap_{i=1}^k\mathcal L_i^{r^\circ},
 \qquad q_{r^\circ}=k.
\tag{4}
\]

For an arbitrary domain \(\mathcal U\), lists
\(\mathcal L_1,\ldots,\mathcal L_k\subseteq\mathcal U\),
\(\varepsilon_s>0\), and \(B\ge0\), define
\[
 s(u)=|\{i:u\in\mathcal L_i\}|\quad
 (u\in\textstyle\bigcup_i\mathcal L_i),
 \qquad s(\perp)=B,
\]
and sample from \(\bigcup_i\mathcal L_i\cup\{\perp\}\) with weights
\[
 \exp(\varepsilon_s s(u)).
 \tag{5}
\]
The exponent has no factor \(1/2\), and \(\perp\) is a separate support
point of score \(B\).  If every list has size at most \(L\) and
\[
 B\geq10\log(L/\delta_s)/\varepsilon_s,
\]
the Lyu~\citep{lyu2025} Sparse Sample interface is
\((2\varepsilon_s,\delta_s)\)-DP for addition, removal, or replacement of
one entire list coordinate.  In the present construction
\(\mathcal U=H_C\), \(\varepsilon_s=\varepsilon_{\rm SS}\),
\(\delta_s=\delta_{\rm SS}\), and the selected tuple is
\((\mathcal L_1^{\widehat r},\ldots,\mathcal L_k^{\widehat r})\).

The AboveThreshold interface draws independent
\(Z_r\sim\operatorname{Lap}(1/\eta)\) and reports \texttt{Above} exactly
when
\[
 q_r+Z_r\geq k/2;
 \tag{6}
\]
otherwise it reports \texttt{Below} and stops at the first \texttt{Above}.
For sensitivity-one adaptive queries and at most \(K\) counted crossings,
its privacy cost is
\[
 \left(\eta\,O(\sqrt{K\log(1/\delta_a)}+\log(1/\delta_a)),
 \delta_a\right),
\]
with the same Above/Below symmetry and concurrent composition as in
Lyu~\citep{lyu2025}.  Here it is one first-crossing process with \(K=1\),
threshold \(k/2\), and at most \(d+1\) stage queries.

The calibration identities and margin are
\[
 (d+1)e^{-\eta\tau_{\rm AT}}=\beta_{\rm AT},
 \qquad
 (kL+1)e^{-\varepsilon_{\rm SS}\tau_{\rm SS}}=\beta_{\rm SS},
 \qquad
 \frac{k}{2}-\tau_{\rm AT}\geq B+\tau_{\rm SS}+2.
 \tag{7}
\]
On a realizable source-good path, the stage scores are
\[
 s_r(\bar h)=|\{i\in[k]:\bar h\in\mathcal L_i^r\}|,
 \qquad q_r=\max_{\bar h\in H_C}s_r(\bar h),
\]
and the DDim descent result supplies a stage with an actual common
function of score \(k\).  An actual Sparse Sample output has a nonempty
finite producer-block occurrence set; the marked lift splits this set and
the projection sums the mark back to the released quotient law \(K_C\).
\end{proposition}

\begin{proof}
The two displayed mechanism laws and their list-coordinate privacy bounds
are Lyu~\citep{lyu2025}'s Algorithm~1/Lemma~3.1 Sparse Sample interface and
Algorithm~2/Lemma~3.2 AboveThreshold interface, with the current domain,
threshold, and privacy parameters substituted.
Lemma~\ref{lem:step-001-calibration} supplies the list cap, the exact
ceiling-defined failure score, and the first two identities in (7), while
Proposition~\ref{prop:step-001-teacher} supplies the last inequality with
all ceilings retained.  Proposition~\ref{prop:step-001-totalization}
defines the stopped transcript and every deterministic fallback, so these
interfaces are defined on arbitrary quotient inputs.

Lemma~\ref{lem:step-001-calibration} bounds every current exact list by
$L$, and Proposition~\ref{prop:step-001-totalization} replaces every
invalid list by the empty list. Hence the selected union has at most
$kL$ elements and the law in (5) has a positive finite normalizer. The
AboveThreshold statement here is conditional on sensitivity-one adaptive
queries; that list-interface condition is proved before the mechanism
accuracy result consumes it. Lemma~\ref{lem:step-009-legal-state}
identifies the scores and exact list membership on the realizable path, and
Proposition~\ref{prop:step-009-termination} supplies the score-\(k\) stage
without an additional confidence or sample charge.  Finally,
Lemma~\ref{lem:step-004-occurrence} gives the finite producer set and
Propositions~\ref{prop:step-004-lift} and
\ref{prop:step-004-projection} give its marked kernel and exact released
projection.  Raw-record locality and raw privacy are established separately
by the later record-locality propositions; the present result is the
quotient list/mechanism interface used for utility.
\end{proof}



\begin{lemma}[Sensitivity-one occurrence scores and finite mechanism support]\label{lem:step-010-score-support}
Under Lemma~\ref{lem:step-001-calibration}, Proposition~\ref{prop:step-001-totalization}, and Lemma~\ref{lem:step-009-legal-state}, fix any stage
\(r\in\{0,\ldots,d\}\). For a tuple
\(\mathbf L=(L_1,\ldots,L_k)\) of current sanitized lists, define

\[
 s_{\mathbf L}(\bar h)
 :=\sum_{i=1}^k\mathbf1\{\bar h\in L_i\},
 \qquad
 q(\mathbf L):=\max_{\bar h\in H_C}s_{\mathbf L}(\bar h).
 \tag{8}
\]

Then:

1. if \(\mathbf L\) and \(\mathbf L'\) differ in at most one entire list
   coordinate, then

   \[
     |q(\mathbf L)-q(\mathbf L')|\leq1;
     \tag{9}
   \]
2. every current list has size at most \(L\), so

   \[
     U(\mathbf L):=\bigcup_{i=1}^kL_i,
     \qquad |U(\mathbf L)|\leq kL;
     \tag{10}
   \]
3. the Sparse Sample support
   \(U(\mathbf L)\cup\{\perp\}\) has size at most \(kL+1\), is finite even
   when \(H_C\) is infinite, and has a positive finite normalizer under
   (5); and
4. if all lists are empty, then \(q(\mathbf L)=0\).

These are list-interface statements only. In particular, (9) does not
claim that one raw record replacement changes only one current list.


\end{lemma}

\begin{proof}
If \(\mathbf L\) and \(\mathbf L'\) differ only at coordinate
\(i_0\), then for every \(\bar h\in H_C\),

\[
 |s_{\mathbf L}(\bar h)-s_{\mathbf L'}(\bar h)|
 =|\mathbf1\{\bar h\in L_{i_0}\}
   -\mathbf1\{\bar h\in L'_{i_0}\}|
 \leq1.
 \tag{11}
\]

Consequently,

\[
 q(\mathbf L)
 =\max_{\bar h}s_{\mathbf L}(\bar h)
 \leq\max_{\bar h}(s_{\mathbf L'}(\bar h)+1)
 =q(\mathbf L')+1.
\]

Interchanging the tuples proves (9). This proof remains valid for ordered
list encodings and repeated positions because the score uses set
membership, not positional multiplicity.

At an adaptive transcript prefix consisting of \(r\) preceding \texttt{Below}
reports, both compared executions request the same stage-\(r\) query. Thus
(9) is the required sensitivity-one statement for every next query
conditioned on a common transcript prefix. It is still only a statement at
the list-tuple interface and makes no raw-record locality claim.

Lemma~\ref{lem:step-001-calibration} bounds every exact current
list by \(L\), while Proposition~\ref{prop:step-001-totalization} replaces every invalid,
nonfinite, non-\(H_C\), or oversized purported list by the empty list. Thus
the same cap holds on every totalized path. The finite union bound gives
(10), and adjoining one distinguished symbol gives support size at most
\(kL+1\). Every score is finite and belongs to \([0,k]\), so every
exponential weight in (5) is finite. The \(\perp\) weight
\(e^{\varepsilon_{\rm SS}B}\) is strictly positive; hence the finite sum of
weights is positive and finite, proving the normalizer claim. If every
list is empty, each summand in (8) is zero for every \(\bar h\), so the
maximum is zero. \(\square\)


\end{proof}

\begin{proposition}[One AboveThreshold transcript selects a score-separated
stage]\label{prop:step-010-abovethreshold}
Under Assumption~\ref{assump:approximate-dp-regime}, Proposition~\ref{prop:step-001-teacher}, Proposition~\ref{prop:step-009-termination}, and
Lemma~\ref{lem:step-010-score-support}, fix an arbitrary realizable
source-good path on which (4) holds. Realize in advance independent

\[
 Z_0,\ldots,Z_d\sim\operatorname{Lap}(1/\eta)
\]

and run the one first-crossing transcript (6). Define

\[
 E_{\rm AT}:=\bigcap_{r=0}^d\{|Z_r|\leq\tau_{\rm AT}\}.
 \tag{12}
\]

Then

\[
 \Pr(E_{\rm AT}^c)\leq\beta_{\rm AT}.
 \tag{13}
\]

On \(E_{\rm AT}\):

1. every reported \texttt{Below} stage satisfies \(q_r<k\);
2. every stage with \(q_r=k\) is reported \texttt{Above} if reached;
3. the stopped transcript has a first legal selected stage
   \(\widehat r\in\{0,\ldots,d\}\); and
4. its true selected score obeys

   \[
     q_{\widehat r}
     \geq\frac{k}{2}-\tau_{\rm AT}
     \geq B+\tau_{\rm SS}+2.
     \tag{14}
   \]

Thus an early \texttt{Above} at a stage with score below \(k\) is a safe false
positive rather than a mechanism failure: (14) is still sufficient for the
Sparse Sample call. An all-empty stage cannot be selected on
\(E_{\rm AT}\). Immediate score-\(k\) success is selected at stage \(0\),
and a first score-\(k\) stage at \(d\) is either reached and selected or is
preceded by an already safe false positive.


\end{proposition}

\begin{proof}
A random variable
\(Z\sim\operatorname{Lap}(1/\eta)\) has density
\((\eta/2)e^{-\eta|z|}\). Direct integration gives, for \(t\geq0\),

\[
 \Pr(|Z|>t)
 =2\int_t^\infty\frac{\eta}{2}e^{-\eta z}\,dz
 =e^{-\eta t}.
 \tag{15}
\]

There are at most \(d+1\) queries. Sampling all possible noises in advance
does not alter the adaptive stopped transcript, and a union bound with (7)
gives

\[
 \Pr(E_{\rm AT}^c)
 \leq(d+1)e^{-\eta\tau_{\rm AT}}
 =\beta_{\rm AT},
\]

which proves (13). No independence among stages is needed for this union
bound, although the source tests are independent.

The teacher margin (2) has a strictly positive right-hand side, so

\[
 \tau_{\rm AT}<k/2.
 \tag{16}
\]

If the process reports \texttt{Below} at stage \(r\), the exact threshold rule (6)
gives \(q_r+Z_r<k/2\). On \(E_{\rm AT}\),

\[
 q_r<k/2-Z_r\leq k/2+\tau_{\rm AT}<k,
 \tag{17}
\]

proving the first assertion. If \(q_r=k\), then on the same event

\[
 q_r+Z_r\geq k-\tau_{\rm AT}>k/2,
 \tag{18}
\]

so that stage is reported \texttt{Above}. Proposition~\ref{prop:step-009-termination} supplies at least one such stage
by stage \(d\). Therefore the stopped transcript has a first legal \texttt{Above}
no later than a score-\(k\) stage, proving existence of \(\widehat r\) and
excluding no-success exhaustion.

At the selected stage, the \texttt{Above} inequality and
\(Z_{\widehat r}\leq\tau_{\rm AT}\) give

\[
 q_{\widehat r}
 \geq k/2-Z_{\widehat r}
 \geq k/2-\tau_{\rm AT}.
\]

Combining this with (2) proves (14). If all lists at a stage are empty,
Lemma~\ref{lem:step-010-score-support} gives \(q_r=0\), and (16) yields

\[
 q_r+Z_r\leq\tau_{\rm AT}<k/2;
\]

hence that stage must report \texttt{Below} on \(E_{\rm AT}\). This also proves
the empty-stage boundary.

For endpoint handling, let \(r_{\min}\) be the least stage with score
\(k\), which exists by (4). If \(r_{\min}=0\), (18) selects the first query.
If \(r_{\min}=d\), then either a previous query reports \texttt{Above}, in which
case (14) makes it a safe selected stage, or all previous reports are
\texttt{Below} and (18) selects stage \(d\). Thus the transcript uses one sequence
of at most \(d+1\) queries and no per-stage success event. \(\square\)


\end{proof}

\begin{lemma}[Sparse Sample score tail and exclusion of the failure symbol]\label{lem:step-010-sparse-output}
Under Lyu~\citep{lyu2025} Algorithm 1 in the form (5), Lemma~\ref{lem:step-001-calibration}, and
Proposition~\ref{prop:step-010-abovethreshold}, condition on an arbitrary
source-good path, a transcript in \(E_{\rm AT}\), its selected stage
\(\widehat r\), and the selected current list tuple. Put

\[
 U_{\widehat r}:=\bigcup_{i=1}^k\mathcal L_i^{\widehat r},
 \qquad q:=q_{\widehat r},
\]

and let \(W\in U_{\widehat r}\cup\{\perp\}\) be the Sparse Sample outcome,
with \(s(\perp)=B\). Define

\[
 E_{\rm SS}:=\{s(W)\geq q-\tau_{\rm SS}\}.
 \tag{19}
\]

Then, uniformly over every such selected tuple,

\[
 \Pr(E_{\rm SS}^c\mid
       \text{selected lists and AboveThreshold transcript})
 \leq\beta_{\rm SS}.
 \tag{20}
\]

Moreover,

\[
 \Pr(W=\perp\mid
       \text{selected lists and AboveThreshold transcript})
 \leq
 e^{-2\varepsilon_{\rm SS}}
 \frac{\beta_{\rm SS}}{kL+1}
 \leq\beta_{\rm SS},
 \tag{21}
\]

and on \(E_{\rm SS}\),

\[
 W\in U_{\widehat r},
 \qquad
 W\in\mathcal L_i^{\widehat r}
 \text{ for at least one }i\in[k].
 \tag{22}
\]

Thus \(W\) is an actual current list member. It is an element of \(H_C\)
but need not be in \(\bar C\).


\end{lemma}

\begin{proof}
Equation (14) gives

\[
 q\geq B+\tau_{\rm SS}+2>0.
 \tag{23}
\]

Therefore \(U_{\widehat r}\) is nonempty and contains an attained
maximum-score item \(\bar h_*\) with \(s(\bar h_*)=q\). By
Lemma~\ref{lem:step-010-score-support}, the complete support has at most
\(kL+1\) points. Let

\[
 \mathcal B
 :=\{z\in U_{\widehat r}\cup\{\perp\}:
          s(z)<q-\tau_{\rm SS}\}.
\]

Under the exact law (5), the denominator contains the term
\(e^{\varepsilon_{\rm SS}q}\) of \(\bar h_*\), while every point in
\(\mathcal B\) has weight at most
\(e^{\varepsilon_{\rm SS}(q-\tau_{\rm SS})}\). Hence

\[
\begin{aligned}
 \Pr(E_{\rm SS}^c\mid\text{selected state})
 &=\frac{\sum_{z\in\mathcal B}
             e^{\varepsilon_{\rm SS}s(z)}}
          {e^{\varepsilon_{\rm SS}B}
           +\sum_{u\in U_{\widehat r}}
              e^{\varepsilon_{\rm SS}s(u)}}\\
 &\leq |\mathcal B|e^{-\varepsilon_{\rm SS}\tau_{\rm SS}}\\
 &\leq(kL+1)e^{-\varepsilon_{\rm SS}\tau_{\rm SS}}\\
 &=\beta_{\rm SS},
\end{aligned}
 \tag{24}
\]

where the last equality is (7). This is the full finite-effective-domain
tail; it charges low-score actual items as well as \(\perp\) and does not
invoke the source's commented utility statement.

For the exact failure-symbol mass, the same maximum-score denominator term
gives

\[
\begin{aligned}
 \Pr(W=\perp\mid\text{selected state})
 &=\frac{e^{\varepsilon_{\rm SS}B}}
         {e^{\varepsilon_{\rm SS}B}
          +\sum_{u\in U_{\widehat r}}
             e^{\varepsilon_{\rm SS}s(u)}}\\
 &\leq e^{-\varepsilon_{\rm SS}(q-B)}\\
 &\leq e^{-\varepsilon_{\rm SS}(\tau_{\rm SS}+2)}\\
 &=e^{-2\varepsilon_{\rm SS}}
   \frac{\beta_{\rm SS}}{kL+1},
\end{aligned}
 \tag{25}
\]

which proves (21) without dropping the extra teacher slack.

Finally, (23) gives

\[
 q-\tau_{\rm SS}\geq B+2>B=s(\perp).
 \tag{26}
\]

Thus \(E_{\rm SS}\) is incompatible with \(W=\perp\). By the exact support
of (5), \(W\ne\perp\) implies \(W\in U_{\widehat r}\), proving (22). The
current lists contain actual quotient functions in \(H_C\); an essential
SOA may be improper, so no conclusion \(W\in\bar C\) is made. \(\square\)


\end{proof}

\begin{proposition}[Mechanism-good event, actual selected item, and exact
fallback ledger]\label{prop:step-010-mechanism-good}
Under Assumption~\ref{assump:approximate-dp-regime}, 
Propositions~\ref{prop:step-001-totalization},
\ref{prop:step-004-lift}, \ref{prop:step-004-projection},
\ref{prop:step-006-good-event}, and
\ref{prop:step-010-abovethreshold}, and
Lemmas~\ref{lem:step-004-occurrence} and
\ref{lem:step-010-sparse-output}, define on a realizable source-good path

\[
 E_{\rm mech}:=E_{\rm AT}\cap E_{\rm SS},
 \tag{27}
\]
where \(E_{\rm SS}\) is declared false if no legal selected Sparse Sample
call occurs. Together with the source event produced by
Proposition~\ref{prop:step-006-good-event}, define the core event
\[
E_{\rm core}:=E_{\rm good}\cap E_{\rm mech}.
\tag{27a}
\]

Then, conditional on any fixed realizable sample and partition for which the
\(E_{\rm good}\) antecedent holds,

\[
 \Pr(E_{\rm mech}^c\mid\bar S,\mathcal P,E_{\rm good})
 \leq\beta_{\rm AT}+\beta_{\rm SS}
 =\frac{\beta}{2}.
 \tag{28}
\]

On \(E_{\rm mech}\), there are a selected stage
\(\widehat r\in\{0,\ldots,d\}\), a producer \(i\in[k]\), and a terminal
quotient output \(\bar H\in H_C\) such that

\[
 \bar H\in\mathcal L_i^{\widehat r}
 \subseteq\mathcal G_i
 :=\bigcup_{r=0}^d\mathcal L_i^r.
 \tag{29}
\]

The transcript is in the actual-output status event, not a
fallback status. Under the marked lift its occurrence set is nonempty,
\(J\in[k]\), and \(\bar H\in\mathcal G_J\); summing out \(J\) gives exactly
the same released \(K_C\) law. These statements remain status-correct if
the actual selected value happens to equal the default \(\bar c_0\).


\end{proposition}

\begin{proof}
Proposition~\ref{prop:step-010-abovethreshold} gives (13) and,
on \(E_{\rm AT}\), a legal selected tuple satisfying (14). Conditional on
every possible selected tuple and preceding AboveThreshold transcript,
Lemma~\ref{lem:step-010-sparse-output} gives (20). Therefore the tower
property and a union bound yield

\[
\begin{aligned}
 \Pr(E_{\rm mech}^c\mid\bar S,\mathcal P,E_{\rm good})
 &\leq\Pr(E_{\rm AT}^c)
   +\Pr(E_{\rm AT}\cap E_{\rm SS}^c
        \mid\bar S,\mathcal P,E_{\rm good})\\
 &\leq\beta_{\rm AT}+\beta_{\rm SS}.
\end{aligned}
 \tag{30}
\]

This argument is uniform over the adaptively selected stage; it does not
multiply \(\beta_{\rm SS}\) by \(d+1\). Equation (1) and the 
allocation give the final equality in (28).

On \(E_{\rm mech}\), equation (22) supplies a producer \(i\) and proves the
first inclusion in (29); the second is the definition of the producer's
all-stage union. The valid selected tuple is finite, its normalizer
is positive and finite by Lemma~\ref{lem:step-010-score-support}, and the
outcome lies in the exact categorical support. Consequently none of the
following totalized terminal branches occurs on \(E_{\rm mech}\):

1. an all-empty selected-stage tuple (excluded already by (14));
2. no \texttt{Above} or stage exhaustion (excluded by (18)); an invalid report or
   a named stage outside \(\{0,\ldots,d\}\) has probability zero under the
   exact finite threshold transcript (6);
3. a nonfinite, non-\(H_C\), oversized, or otherwise invalid selected list
   (the  exact current lists pass the sanitizer and obey the cap);
4. an invalid or zero Sparse Sample normalizer (excluded by finiteness and
   the positive \(\perp\) weight);
5. the failure symbol \(\perp\) (excluded by (26));
6. an out-of-support outcome or residual mechanism failure (the exact
   finite categorical law has support only
   \(U_{\widehat r}\cup\{\perp\}\)); or
7. the no-success or default terminal output as a status path.

Thus Proposition~\ref{prop:step-001-totalization} returns the same
actual \(\bar H=W\) rather than invoking \(\bar c_0\) as fallback. If
\(W=\bar c_0\) as a function, its categorical actual status, not its value,
still distinguishes it from fallback.

Lemma~\ref{lem:step-004-occurrence} now applies to this actual
transcript and gives a nonempty finite producer occurrence set. The marked
kernel chooses \(J\) only from that set, so \(J\in[k]\) and
\(\bar H\in\mathcal G_J\). Proposition~\ref{prop:step-004-projection} then removes the analysis-only
mark with exact equality of the released marginal. No privacy statement for
\((\bar H,J)\) is made.

The boundary cases do not create a hidden branch or factor:

\noindent \ If the first score-\(k\) stage is \(0\), it is selected immediately on
  \(E_{\rm AT}\). If it first occurs at \(d\), the process either safely
  selects an earlier false positive or selects stage \(d\).
\noindent \ If a queried stage has all lists empty, it reports \texttt{Below} on
  \(E_{\rm AT}\). Empty lists in only some blocks are harmless; the selected
  score in (14) still supplies an actual maximum-score item.
\noindent \ The candidate dictionary is defined at \(t=2\), and all sensitivity,
  support, and marking formulas remain well-formed there. For the realized
  least feasible teacher, however, \(k=2\) is impossible: (2) would require
  \(1-\tau_{\rm AT}\geq B+\tau_{\rm SS}+2\), whose left side is less than
  (1) and whose right side is greater than (2). Thus no proof line
  silently divides by \(k-1\) or assumes a nonexistent feasible \(k=2\)
  mechanism path.
\noindent \ The output may be improper. Actual membership means membership in a
  current \(H_C\)-valued essential list, not membership in \(\bar C\).

There is one master sample, one fixed block partition reused at all stages,
one AboveThreshold transcript, and one Sparse Sample call. The only stage
multiplicity in this argument is the \(d+1\) Laplace union already present
inside \(\tau_{\rm AT}\); no sample factor, privacy composition, or
additional confidence share is introduced. 
Lemma~\ref{lem:step-001-calibration} and
Proposition~\ref{prop:step-001-teacher} provide the exact source parameters,
list cap, confidence identities, and teacher margin. Proposition~\ref{prop:step-001-totalization} provides the one-transcript,
one-call procedure and every fallback branch. Lemma
~\ref{lem:step-010-score-support} proves directly that the maximum
occurrence query is sensitivity one at the list interface and that every
Sparse Sample effective support is finite of size at most \(kL+1\), even
though \(H_C\) can be infinite.

Proposition~\ref{prop:step-009-termination} supplies a literal
score-\(k\) actual function at some current stage, conditional on its
 realizable source-good antecedent. Proposition
~\ref{prop:step-010-abovethreshold} realizes all \(d+1\) possible Laplace
noises within one stopped process. Its simultaneous event costs exactly one
\(\beta_{\rm AT}\) share, makes every \texttt{Below} report structurally accurate
enough \((q_r<k)\), detects a score-\(k\) stage, and gives the selected-stage
lower bound

\[
 q_{\widehat r}\geq k/2-\tau_{\rm AT}
 \geq B+\tau_{\rm SS}+2.
\]

Thus empty stages cannot be selected on the good-noise event, while an
early false positive remains safe.

Lemma~\ref{lem:step-010-sparse-output} then applies the exact source
exponential weights to the one selected finite support. Its maximum-score
denominator term and support size give

\[
 \Pr[s(W)<q_{\widehat r}-\tau_{\rm SS}]
 \leq(kL+1)e^{-\varepsilon_{\rm SS}\tau_{\rm SS}}
 =\beta_{\rm SS},
\]

and the teacher margin puts the failure symbol's score strictly below the
 output-score threshold. Hence the good Sparse Sample outcome is
literally in one selected current list, not merely equal to a trace or a
fallback value.

The conditioning argument above composes these two finite tails on the
single selected list tuple and obtains the
mechanism ledger

\[
 \Pr(E_{\rm mech}^c\mid\bar S,\mathcal P,E_{\rm good})
 \leq\beta_{\rm AT}+\beta_{\rm SS}=\beta/2.
\]

The totalized branches are all covered by (29). Lemma~\ref{lem:step-004-occurrence} turns actual selected-stage membership
into a nonempty all-stage producer occurrence set, while Proposition~\ref{prop:step-004-projection} confirms that the analysis-only
mark changes no released mass.

Hence one adaptive AboveThreshold transcript and one Sparse Sample call
produce, on \(E_{\rm mech}\), an actual current quotient-list member rather
than fallback.



\end{proof}

\subsection{All-Input Privacy and Raw Transfer}\label{app:step_011}

The two mechanism charges fixed before sampling are
\[
 (\varepsilon_{\rm AT},\delta_{\rm AT})
   =(\varepsilon/4,\delta/2),
 \qquad
 (\varepsilon_{\rm sp},\delta_{\rm SS})
   =(\varepsilon/4,\delta/2),
\tag{1}
\]
where \(\varepsilon_{\rm sp}=2\varepsilon_{\rm SS}\) and
\(\varepsilon_{\rm SS}=\varepsilon/8\).  The raw learner kernel supplied
by Proposition~\ref{prop:step-003-raw-pullback} is
\[
 A_N(s,E)=K_C(T_N(s),E),
 \qquad E\in\mathcal H_C.
\tag{2}
\]
Every totalized current list has the all-input cap
\[
 |\mathcal L_i^r|\leq L
 \qquad(i\in[k],\ 0\leq r\leq d),
\tag{3}
\]
and is determined blockwise from
\[
 H_i^r
 =\{\bar h\in\bar C:
    \operatorname{err}_{\bar S_i}(\bar h)
    \leq\rho^{r+1}\gamma\},
 \qquad 0\leq r\leq d.
\tag{4}
\]
For finite lists \(L_1,\ldots,L_k\subseteq H_C\), define
\[
 q(L_1,\ldots,L_k)
 :=\max_{\bar h\in H_C}
   \sum_{i=1}^k\mathbf1\{\bar h\in L_i\}.
\tag{5}
\]
Lemma~\ref{lem:step-010-score-support} shows that this query changes by at
most one under replacement of one list coordinate and that the Sparse
Sample support has size at most \(kL+1\), including \(\perp\).

The Sparse Sample privacy interface of Lyu~\citep{lyu2025} is
\[
 B\geq\frac{10\log(L/\delta_s)}{\varepsilon_s}
\tag{6}
\]
for \((2\varepsilon_s,\delta_s)\)-DP under addition, removal, or
replacement of one list coordinate.  Here
\[
 \varepsilon_s=\varepsilon_{\rm SS}:=\varepsilon/8,
 \qquad \delta_s=\delta_{\rm SS}:=\delta/2,
 \qquad
 B=\left\lceil
 \frac{10\log(L/\delta_{\rm SS})}{\varepsilon_{\rm SS}}
 \right\rceil,
\tag{7}
\]
so the conditional selected-stage law has charge
\[
 (2\varepsilon_{\rm SS},\delta_{\rm SS})
 =(\varepsilon/4,\delta/2).
\tag{8}
\]
For sensitivity-one adaptive queries, the stopped AboveThreshold interface
in the same source has privacy parameters
\[
 \left(
 \eta\,O\!\left(
   \sqrt{K\log(1/\delta_a)}+\log(1/\delta_a)
 \right),
 \delta_a
 \right).
\tag{9}
\]
With \(K=1\), \(\delta_a=\delta_{\rm AT}:=\delta/2\), and the fixed
universal \(c_{\rm AT}\), its first coordinate is at most
\[
 c_{\rm AT}\eta
 \left(\sqrt{\log(1/\delta_{\rm AT})}
       +\log(1/\delta_{\rm AT})\right).
\tag{10}
\]
Finally,
\[
 g_\delta:=\log(4/\delta),
 \qquad
 \eta:=\frac{\varepsilon}
 {4c_{\rm AT}(\sqrt{g_\delta}+g_\delta)},
\tag{11}
\]
and \(\log(1/\delta_{\rm AT})=\log(2/\delta)\le g_\delta\), so (10)
is at most \(\varepsilon/4\).  The results below establish the missing
record-to-one-list premise and then compose these two interfaces without
using any utility event.

\begin{lemma}[Recordwise quotient adjacency and fixed-block localization]\label{lem:step-011-record-locality}
Under Assumption~\ref{assump:countable-evaluation-quotient} and
Proposition~\ref{prop:step-002-record-map}, let \(N\in\mathbb N_0\). If
\(s,s'\in Z_X^N\) are replace-one
neighbors, then \(T_N(s)\) and \(T_N(s')\) are equal or replace-one neighbors
in \(Z_Q^N\). If \(\bar s,\bar s'\in Z_Q^N\) are equal or replace-one
neighbors and a fixed partition \(\pi\) assigns the \(N\) coordinate indices
to \(k\) ordered blocks, then equal quotient tuples give equal ordered block
tuples; otherwise the block tuples are equal except possibly at the unique
block containing the changed coordinate. For \(N=0\), all these tuples are
the same empty tuple.


\end{lemma}

\begin{proof}
Write

\[
 s=((x_j,y_j))_{j=1}^N,
 \qquad
 s'=((x'_j,y'_j))_{j=1}^N.
\]

If the raw tuples are equal, their quotient images are equal. Otherwise,
there is one index \(j_*\) such that
\((x_j,y_j)=(x'_j,y'_j)\) for every \(j\ne j_*\). Since \(T_N\) acts
coordinatewise,

\[
 T_N(s)_j=(\kappa(x_j),y_j)
          =(\kappa(x'_j),y'_j)=T_N(s')_j
 \qquad(j\ne j_*).
\tag{12}
\]

In this latter case, at \(j_*\) the two quotient records may also coincide,
for example when
the raw points lie in the same evaluation cell and the labels agree. Thus
the quotient tuples are either equal or differ at precisely that one
coordinate. No realizability relation between \(x_j\) and \(y_j\) was used.

Now fix quotient tuples \(\bar s,\bar s'\) and an index partition \(\pi\). If
\(\bar s=\bar s'\), every ordered block projection is identical, so no
changed block need be named. Otherwise, let \(j_*\) be their unique differing
coordinate. Every coordinate belongs to exactly one block; call the block
containing \(j_*\) by \(i_*=i_*(\pi,j_*)\). The
ordered block projection for every \(i\ne i_*\) uses only unchanged
coordinates and is therefore identical under \(\bar s\) and \(\bar s'\).
Only block \(i_*\) can differ, and within it only one ordered record differs.
This conclusion is deterministic conditional on \(\pi\), holds for every
possible changed label, and is unaffected by duplicate records. If \(N=0\),
there is only the empty raw tuple, its empty quotient image, and the unique
empty block state; the assertions are identities. \(\square\)


\end{proof}

\begin{proposition}[All-stage one-list locality and sensitivity at every common
prefix]\label{prop:step-011-list-locality}
Under Proposition~\ref{prop:step-003-coding}, Lemma~\ref{lem:step-008-stage-map},  Propositions
~\ref{prop:step-008-essential-lists} and
~\ref{prop:step-008-list-envelope}, Lemma
~\ref{lem:step-010-score-support}, and Lemma
~\ref{lem:step-011-record-locality}, suppose \(d\geq1\). Fix a partition
\(\pi\) and quotient neighbors \(\bar s\sim\bar s'\). There is a block
\(i_*\in[k]\) such that, simultaneously for every \(0\leq r\leq d\),

\[
 \mathcal L_i^r(\bar s,\pi)
 =\mathcal L_i^r(\bar s',\pi)
 \qquad(i\ne i_*),
\tag{13}
\]

while both possibly different \(i_*\)-lists are finite subsets of \(H_C\)
of size at most \(L\). Consequently the two stage-\(r\) list tuples differ
by replacement of at most one entire list coordinate, and

\[
 |q_r(\bar s,\pi)-q_r(\bar s',\pi)|\leq1,
 \qquad 0\leq r\leq d,
\tag{14}
\]

where

\[
 q_r(\bar z,\pi)
 :=\max_{\bar h\in H_C}
   \sum_{i=1}^k
   \mathbf1\{\bar h\in\mathcal L_i^r(\bar z,\pi)\}.
\tag{15}
\]

At every common prefix of the stopped AboveThreshold transcript, both
executions request the same next stage \(r\), so (14) is the exact
sensitivity-one adaptive-query condition. The statement includes arbitrary
nonrealizable labels, empty restrictions and lists, pointwise sanitation,
fixed list orderings, repeated encodings, and duplicate functions across
different block lists.


\end{proposition}

\begin{proof}
If \(\bar s=\bar s'\), all block data, sanitized lists, and scores
are equal; since the positive branch has \(k\geq2\), choose any
\(i_*\in[k]\), and (13)-(14) follow. Hence assume that the quotient inputs
are strict replace-one neighbors. By
Lemma~\ref{lem:step-011-record-locality}, only the ordered data of their
unique changed block \(i_*\) can differ. The restriction definition in
Lemma~\ref{lem:step-008-stage-map} shows that \(H_i^r\) is a function only
of the ordered data in block \(i\), the fixed public stage, and the fixed
public thresholds. The binding construction fixes a
decomposition and a list ordering by a lookup on this local quotient-block
state. Therefore identical local states use identical restrictions,
decomposition choices, essential-function sets, and orderings. This proves
(13) for every \(i\ne i_*\) and every stage simultaneously.

If a local restriction is empty, Proposition
~\ref{prop:step-008-essential-lists} assigns that coordinate the empty
list. If a purported local list is nonfinite, outside \(H_C\), oversized, or
otherwise invalid, the  totalization sanitizes that one coordinate
to the empty list before either mechanism reads it. This sanitation is
pointwise in the same local block state and therefore cannot change a
second coordinate. For a nonempty valid restriction, the cited source
construction supplies its finite actual-function set; Proposition
~\ref{prop:step-008-list-envelope} gives
\(|\mathcal L_i^r|\le L\). Hence both possible
\(i_*\)-lists are legal and capped even though they may be completely
different. No stability of individual candidates is required.

For completeness, fix \(r\) and abbreviate the two list tuples by
\(\mathbf L=(L_1,\ldots,L_k)\) and
\(\mathbf L'=(L'_1,\ldots,L'_k)\). For every \(\bar h\in H_C\), (13) gives

\[
 \left|
  \sum_i\mathbf1\{\bar h\in L_i\}
  -\sum_i\mathbf1\{\bar h\in L'_i\}
 \right|
 =|\mathbf1\{\bar h\in L_{i_*}\}
   -\mathbf1\{\bar h\in L'_{i_*}\}|
 \leq1.
\tag{16}
\]

Taking maxima and applying (16) in each direction proves (14). The maximum
is well-defined even if \(H_C\) is infinite: if the union is nonempty it is
attained on the finite union of lists, and if all lists are empty its value
is zero. An algorithmic ordering is only an encoding of each finite set.
Repeated positions are setified by the  list interface, and the
score counts membership in a block list once. The same function appearing
in several different block lists is intentionally counted once per block
and does not affect (16).

Before the first \texttt{Above}, a common transcript prefix consists of the same
number \(r\) of preceding \texttt{Below} reports, so both executions request the
same setting-defined query \(q_r\). If a prefix is exhausted or invalid, no
further data query is made. Thus every next query at every common prefix is
either the same sensitivity-one function in (14) or a constant no-query
continuation. This argument is pointwise for all quotient inputs and never
uses \(E_{\rm good}\), realizability, or a mechanism-accuracy event.
\(\square\)


\end{proof}

\begin{proposition}[All-input privacy of the stopped AboveThreshold prefix]\label{prop:step-011-at-privacy}
Under Assumption~\ref{assump:approximate-dp-regime},
Lemma~\ref{lem:step-001-calibration},
Propositions~\ref{prop:step-010-interfaces} and
\ref{prop:step-011-list-locality}, and Lyu~\citep{lyu2025}
Lemma~3.2, fix a
data-independent partition \(\pi\). Let \(M_{\rm AT}^{\pi}\) be the finite
transcript kernel that queries \(q_0,\ldots,q_d\) in order, stops at its
first \texttt{Above}, and otherwise records exhaustion after stage \(d\). Include
the selected stage or the no-selection tag in the transcript. Then, on all
quotient inputs and for every measurable transcript event \(F\),

\[
 M_{\rm AT}^{\pi}(\bar s,F)
 \leq e^{\varepsilon/4}
       M_{\rm AT}^{\pi}(\bar s',F)+\delta/2
 \qquad(\bar s\sim\bar s').
\tag{17}
\]

The same inequality holds with \(\bar s,\bar s'\) interchanged.


\end{proposition}

\begin{proof}
At every common transcript prefix, Proposition
~\ref{prop:step-011-list-locality} gives exactly the sensitivity-one
premise of Lyu~\citep{lyu2025} Lemma 3.2. The process counts at most one \texttt{Above}. Its
termination after at most \(d+1\) queries and its retention of only the
first-crossing prefix are deterministic postprocessings of the source
transcript; neither operation increases privacy loss.

Put \(\delta_{\rm AT}=\delta/2\). Lemma~\ref{lem:step-001-calibration}
sets
\[
g_\delta=\log(4/\delta),
\qquad
\eta=\frac{\varepsilon}
 {4c_{\rm AT}(\sqrt{g_\delta}+g_\delta)}.
\]
Moreover,

\[
 \log(1/\delta_{\rm AT})=\log(2/\delta)
 \leq \log(4/\delta)=g_\delta
\]

Proposition~\ref{prop:step-010-interfaces} gives the \(K=1\)
AboveThreshold cost bound, so

\[
\begin{aligned}
 \varepsilon_{\rm AT}
 &\leq c_{\rm AT}\eta
   \left(\sqrt{\log(1/\delta_{\rm AT})}
         +\log(1/\delta_{\rm AT})\right)\\
 &\leq c_{\rm AT}
   \frac{\varepsilon}
        {4c_{\rm AT}(\sqrt{g_\delta}+g_\delta)}
   (\sqrt{g_\delta}+g_\delta)
 =\frac{\varepsilon}{4}.
\end{aligned}
\tag{18}
\]

Lyu~\citep{lyu2025} Lemma 3.2 therefore gives (17). Its DP statement is for every ordered
neighbor pair and every transcript event. Because replacement adjacency is
symmetric, applying the same statement to the ordered pair
\((\bar s',\bar s)\) gives the reverse inequality; it is not obtained by
algebraically reversing (17). Empty-list scores, arbitrary labels, and an
eventual no-selection transcript are ordinary inputs and outputs of this
same source process. No accuracy condition on its Laplace noises was
conditioned upon. \(\square\)


\end{proof}

\begin{proposition}[Uniform selected-stage Sparse Sample privacy and totalized
no-call continuations]\label{prop:step-011-sparse-privacy}
Under Assumption~\ref{assump:approximate-dp-regime},
Lemma~\ref{lem:step-001-calibration},
Propositions~\ref{prop:step-010-interfaces} and
\ref{prop:step-011-list-locality}, and Lyu~\citep{lyu2025}
Lemma~3.1, fix a
partition \(\pi\) and any possible AboveThreshold history \(t\). Define a
total conditional \(H_C\)-valued kernel \(M_{2,t}^{\pi}\) as follows:

1. if \(t\) selects a stage \(r\in\{0,\ldots,d\}\), apply Sparse Sample to
   the sanitized tuple
   \((\mathcal L_1^r,\ldots,\mathcal L_k^r)\) with
   \(\varepsilon_{\rm SS}=\varepsilon/8\),
   \(\delta_{\rm SS}=\delta/2\), and
   \[
   B=\left\lceil
      \frac{10\log(L/\delta_{\rm SS})}{\varepsilon_{\rm SS}}
     \right\rceil,
   \]
   then map
   \(\perp\) to \(\bar c_0\) and retain every actual \(H_C\)-item; and
2. if \(t\) records no selection, exhaustion, an out-of-range stage, or an
   invalid/no-call tag, return \(\bar c_0\) deterministically.

Then for every history \(t\), every quotient neighboring pair, and every
\(E\in\mathcal H_C\),

\[
 M_{2,t}^{\pi}(\bar s,E)
 \leq e^{\varepsilon/4}
       M_{2,t}^{\pi}(\bar s',E)+\delta/2.
\tag{19}
\]

This uniform statement does not require the selected stages in two actual
neighboring executions to agree.


\end{proposition}

\begin{proof}
First suppose that \(t\) selects the common history-indexed stage
\(r\). Proposition~\ref{prop:step-011-list-locality} says that the two
stage-\(r\) input tuples differ in at most one entire list coordinate, and
Proposition~\ref{prop:step-008-list-envelope} bounds every list by \(L\).
Lemma~\ref{lem:step-001-calibration} defines the displayed ceiling, which
implies the exact Sparse Sample threshold in
Proposition~\ref{prop:step-010-interfaces}. Lyu~\citep{lyu2025} Lemma~3.1
therefore makes the Sparse Sample law
on the fixed global space \(H_C\cup\{\perp\}\)

\[
 (2\varepsilon_{\rm SS},\delta_{\rm SS})
 =(\varepsilon/4,\delta/2)\text{-DP}.
\tag{20}
\]

The deterministic map

\[
 \phi(z):=
 \begin{cases}
 z,&z\in H_C,\\
 \bar c_0,&z=\perp
 \end{cases}
\tag{21}
\]

is measurable. For any \(E\in\mathcal H_C\), the probability after
totalization is the pre-totalization probability of
\(\phi^{-1}(E)\). Applying (20) to this measurable preimage proves (19).
This remains valid when an actual categorical item equals \(\bar c_0\):
the preimage then contains both that actual item and \(\perp\), and
postprocessing still applies to the union.

All stage lists have already been sanitized before the occurrence query and
before Sparse Sample. Thus a private-data-dependent invalid list does not
choose between an unaccounted mechanism and a constant branch: its one
coordinate is replaced by the empty list, and the same source mechanism is
run on the resulting legal tuple. If all selected lists are empty, the
source support is only \(\{\perp\}\), its normalizer is the positive weight
\(e^{\varepsilon_{\rm SS}B}\), and (21) makes the continuation constant.
Partially empty tuples require no special case. A zero normalizer or an
out-of-support categorical outcome has probability zero under the exact
source law; its predeclared fallback is still the same measurable
postprocessing.

If \(t\) is a no-call, exhausted, out-of-range, or invalid history, the
conditional law is \(\delta_{\bar c_0}\) for every input. It is
\((0,0)\)-DP and hence satisfies (19). We define this constant continuation
even for histories having probability zero so that the adaptive family is
total on the whole finite transcript space.

The proof is indexed by a \emph{common value} \(t\) of the first transcript. If
one actual execution selects stage \(r\) and the other selects stage \(r'\),
we never compare \(\mathcal L^r(\bar s)\) directly with
\(\mathcal L^{r'}(\bar s')\). The private first transcript accounts for the
different-history probabilities; for each possible shared history, the
conditional second kernel satisfies (19). This is precisely the interface
used by adaptive composition. No actual-output or good-noise event has been
used. \(\square\)


\end{proof}

\begin{lemma}[Exact finite-transcript adaptive composition and
postprocessing]\label{lem:step-011-adaptive-composition}
Let \(M_1\) be a kernel from a database space to a finite discrete
transcript space \((\mathsf T,2^{\mathsf T})\). For every
\(t\in\mathsf T\), let \(M_{2,t}\)
be a kernel from the same database space to a measurable space
\((\mathsf Y,\mathcal Y)\). Let \(\varepsilon_1,\varepsilon_2\geq0\) and
\(0\leq\delta_1,\delta_2\leq1\). Suppose, for every ordered neighboring pair,
\(M_1\) is \((\varepsilon_1,\delta_1)\)-DP and every \(M_{2,t}\) is
\((\varepsilon_2,\delta_2)\)-DP, uniformly in \(t\). Then the adaptive joint
law

\[
 \mathsf J_x(dt,dy):=M_1(x,dt)M_{2,t}(x,dy)
\tag{22}
\]

is
\((\varepsilon_1+\varepsilon_2,\delta_1+\delta_2)\)-DP. For every
measurable output space \((\mathsf O,\mathcal O)\) and measurable map
\(g:\mathsf T\times\mathsf Y\to\mathsf O\), the
postprocessed law \(g_\#\mathsf J_x\) has the same guarantee.


\end{lemma}

\begin{proof}
Fix an ordered neighbor pair \(x\sim x'\). Write \(P,Q\) for
the two first-transcript laws. The one-direction DP inequality is

\[
 P(F)\leq e^{\varepsilon_1}Q(F)+\delta_1
 \qquad(F\subseteq\mathsf T).
\tag{23}
\]

We record an exact submeasure form of (23). With
\(\lambda=P+Q\), let \(p=dP/d\lambda\), \(q=dQ/d\lambda\), and define

\[
 dP^0:=\min\{p,e^{\varepsilon_1}q\}\,d\lambda.
\tag{24}
\]

Then \(P^0\leq P\), \(P^0\leq e^{\varepsilon_1}Q\), and

\[
\begin{aligned}
 P(\mathsf T)-P^0(\mathsf T)
 &=\int(p-e^{\varepsilon_1}q)_+\,d\lambda\\
 &=\sup_{F\subseteq\mathsf T}
   \{P(F)-e^{\varepsilon_1}Q(F)\}
 \leq\delta_1.
\end{aligned}
\tag{25}
\]

For each \(t\in\mathsf T\), apply the same construction to the two
conditional second-output laws

\[
 P_t:=M_{2,t}(x,\cdot),
 \qquad Q_t:=M_{2,t}(x',\cdot).
\]

It yields a subprobability \(P_t^0\leq P_t\) satisfying

\[
 P_t^0\leq e^{\varepsilon_2}Q_t,
 \qquad
 P_t^0(\mathsf Y)\geq1-\delta_2.
\tag{26}
\]

Because \(\mathsf T\) is finite, these finitely many choices form a
measurable history-indexed family without a selection issue. Define the
joint submeasure

\[
 \mathsf J_x^0(dt,dy):=P^0(dt)P_t^0(dy).
\tag{27}
\]

It obeys \(\mathsf J_x^0\leq\mathsf J_x\), and (25)-(26) give

\[
 \mathsf J_x^0(\mathsf T\times\mathsf Y)
 \geq(1-\delta_1)(1-\delta_2)
 \geq1-\delta_1-\delta_2.
\tag{28}
\]

For any measurable \(G\subseteq\mathsf T\times\mathsf Y\), with section
\(G_t:=\{y:(t,y)\in G\}\), (24) and (26) imply

\[
\begin{aligned}
 \mathsf J_x^0(G)
 &=\sum_{t\in\mathsf T}P^0(\{t\})P_t^0(G_t)\\
 &\leq e^{\varepsilon_1+\varepsilon_2}
   \sum_{t\in\mathsf T}Q(\{t\})Q_t(G_t)\\
 &=e^{\varepsilon_1+\varepsilon_2}\mathsf J_{x'}(G).
\end{aligned}
\tag{29}
\]

Since the missing mass of the submeasure in (27) is at most
\(\delta_1+\delta_2\),

\[
\begin{aligned}
 \mathsf J_x(G)
 &\leq \mathsf J_x^0(G)
       +(\mathsf J_x-\mathsf J_x^0)
          (\mathsf T\times\mathsf Y)\\
 &\leq e^{\varepsilon_1+\varepsilon_2}
       \mathsf J_{x'}(G)+\delta_1+\delta_2.
\end{aligned}
\tag{30}
\]

This derivation explains why the additive term is
\(\delta_1+\delta_2\), with no erroneous exponential multiplier. Repeating
it for the reversed ordered pair supplies the reverse DP inequality.

Finally, for a measurable output event \(E\in\mathcal O\),

\[
 (g_\#\mathsf J_x)(E)
 =\mathsf J_x(g^{-1}(E)).
\]

Using \(G=g^{-1}(E)\) in (30) proves the postprocessing statement for every
measurable event. \(\square\)


\end{proof}

\begin{proposition}[All-input quotient privacy after adaptive composition and
partition mixing]\label{prop:step-011-quotient-dp}
Under Assumption~\ref{assump:approximate-dp-regime},  Propositions
~\ref{prop:step-003-coding} and
~\ref{prop:step-003-quotient-kernel}, Propositions
~\ref{prop:step-011-at-privacy} and
~\ref{prop:step-011-sparse-privacy}, and Lemma
~\ref{lem:step-011-adaptive-composition}, the positive-branch released
quotient kernel satisfies, for all quotient replace-one neighbors and all
\(E\in\mathcal H_C\),

\[
 K_C(\bar s,E)
 \leq e^{\varepsilon/2}K_C(\bar s',E)+\delta.
\tag{31}
\]

This holds on every input, without realizability or conditioning on any
generated event.


\end{proposition}

\begin{proof}
Fix a partition \(\pi\). Use the finite first transcript

\[
 \mathsf T_{\rm AT}
 =\{t_r:0\leq r\leq d\}\cup\{t_{\rm none}\}
\tag{32}
\]

where \(t_r\) records \(r\) preceding \texttt{Below} reports followed by \texttt{Above}
at stage \(r\), and \(t_{\rm none}\) records \(d+1\) \texttt{Below} reports. Any
predeclared zero-probability invalid transcript tags can be adjoined to this
finite space with the constant continuation of Proposition
~\ref{prop:step-011-sparse-privacy}.

Proposition~\ref{prop:step-011-at-privacy} gives

\[
 (\varepsilon_1,\delta_1)
 =(\varepsilon/4,\delta/2)
\]

for the first kernel. Proposition
~\ref{prop:step-011-sparse-privacy} gives, uniformly for every \(t\),

\[
 (\varepsilon_2,\delta_2)
 =(\varepsilon/4,\delta/2)
\]

for the totalized second kernel; on no-call histories this is a conservative
upper bound on its exact \((0,0)\) cost. Lemma
~\ref{lem:step-011-adaptive-composition} therefore makes the augmented
history/output law

\[
 (\varepsilon_1+\varepsilon_2,
   \delta_1+\delta_2)
 =(\varepsilon/2,\delta)\text{-DP}.
\tag{33}
\]

The released quotient output is a measurable postprocessing: retain the
Sparse Sample \(H_C\)-item on a selected actual outcome and otherwise return
\(\bar c_0\). Proposition~\ref{prop:step-003-coding} supplies the
measurable status and output coordinates. Lemma
~\ref{lem:step-011-adaptive-composition} therefore proves (31) for the
conditional released kernel \(K_C^{\pi}\).

The terminal routes before removing \(\pi\) are as follows:

1. An empty restriction is represented by an empty list before either
   mechanism. An all-empty stage still has the legal query value zero.
2. An \texttt{Above} on an all-empty or partially empty tuple invokes the same
   private Sparse Sample kernel; an all-empty union returns only \(\perp\),
   which postprocesses to \(\bar c_0\).
3. If no \texttt{Above} occurs by stage \(d\), the private first transcript is
   followed by the constant no-call output \(\bar c_0\).
4. A selected Sparse Sample \(\perp\) outcome maps to \(\bar c_0\). An
   actual selected function equal to \(\bar c_0\) may collide in released
   value with this fallback, but collision under a deterministic map cannot
   increase privacy loss.
5. A nonfinite, non-\(H_C\), oversized, or otherwise invalid local list is
   sanitized coordinatewise before querying. Proposition
   \ref{prop:step-008-list-envelope} guarantees
   that valid nonempty restrictions in fact produce finite capped lists on
   every input, so the remaining structural-invalid tags are unreachable;
   their predeclared constant continuation is harmless.
6. An out-of-range transcript, zero normalizer, or out-of-support mechanism
   outcome has probability zero under the exact finite mechanisms. Its
   totalized value is nevertheless defined and is included in the same
   measurable postprocessing.

Thus no private-data-dependent branch outside the two composed mechanisms
is released. In particular, none of these cases assumes that the output
is an actual list item, that the first crossing is accurate, or that a
source-good path occurred.

It remains to average over the uniform data-independent partition. Let
\(\Pi_N\) be its finite set of index partitions and let \(\nu\) be its law,
which is the same for \(\bar s\) and \(\bar s'\). For every fixed
\(\pi\in\Pi_N\), the preceding argument gives

\[
 K_C^{\pi}(\bar s,E)
 \leq e^{\varepsilon/2}K_C^{\pi}(\bar s',E)+\delta.
\tag{34}
\]

Consequently,

\[
\begin{aligned}
 K_C(\bar s,E)
 &=\sum_{\pi\in\Pi_N}\nu(\pi)K_C^{\pi}(\bar s,E)\\
 &\leq e^{\varepsilon/2}
       \sum_{\pi\in\Pi_N}\nu(\pi)K_C^{\pi}(\bar s',E)
       +\delta\sum_{\pi\in\Pi_N}\nu(\pi)\\
 &=e^{\varepsilon/2}K_C(\bar s',E)+\delta.
\end{aligned}
\tag{35}
\]

The additive term is paid once, not once per partition. Mechanism coins are
already integrated inside each conditional kernel. The proof applies to
every quotient input atom, including arbitrary labels and duplicate
records. Since \(E\) was arbitrary in \(\mathcal H_C\), (31) holds for all
measurable outputs, whether proper, improper, default-valued, or outside all
current lists. \(\square\)


\end{proof}

\begin{proposition}[Exact raw pullback, privacy monotonicity, and null boundary]\label{prop:step-011-raw-dp}
Under Assumption~\ref{assump:approximate-dp-regime}, Proposition
~\ref{prop:step-003-raw-pullback}, Lemma
~\ref{lem:step-011-record-locality}, and Proposition
~\ref{prop:step-011-quotient-dp}, the setting-defined raw learner kernel
\(A_N\) is \((\varepsilon/2,\delta)\)-DP and hence
\((\varepsilon,\delta)\)-DP for every raw replace-one pair and every
\(E\in\mathcal H_C\). If \(d=0\), the setting uses \(N=0\) and the law is
exactly \((0,0)\)-DP.


\end{proposition}

\begin{proof}
Proposition~\ref{prop:step-003-raw-pullback} gives the
exact kernel identity (2). Fix raw neighbors \(s\sim s'\) and an arbitrary
\(E\in\mathcal H_C\). Lemma~\ref{lem:step-011-record-locality} leaves two
cases.

If \(T_N(s)=T_N(s')\), then exact pullback gives

\[
 A_N(s,E)=K_C(T_N(s),E)
         =K_C(T_N(s'),E)=A_N(s',E).
\tag{36}
\]

This covers a raw point replacement within one quotient cell when the
label is unchanged. If the labels differ, the quotient records differ and
the next case applies.

Otherwise the quotient images are distinct, and Lemma
~\ref{lem:step-011-record-locality} makes them strict replace-one neighbors.
Proposition
~\ref{prop:step-011-quotient-dp} gives

\[
\begin{aligned}
 A_N(s,E)
 &=K_C(T_N(s),E)\\
 &\leq e^{\varepsilon/2}K_C(T_N(s'),E)+\delta\\
 &=e^{\varepsilon/2}A_N(s',E)+\delta.
\end{aligned}
\tag{37}
\]

Equations (36)-(37) are pointwise for arbitrary labels; neither case uses a
data distribution or target concept. Adjacency is symmetric, so applying
the same two-case proof to the ordered pair \((s',s)\) gives the reverse
inequality. Again, the reverse is not inferred by rearranging an
approximate-DP inequality.

Because \(0<\varepsilon\),
\(e^{\varepsilon/2}\leq e^\varepsilon\). Therefore (37) implies

\[
 A_N(s,E)\leq e^\varepsilon A_N(s',E)+\delta,
\tag{38}
\]

which is exactly the setting's raw replacement-
\((\varepsilon,\delta)\)-DP definition for the released \(H_C\)-valued
output. No decoder is needed for (38); applying the fixed decoder after
release would itself be postprocessing.

For \(d=0\), Proposition
~\ref{prop:step-003-quotient-kernel} supplies \(N=0\) and the exact law
\(\delta_{\bar c_0}\) on the unique empty quotient input. Proposition~\ref{prop:step-003-raw-pullback} gives the same Dirac law on the
unique empty raw input. Thus all event probabilities are identical and the
kernel is \((0,0)\)-DP; no partition, AboveThreshold query, or Sparse Sample
call exists.

For the formal \(k=2\) positive-branch boundary, the privacy proof remains
well-formed: Lyu~\citep{lyu2025} Sparse Sample permits every \(k\geq1\), the list cap and
one-coordinate relation remain valid, AboveThreshold counts one crossing,
and no line divides by \(k-1\) or uses the teacher utility margin. Whether
the  least-feasible utility calibration actually selects \(k=2\)
is irrelevant to privacy. First-stage selection, stage-\(d\) selection,
different selected stages, no selection, and every fallback have already
been included in the finite transcript composition. This closes every
claimed boundary without changing the mechanism. 
Lemma~\ref{lem:step-011-record-locality} first proves the exact pointwise
transport: one raw replacement becomes either no
change or one quotient replacement, and, conditional on a fixed
data-independent index partition, only one ordered block can change.
Proposition~\ref{prop:step-011-list-locality} combines that fact with the
 block-local restriction lists, fixed choices, empty-list
totalization, and cap. It proves simultaneously over all \(d+1\) stages
that only one entire list coordinate can change. The direct membership
calculation then gives sensitivity one for every next maximum-occurrence
query at every common transcript prefix. This derivation holds for arbitrary
labels and does not consume eventwise inclusion from the restriction-list
proposition.

Proposition~\ref{prop:step-011-at-privacy} applies Lyu~\citep{lyu2025} Lemma 3.2
with \(K=1\), \(\delta_{\rm AT}=\delta/2\), and the exact  \(\eta\).
The displayed calculation (18) gives the first allocation
\((\varepsilon/4,\delta/2)\) without multiplying privacy by \(d+1\).
Proposition~\ref{prop:step-011-sparse-privacy} applies Lyu~\citep{lyu2025} Lemma 3.1
at every possible common selected-stage history with
\(\varepsilon_{\rm SS}=\varepsilon/8\),
\(\delta_{\rm SS}=\delta/2\), the exact ceiling \(B\), and one-list
replacement. It gives the second allocation
\((\varepsilon/4,\delta/2)\). No-call and invalid-history continuations are
constant, while \(\perp\) and all other fallbacks are postprocessed to
\(\bar c_0\).

Lemma~\ref{lem:step-011-adaptive-composition} proves, with an explicit
submeasure argument, that these history-dependent kernels compose to

\[
 (\varepsilon/4+\varepsilon/4,
   \delta/2+\delta/2)
 =(\varepsilon/2,\delta)\text{-DP}.
\]

This proof handles differing stopping times and selected stages because it
compares the second kernel only at a common transcript value and charges
the transcript mismatch to the first private kernel. Proposition
~\ref{prop:step-011-quotient-dp} then applies measurable postprocessing and
averages over the common data-independent partition, paying neither a new
privacy term nor a partition multiplier. It yields the all-input quotient
inequality (31) for every \(E\in\mathcal H_C\).

Finally, Proposition~\ref{prop:step-003-raw-pullback} and the quotient
privacy inequality just proved transfer (31) exactly through \(T_N\).
Equal quotient images give equality of laws; neighboring images use (31).
Thus the raw learner is in fact \((\varepsilon/2,\delta)\)-DP and therefore
is \((\varepsilon,\delta)\)-DP. The exact
\(d=0,N=0\) branch is \((0,0)\)-DP. This completes the all-input privacy
and raw-transfer argument.



\end{proof}

\subsection{Exact SOA Identity and Empirical Utility}\label{app:step_012}

Fix the positive branch \(d\ge1\), an indexed quotient master sample, and
its partition:
\[
 \bar S=((q_j,y_j))_{j=1}^{n_0},
 \qquad \mathcal P=(B_1,\ldots,B_k),
 \qquad |B_i|=m,
 \qquad n_0=km.
\tag{1}
\]
For \(f\in H_C\), define the master and block empirical errors
\[
 e_{\bar S}(f)
 :=\frac1{n_0}\sum_{j=1}^{n_0}
       \mathbf1\{f(q_j)\ne y_j\},
 \qquad
 e_i(f)
 :=\frac1m\sum_{j\in B_i}
       \mathbf1\{f(q_j)\ne y_j\}.
\tag{2}
\]
Retain
\[
 \gamma:=\frac{\alpha}{16},
 \qquad
 a_d:=\frac1{5d},
 \qquad
 \rho:=1-\frac1{2d},
\tag{3}
\]
and
\[
 H_i^r:=\{\bar c\in\bar C:
              e_i(\bar c)\leq\rho^{r+1}\gamma\},
 \qquad
 p_r:=2^rn_0d,
 \qquad 0\leq r\leq d.
\tag{4}
\]
Write \(E_{\mathrm{core}}:=E_{\mathrm{good}}\cap E_{\mathrm{mech}}\)
only for the joint path condition produced by the trace and mechanism
results.  The statements proved below are that an actual selected output
has the exact form
\[
 \bar H=\operatorname{SOA}_{\mathcal G}
 \quad\text{for a nonempty }\mathcal G\subseteq\bar C
 \text{ that is both }n_0\text{- and }(d+1)\text{-irreducible},
\tag{5}
\]
that
\[
 \bar H\in
 \widehat C_{d+1}
 :=\left\{\operatorname{SOA}_{\mathcal A}:
       \varnothing\ne\mathcal A\subseteq\bar C
       \text{ is }(d+1)\text{-irreducible}\right\},
 \qquad
 \operatorname{LD}(\widehat C_{d+1})\leq d,
\tag{6}
\]
and that
\[
 e_{\bar S}(\bar H)\leq2\gamma=\frac{\alpha}{8}.
\tag{7}
\]
The family in (6) is fixed by \(\bar C\) and may be improper.

On \(E_{\mathrm{good}}\), Proposition
~\ref{prop:step-006-good-event} gives, for every actual
\(\bar c\in\bar C\) and every \(i\in[k]\),
\[
\begin{cases}
 (1-a_d)e_{\bar S}(\bar c)
       \leq e_i(\bar c)
       \leq(1+a_d)e_{\bar S}(\bar c),
       &e_{\bar S}(\bar c)>\gamma/3,\\
 0\leq e_i(\bar c)\leq\gamma/2,
       &e_{\bar S}(\bar c)\leq\gamma/3.
\end{cases}
\tag{8}
\]
For every \(\bar f\in\mathcal L_i^r\), Proposition
~\ref{prop:step-008-leaf-scale} supplies a maximal leaf \(\mathcal G\)
with
\[
 \bar f=\operatorname{SOA}_{\mathcal G},
 \qquad
 \mathcal G\subseteq H_i^r\subseteq\bar C,
 \qquad
 \mathcal G\text{ is }p_r2^{d-t}\text{-irreducible},
\tag{9}
\]
where
\(t=\operatorname{DDim}_{p_r,d}(H_i^r)
   =\operatorname{LD}(\mathcal G)\), and
\[
 p_r2^{d-t}\geq p_0=n_0d\geq\max\{n_0,d+1\}.
\tag{10}
\]
On \(E_{\mathrm{core}}\), Proposition
~\ref{prop:step-010-mechanism-good} gives a selected stage and producer
with
\[
 \bar H\in\mathcal L_i^{\widehat r}.
\tag{11}
\]
The exact source/current identities from
Lemma~\ref{lem:step-008-stage-map} are
\[
 H_i^r=H_{i,\mathrm{src}}^{r+1},
 \qquad
 p_r=2^rn_0d=\frac{p_{r+1,\mathrm{src}}}{2},
\tag{12}
\]
and, only when \(r<d\),
\[
 p_{r+1}=2p_r.
\tag{13}
\]

For a nonempty class \(\mathcal A\) with
\(t=\operatorname{DDim}_{p,d}(\mathcal A)\), Lyu's essential-function
definition~\citep{lyu2025} gives, in every optimal decomposition, a leaf
\(\mathcal A_\ell\) satisfying
\[
 \operatorname{LD}(\mathcal A_\ell)=t,
 \qquad
 \operatorname{SOA}_{\mathcal A_\ell}=f
 \quad\text{pointwise on all of }Q_C.
\tag{14}
\]
The empirical restriction argument in the proof of Lyu's Theorem~3 is
instantiated here with
\[
 \mathcal X=Q_C,
 \quad S=\bar S,
 \quad n=n_0,
 \quad h=\bar H,
 \quad u=\gamma.
\tag{15}
\]
It says that if a nonempty \(\mathcal G\) is \(n\)-irreducible,
\(h=\operatorname{SOA}_{\mathcal G}\), and every \(g\in\mathcal G\)
has empirical error at most \(2u\), then \(h\) has empirical error at
most \(2u\); the contradiction is reproduced below before the interface is
used. Finally, Lyu~\citep[Lemma~4.2]{lyu2025} defines
\[
 \widehat{\mathcal H}_{d+1}
 :=\{\operatorname{SOA}_{\mathcal A}:
       \varnothing\ne\mathcal A\subseteq\mathcal H
       \text{ is }(d+1)\text{-irreducible}\}
\tag{16}
\]
and concludes
\[
 \operatorname{LD}(\widehat{\mathcal H}_{d+1})\leq d
\tag{17}
\]
whenever \(\operatorname{LD}(\mathcal H)\leq d\).  We apply it only with
\(\mathcal H=\bar C\).

\begin{proposition}[Actual selected item has an exact current maximal-leaf
SOA]\label{prop:step-012-selected-leaf}
Under Assumption~\ref{assump:finite-littlestone},
Lemma~\ref{lem:step-008-stage-map},
Propositions~\ref{prop:step-006-good-event},
\ref{prop:step-008-essential-lists},
\ref{prop:step-008-leaf-scale}, and
\ref{prop:step-010-mechanism-good}, suppose \(d\ge1\) and the fixed path
lies in \(E_{\mathrm{core}}=E_{\mathrm{good}}\cap E_{\mathrm{mech}}\).
Then there are
\(\widehat r\in\{0,\ldots,d\}\), \(i\in[k]\), and, in every optimal
\((p_{\widehat r},d)\)-decomposition of
\(H_i^{\widehat r}\), a nonempty leaf \(\mathcal G\) such that, with

\[
 t:=\operatorname{DDim}_{p_{\widehat r},d}(H_i^{\widehat r})
   =\operatorname{LD}(\mathcal G),
\tag{18}
\]

one has

\[
 \bar H\in\mathcal L_i^{\widehat r},
 \qquad
 \mathcal G\subseteq H_i^{\widehat r}\subseteq\bar C,
 \qquad
 \bar H=\operatorname{SOA}_{\mathcal G}
 \quad\text{pointwise on }Q_C.
\tag{19}
\]

The leaf \(\mathcal G\) is
\(p_{\widehat r}2^{d-t}\)-irreducible, and

\[
 p_{\widehat r}2^{d-t}
 \ge p_0=n_0d
 \ge\max\{n_0,d+1\}.
\tag{20}
\]

Thus \(\mathcal G\) is both \(n_0\)-irreducible and
\((d+1)\)-irreducible. The exact source/current identities are

\[
 H_i^{\widehat r}=H_{i,\mathrm{src}}^{\widehat r+1},
 \qquad
 p_{\widehat r}
 =2^{\widehat r}n_0d
 =\frac{p_{\widehat r+1,\mathrm{src}}}{2}.
\tag{21}
\]

If \(\widehat r<d\), the next current scale satisfies
\(p_{\widehat r+1}=2p_{\widehat r}\). If \(\widehat r=d\), there is no
current transition beyond the selected endpoint. The conclusion records
actual selected status even if \(\bar H=\bar c_0\) as a function.


\end{proposition}

\begin{proof}
Proposition~\ref{prop:step-010-mechanism-good} gives the selected stage,
producer, actual-status output, and first relation in (19). Because the
output is a literal member of \(\mathcal L_i^{\widehat r}\), this list is
nonempty. By the piecewise definition in Proposition~\ref{prop:step-008-essential-lists}, the corresponding outer
restriction \(H_i^{\widehat r}\) is therefore nonempty and the selected
function is genuinely \((p_{\widehat r},d)\)-essential to it. No DDim or
SOA is evaluated on an empty class.

Apply Proposition~\ref{prop:step-008-leaf-scale} to this exact
list membership. It gives, in every optimal current decomposition, a
maximal leaf \(\mathcal G\), (18), the subset relation in (19), the
pointwise identity in (19), and the full scale (20). That proposition also
proves that irreducibility at the scale in (20) implies irreducibility at
each smaller positive integer scale, giving \(n_0\) and \(d+1\).

Finally, Lemma~\ref{lem:step-008-stage-map} gives (21). Its
factor-two statement gives \(p_{\widehat r+1}=2p_{\widehat r}\) only when
\(\widehat r<d\). The selected list throughout this argument remains the
current half-source-scale list at \(p_{\widehat r}\); it is not relabeled
as a literal source-stage list. At \(\widehat r=d\), (21) uses the source
endpoint \(d+1\) and does not invent a current stage \(d+1\).

Actual versus fallback status comes from the  transcript, not from
the function value. Hence equality \(\bar H=\bar c_0\) does not change any
of the preceding memberships or identities. \(\square\)


\end{proof}

\begin{lemma}[Every member of the selected leaf has small master error]\label{lem:step-012-leaf-error}
Under Assumption~\ref{assump:finite-littlestone},
Propositions~\ref{prop:step-006-good-event},
\ref{prop:step-010-mechanism-good}, and
\ref{prop:step-012-selected-leaf}, suppose \(d\ge1\) and
\(E_{\mathrm{core}}=E_{\mathrm{good}}\cap E_{\mathrm{mech}}\) holds. Then
every \(g\in\mathcal G\) satisfies

\[
 e_{\bar S}(g)\le2\gamma.
\tag{22}
\]

This is a statement about actual members \(g\in\bar C\); it makes no
application of \(E_{\mathrm{good}}\) to the potentially improper
\(\bar H=\operatorname{SOA}_{\mathcal G}\).


\end{lemma}

\begin{proof}
Fix \(g\in\mathcal G\). By (19) and the selected current restriction
defined in Lemma~\ref{lem:step-008-stage-map},

\[
 e_i(g)\le\rho^{\widehat r+1}\gamma\le\gamma,
\tag{23}
\]

because \(0<\rho<1\) and \(\widehat r+1\ge1\). Also
\(g\in\mathcal G\subseteq H_i^{\widehat r}\subseteq\bar C\), so the exact
event produced by Proposition~\ref{prop:step-006-good-event} applies to this
same actual function.

If \(e_{\bar S}(g)\le\gamma/3\), then immediately

\[
 e_{\bar S}(g)\le\gamma/3\le2\gamma.
\tag{24}
\]

This branch includes equality at \(\gamma/3\). If instead
\(e_{\bar S}(g)>\gamma/3\), the lower half of the high-error relative
clause in Proposition~\ref{prop:step-006-good-event}, together with (23),
gives

\[
 (1-a_d)e_{\bar S}(g)
 \le e_i(g)
 \le\rho^{\widehat r+1}\gamma
 \le\gamma.
\tag{25}
\]

Since \(d\ge1\),

\[
 a_d=\frac1{5d}\le\frac15,
 \qquad
 1-a_d\ge\frac45,
 \qquad
 \frac1{1-a_d}\le\frac54\le2.
\tag{26}
\]

Dividing (25) by the positive \(1-a_d\) and using (26) yields

\[
 e_{\bar S}(g)
 \le\frac{\gamma}{1-a_d}
 \le\frac54\gamma
 \le2\gamma.
\tag{27}
\]

The two exhaustive event branches prove (22). The argument is uniform in
the selected stage, including \(0\) and \(d\), and in every leaf member;
no union, independence, or new confidence factor is used. \(\square\)


\end{proof}

\begin{lemma}[Indexed irreducible-SOA empirical contradiction]\label{lem:step-012-empirical}
Under Assumption~\ref{assump:finite-littlestone},
Propositions~\ref{prop:step-006-good-event},
\ref{prop:step-010-mechanism-good}, and
\ref{prop:step-012-selected-leaf},
Lemma~\ref{lem:step-012-leaf-error}, and the empirical argument in
Lyu~\citep{lyu2025}'s proof of Theorem 3, suppose \(d\ge1\) and
\(E_{\mathrm{core}}=E_{\mathrm{good}}\cap E_{\mathrm{mech}}\) holds. Then

\[
 e_{\bar S}(\bar H)\le2\gamma=\frac{\alpha}{8}.
\tag{28}
\]

The empirical error in (28) uses all \(n_0\) indexed master records with
normalization \(1/n_0\), including repetitions.


\end{lemma}

\begin{proof}
Suppose, for contradiction, that

\[
 e_{\bar S}(\bar H)>2\gamma.
\tag{29}
\]

Restrict the nonempty leaf \(\mathcal G\) along the complete ordered
sequence of master inputs, labeled by the exact SOA function:

\[
 \mathcal R
 :=\mathcal G|_{(q_1,\bar H(q_1)),\ldots,
                       (q_{n_0},\bar H(q_{n_0}))}.
\tag{30}
\]

If \(g\in\mathcal R\), then
\(g(q_j)=\bar H(q_j)\) for every indexed coordinate \(j\). Consequently,
for every \(j\), including every repeated occurrence of a quotient point,

\[
 \mathbf1\{g(q_j)\ne y_j\}
 =\mathbf1\{\bar H(q_j)\ne y_j\}.
\tag{31}
\]

Summing (31) and dividing by the same \(n_0\) gives

\[
 e_{\bar S}(g)=e_{\bar S}(\bar H)>2\gamma,
\tag{32}
\]

contradicting Lemma~\ref{lem:step-012-leaf-error}. Hence

\[
 \mathcal R=\varnothing.
\tag{33}
\]

On the other hand, Proposition~\ref{prop:step-012-selected-leaf} gives
\(\bar H=\operatorname{SOA}_{\mathcal G}\) pointwise and says that
\(\mathcal G\) is \(n_0\)-irreducible. Applying the definition of
\(n_0\)-irreducibility to the exact point sequence
\(q_1,\ldots,q_{n_0}\) gives

\[
 \operatorname{LD}(\mathcal R)
 =\operatorname{LD}(\mathcal G)=t\ge0.
\tag{34}
\]

Thus \(\mathcal R\) must be nonempty: an empty binary class has no
depth-zero Littlestone witness and cannot have the nonnegative dimension
of the nonempty leaf \(\mathcal G\). This contradicts (33). Repeated
sample points cause no difficulty, because irreducibility quantifies over
sequences of points and permits repetitions; moreover the labels in (30)
are the single function values \(\bar H(q_j)\).

Therefore (29) is false and \(e_{\bar S}(\bar H)\le2\gamma\). Finally,

\[
 2\gamma=2\left(\frac{\alpha}{16}\right)=\frac{\alpha}{8},
\tag{35}
\]

which proves (28). The contradiction began with a strict \(>\), so the
exported conclusion correctly permits exact equality at
\(\alpha/8\). \(\square\)


\end{proof}

\begin{proposition}[Membership in the fixed improper irreducible-SOA
family]\label{prop:step-012-fixed-family}
Under Assumption~\ref{assump:finite-littlestone},
Lemma~\ref{lem:step-008-stage-map},
Propositions~\ref{prop:step-006-good-event},
\ref{prop:step-010-mechanism-good}, and
\ref{prop:step-012-selected-leaf}, and Lyu~\citep{lyu2025}
Lemma~4.2, suppose \(d\ge1\) and
\(E_{\mathrm{core}}=E_{\mathrm{good}}\cap E_{\mathrm{mech}}\) holds. Define
the static family

\[
 \widehat C_{d+1}
 :=\left\{\operatorname{SOA}_{\mathcal A}:
       \varnothing\ne\mathcal A\subseteq\bar C
       \text{ is }(d+1)\text{-irreducible}\right\}.
\tag{36}
\]

Then

\[
 \bar H\in\widehat C_{d+1},
 \qquad
 \operatorname{LD}(\widehat C_{d+1})\le d.
\tag{37}
\]

The first membership in (37) is literal equality of functions on \(Q_C\).
It does not assert \(\bar H\in\bar C\), and the second conclusion does not
assert any measurability, empirical, population, or privacy property.


\end{proposition}

\begin{proof}
Proposition~\ref{prop:step-012-selected-leaf} gives a nonempty
\(\mathcal G\subseteq H_i^{\widehat r}\subseteq\bar C\), proves that it
is \((d+1)\)-irreducible, and gives the pointwise identity
\(\bar H=\operatorname{SOA}_{\mathcal G}\). Thus the same function
\(\bar H\) is one of the elements in the sample-independent definition
(36). This proves the first part of (37) without selecting a trace
representative or a proper concept.

Lemma~\ref{lem:step-008-stage-map} proves
\(\operatorname{LD}(\bar C)\le d\). Every hypothesis of Lyu~\citep{lyu2025}
Lemma~4.2, applied to the static family in (36), is now discharged with
\(\mathcal H=\bar C\). Its conclusion is precisely the second part of
(37). \(\square\)


\end{proof}

\begin{proposition}[Endpoint, normalization, status, and improperness
boundary cases]\label{prop:step-012-boundaries}
Under Assumption~\ref{assump:finite-littlestone},
Lemma~\ref{lem:step-012-empirical}, and
Propositions~\ref{prop:step-006-good-event},
\ref{prop:step-010-mechanism-good},
\ref{prop:step-012-selected-leaf}, and
\ref{prop:step-012-fixed-family}, suppose \(d\ge1\) and
\(E_{\mathrm{core}}=E_{\mathrm{good}}\cap E_{\mathrm{mech}}\) holds. Then
the exact SOA, fixed-family, and
empirical-error conclusions remain valid in all of the following boundary
and identity regimes.

1. If \(\widehat r=0\), the selected current class is source stage \(1\),
   \(H_i^0=H_{i,\mathrm{src}}^1\), and the selected-list parameter is
   \(p_0=n_0d=p_{1,\mathrm{src}}/2\). The proof uses this exact current
   list, not a same-scale or literal source-stage substitute.
2. If \(\widehat r=d\), the selected current class is the explicit source
   endpoint \(H_i^d=H_{i,\mathrm{src}}^{d+1}\), with
   \(p_d=p_{d+1,\mathrm{src}}/2\). No current \(H_i^{d+1}\), no next-stage
   essential list, and no transition beyond \(d\) is used.
3. If \(d=1\), the only possible selected stages are \(0\) and \(1\),
   \(p_0=n_0\), \(p_1=2n_0\), and
   \(p_{\widehat r}2^{1-t}\ge n_0\ge2=d+1\). Thus both the full-sample and
   fixed-family irreducibility uses remain legal, including \(t=0\).
4. The proof is unchanged at \(v=1\) and at \(v=d\). The parameter \(v\)
   does not enter (18)-(37); it only affects already fixed upstream sample
   quantities. No additional \(v\)- or \(d\)-power is introduced here.
5. A singleton current list, one function repeated in several lists, or
   duplicate encodings of the same list item still supplies one actual
   function \(\bar H\) and one producer \(i\). Essentiality and SOA
   equality concern that function, not its multiplicity. If the selected
   function equals \(\bar c_0\),  actual transcript status keeps it
   actual; equality of values never turns it into fallback.
6. An empty selected tuple, an empty selected list, \(\perp\), and every
   totalized fallback path are absent on the joint source/mechanism-good
   path. Specifically, Proposition~\ref{prop:step-010-mechanism-good}
   supplies literal selected-list membership, which implies that the chosen
   list and current restriction are nonempty.
   The source event \(E_{\mathrm{good}}\) supplies the error transfer, but
   is not by itself misrepresented as a mechanism-status event.
7. If the selected maximal leaf has \(t=0\), it is still nonempty and the
   same scale (20), indexed restriction (30), and fixed-family argument
   apply. DDim zero is never assigned to an empty outer restriction.
8. Repeated quotient points or repeated labeled records remain separate
   indices in both sums in (2). Equality of functions on all \(Q_C\) gives
   coordinatewise equality (31), so no deduplication changes the
   \(1/n_0\) normalization or shortens the \(n_0\)-point irreducibility
   sequence.
9. The threshold identity is exactly
   \(2\gamma=\alpha/8\). The proof excludes only
   \(e_{\bar S}(\bar H)>\alpha/8\) and therefore makes no unjustified
   strict-error claim at equality.
10. The output is guaranteed to lie in \(H_C\) and in the fixed family
    \(\widehat C_{d+1}\). It may lie outside \(\bar C\); no properness of
    the quotient output or decoded raw hypothesis is asserted.


\end{proposition}

\begin{proof}
Items 1 and 2 are the endpoint cases of (21), with the
transition range recorded in
Proposition~\ref{prop:step-012-selected-leaf}. Item 3 substitutes \(d=1\)
into (18)-(21); the scale inequality is the  exact bound (20), not
an asymptotic estimate. Item 4 follows because no derivation after the
fixed upstream dictionary contains \(v\).

For Item 5, mechanism-good status uses set membership rather than positional
multiplicity, while Definition~4.3 and (19) identify actual functions
pointwise. Item 6 follows from the selected-list status in
Proposition~\ref{prop:step-010-mechanism-good}; the separate role of
Proposition~\ref{prop:step-006-good-event} remains visible. Item 7 follows
from the nonempty leaf
in Proposition~\ref{prop:step-012-selected-leaf} and does not require a
positive \(t\).

Items 8 and 9 are the explicit indexed calculation (31)-(35). Item 10 is
exactly Proposition~\ref{prop:step-012-fixed-family} together with the
actual-output boundary in (19). None of these specializations adds a
privacy, probability, population-risk, raw-risk, properness, or sample-size
conclusion. 
On the joint source-good plus mechanism-good path, Proposition~\ref{prop:step-010-mechanism-good} first supplies a literal
actual current-list member and its selected stage. Proposition
~\ref{prop:step-012-selected-leaf} combines that status with 
Definition 4.3 and leaf-scale interfaces to obtain one nonempty
\(\mathcal G\subseteq H_i^{\widehat r}\subseteq\bar C\) satisfying the
exact pointwise identity

\[
 \bar H=\operatorname{SOA}_{\mathcal G}
\]

and the exact scale

 \[
 \mathcal G\text{ is }p_{\widehat r}2^{d-t}\text{-irreducible},
 \qquad
 p_{\widehat r}2^{d-t}\ge n_0d
 \ge\max\{n_0,d+1\}.
 \]

The same proposition preserves the selected current/source stage identity
and the half-scale convention, with a factor-two next-current pairing only
when \(\widehat r<d\).

Lemma~\ref{lem:step-012-leaf-error} next uses the exact two-clause
\(E_{\mathrm{good}}\) interface on each actual
\(g\in\mathcal G\). Membership in the selected restriction gives
\(e_i(g)\le\rho^{\widehat r+1}\gamma\le\gamma\); the low event branch is
already below \(2\gamma\), while the high branch gives

\[
 e_{\bar S}(g)
 \le\frac{\gamma}{1-a_d}
 \le\frac54\gamma
 \le2\gamma.
\]

Lemma~\ref{lem:step-012-empirical} then applies Lyu~\citep{lyu2025} Theorem 3
contradiction to the same actual function and the same indexed master
sample. If \(e_{\bar S}(\bar H)>2\gamma\), every member of
\(\mathcal G\) is excluded from the restriction along the complete
\(\bar H=\operatorname{SOA}_{\mathcal G}\)-labeled \(n_0\)-tuple. The
restriction is empty, contradicting \(n_0\)-irreducibility. Hence

\[
 e_{\bar S}(\bar H)\le2\gamma=\alpha/8.
\]

Finally, Proposition~\ref{prop:step-012-fixed-family} uses the same leaf,
not a surrogate class or trace, to prove

\[
 \bar H\in\widehat C_{d+1},
 \qquad
 \operatorname{LD}(\widehat C_{d+1})\le d.
\]

The endpoint and object-integrity calculations above verify that the two exported
conclusions retain literal identity, normalization, actual status, endpoint
scope, equality allowance, and improperness in every boundary
regime. Together, these conclusions give the exact SOA utility guarantee
and the empirical-error/actual-SOA interface used in the PAC conversion.



\end{proof}

\subsection{Marked Holdout and PAC Conversion}\label{app:step_013}


\begin{proposition}[Core-event and marked-holdout interfaces]
\label{prop:step-013-interfaces}
Under Assumptions~\ref{assump:finite-littlestone},
\ref{assump:realizable-iid}, and
\ref{assump:approximate-dp-regime},
Lemmas~\ref{lem:step-004-occurrence},
\ref{lem:step-007-fixed-point}, and
\ref{lem:step-012-empirical}, and
Propositions~\ref{prop:step-004-lift},
\ref{prop:step-004-projection},
\ref{prop:step-007-boundaries},
\ref{prop:step-007-tower},
\ref{prop:step-008-list-envelope}, and
\ref{prop:step-010-mechanism-good}, suppose \(d\ge1\). Then the
analysis-only mark is assigned pathwise to the nonempty producer set of
an actual terminal output, while nonactual paths receive mark \(0\). Its
marked kernel has the exact released marginal
\[
\sum_{j=0}^k\widetilde K_C(\bar s,B\times\{j\})
=K_C(\bar s,B),
\qquad B\in\mathcal H_C.
\tag{1}
\]
This identity makes no privacy assertion for the unreleased pair
\((\bar H,J)\).

Write

\[
 b:=\log(4k/\beta),\qquad a:=v+b,
 \qquad Q:=e+\frac{ekd^2a}{\alpha v},
 \qquad n_0:=km,
 \tag{2}
\]

and

\[
 C_{\rm fp}:=2+\log(1+C_{\rm blk}).
 \tag{3}
\]

Then the fixed-point and exact-ceiling bounds are

\[
 \log\frac{en_0}{v}\leq C_{\rm fp}\log Q,
 \qquad
 m\geq C_{\rm blk}\frac{d^2}{\alpha}a\log Q,
 \tag{4}
\]

where the second inequality is the lower side of the exact ceiling, and the
universal block calibration is

\[
 \frac{C_{\rm blk}}{3600}
 \geq4+\log(1+C_{\rm blk})=C_{\rm fp}+2.
 \tag{5}
\]

\[
 \Pr(E_{\rm good}^c)\leq\beta_{\rm tr}=\beta/4.
 \tag{6}
\]
The complementary \(d=0\) branch is the exact \(N=0\) singleton law of
Proposition~\ref{prop:step-007-boundaries}.

For the actual current essential lists, including the totalized empty list
for an empty outer restriction, define

\[
 \mathcal G_i:=\bigcup_{r=0}^d\mathcal L_i^r,
 \qquad
 |\mathcal G_i|\leq(d+1)L,
 \tag{7}
\]

with

\[
 L=(2^dn_0d)^d2^{d^2},
 \qquad
 \log L=d\log n_0+d\log d+2d^2\log2.
 \tag{8}
\]

The elements in (7) are actual functions in \(H_C\), not empirical-trace
representatives. For every fixed realizable sample and partition satisfying
\(E_{\rm good}\), the measurable event \(E_{\rm mech}\) satisfies

\[
 \Pr(E_{\rm mech}^c\mid\bar S,\mathcal P,E_{\rm good})
 \leq\beta_{\rm AT}+\beta_{\rm SS}=\beta/2.
 \tag{9}
\]

On \(E_{\rm mech}\), the terminal output is an actual selected-stage list
item, not fallback; its  marked occurrence is positive and obeys

\[
 J=i\quad\Longrightarrow\quad
 \bar H\in\mathcal G_i(\bar S_i).
 \tag{10}
\]

Equation (10) is pathwise; it asserts no independence conditional on
\(J=i\). On the previously produced event
\(E_{\rm core}=E_{\rm good}\cap E_{\rm mech}\), the same literal actual
output satisfies

\[
 \widehat{\operatorname{err}}_{\bar S}(\bar H)
 :=\frac1{n_0}\sum_{(q,y)\in\bar S}
       \mathbf1\{\bar H(q)\ne y\}
 \leq\frac{\alpha}{8}.
 \tag{12}
\]

This is a deterministic full-master empirical certificate on the core path.

For the downstream holdout argument, fix arbitrary \(D\) and \(c\in C\),
write \(\bar D=\kappa_\#D\) and \(\bar c\) for the quotient target, and let
\[
 \bar S=(Z_1,\ldots,Z_{n_0})
 \sim P_{\bar D,\bar c}^{n_0},
 \qquad Z_j=(Q_j,\bar c(Q_j)).
 \tag{13}
\]
After drawing this iid quotient master sample, draw the setting's
data-independent uniform indexed partition \(\mathcal P\), the learner
transcript, and the analysis-only occurrence mark. Denote the resulting full
marked law by \(\widetilde{\mathbb P}_{\bar D,\bar c}\). For a realized
partition \(\pi\), let \(\bar S_i^\pi\) be the ordered \(m\)-record producer
block and \(\bar S_{-i}^\pi\) its ordered
\(M=(k-1)m\)-record complement, and define
\[
 \widehat{\operatorname{err}}_{-i}^\pi(h)
 :=\frac1M\sum_{(q,y)\in\bar S_{-i}^\pi}
       \mathbf1\{h(q)\ne y\},
 \qquad
 r(h):=\operatorname{err}_{\bar D}(h,\bar c).
 \tag{14}
\]
The superscript \(\pi\) is suppressed when the realized partition is clear.
Finally, let
\[
 F_\alpha:=\{r(\bar H)>\alpha\}.
 \tag{15}
\]
These definitions name the common probability space and assert no
independence beyond Assumption~\ref{assump:realizable-iid}.
\end{proposition}

\begin{proof}
Lemma~\ref{lem:step-004-occurrence} and
Propositions~\ref{prop:step-004-lift} and
\ref{prop:step-004-projection} give the mark and (1).
Lemma~\ref{lem:step-007-fixed-point} and
Proposition~\ref{prop:step-007-boundaries} give (2)--(5), while
Proposition~\ref{prop:step-007-tower} gives (6).
Proposition~\ref{prop:step-008-list-envelope} gives (7)--(8), and
Proposition~\ref{prop:step-010-mechanism-good} gives (9)--(10) and defines
\(E_{\rm core}=E_{\rm good}\cap E_{\rm mech}\), while
Lemma~\ref{lem:step-012-empirical} gives (12).
Assumption~\ref{assump:realizable-iid} gives the iid quotient law (13); the
data-independent partition and Proposition~\ref{prop:step-004-lift} then
define the full marked experiment, and (14)--(15) are the stated complement
loss, population risk, and failure-event definitions. They add no
independence assertion. Together, the cited results prove all displayed
interfaces without requiring an additional event or tail bound. \(\square\)
\end{proof}

\begin{lemma}[Measurable block-local all-stage candidate family]\label{lem:step-013-block-family}
Under Lemma~\ref{lem:step-004-occurrence} and
Propositions~\ref{prop:step-008-list-envelope} and
\ref{prop:step-013-interfaces}, suppose \(d\geq1\). For every
producer \(i\in[k]\) and local labeled block \(s_i\in Z_Q^m\), let
\(\mathcal L_i^r(s_i)\) be the actual ordered current essential list at stage
\(r\), with its empty-list totalization, and set

\[
 \mathcal G_i(s_i)
 :=\bigcup_{r=0}^d\mathcal L_i^r(s_i).
 \tag{16}
\]

Then:

1. \(\mathcal G_i(s_i)\) depends only on the \(m\) labeled records in
   \(s_i\) and the public parameter tuple, not on \(s_{-i}\), the selected
   stage, mechanism coins, output, or mark;
2. its concatenated finite-list encoding and membership graph
   \[
     \{(s_i,h)\in Z_Q^m\times H_C:h\in\mathcal G_i(s_i)\}
     \tag{17}
   \]
   are measurable;
3. for every local block state,
   \[
     |\mathcal G_i(s_i)|\leq(d+1)L;
     \tag{18}
   \]
4. for every partition \(\pi\), the event
   \[
    A_i^\pi:=
    \bigcup_{h\in\mathcal G_i(\bar S_i^\pi):\ r(h)>\alpha}
      \{\widehat{\operatorname{err}}_{-i}^\pi(h)\leq\alpha/4\}
    \tag{19}
   \]
   is measurable. This is a finite random union, not an uncountable
   supremum event.


\end{lemma}

\begin{proof}
The restriction list at stage \(r\) is computed from the
block empirical errors of \(s_i\) alone. Its decomposition/list convention
was fixed before sampling as a deterministic lookup on the local quotient
state. Consequently each actual list \(\mathcal L_i^r(s_i)\), and therefore
their all-stage union (16), is a function only of \(s_i\) and public
parameters. Taking the union over all stages rather than only the selected
stage is what removes any dependence of the candidate family on the
adaptive mechanism transcript.

The local input space \(Z_Q^m\) is countable discrete. Thus the pre-fixed
finite-list coordinate \(s_i\mapsto\mathcal L_i^r(s_i)\) is measurable for
every \(r\). Concatenating the \(d+1\) ordered list encodings is measurable,
and dynamic membership in the concatenation is the finite-list membership
relation already proved measurable by Lemma~\ref{lem:step-004-occurrence}.
This proves (17). Removing repeated entries across stages is unnecessary
for measurability; it only converts the encoding to the mathematical set in
(16).

Proposition~\ref{prop:step-008-list-envelope} gives
\(|\mathcal L_i^r(s_i)|\leq L\) at every stage, including an empty
totalized list. Subadditivity over exactly \(d+1\) stages proves (18).

It remains to check the section in (19). Because \(Q_C\) is finite or
countable on the  quotient interface, enumerate it as
\((q_\ell)\). For \(h\in H_C\),

\[
 r(h)=\sum_{\ell}\bar D(q_\ell)
       \mathbf1\{h(q_\ell)\ne\bar c(q_\ell)\}
 \tag{20}
\]

is the monotone limit of measurable finite coordinate sums. Hence
\(h\mapsto r(h)\) is \(\mathcal H_C\)-measurable. The complement empirical
error in (14) is a finite measurable sum in \((h,\bar S_{-i})\). Combining
these two facts with the measurable finite-list graph (17), (19) is a
finite measurable disjunction over the entries of the concatenated list.
If all lists are empty, that disjunction is empty. If a function occurs at
several stages, retaining several encoding positions can only repeat the
same measurable event and never enlarges the cardinality bound beyond
(18). \(\square\)


\end{proof}

\begin{proposition}[Full-path marked inclusion and Holdout]\label{prop:step-013-pathwise}
Under Propositions~\ref{prop:step-004-lift},
\ref{prop:step-010-mechanism-good}, and
\ref{prop:step-013-interfaces}, Lemma~\ref{lem:step-012-empirical}, and
Lemma~\ref{lem:step-013-block-family}, suppose \(d\geq1\). For every
complete realization of the sample, partition, learner transcript,
mechanism coins, and analysis mark, and every \(i\in[k]\),

\[
 E_{\rm core}\cap\{J=i\}
 \quad\Longrightarrow\quad
 \bar H\in\mathcal G_i(\bar S_i)
 \quad\text{and}\quad
 \widehat{\operatorname{err}}_{-i}(\bar H)
 \leq\frac{k\alpha}{8(k-1)}\leq\frac\alpha4.
 \tag{21}
\]

Consequently the following inclusion holds pathwise:

\[
 \boxed{
 E_{\rm core}\cap\{J=i\}\cap F_\alpha
 \subseteq
 \bigcup_{h\in\mathcal G_i(\bar S_i):\ r(h)>\alpha}
 \{\widehat{\operatorname{err}}_{-i}(h)\leq\alpha/4\}.
 }
 \tag{22}
\]

No probability conditioning, independence statement, or distributional
claim about \(J\) is used in (21)--(22).


\end{proposition}

\begin{proof}
On \(E_{\rm core}\), Proposition~\ref{prop:step-010-mechanism-good} gives actual rather than
fallback status. The occurrence mark is therefore positive, and
if its realized value is \(J=i\), its support rule (10) gives the literal
membership

\[
 \bar H\in\mathcal G_i(\bar S_i).
 \tag{23}
\]

Lemma~\ref{lem:step-012-empirical} applies to this same function
and the same full indexed master sample, giving

\[
 \sum_{(q,y)\in\bar S}
   \mathbf1\{\bar H(q)\ne y\}
 \leq n_0\frac\alpha8.
 \tag{24}
\]

Every summand is nonnegative, so deleting the \(m\) producer-block terms
can only reduce the numerator. Since \(n_0=km\) and
\(n_0-m=(k-1)m\),

\[
\begin{aligned}
 \widehat{\operatorname{err}}_{-i}(\bar H)
 &\leq\frac{n_0}{n_0-m}\frac\alpha8\\
 &=\frac{k\alpha}{8(k-1)}
 \leq\frac\alpha4,
\end{aligned}
 \tag{25}
\]

where the last inequality is exactly \(k/(k-1)\leq2\) for \(k\geq2\).
This is (Holdout). It compares the same literal \(\bar H\), the same target
labels \(y=\bar c(q)\), and the same zero-one loss in the full and
complement samples.

If \(F_\alpha\) also occurs, then the function \(h=\bar H\) supplied by
(23) satisfies \(r(h)>\alpha\), and (25) supplies its right-hand event in
(22). Thus that literal function witnesses the finite union, proving the
inclusion.

The proof has not conditioned on \(J=i\); \(\{J=i\}\) is simply one
measurable event on the already realized full path. In particular, it does
not claim that \(\bar S_{-i}\) is independent of \(J\), the selected stage,
the transcript, or \(\bar H\). The independent fixed-candidate integration
is applied only after (22). \(\square\)


\end{proof}

\begin{lemma}[Direct one-sided iid Bernoulli lower tail]\label{lem:step-013-lower-tail}
Under Assumption~\ref{assump:realizable-iid}, let \(h\in H_C\) be fixed,
put \(p=r(h)>\alpha\), and let \(M=(k-1)m\geq1\) fresh quotient examples
be iid from \(P_{\bar D,\bar c}\). If \(Y\) is the number of errors of
\(h\) on those examples, then

\[
\begin{aligned}
 \Pr\left[\frac YM\leq\frac\alpha4\right]
 &\leq \Pr[Y\leq Mp/4]\\
 &\leq e^{-9Mp/32}\\
 &\leq e^{-9\alpha(k-1)m/32}.
\end{aligned}
 \tag{26}
\]


\end{lemma}

\begin{proof}
The error indicators are iid Bernoulli\((p)\), so
\(Y\sim\operatorname{Bin}(M,p)\). Because \(p>\alpha\), the first event in
(26) is contained in \(\{Y\leq Mp/4\}\).

For \(\lambda=\log4>0\), exponential Markov applied to
\(e^{-\lambda Y}\) gives

\[
\begin{aligned}
 \Pr[Y\leq Mp/4]
 &\leq e^{\lambda Mp/4}\mathbb E[e^{-\lambda Y}]\\
 &=e^{\lambda Mp/4}(1-p+pe^{-\lambda})^M\\
 &=e^{\lambda Mp/4}(1-3p/4)^M\\
 &\leq\exp\left\{Mp\left(\frac{\lambda}{4}-\frac34\right)\right\},
\end{aligned}
 \tag{27}
\]

where \(1-x\leq e^{-x}\) was used in the last line. The elementary series
bound

\[
 e^{3/2}>1+\frac32+\frac{(3/2)^2}{2}
              +\frac{(3/2)^3}{6}>4
\]

implies \(\log4<3/2\). Hence the exponent in (27) is at most
\(-3Mp/8\), which is no larger than \(-9Mp/32\). This proves the second
line of (26). The last line follows from \(p>\alpha\) and
\(M=(k-1)m\). Every constant is explicit, and no cited Chernoff statement
or asymptotic estimate is used. \(\square\)


\end{proof}

\begin{proposition}[Finite producer integration after the pathwise inclusion]\label{prop:step-013-finite-integration}
Under Assumption~\ref{assump:realizable-iid},
Lemmas~\ref{lem:step-013-block-family} and
\ref{lem:step-013-lower-tail}, and
Propositions~\ref{prop:step-013-interfaces} and
\ref{prop:step-013-pathwise}, suppose \(d\geq1\). Then

\[
 \widetilde{\mathbb P}_{\bar D,\bar c}
 \left(E_{\rm core}\cap F_\alpha\right)
 \leq k(d+1)L
       \exp\left(-\frac{9\alpha(k-1)m}{32}\right).
 \tag{28}
\]


\end{proposition}

\begin{proof}
Proposition~\ref{prop:step-013-pathwise} has already proved
(22) on the full sample/partition/mechanism/mark path. We now, and only
now, condition on the data-independent partition and producer block.

Fix a deterministic indexed partition \(\pi\) and \(i\in[k]\). Because
the master coordinates in (13) are iid and \(\pi\) is fixed independently
of them, the random vectors \(\bar S_i^\pi\) and
\(\bar S_{-i}^\pi\) are independent; the complement consists of exactly
\(M=(k-1)m\) iid quotient examples. This is conditioning on the partition,
not conditioning on the complete master sample. No independence among
blocks conditional on the complete sample is asserted.

Next condition on \(\bar S_i^\pi=s_i\). By
Lemma~\ref{lem:step-013-block-family}, the set
\(\mathcal G_i(s_i)\) is now a fixed finite set of at most \((d+1)L\)
literal functions and does not depend on the complement. For every member
\(h\) of this set with \(r(h)>\alpha\),
Lemma~\ref{lem:step-013-lower-tail} applies to the independent complement
and gives

\[
 \Pr\left[
   \widehat{\operatorname{err}}_{-i}^\pi(h)\leq\alpha/4
   \mid\mathcal P=\pi,\bar S_i^\pi=s_i
 \right]
 \leq e^{-9\alpha(k-1)m/32}.
 \tag{29}
\]

Here the product regular conditional law is chosen to be the iid complement
law also on producer states of probability zero; this is legitimate on the
countable discrete block space and makes (29) pointwise in \(s_i\).
The function in (29) is fixed by \(s_i\). It is not obtained by
conditioning on \(J=i\), the selected stage, \(\bar H\), or any mechanism
event. Applying the finite union bound to (19) therefore yields

\[
 \Pr[A_i^\pi\mid\mathcal P=\pi,\bar S_i^\pi=s_i]
 \leq(d+1)L e^{-9\alpha(k-1)m/32}.
 \tag{30}
\]

Empty lists give an empty union, and repeated candidates only reduce the
number of distinct terms. The bound in (30) is uniform in \(s_i\) and
\(\pi\). Measurability from
Lemma~\ref{lem:step-013-block-family} permits the tower integral, first
over \(\bar S_i^\pi\) and then over the finite partition law. Writing
\(A_i:=A_i^{\mathcal P}\), this gives

\[
 \widetilde{\mathbb P}_{\bar D,\bar c}(A_i)
 \leq(d+1)L e^{-9\alpha(k-1)m/32}.
 \tag{31}
\]

The event \(A_i\) depends only on the sample and partition, so integrating
the learner and mark kernels contributes unit mass. On \(E_{\rm core}\),
 actual status excludes mark \(0\). The positive mark events are
disjoint, and (22) gives

\[
\begin{aligned}
 \widetilde{\mathbb P}_{\bar D,\bar c}
   (E_{\rm core}\cap F_\alpha)
 &=\sum_{i=1}^k
   \widetilde{\mathbb P}_{\bar D,\bar c}
   (E_{\rm core}\cap\{J=i\}\cap F_\alpha)\\
 &\leq\sum_{i=1}^k
   \widetilde{\mathbb P}_{\bar D,\bar c}(A_i)\\
 &\leq k(d+1)L e^{-9\alpha(k-1)m/32}.
\end{aligned}
 \tag{32}
\]

This proves (28). The mark is used only to obtain the finite pathwise
decomposition in the first line of (32); it is never used as a source of
conditional independence. \(\square\)


\end{proof}

\begin{lemma}[Exact finite-multiplicity domination by the  block
calibration]\label{lem:step-013-beta-gen}
Under Lemma~\ref{lem:step-007-fixed-point} and
Propositions~\ref{prop:step-007-boundaries},
\ref{prop:step-008-list-envelope}, and
\ref{prop:step-013-interfaces}, suppose \(d\geq1\). Then

\[
 k(d+1)L
 \exp\left(-\frac{9\alpha(k-1)m}{32}\right)
 \leq\beta_{\rm gen}=\frac\beta4.
 \tag{33}
\]


\end{lemma}

\begin{proof}
Retain \(a,b,Q,C_{\rm fp}\) from (2)--(3). The  positive
branch gives \(1\leq v\leq d\), \(k\geq2\), and \(a=v+b\geq1\). Since
\(a\geq v\) and \(0<\alpha<1/4\),

\[
 Q=e+\frac{ekd^2a}{\alpha v}>e+8ed^2>d,
 \tag{34}
\]

so \(\log Q>1\), \(\log d\leq\log Q\), and
\(\log2\leq\log Q\). From the  fixed point (4),

\[
\begin{aligned}
 \log n_0
 &=\log(en_0/v)+\log v-1\\
 &\leq C_{\rm fp}\log Q+\log d\\
 &\leq(C_{\rm fp}+1)\log Q.
\end{aligned}
 \tag{35}
\]

Substituting (35) into the exact  list formula (8) gives

\[
\begin{aligned}
 \log L
 &=d\log n_0+d\log d+2d^2\log2\\
 &\leq(C_{\rm fp}+1)d\log Q+d\log Q+2d^2\log Q\\
 &\leq(C_{\rm fp}+4)d^2\log Q.
\end{aligned}
 \tag{36}
\]

The exact confidence multiplicity is

\[
\begin{aligned}
 \log\frac{k(d+1)L}{\beta_{\rm gen}}
 &=\log\frac{4k}{\beta}+\log(d+1)+\log L\\
 &=b+\log(d+1)+\log L.
\end{aligned}
 \tag{37}
\]

Here \(b\leq a\), \(d+1\leq2d\), \(d^2a\geq1\), and \(\log Q>1\).
Consequently

\[
 b\leq d^2a\log Q,
 \qquad
 \log(d+1)\leq2\log Q\leq2d^2a\log Q,
 \tag{38}
\]

and (36)--(38) yield the fully displayed multiplicity bound

\[
 \log\frac{k(d+1)L}{\beta_{\rm gen}}
 \leq(C_{\rm fp}+7)d^2a\log Q.
 \tag{39}
\]

The same confidence-budget result block calibration pays this term. Indeed, (5)
implies

\[
 \frac{9C_{\rm blk}}{32}
 =\frac{2025}{2}\frac{C_{\rm blk}}{3600}
 \geq\frac{2025}{2}(C_{\rm fp}+2)
 \geq C_{\rm fp}+7,
 \tag{40}
\]

where the last difference is positive because \(C_{\rm fp}>0\). Using the
lower ceiling inequality from Proposition~\ref{prop:step-007-boundaries}, displayed in (4), and
\(k-1\geq1\),

\[
\begin{aligned}
 \frac{9\alpha(k-1)m}{32}
 &\geq\frac{9C_{\rm blk}}{32}(k-1)d^2a\log Q\\
 &\geq(C_{\rm fp}+7)d^2a\log Q\\
 &\geq\log\frac{k(d+1)L}{\beta_{\rm gen}}.
\end{aligned}
 \tag{41}
\]

Exponentiating (41) proves (33). Thus the exact  trace fixed point
and the already fixed universal block constant pay the producer, stage,
list, and confidence multiplicities. No new parameter-dependent constant is
chosen and no term is absorbed by prose. \(\square\)


\end{proof}

\begin{proposition}[Unconditional released quotient PAC ledger and boundary
integrity]\label{prop:step-013-pac}
Under Assumption~\ref{assump:realizable-iid}, 
Propositions~\ref{prop:step-004-projection},
\ref{prop:step-007-tower},
\ref{prop:step-007-boundaries},
\ref{prop:step-010-mechanism-good},
\ref{prop:step-013-interfaces}, Lemma~\ref{lem:step-012-empirical},
Proposition~\ref{prop:step-013-finite-integration}, and
Lemma~\ref{lem:step-013-beta-gen}, the released quotient law obeys, for
every allowed \(D\) and \(c\in C\),

\[
 \Pr_{\bar S\sim P_{\bar D,\bar c}^{n_0},\,
          \bar H\sim K_C(\bar S,\cdot)}
 \left[\operatorname{err}_{\bar D}(\bar H,\bar c)>\alpha\right]
 \leq\beta
 \tag{42}
\]

on the positive branch. Equivalently, the released quotient output has
population error at most \(\alpha\) with probability at least \(1-\beta\).
The  \(d=0\) branch satisfies the same conclusion with zero failure
probability and no data.


\end{proposition}

\begin{proof}
First suppose \(d\geq1\). Since \(E_{\rm good}\) is determined
by the sample and partition, integrating the uniform conditional bound (9)
over its successful sections gives

\[
\begin{aligned}
 \widetilde{\mathbb P}_{\bar D,\bar c}
   (E_{\rm good}\cap E_{\rm mech}^c)
 &=\mathbb E\left[
   \mathbf1_{E_{\rm good}}
   \Pr(E_{\rm mech}^c\mid\bar S,\mathcal P,E_{\rm good})
   \right]\\
 &\leq\beta_{\rm AT}+\beta_{\rm SS}.
\end{aligned}
 \tag{43}
\]

There is no additional empirical-certificate failure on
\(E_{\rm core}=E_{\rm good}\cap E_{\rm mech}\): Lemma~\ref{lem:step-012-empirical} is deterministic on that path.
Proposition~\ref{prop:step-007-tower}, (43),
Proposition~\ref{prop:step-013-finite-integration}, and
Lemma~\ref{lem:step-013-beta-gen} therefore give

\[
\begin{aligned}
 \widetilde{\mathbb P}_{\bar D,\bar c}(F_\alpha)
 &\leq
 \widetilde{\mathbb P}_{\bar D,\bar c}(E_{\rm good}^c)
 +\widetilde{\mathbb P}_{\bar D,\bar c}
    (E_{\rm good}\cap E_{\rm mech}^c)
 +\widetilde{\mathbb P}_{\bar D,\bar c}
    (E_{\rm core}\cap F_\alpha)\\
 &\leq\beta_{\rm tr}+\beta_{\rm AT}
          +\beta_{\rm SS}+\beta_{\rm gen}\\
 &=\frac\beta4+\frac\beta4+\frac\beta4+\frac\beta4=\beta.
\end{aligned}
 \tag{44}
\]

Let

\[
 B_\alpha:=\{h\in H_C:r(h)>\alpha\}.
\]

It is measurable by (20), and \(F_\alpha\) is its pullback by the terminal
output coordinate. The exact projection (1), integrated over the iid
sample, says

\[
 \widetilde{\mathbb P}_{\bar D,\bar c}
  [\bar H\in B_\alpha]
 =\Pr_{\bar S,\,\bar H\sim K_C(\bar S,\cdot)}
  [\bar H\in B_\alpha].
 \tag{45}
\]

Combining (44)--(45) proves (42). This projects out the mark; it neither
releases \(J\) nor assigns it a privacy property.

The remaining boundary cases preserve the same argument:

1. Proposition~\ref{prop:step-007-boundaries} gives the setting's
   \(d=0\) branch: \(\bar C\) has the unique released member, \(N=0\), and
   its quotient population error against the unique target is zero. The
   teacher, blocks, \(\mathcal G_i\), positive marks, and every denominator
   containing \(k-1\) are bypassed.
2. At \(k=2\), the complement size is \(M=m\), and (25) is exactly
   \(2(\alpha/8)=\alpha/4\). No block-independence claim conditional on the
   complete sample is introduced.
3. At \(d=1\), \(\mathcal G_i\) is the union of the two actual endpoint
   lists \(r=0,1\), and (18) is \(|\mathcal G_i|\leq2L\). Selection at
   stage \(0\) or stage \(d\) is included because (16) uses all stages.
4. At \(v=1\) and \(v=d\), the  fixed point and inequalities
   (34)--(41) remain valid; the argument introduces no extra structural power
   or rate claim.
5. Empty stage lists give empty finite sections. Repeated functions within
   the all-stage encoding or across producer blocks only repeat union terms
   and cannot increase the bound. Repeated sample values remain independent
   indexed iid draws before realization and retain their literal
   multiplicities in every empirical sum.
6. A fallback path lies in the complement of \(E_{\rm mech}\) whenever
   \(E_{\rm good}\) holds and is paid in (43); outside
   \(E_{\rm good}\) it is paid by \(\beta_{\rm tr}\). Its mark is \(0\),
   and it is never inserted into a positive producer union. An actual
   selected function equal in value to \(\bar c_0\) remains actual by
   status and is treated normally.
7. The space \(H_C\) may be infinite or uncountable as a set. The proof
   uses only the finite families in (16), measurable coordinate risks, and
   finite unions; it never forms an uncountable supremum over \(H_C\).

Thus all fallback/core, multiplicity, empty-list, selected-endpoint,
smallest-parameter, repeated-record, and infinite-output-space regimes obey
the same unconditional quotient PAC conclusion. 
The restriction-list propositions first supply the actual block-local lists.
Lemma
~\ref{lem:step-013-block-family} takes their union over all \(d+1\) stages,
not over a selected stage, and proves that this family depends only on the
producer block, is measurably encoded, and has size at most \((d+1)L\).

Proposition~\ref{prop:step-013-pathwise} then performs the high-risk step in
the required order. On every complete path,  actual marking gives
\(\bar H\in\mathcal G_i(\bar S_i)\) when \(J=i\), and Lemma
\ref{lem:step-012-empirical} gives the full-master error of that same
function. Deleting the producer
block yields

\[
 \widehat{\operatorname{err}}_{-i}(\bar H)
 \leq\frac{k}{k-1}\frac\alpha8\leq\frac\alpha4.
\]

Together with population failure, this proves the full-path inclusion (22)
before any conditional independence is invoked.

Lemma~\ref{lem:step-013-lower-tail} proves directly that a fixed bad
candidate has complement lower-tail probability at most
\(e^{-9\alpha(k-1)m/32}\). Only after (22) is available does
Proposition~\ref{prop:step-013-finite-integration} fix the independent
partition, condition on \(\bar S_i\), apply that fixed-candidate bound to
the now deterministic finite family, and integrate. It obtains exactly

\[
 \widetilde{\mathbb P}(E_{\rm core}\cap F_\alpha)
 \leq k(d+1)L e^{-9\alpha(k-1)m/32}.
\]

Lemma~\ref{lem:step-013-beta-gen} uses the confidence-budget fixed point and
lower ceiling inequality, the exact restriction-list formula for \(\log L\), and
the same universal block calibration to prove that this term is at most
\(\beta_{\rm gen}\), with every logarithmic domination displayed in
(34)--(41).

Finally, Proposition~\ref{prop:step-013-finite-integration} and
Lemma~\ref{lem:step-013-beta-gen} integrate the generalization charge, and
the mechanism ledger assembles

\[
 \beta_{\rm tr}+\beta_{\rm AT}+\beta_{\rm SS}+\beta_{\rm gen}=\beta.
\]

Proposition~\ref{prop:step-004-projection} transfers the marked
bound exactly to the released quotient law. This proves the exact PAC
conversion without conditioning on \(J\) for independence, releasing the
mark, using an uncountable union, or recharging privacy.



\end{proof}

\subsection{VC-Arm Rate Specialization}\label{app:step_014}

\begin{proposition}[Exact quotient-to-raw private PAC interface]\label{prop:step-014-interface}
Under Assumptions~\ref{assump:finite-littlestone},
\ref{assump:countable-evaluation-quotient},
\ref{assump:realizable-iid}, and
\ref{assump:approximate-dp-regime}, 
Propositions~\ref{prop:step-002-iid-pushforward},
\ref{prop:step-002-risk},
\ref{prop:step-011-raw-dp}, and
\ref{prop:step-013-pac}, suppose \(d\geq1\) and use the exact sample size
\(N=n_0\). Then \(K_C=K_C^{\mathrm{VC\text{-}Lyu}}\) is the released
quotient Markov kernel and
\[
 A_N(s,E):=K_C(T_N(s),E),
 \qquad s\in Z_X^N,\quad E\in\mathcal H_C,
\tag{1}
\]
is its raw-input Markov-kernel pullback. For every raw replace-one pair
\(s\sim s'\), including nonrealizable inputs and arbitrary labels,
\[
 A_N(s,E)\leq e^\varepsilon A_N(s',E)+\delta
 \qquad(E\in\mathcal H_C).
\tag{2}
\]
Moreover, for every probability measure \(D\) on \((X,\Sigma)\) and every
\(c\in C\),
\[
\begin{aligned}
&\Pr_{S\sim P_{D,c}^N,\,\bar H\sim A_N(S,\cdot)}
 \left[
 \operatorname{err}_D(\operatorname{Dec}_C(\bar H),c)>\alpha
 \right]\\
&\quad =
\Pr_{\bar S\sim P_{\bar D,\bar c}^N,\,
      \bar H\sim K_C(\bar S,\cdot)}
 \left[
 \operatorname{err}_{\bar D}(\bar H,\bar c)>\alpha
 \right]
\leq\beta.
\end{aligned}
\tag{3}
\]
The equality in (3) is exact for every possibly improper output.


\end{proposition}

\begin{proof}
Proposition~\ref{prop:step-011-raw-dp}, together with
the kernel inputs discharged in its  proof, gives both Markov
kernels in (1), the exact pullback identity, and (2). Its conclusion is
deterministic and pointwise over every raw labeled input; it does not use
realizability or any utility event.

Fix \(D,c\). Proposition~\ref{prop:step-002-risk} makes
\[
 B_{\alpha,D,c}
 :=\{\bar h\in H_C:
       \operatorname{err}_{\bar D}(\bar h,\bar c)>\alpha\}
\tag{4}
\]
measurable and gives, pointwise for every \(\bar h\in H_C\),
\[
 \mathbf1\!\left\{
 \operatorname{err}_D(\operatorname{Dec}_C(\bar h),c)>\alpha
 \right\}
 =\mathbf1\{\bar h\in B_{\alpha,D,c}\}.
\tag{5}
\]
Using (1), the raw failure probability is therefore
\[
 \int_{Z_X^N}
 K_C(T_N(s),B_{\alpha,D,c})\,P_{D,c}^N(ds).
\tag{6}
\]
The integrand is measurable because \(K_C\) is a kernel and \(T_N\) is
measurable. By the defining change-of-variables identity for the pushforward
measure and Proposition~\ref{prop:step-002-iid-pushforward},
\[
\begin{aligned}
 (6)
 &=\int_{Z_Q^N}
 K_C(\bar s,B_{\alpha,D,c})\,
 (T_N)_\#P_{D,c}^N(d\bar s)\\
 &=\int_{Z_Q^N}
 K_C(\bar s,B_{\alpha,D,c})\,
 P_{\bar D,\bar c}^N(d\bar s).
\end{aligned}
\tag{7}
\]
This is exactly the quotient failure probability in (3), including all
partition and mechanism randomness already integrated into \(K_C\).
Proposition~\ref{prop:step-013-pac} bounds it by \(\beta\) after
the analysis-only mark has been projected out. Thus (3) preserves the
released law, target function, and binary population-risk metric with zero
residual. \(\square\)


\end{proof}

\begin{lemma}[Ceiling-aware elimination of every VC-arm auxiliary]\label{lem:step-014-elimination}
Under Assumptions~\ref{assump:finite-littlestone}
and~\ref{assump:approximate-dp-regime}, Lemma~\ref{lem:step-001-envelope}, Proposition~\ref{prop:step-001-teacher}, and 
Lemmas~\ref{lem:step-007-fixed-point}
and~\ref{lem:step-007-sample-envelope}, suppose \(d\geq1\). Put
\[
 \ell:=\log\frac{64}{\delta\beta},
 \qquad
 A_{\log}:=80+\log(1+C_{\rm blk}),
 \qquad
 A_{\rm def}:=256(A_{\log}+c_{\rm AT}+1),
\tag{8}
\]
\[
 C_{\rm teach}:=2^{12}A_{\rm def}^2,
 \qquad
 H:=A_{\log}(1+\log C_{\rm teach}),
\tag{9}
\]
and
\[
 K_{\rm fp}:=
 2(1+C_{\rm blk})C_{\rm teach}H(H+3).
\tag{10}
\]
All constants in (8)-(10) are universal. For the exact  tuple
\[
 a(t)=v+\log(4t/\beta),\qquad
 Q(t)=e+\frac{etd^2a(t)}{\alpha v},
\tag{11}
\]
\[
 m(t)=\left\lceil
 C_{\rm blk}\frac{d^2}{\alpha}a(t)\log Q(t)
 \right\rceil,\qquad
 N=n_0=km(k),
\tag{12}
\]
one has, with every ceiling retained,
\[
 N\leq
 K_{\rm fp}
 \frac{d^4\ell\Lambda^3}{\varepsilon\alpha}
 (v+\Lambda).
\tag{13}
\]
Thus \(k,a(k),Q(k),m\), and \(N=km\) have all been replaced by displayed
parameters and universal constants.


\end{lemma}

\begin{proof}
Proposition~\ref{prop:step-001-teacher} defines
\[
 \bar k
 =
 \left\lceil
 C_{\rm teach}\frac{d^2\ell\Lambda^2}{\varepsilon}
 \right\rceil
\tag{14}
\]
before sampling and proves
\[
 2\leq k\leq\bar k
 \leq
 2C_{\rm teach}\frac{d^2\ell\Lambda^2}{\varepsilon}.
\tag{15}
\]
The final inequality is the exact ceiling bound
\(\lceil y\rceil\leq y+1\leq2y\), applied only after the  proof
establishes
\(y=C_{\rm teach}d^2\ell\Lambda^2/\varepsilon\geq1\).

We next retain the independent ceiling in \(m\). With
\[
 x:=C_{\rm blk}\frac{d^2}{\alpha}a(k)\log Q(k),
\]
write exactly \(m=x+\theta\), where \(0\leq\theta<1\). Lemma~\ref{lem:step-007-fixed-point} proves
\(d^2a(k)\log Q(k)/\alpha>4\), hence
\[
\begin{aligned}
 m
 &\leq
 C_{\rm blk}\frac{d^2}{\alpha}a(k)\log Q(k)+1\\
 &\leq
 (1+C_{\rm blk})\frac{d^2}{\alpha}a(k)\log Q(k).
\end{aligned}
\tag{16}
\]
The additive integer remainder has therefore been paid explicitly, rather
than hidden in asymptotic notation.

For integers \(2\leq s\leq t\), (11) gives
\(a(s)\leq a(t)\), then \(sa(s)\leq ta(t)\), and therefore
\(Q(s)\leq Q(t)\). Applying this deterministic monotonicity only after
\(k\leq\bar k\) is known, (15)-(16) give
\[
 N=km
 \leq
 (1+C_{\rm blk})\bar k\frac{d^2}{\alpha}
 a(\bar k)\log Q(\bar k).
\tag{17}
\]
No monotonicity of the teacher feasible set is asserted.

Lemma~\ref{lem:step-001-envelope}, at the fixed universal
\(C=C_{\rm teach}\), gives
\[
 \log\bar k\leq H\Lambda,
 \qquad
 \log Q(\bar k)\leq H\Lambda.
\tag{18}
\]
Writing \(b=\log(1/\beta)\), the parameter ranges give
\(b\leq\Lambda\), \(\Lambda\geq1\), and \(\log4<2\), so
\[
\begin{aligned}
 a(\bar k)
 &=v+\log4+\log\bar k+b\\
 &\leq v+(H+3)\Lambda.
\end{aligned}
\tag{19}
\]
Substituting (15), (18), and (19) into (17) yields the fully visible chain
\[
 N
 \leq
 2(1+C_{\rm blk})C_{\rm teach}H
 \frac{d^4\ell\Lambda^3}{\varepsilon\alpha}
 \bigl[v+(H+3)\Lambda\bigr].
\tag{20}
\]
Finally,
\[
 v+(H+3)\Lambda\leq(H+3)(v+\Lambda),
\tag{21}
\]
because \(H+3\geq1\). Equations (10), (20), and (21) prove (13).
This line-by-line chain eliminates \(k\) through (15), \(a(k)\) through
(19), \(Q(k)\) through its logarithm in (18), \(m\) through the exact
ceiling calculation (16), and \(N=km\) through (17)-(20). \(\square\)


\end{proof}

\begin{proposition}[Exact normalized VC rate and structural boundary profile]\label{prop:step-014-rate}
Under Assumption~\ref{assump:approximate-dp-regime} and
Lemma~\ref{lem:step-014-elimination}, suppose \(d\geq1\). Define
\[
 K_{\mathrm V}:=\max\{1,4K_{\rm fp}\},
 \qquad q:=4.
\tag{22}
\]
Then \(K_{\mathrm V}\geq1\) is universal, \(q\in\mathbb N_0\), and
\[
 N\leq
 K_{\mathrm V}\Lambda(d,v,\alpha,\beta,\varepsilon,\delta)^q
 R_{\mathrm{VC}}(d,v,\alpha,\beta,\varepsilon,\delta).
\tag{23}
\]
The bound is pointwise valid for every fixed \(0<\delta<1\), has only
displayed-parameter logarithms inside \(\Lambda\), has the \(d^4\) profile
at \(v=1\), and at \(v=d\) has the exact polynomial numerator
\(d^5+d^4\log(1/\beta)\).


\end{proposition}

\begin{proof}
For clarity only inside this proof, set
\[
 b:=\log(1/\beta),\qquad
 s:=\log(1/(\delta\beta)),\qquad
 \ell:=\log(64/(\delta\beta)).
\tag{24}
\]
Since \(0<\beta<1/4\) and \(0<\delta<1\),
\[
 b>\log4>1,\qquad s>\log4>1.
\tag{25}
\]
Moreover \(\log64=3\log4\), so
\[
 \ell=s+\log64\leq4s.
\tag{26}
\]
Because \(v\geq1\) and \(\Lambda\geq1\),
\[
 v+\Lambda\leq\Lambda(v+1)\leq\Lambda(v+b).
\tag{27}
\]
Applying (26)-(27) to (13) gives
\[
 N
 \leq
 4K_{\rm fp}\Lambda^4
 \frac{d^4(v+b)s}{\varepsilon\alpha}.
\tag{28}
\]
This is the only logarithmic absorption in the rate bridge, and both
inequalities proving it are displayed.

The setting's exact rate is
\[
 R_{\mathrm{VC}}
 =
 \frac{d^4(v+b)s}{\varepsilon\alpha}
 +\frac{d+b}{\alpha}.
\tag{29}
\]
The second term has not been silently discarded. Indeed
\(d\geq1\), \(v\geq1\), \(s>1\), and \(\varepsilon\leq1\) give
\[
 \frac{d^4(v+b)s}{\varepsilon\alpha}
 \geq
 \frac{d^4v+d^4b}{\alpha}
 \geq
 \frac{d+b}{\alpha}.
\tag{30}
\]
Thus the two terms in (29) are explicitly comparable, while in the
direction needed for (23),
\[
 \frac{d^4(v+b)s}{\varepsilon\alpha}
 \leq R_{\mathrm{VC}}.
\tag{31}
\]
Equations (22), (28), and (31) prove (23).

At \(v=1\), (28) reads
\[
 N\leq
 K_{\mathrm V}\Lambda^4
 \frac{d^4(1+b)s}{\varepsilon\alpha},
\tag{32}
\]
so no additional positive power of \(v\) or \(d\) is hidden. At \(v=d\),
the first numerator in (29) is exactly
\[
 d^4(d+b)=d^5+d^4b,
\tag{33}
\]
which is the claimed \(d^5\) structural scale, with its confidence term
still explicit. These are specializations of the same VC arm and do not
invoke or compare another arm.

Finally, (25)-(31) require only the fixed allowed range \(0<\delta<1\).
They remain finite and valid for an arbitrary such \(\delta\), whether
moderate or tending to zero. No limiting claim has been used to prove the
fixed-sample theorem. \(\square\)


\end{proof}

\begin{proposition}[Exact zero-dimensional no-data baseline]\label{prop:step-014-zero}
Under Assumptions~\ref{assump:finite-littlestone},
\ref{assump:countable-evaluation-quotient},
\ref{assump:realizable-iid}, and
\ref{assump:approximate-dp-regime}, and 
Propositions~\ref{prop:step-001-zero},
\ref{prop:step-002-risk},
\ref{prop:step-011-raw-dp}, and
\ref{prop:step-013-pac}, if \(d=0\), then \(v=0\), \(N=0\), the quotient
law and its raw pullback are the exact Dirac law at the unique
\(\bar c_0\in\bar C\), the raw law is \((0,0)\)-DP, and for every \(D\)
and \(c\in C\),
\[
 \operatorname{err}_D(\operatorname{Dec}_C(\bar c_0),c)=0
\tag{34}
\]
deterministically. Consequently the PAC failure probability is exactly
zero.


\end{proposition}

\begin{proof}
Proposition~\ref{prop:step-001-zero} proves that two
distinct concepts would shatter a depth-one Littlestone tree. Hence
\(d=0\) makes the nonempty class \(C\), and therefore \(\bar C\), a
singleton; it also implies \(v=0\). The  construction bypasses
the teacher, blocks, stages, lists, partition, mechanisms, and every
expression containing \(v^{-1}\), and sets \(N=0\).

Proposition~\ref{prop:step-011-raw-dp} identifies both quotient
and raw empty-input laws as the same Dirac kernel and proves exact
\((0,0)\)-DP. Since the unique released quotient concept is the quotient
factor of the unique target, Proposition~\ref{prop:step-002-risk} gives (34). Proposition~\ref{prop:step-013-pac} records the same released-law conclusion
with zero failure probability. This is an exact baseline, not a limit of
the positive-dimensional formula. \(\square\)


\end{proof}

\begin{proposition}[Conditional normalized VC-arm theorem and exact
small-\(\delta\) schedule]\label{prop:step-014-vc-arm}
Under Assumptions~\ref{assump:finite-littlestone},
\ref{assump:countable-evaluation-quotient},
\ref{assump:realizable-iid}, and
\ref{assump:approximate-dp-regime}, and
Propositions~\ref{prop:step-014-interface},
\ref{prop:step-014-rate}, and
\ref{prop:step-014-zero}, the quotient-first totalized VC-sensitive law has
the following exact common interface.

If \(d=0\), it uses \(N=0\), is a quotient/raw Markov-kernel Dirac law, is
\((0,0)\)-DP, has identically zero decoded population risk, and therefore
\(m_C(\alpha,\beta;\varepsilon,\delta)=0\). If
\(d\geq1\), it uses the exact  \(N=n_0=km\), its raw pullback is
\((\varepsilon,\delta)\)-DP on every raw neighboring labeled input, and
\[
\sup_D\sup_{c\in C}
\Pr_{S\sim P_{D,c}^N,\,\bar H\sim A_N(S,\cdot)}
\left[
\operatorname{err}_D(\operatorname{Dec}_C(\bar H),c)>\alpha
\right]
\leq\beta,
\tag{35}
\]
while the universal constants in (22) and the definition of \(m_C\) give
\[
 m_C(\alpha,\beta;\varepsilon,\delta)
 \leq N\leq K_{\mathrm V}\Lambda^qR_{\mathrm{VC}}.
\tag{36}
\]
Along any parameter sequence satisfying exactly
\[
 \delta K_{\mathrm V}\Lambda^qR_{\mathrm{VC}}\longrightarrow0,
\tag{37}
\]
the positive-dimensional sample sizes satisfy \(N\delta\to0\). The
fixed-parameter theorem, including privacy and PAC utility, remains valid
for every \(0<\delta<1\); no \(N\delta\) limit is claimed for fixed positive
\(\delta\) absent (37).

This is an explicitly conditional theorem for finite-or-countable
measurable evaluation quotients. It makes no claim for uncountable
evaluation quotients, proves no universal polynomial in
\((v,\log d)\) or \((v,\log^*d)\), and makes no comparison with the old-Lyu~\citep{lyu2025}
or finite-class arms.


\end{proposition}

\begin{proof}
On the positive branch,
Proposition~\ref{prop:step-014-interface} supplies the exact quotient/raw
kernel identity, all-input privacy, exact decoder-risk transfer, and (35).
Proposition~\ref{prop:step-014-rate} supplies (36) with the universal
\(K_{\mathrm V}=\max\{1,4K_{\rm fp}\}\) and integer \(q=4\), including every ceiling
and structural boundary checks; the first inequality in (36) follows
directly because Proposition~\ref{prop:step-014-interface} supplies an
admissible learner at sample size \(N\). On the null branch,
Proposition~\ref{prop:step-014-zero} supplies the exact no-data conclusion.

For a sequence satisfying (37), nonnegativity and (36) give the displayed
squeeze
\[
 0\leq N\delta
 \leq\delta K_{\mathrm V}\Lambda^qR_{\mathrm{VC}}
 \longrightarrow0.
\tag{38}
\]
This is the entire small-\(\delta\) argument. It does not infer (37) from
\(\delta\to0\), and it does not infer \(N\delta\to0\) for a fixed positive
\(\delta\). The scope statements merely record the unchanged primitive
restriction and the arms outside this result. 
Proposition~\ref{prop:step-014-interface} first composes the quotient
factorization, privacy, and holdout results; its exact integral identity keeps
the released quotient law, raw pullback, decoder, target, and binary
population risk aligned. Lemma~\ref{lem:step-014-elimination} then composes
the parameter closure and trace fixed point; equations (14)--(21) pay both ceilings
and eliminate \(k,a(k),Q(k),m,N=km\) with the explicit universal constant
\(K_{\rm fp}\). Proposition~\ref{prop:step-014-rate} proves each logarithmic
domination, retains the exact two-term \(R_{\mathrm{VC}}\), fixes
\(K_{\mathrm V}=\max\{1,4K_{\rm fp}\}\) and \(q=4\), and checks \(v=1\), \(v=d\), and
every allowed fixed \(\delta\). Proposition~\ref{prop:step-014-zero}
separately preserves the exact \(d=0\) baseline.

Together these named results establish the normalized VC-arm theorem. Its
final squeeze proves \(N\delta\to0\) only
under the displayed schedule, while fixed positive \(\delta\) remains covered
by the finite-parameter statement.



\end{proof}

\subsection{Old-Lyu Comparison Arm}\label{app:step_015}


\begin{proposition}[Old-arm quotient and source interfaces]\label{prop:step-015-interfaces}
Under Assumptions~\ref{assump:finite-littlestone},
\ref{assump:countable-evaluation-quotient}, and
\ref{assump:approximate-dp-regime},
Lemmas~\ref{lem:step-003-countable-promotion},
\ref{lem:step-004-occurrence}, and
Propositions~\ref{prop:step-002-record-map},
\ref{prop:step-002-iid-pushforward},
\ref{prop:step-002-risk},
\ref{prop:step-003-quotient-kernel},
\ref{prop:step-003-raw-pullback},
\ref{prop:step-004-lift}, and
\ref{prop:step-004-projection}, suppose \(d\ge1\). The following
transport and source interfaces are available for the old-arm
construction on \(Q_C\).

The map \(T_N\) is measurable and sends raw replace-one neighbors to equal
or replace-one quotient samples; realizable iid samples satisfy
\(T_N(S)\sim P_{\bar D,\bar c}^N\); and, for every \(\bar h\in H_C\),
\[
 \operatorname{err}_D(\operatorname{Dec}_C(\bar h),c)
 =\operatorname{err}_{\bar D}(\bar h,\bar c).
 \tag{D1}
\]
The countable quotient promotes every pointwise totalized law to a Markov
kernel. Whenever a totalized finite-list transcript has a nonempty producer
occurrence set, Lemma~\ref{lem:step-004-occurrence} and
Propositions~\ref{prop:step-004-lift} and
\ref{prop:step-004-projection} supply an analysis-only mark whose marginal is
exactly the released quotient law.

For a fixed binary population of size \(N\), mean \(\mu\), and a uniformly
sampled subset of size \(m\), the without-replacement bound from
Lyu~\citep{lyu2025} is, for \(0<\xi<1\),
\[
 \Pr(|\widehat\mu-\mu|>\xi\mu)
 \le2\exp(-\xi^2m\mu/3).
 \tag{C1}
\]
If \(\operatorname{LD}(\mathcal H)\le d\), valid \((p,d)\)-decompositions
exist, a valid tree has at most \(p^d2^{d^2}\) leaves, and a
maximum-dimensional leaf of dimension \(t\) is
\(p2^{d-t}\)-irreducible.  For \(\mathcal G\subseteq\mathcal H\),
\[
 \operatorname{DDim}_{2p,d}(\mathcal G)
 \le\operatorname{DDim}_{p,d}(\mathcal H).
 \tag{C2}
\]
When both sides equal \(t\), every dimension-\(t\) SOA from an optimal
\((2p,d)\)-decomposition of \(\mathcal G\) occurs in every optimal
\((p,d)\)-decomposition of \(\mathcal H\), and there are at most
\(p^d2^{d^2}\) essential hypotheses; at DDim zero the essential set is the
whole nonempty class.

The fixed family
\[
 \widehat{\bar C}_{d+1}:=
 \{\operatorname{SOA}_{\mathcal G}:\mathcal G\subseteq\bar C
       \text{ is }(d+1)\text{-irreducible}\}
 \tag{C3}
\]
has Littlestone dimension at most \(d\). For finite lists
\(\mathcal L_1,\ldots,\mathcal L_k\) of size at most \(L\), the Sparse
Sample law uses score \(s(h)=|\{i:h\in\mathcal L_i\}|\), a distinct failure
symbol of score \(B\), and satisfies
\[
 B\geq\frac{10\log(L/\delta_s)}{\varepsilon_s}
 \quad\Longrightarrow\quad
 (2\varepsilon_s,\delta_s)\text{-DP}
 \tag{C4}
\]
under one-list replacement. The stopped one-crossing AboveThreshold
interface for sensitivity-one queries has
\[
 \bigl(c_{\mathrm{AT}}\eta
 [\sqrt{\log(1/\delta_a)}+\log(1/\delta_a)],
 \delta_a\bigr)\text{-DP},
 \tag{C5}
\]
for a suitable universal \(c_{\mathrm{AT}}\ge1\). The accuracy event for
the Laplace noises of the at most \(d+1\) stage queries is derived directly
from their elementary tails when it is first used; it is not part of this
source-privacy interface.

The empirical implication used from Lyu~\citep{lyu2025}, Theorem~3 is the
following precise restriction argument: if a nonempty class \(\mathcal G\)
is irreducible for at least the indexed sample length,
\(h=\operatorname{SOA}_{\mathcal G}\), and every \(g\in\mathcal G\) has
empirical error at most \(u\), then \(h\) also has empirical error at most
\(u\). Indeed, a larger error for \(h\) would make the restriction of
\(\mathcal G\) along the full SOA-labeled sample empty, contradicting
irreducibility. No privacy or probability-one conclusion is part of this
interface.

Finally, Sauer--Shelah~\citep{sauer1972} gives at most
\(\sum_{j=0}^d\binom nj\le(en/d)^d\) traces when \(n\ge d\ge1\).
The finite-population and Bernoulli lower-tail inequalities needed later
are proved locally at their first use, rather than asserted by this
source wrapper or imported from an external authority.
For adaptive composition, let \(M_1\) be a kernel from a database space to a finite transcript
space \(\mathsf T\), and, for each \(t\in\mathsf T\), let \(M_{2,t}\) be a
kernel from the same database space to a measurable space
\((\mathsf Y,\mathcal Y)\). If \(M_1\) is
\((\varepsilon_1,\delta_1)\)-DP and the history-indexed kernels
\(M_{2,t}\) are uniformly \((\varepsilon_2,\delta_2)\)-DP, then their
adaptive joint law is
\((\varepsilon_1+\varepsilon_2,\delta_1+\delta_2)\)-DP. Measurable
postprocessing and mixing over common data-independent randomness preserve
these parameters.
\end{proposition}

\begin{proof}
Propositions~\ref{prop:step-002-record-map},
\ref{prop:step-002-iid-pushforward}, and
\ref{prop:step-002-risk} give the transport and risk identities in (D1).
Lemma~\ref{lem:step-003-countable-promotion} and
Propositions~\ref{prop:step-003-quotient-kernel} and
\ref{prop:step-003-raw-pullback} give the kernel promotions, while
Lemma~\ref{lem:step-004-occurrence} and
Propositions~\ref{prop:step-004-lift} and
\ref{prop:step-004-projection} give the stated analysis-only mark and
exact-projection interface.

Equation (C1) is Lyu~\citep{lyu2025}, Proposition~2. The decomposition,
leaf, irreducibility, DDim, and essential-set assertions in (C2) are the
current-notation forms of Lyu~\citep{lyu2025}, Definitions~4.2--4.3,
Lemmas~4.1 and~4.3, and Corollary~4.1. Equation (C3) is Lyu~\citep{lyu2025},
Lemma~4.2, applied with \(\operatorname{LD}(\bar C)=d\) from
Lemma~\ref{lem:step-002-ld}. Equations (C4)--(C5) are the privacy
interfaces of Lyu~\citep{lyu2025}, Algorithms~1--2 and Lemmas~3.1--3.2.
The displayed trace bound is Sauer--Shelah~\citep{sauer1972}. The
finite-population and Bernoulli lower-tail inequalities needed later are
proved locally at their first use; this source wrapper does not depend on
those later results.

It remains to verify the stated adaptive-composition interface. Fix an
ordered neighboring pair \(x,x'\), and write
\[
 P=M_1(x),\qquad Q=M_1(x'),\qquad
 P_t=M_{2,t}(x),\qquad Q_t=M_{2,t}(x').
\]
Because \(\mathsf T\) is finite, the subprobability measure
\[
 P^0(\{t\}):=\min\{P(\{t\}),e^{\varepsilon_1}Q(\{t\})\}
\]
satisfies \(P^0\le P\), \(P^0\le e^{\varepsilon_1}Q\), and
\(P^0(\mathsf T)\ge1-\delta_1\). Indeed, its missing mass is the
positive-part sum
\(\sup_{A\subseteq\mathsf T}\{P(A)-e^{\varepsilon_1}Q(A)\}\), which is at
most \(\delta_1\) by privacy of \(M_1\).

For each \(t\), take densities \(p_t,q_t\) with respect to
\(P_t+Q_t\) and define
\[
 dP_t^0:=\min\{p_t,e^{\varepsilon_2}q_t\}\,d(P_t+Q_t).
\]
The same positive-part identity and the uniform privacy of \(M_{2,t}\)
give
\[
 P_t^0\le P_t,\qquad P_t^0\le e^{\varepsilon_2}Q_t,\qquad
 P_t^0(\mathsf Y)\ge1-\delta_2
 \quad\text{for every }t.
\]
If \(J_x\) is the adaptive joint law, define, for every measurable
\(A\subseteq\mathsf T\times\mathsf Y\) with sections
\(A_t=\{y:(t,y)\in A\}\),
\[
 J_x^0(A):=\sum_{t\in\mathsf T}P^0(\{t\})P_t^0(A_t).
\]
Then \(J_x^0\le J_x\), and
\[
 J_x^0(\mathsf T\times\mathsf Y)
 \ge(1-\delta_2)P^0(\mathsf T)
 \ge1-\delta_1-\delta_2.
\]
Moreover,
\[
 \begin{aligned}
 J_x^0(A)
 &\le e^{\varepsilon_2}
       \sum_{t\in\mathsf T}P^0(\{t\})Q_t(A_t)\\
 &\le e^{\varepsilon_1+\varepsilon_2}
       \sum_{t\in\mathsf T}Q(\{t\})Q_t(A_t)
  =e^{\varepsilon_1+\varepsilon_2}J_{x'}(A).
 \end{aligned}
\]
Consequently
\[
 J_x(A)\le
 e^{\varepsilon_1+\varepsilon_2}J_{x'}(A)
 +\delta_1+\delta_2.
\]
Repeating the construction with \(x,x'\) reversed proves the reverse
neighbor inequality. Postprocessing follows by applying this bound to
inverse images, and integration over the same data-independent mixing law
preserves it. Combining these interfaces gives the stated old-arm package
without changing any source claim.
\(\square\)
\end{proof}





\begin{proposition}[Exact old-arm singleton branch]\label{prop:step-015-zero}
Under Assumptions~\ref{assump:finite-littlestone} and
\ref{assump:countable-evaluation-quotient}, if \(d=0\), then \(\bar C\)
has a unique element \(\bar c_0\). The old arm with
\(N_{\mathrm{old}}=0\), quotient law
\[
K_{\mathrm o}(\varnothing,E)=\mathbf1_E(\bar c_0),
\tag{1.1}
\]
and raw pullback along \(T_0\) is a measurable \((0,0)\)-DP Markov kernel.
For every \(D\) and \(c\in C\), its decoded population error is zero.


\end{proposition}

\begin{proof}
Lemma~\ref{lem:step-002-ld} gives
\(\operatorname{LD}(\bar C)=0\). If two binary functions in \(\bar C\)
differed at a quotient point, that point would form a depth-one
Littlestone tree; hence \(\bar C\) is a singleton. Equation (1.1) is a
Dirac probability law on the standard-Borel output space and is constant
on the unique empty input. It is therefore a Markov kernel and is
\((0,0)\)-DP. Proposition~\ref{prop:step-002-factorization}
identifies the unique \(\bar c_0\) with every raw target, and (D1) gives
zero decoded error. No data, partition, list, mechanism, or asymptotic
condition is present. \(\square\)


\end{proof}

\begin{lemma}[Independent old-arm scalar dictionary]\label{lem:step-015-dictionary}
Under Assumptions~\ref{assump:finite-littlestone} and
\ref{assump:approximate-dp-regime} and
Proposition~\ref{prop:step-015-interfaces}, suppose \(d\ge1\). Fix a universal
\(c_{\mathrm{AT}}\ge1\) for which the bound (C5) holds, and
fix the universal old block constant
\[
C_{\mathrm o}:=2^{20}.
\tag{2.1}
\]
Define
\[
\gamma_{\mathrm o}:=\frac{\alpha}{16},\qquad
\rho:=1-\frac1{2d},\qquad \xi_d:=\frac1{5d},
\tag{2.2}
\]
\[
\beta_{\mathrm{tr,o}}=\beta_{\mathrm{AT,o}}
=\beta_{\mathrm{SS,o}}=\beta_{\mathrm{gen,o}}:=\frac\beta4,
\qquad
\delta_{\mathrm{AT,o}}=\delta_{\mathrm{SS,o}}:=\frac\delta2,
\tag{2.3}
\]
\[
g_\delta:=\log(4/\delta),\qquad
\eta_{\mathrm o}:=
\frac{\varepsilon}
{4c_{\mathrm{AT}}(\sqrt{g_\delta}+g_\delta)},\qquad
\varepsilon_{\mathrm{SS,o}}:=\frac\varepsilon8.
\tag{2.4}
\]
For every integer \(t\ge2\), define
\[
a_{\mathrm o}(t):=d+\log(4t/\beta),\qquad
Q_{\mathrm o}(t):=
e+\frac{etd^2a_{\mathrm o}(t)}{\alpha d},
\tag{2.5}
\]
\[
m_{\mathrm o}(t):=
\left\lceil C_{\mathrm o}\frac{d^2}{\alpha}
a_{\mathrm o}(t)\log Q_{\mathrm o}(t)\right\rceil,\qquad
n_{\mathrm o}(t):=t\,m_{\mathrm o}(t),
\tag{2.6}
\]
\[
p_{\mathrm o,r}(t):=2^rn_{\mathrm o}(t)d
\quad(0\le r\le d),\qquad
L_{\mathrm o}(t):=p_{\mathrm o,d}(t)^d2^{d^2},
\tag{2.7}
\]
\[
B_{\mathrm o}(t):=
\left\lceil
\frac{10\log(L_{\mathrm o}(t)/\delta_{\mathrm{SS,o}})}
{\varepsilon_{\mathrm{SS,o}}}\right\rceil,
\tag{2.8}
\]
\[
\tau_{\mathrm{AT,o}}:=
\eta_{\mathrm o}^{-1}
\log((d+1)/\beta_{\mathrm{AT,o}}),\qquad
\tau_{\mathrm{SS,o}}(t):=
\varepsilon_{\mathrm{SS,o}}^{-1}
\log((tL_{\mathrm o}(t)+1)/\beta_{\mathrm{SS,o}}).
\tag{2.9}
\]
Every displayed quantity is finite and positive; \(m_{\mathrm o}(t)\),
\(n_{\mathrm o}(t)\), \(p_{\mathrm o,r}(t)\), \(L_{\mathrm o}(t)\), and
\(B_{\mathrm o}(t)\) are integers. Moreover
\[
B_{\mathrm o}(t)\ge
\frac{10\log(L_{\mathrm o}(t)/\delta_{\mathrm{SS,o}})}
{\varepsilon_{\mathrm{SS,o}}},
\qquad
|\mathcal L|\le p_{\mathrm o,r}(t)^d2^{d^2}
\le L_{\mathrm o}(t)
\tag{2.10}
\]
for every valid \((p_{\mathrm o,r}(t),d)\)-essential list
\(\mathcal L\).


\end{lemma}

\begin{proof}
The parameter ranges give \(g_\delta>\log4>0\),
\(a_{\mathrm o}(t)>0\), and \(Q_{\mathrm o}(t)>e\). Hence all logarithms
and reciprocals in (2.2)-(2.9) are legal and finite. The stated
integrality follows from the ceilings and integer products and powers.
The first inequality in (2.10) is exactly (2.8). For the second,
Lyu~\citep[Lemma~4.1 and Corollary~4.1]{lyu2025} gives
\(|\mathcal L|\le p_{\mathrm o,r}(t)^d2^{d^2}\); monotonicity
\(p_{\mathrm o,r}(t)\le p_{\mathrm o,d}(t)\) gives the final cap.
Thus the privacy threshold (C4) and every finite list cap are fixed
before any input is observed. \(\square\)


\end{proof}

\begin{lemma}[Ceiling-aware old teacher envelope]\label{lem:step-015-envelope}
Under Lemma~\ref{lem:step-015-dictionary}, put
\[
\ell:=\log(64/(\delta\beta)),\qquad
R_{\mathrm T,o}:=\frac{d^2\ell\Lambda^2}{\varepsilon}.
\tag{3.1}
\]
There is a universal constant \(A_{\mathrm o}\ge1\), depending only on
the two fixed universal constants \(c_{\mathrm{AT}}\) and \(C_{\mathrm o}\),
such that, for every universal \(C\ge1\) and
\[
t_C:=\lceil C R_{\mathrm T,o}\rceil,
\tag{3.2}
\]
one has
\[
\begin{aligned}
\log t_C+\log m_{\mathrm o}(t_C)+\log Q_{\mathrm o}(t_C)
&\le A_{\mathrm o}(1+\log C)\Lambda^2,\\
\log L_{\mathrm o}(t_C)
&\le A_{\mathrm o}(1+\log C)d^2\Lambda^2,
\end{aligned}
\tag{3.3}
\]
and
\[
B_{\mathrm o}(t_C)+\tau_{\mathrm{SS,o}}(t_C)
+\tau_{\mathrm{AT,o}}
\le A_{\mathrm o}(1+\log C)R_{\mathrm T,o}.
\tag{3.4}
\]


\end{lemma}

\begin{proof}
First note that \(\ell>1\), \(R_{\mathrm T,o}>1\), and
\[
t_C\le2CR_{\mathrm T,o}.
\tag{3.5}
\]
The definition of \(\Lambda\) gives
\[
\log d,\ \log(1/\alpha),\ \log(1/\beta),\
\log(1/\varepsilon),\ \log\Lambda\le\Lambda.
\tag{3.6}
\]
Also, with \(u=\log(e/\delta)\),
\[
\ell\le \log64+u+\log(1/\beta),\qquad
\log\ell\le 3\Lambda.
\tag{3.7}
\]
Indeed \(\log u\le\log(e+\log(e/\delta))\le\Lambda\), and the logarithm
of the sum in (3.7) is bounded by a universal constant plus the maximum of
\(\log u,\log\log(1/\beta),1\).
Equations (3.5)-(3.7) imply
\[
\log t_C\le A_1(1+\log C)\Lambda
\tag{3.8}
\]
for a numerical \(A_1\).

Because \(d\ge1\) and \(\Lambda\ge1\), (2.5) and (3.8) give
\[
a_{\mathrm o}(t_C)
\le A_2(1+\log C)d\Lambda,
\qquad
\log Q_{\mathrm o}(t_C)
\le A_3(1+\log C)\Lambda.
\tag{3.9}
\]
To see the second inequality without suppressing a parameter, use
\[
Q_{\mathrm o}(t_C)
\le2e\,t_Cd\,a_{\mathrm o}(t_C)/\alpha
\]
(the product on the right exceeds \(e\)) and take logarithms, invoking
(3.6), (3.8), and the first inequality in (3.9).

The unrounded expression in (2.6) exceeds one. Consequently
\[
m_{\mathrm o}(t_C)
\le2C_{\mathrm o}\frac{d^2}{\alpha}
a_{\mathrm o}(t_C)\log Q_{\mathrm o}(t_C),
\tag{3.10}
\]
and taking logarithms proves the remaining first-line term of (3.3).
Furthermore
\[
\begin{aligned}
\log L_{\mathrm o}(t_C)
&=d\log(2^dn_{\mathrm o}(t_C)d)+d^2\log2\\
&=2d^2\log2+d\log t_C+d\log m_{\mathrm o}(t_C)+d\log d\\
&\le A_4(1+\log C)d^2\Lambda,
\end{aligned}
\tag{3.11}
\]
which is stronger than the second line of (3.3).

It remains to expose the nonlogarithmic \(\delta\) charge. From
\(\varepsilon_{\mathrm{SS,o}}=\varepsilon/8\) and
\(\delta_{\mathrm{SS,o}}=\delta/2\),
\[
B_{\mathrm o}(t_C)
\le1+\frac{80}{\varepsilon}
[\log L_{\mathrm o}(t_C)+\log(2/\delta)].
\tag{3.12}
\]
Since \(t_CL_{\mathrm o}(t_C)\ge2\),
\[
\tau_{\mathrm{SS,o}}(t_C)
\le\frac8\varepsilon
[\log(8/\beta)+\log t_C+\log L_{\mathrm o}(t_C)].
\tag{3.13}
\]
Finally \(g_\delta>1\), so
\(\sqrt{g_\delta}+g_\delta\le2g_\delta\le2\ell\), while
\[
\log(4(d+1)/\beta)\le4\Lambda.
\]
Thus
\[
\tau_{\mathrm{AT,o}}
\le\frac{32c_{\mathrm{AT}}\ell\Lambda}{\varepsilon}.
\tag{3.14}
\]
Each summand in (3.12)-(3.14), including the ceiling remainder, is at
most a universal multiple of
\(d^2\ell\Lambda^2/\varepsilon\). Combining (3.8)-(3.14) and enlarging
one universal \(A_{\mathrm o}\) proves (3.3)-(3.4). No feasibility
assumption was used to define \(t_C\) or any quantity evaluated at it.
\(\square\)


\end{proof}

\begin{proposition}[Feasible least old teacher and exact margin]\label{prop:step-015-teacher}
Under Lemmas~\ref{lem:step-015-dictionary} and
\ref{lem:step-015-envelope}, let
\[
C_{\mathrm{teach,o}}:=(16A_{\mathrm o})^2,\qquad
\bar k_{\mathrm o}:=
\left\lceil C_{\mathrm{teach,o}}
\frac{d^2\log(64/(\delta\beta))\Lambda^2}{\varepsilon}
\right\rceil.
\tag{4.1}
\]
Then
\[
\frac{\bar k_{\mathrm o}}2-\tau_{\mathrm{AT,o}}
\ge B_{\mathrm o}(\bar k_{\mathrm o})
+\tau_{\mathrm{SS,o}}(\bar k_{\mathrm o})+2.
\tag{4.2}
\]
Consequently the least feasible integer
\[
k_{\mathrm o}:=
\min\left\{t\in\mathbb Z:t\ge2,\quad
\frac t2-\tau_{\mathrm{AT,o}}
\ge B_{\mathrm o}(t)+\tau_{\mathrm{SS,o}}(t)+2\right\}
\tag{4.3}
\]
exists and
\[
2\le k_{\mathrm o}\le\bar k_{\mathrm o}
\le2C_{\mathrm{teach,o}}
\frac{d^2\log(64/(\delta\beta))\Lambda^2}{\varepsilon}.
\tag{4.4}
\]
Define, only after this minimization,
\[
\begin{gathered}
a_{\mathrm o}:=a_{\mathrm o}(k_{\mathrm o}),\quad
Q_{\mathrm o}:=Q_{\mathrm o}(k_{\mathrm o}),\quad
m_{\mathrm o}:=m_{\mathrm o}(k_{\mathrm o}),\\
n_{\mathrm o}:=k_{\mathrm o}m_{\mathrm o}
=N_{\mathrm o}=N_{\mathrm{old}},\quad
p_{\mathrm o,r}:=2^rn_{\mathrm o}d,\quad
L_{\mathrm o}:=L_{\mathrm o}(k_{\mathrm o}),\\
B_{\mathrm o}:=B_{\mathrm o}(k_{\mathrm o}),\qquad
\tau_{\mathrm{SS,o}}:=\tau_{\mathrm{SS,o}}(k_{\mathrm o}).
\end{gathered}
\tag{4.5}
\]
This exact tuple satisfies
\[
\frac{k_{\mathrm o}}2-\tau_{\mathrm{AT,o}}
\ge B_{\mathrm o}+\tau_{\mathrm{SS,o}}+2,
\qquad
p_{\mathrm o,0}=n_{\mathrm o}d
\ge\max\{n_{\mathrm o},d+1\}.
\tag{4.6}
\]


\end{proposition}

\begin{proof}
For \(x\ge1\), \(\log x\le\sqrt x\). Hence, with
\(C=C_{\mathrm{teach,o}}\),
\[
A_{\mathrm o}(1+\log C)
\le A_{\mathrm o}+16A_{\mathrm o}^2
\le\frac C4.
\tag{4.7}
\]
Lemma~\ref{lem:step-015-envelope} therefore bounds the three defects at
\(\bar k_{\mathrm o}\) by
\[
\frac{C_{\mathrm{teach,o}}R_{\mathrm T,o}}4.
\tag{4.8}
\]
On the other hand,
\(\bar k_{\mathrm o}/2\ge
C_{\mathrm{teach,o}}R_{\mathrm T,o}/2\).
Since \(R_{\mathrm T,o}>\log256>5\), the remaining quarter of the
teacher scale exceeds \(2\), proving (4.2). Thus the set in (4.3) is a
nonempty subset of the integers bounded below by two, so it has a least
element and (4.4) follows from
\(\lceil C_{\mathrm{teach,o}}R_{\mathrm T,o}\rceil
\le2C_{\mathrm{teach,o}}R_{\mathrm T,o}\).
Equation (4.6)'s first part is (4.3) evaluated at the least element.
Finally \(n_{\mathrm o}\ge k_{\mathrm o}\ge2\); hence
\(n_{\mathrm o}d\ge n_{\mathrm o}\) and
\(n_{\mathrm o}d\ge2d\ge d+1\). All ceilings remain in (4.5).
\(\square\)


\end{proof}

\begin{proposition}[All-path totalization of the old quotient procedure]\label{prop:step-015-total}
Under Assumptions~\ref{assump:finite-littlestone},
\ref{assump:countable-evaluation-quotient}, and
\ref{assump:approximate-dp-regime},
Lyu~\citep[Claim~4.1, Definitions~4.2--4.3, Lemmas~4.1 and~4.3,
and Corollary~4.1]{lyu2025},
and Proposition~\ref{prop:step-015-teacher}, the positive-dimensional old
quotient procedure can be fixed before data so that it outputs an element
of \(H_C\) on every input in \(Z_Q^{N_{\mathrm o}}\) and every internal
path. It uses only the old tuple (4.5), preserves every exact valid
path, and sends every empty, invalid, failure-symbol, no-success, exhausted,
or residual path to the fixed \(\bar c_0\).


\end{proposition}

\begin{proof}
Fix once and for all the old tuple, the default \(\bar c_0\), a
data-independent uniform partition rule into \(k_{\mathrm o}\) indexed
blocks of size \(m_{\mathrm o}\), and deterministic choices for identical
local states. For block \(i\) and stage \(r=0,\ldots,d\), form the exact
restriction
\[
H_{\mathrm o,i}^r
:=\{h\in\bar C:
e_{\bar S_{\mathrm o,i}}(h)
\le\rho^{r+1}\gamma_{\mathrm o}\}.
\tag{5.1}
\]
For every nonempty restriction,
Lyu~\citep[Claim~4.1 and Definition~4.3]{lyu2025} supplies a
valid optimal \((p_{\mathrm o,r},d)\)-decomposition. Define
\(\mathcal L_{\mathrm o,i}^r\) to be its exact set of
\((p_{\mathrm o,r},d)\)-essential SOA functions, in one fixed
no-repetition order. By (2.10) it is a finite \(H_C\)-list of size at most
\(L_{\mathrm o}\). Empty restrictions receive the empty list. If a
  purported list object is undefined, nonfinite, outside \(H_C\), or over
the cap, record an internal failure token and replace only that local list
by the empty list. These checks do not alter a valid list.

At each stage define the total score
\[
q_{\mathrm o,r}:=
\begin{cases}
\displaystyle\max_{h\in\cup_i\mathcal L_{\mathrm o,i}^r}
|\{i:h\in\mathcal L_{\mathrm o,i}^r\}|,&
\cup_i\mathcal L_{\mathrm o,i}^r\ne\varnothing,\\
0,&\text{otherwise}.
\end{cases}
\tag{5.2}
\]
Feed the \(d+1\) scores in order to one stopped AboveThreshold process at
threshold \(k_{\mathrm o}/2\). At its first legal Above stage, invoke the
exact Sparse Sample law on the \(k_{\mathrm o}\) lists using
\((\varepsilon_{\mathrm{SS,o}},B_{\mathrm o})\). Return its actual union
item if one is drawn. Return \(\bar c_0\) if AboveThreshold never selects a
legal stage, if Sparse Sample draws \(\perp\), or if any transcript
coordinate is invalid or exhausted. Because the failure-symbol weight is
always positive and the effective union is finite, the valid Sparse
Sample normalizer is finite and nonzero. Thus every path terminates in
\(H_C\), while exact restrictions, decompositions, lists, mechanism laws,
and successful outputs are unchanged on valid paths. No event,
realizability condition, privacy conclusion, or utility conclusion enters
this definition. \(\square\)


\end{proof}

\begin{proposition}[Old quotient kernel, raw pullback, and exact risk]\label{prop:step-015-kernel}
Under Assumption~\ref{assump:countable-evaluation-quotient}, 
Propositions~\ref{prop:step-002-record-map},
\ref{prop:step-002-risk}, and
\ref{prop:step-003-raw-pullback}, Lemma~\ref{lem:step-003-countable-promotion}, and
Proposition~\ref{prop:step-015-total}, the totalized pointwise old law
defines a Markov kernel
\[
K_{\mathrm o}=K_C^{\mathrm{old\text{-}Lyu}}:
Z_Q^{N_{\mathrm o}}\leadsto H_C.
\tag{6.1}
\]
Its raw pullback
\[
A_{\mathrm o,N_{\mathrm o}}(s,E)
:=K_{\mathrm o}(T_{N_{\mathrm o}}(s),E)
\tag{6.2}
\]
is a Markov kernel
\(Z_X^{N_{\mathrm o}}\leadsto H_C\). For every \(D,c,\bar h\),
the released decoded risk is exactly (D1).


\end{proposition}

\begin{proof}
Proposition~\ref{prop:step-015-total} assigns a probability law on
\((H_C,\mathcal H_C)\) to every atom of the countable-discrete input
\(Z_Q^{N_{\mathrm o}}\). Lemma~\ref{lem:step-003-countable-promotion} promotes this arbitrary
pointwise family to the Markov kernel (6.1); in particular, data-dependent
finite lists and fallback choices do not require a measurable global
selector. Proposition~\ref{prop:step-002-record-map} makes
\(T_{N_{\mathrm o}}\) measurable, so Proposition~\ref{prop:step-003-raw-pullback} proves (6.2) is a kernel.
The output codomain is the full \(H_C\), and Proposition~\ref{prop:step-002-risk} applies to every such output,
including improper SOA functions and \(\bar c_0\). \(\square\)


\end{proof}

\begin{lemma}[Coarse \(d\)-trace count for the old master sample]\label{lem:step-015-traces}
Under Assumption~\ref{assump:finite-littlestone} and 
Lemmas~\ref{lem:step-002-vc} and \ref{lem:step-002-ld}, suppose \(d\ge1\)
and fix any labeled quotient tuple
\(\bar s=((q_j,y_j))_{j=1}^{n_{\mathrm o}}\). Then its old error-trace
family
\[
\mathcal E_{\bar C}(\bar s)
:=\{(\mathbf1\{h(q_j)\ne y_j\})_{j=1}^{n_{\mathrm o}}:
h\in\bar C\}
\tag{7.1}
\]
is finite and satisfies
\[
|\mathcal E_{\bar C}(\bar s)|
\le\left(\frac{en_{\mathrm o}}d\right)^d.
\tag{7.2}
\]


\end{lemma}

\begin{proof}
 quotient preservation gives
\(\operatorname{VC}(\bar C)=v\le d\). On the indexed tuple, xor with the
fixed label vector \((y_j)_j\) is a bijection between prediction traces and
the error traces (7.1). Repeated quotient points can only identify
prediction traces, never create new ones. Proposition
\ref{prop:step-015-teacher} gives \(n_{\mathrm o}\ge d\), so
Sauer--Shelah~\citep{sauer1972} applied with the coarser upper bound \(d\) yields
\[
|\mathcal E_{\bar C}(\bar s)|
\le\sum_{j=0}^d\binom{n_{\mathrm o}}j
\le(en_{\mathrm o}/d)^d.
\]
No value of \(v\), no finite cardinality of \(C\), and no event from
another arm is used. \(\square\)


\end{proof}

\begin{lemma}[Both fixed-trace without-replacement tails]\label{lem:step-015-tails}
Under Assumption~\ref{assump:approximate-dp-regime} and
Proposition~\ref{prop:step-015-interfaces} and
Lemma~\ref{lem:step-015-dictionary}, condition on a complete indexed old
master sample and fix one binary error trace of mean
\(\mu\in[0,1]\). If \(\widehat\mu_i\) is its mean on one uniformly
partitioned block of size \(m_{\mathrm o}\), then:
\[
\Pr\!\left[
|\widehat\mu_i-\mu|>\xi_d\mu\mid\bar S_{\mathrm o}
\right]
\le2\exp\!\left(-\frac{m_{\mathrm o}\mu}{75d^2}\right)
\le2\exp\!\left(-\frac{m_{\mathrm o}\alpha}{3600d^2}\right)
\tag{8.1}
\]
whenever \(\mu>\gamma_{\mathrm o}/3\), and
\[
\Pr[\widehat\mu_i\ge\gamma_{\mathrm o}/2\mid\bar S_{\mathrm o}]
\le\exp[-m_{\mathrm o}
D(\gamma_{\mathrm o}/2\Vert\mu)]
\le\exp(-m_{\mathrm o}\gamma_{\mathrm o}/30)
\tag{8.2}
\]
whenever \(0\le\mu\le\gamma_{\mathrm o}/3\). The second statement includes
the boundary \(\mu=\gamma_{\mathrm o}/3\); at \(\mu=0\) its failure event
is empty.


\end{lemma}

\begin{proof}
For (8.1), apply Lyu~\citep{lyu2025} Proposition 2, equation (C1), with
\(\xi_d=1/(5d)\). Its exponent is
\(\xi_d^2m_{\mathrm o}\mu/3=m_{\mathrm o}\mu/(75d^2)\).
Since
\(\mu>\gamma_{\mathrm o}/3=\alpha/48\), the second inequality follows.

For (8.2), write the fixed population as
\(x_1,\ldots,x_{n_{\mathrm o}}\in\{0,1\}\). We first prove (C6)
directly. For \(N\ge1\), let \(e_m(w)\) be the degree-\(m\) elementary
symmetric polynomial of a nonnegative vector \(w\in[0,\infty)^N\), with
\(e_0=1\). We establish
\[
 e_m(w_1,\ldots,w_N)
 \le\binom Nm\left(\frac1N\sum_{j=1}^Nw_j\right)^m,
 \qquad 0\le m\le N,
\tag{C6}
\]
where a zeroth power is one. The
cases \(m=0\) and \(m=1\) are equalities in (C6), as is every case when
\(N=1\). Suppose now that \(N\ge2\) and \(2\le m\le N\). Fix two
coordinates \(a,b\) and all remaining coordinates \(z\). Expanding by
whether a monomial contains neither, one, or both of the selected
coordinates, with \(e_j(z):=0\) outside \(0\le j\le N-2\), gives
\[
e_m(a,b,z)=A_{\mathrm{sm}}+(a+b)B_{\mathrm{sm}}+abC_{\mathrm{sm}},
\qquad
A_{\mathrm{sm}}=e_m(z),\quad B_{\mathrm{sm}}=e_{m-1}(z),\quad
C_{\mathrm{sm}}=e_{m-2}(z)\ge0.
\tag{8.3a}
\]
Replacing \(a,b\) by their common average preserves \(a+b\), while
\[
\left(\frac{a+b}{2}\right)^2-ab=\frac{(a-b)^2}{4}\ge0.
\tag{8.3b}
\]
Thus this averaging operation cannot decrease \(e_m\).

Starting from \(w^{(0)}=w\), repeatedly average a maximum and a minimum
coordinate, and let \(\bar w=N^{-1}\sum_jw_j\). The sum and hence
\(\bar w\) stay fixed. If
\[
V_s:=\sum_{j=1}^N(w_j^{(s)}-\bar w)^2,
\]
then one averaging step satisfies
\[
V_s-V_{s+1}
=\frac12\left(\max_jw_j^{(s)}-\min_jw_j^{(s)}\right)^2.
\tag{8.3c}
\]
The nonnegative sequence \((V_s)_s\) decreases and converges, so the
successive differences tend to zero. Hence the displayed maximum--minimum
gap tends to zero. Since \(\bar w\) always lies between the minimum and the
maximum, every coordinate of \(w^{(s)}\) tends to \(\bar w\). Polynomial
continuity and the monotonicity just proved now give
\[
e_m(w)=e_m(w^{(0)})
\le\lim_{s\to\infty}e_m(w^{(s)})
=\binom Nm\bar w^m,
\tag{8.3d}
\]
which proves (C6), including \(m=N\); the already separated \(m=0,1\)
cases cover the other endpoints. Here
\(N=n_{\mathrm o}=k_{\mathrm o}m_{\mathrm o}\),
\(m=m_{\mathrm o}\), \(m_{\mathrm o}\ge1\), and \(k_{\mathrm o}\ge2\),
so \(1\le m<N\).

For every \(\lambda>0\), apply the proved inequality (C6) to
\(w_j=e^{\lambda x_j}\). Since
\(N^{-1}\sum_jw_j=1-\mu+\mu e^\lambda\), it gives
\[
\begin{aligned}
\mathbb E\exp\{\lambda m_{\mathrm o}\widehat\mu_i\}
&=\frac{e_{m_{\mathrm o}}
(e^{\lambda x_1},\ldots,e^{\lambda x_{n_{\mathrm o}}})}
{\binom{n_{\mathrm o}}{m_{\mathrm o}}}\\
&\le(1-\mu+\mu e^\lambda)^{m_{\mathrm o}}.
\end{aligned}
\tag{8.3}
\]
Markov's inequality and minimization over \(\lambda>0\) therefore yield,
for \(q>\mu\),
\[
\Pr(\widehat\mu_i\ge q)
\le\inf_{\lambda>0}
\exp\{m_{\mathrm o}[\log(1-\mu+\mu e^\lambda)-\lambda q]\}
=e^{-m_{\mathrm o}D(q\Vert\mu)}.
\tag{8.4}
\]
For \(\mu=0\) the event is impossible. Otherwise, for fixed
\(q=\gamma_{\mathrm o}/2\), the derivative of
\(D(q\Vert\mu)\) in \(\mu<q\) is
\((\mu-q)/[\mu(1-\mu)]<0\). Hence the minimum over
\(\mu\le\gamma_{\mathrm o}/3\) occurs at \(\gamma_{\mathrm o}/3\).
Using \(\log(1-u)\ge-u/(1-u)\) for \(0\le u<1\),
\[
\begin{aligned}
D(\gamma_{\mathrm o}/2\Vert\gamma_{\mathrm o}/3)
&=\frac{\gamma_{\mathrm o}}2\log\frac32
+\left(1-\frac{\gamma_{\mathrm o}}2\right)
\log\frac{1-\gamma_{\mathrm o}/2}{1-\gamma_{\mathrm o}/3}\\
&\ge\frac{\gamma_{\mathrm o}}2\log\frac32
-\frac{\gamma_{\mathrm o}}6
\ge\frac{\gamma_{\mathrm o}}{30},
\end{aligned}
\tag{8.5}
\]
where the last line uses \(\log(3/2)\ge2/5\). Substitution in (8.4)
proves (8.2). This local one-sided proof is why Proposition 2 is not
misapplied at or near zero. \(\square\)


\end{proof}

\begin{proposition}[Generated old trace event and confidence charge]\label{prop:step-015-good}
Under Assumptions~\ref{assump:finite-littlestone} and
\ref{assump:approximate-dp-regime},
Proposition~\ref{prop:step-015-teacher}, and
Lemmas~\ref{lem:step-015-traces} and \ref{lem:step-015-tails}, define
\(E_{\mathrm{good,o}}\) to be the event that, simultaneously for every
trace in \(\mathcal E_{\bar C}(\bar S_{\mathrm o})\) and every block \(i\),
\[
\widehat\mu_i\in
\begin{cases}
[(1-\xi_d)\mu,(1+\xi_d)\mu],
&\mu>\gamma_{\mathrm o}/3,\\
[0,\gamma_{\mathrm o}/2],
&0\le\mu\le\gamma_{\mathrm o}/3.
\end{cases}
\tag{9.1}
\]
Then
\[
\Pr(E_{\mathrm{good,o}}^c\mid\bar S_{\mathrm o})
\le4k_{\mathrm o}
\left(\frac{en_{\mathrm o}}d\right)^d
\exp\!\left(-\frac{m_{\mathrm o}\alpha}{3600d^2}\right),
\tag{9.2}
\]
and the exact old calibration gives
\[
\Pr(E_{\mathrm{good,o}}^c)\le\beta_{\mathrm{tr,o}}=\beta/4.
\tag{9.3}
\]


\end{proposition}

\begin{proof}
The two cases in (9.1) are disjoint and exhaustive, with equality at
\(\gamma_{\mathrm o}/3\) assigned to the low case. Conditional on the
master sample, each block is marginally a uniform subset; no conditional
independence among different blocks is asserted. Lemma
\ref{lem:step-015-tails}, followed by the finite union over the
\(k_{\mathrm o}\) blocks and the traces counted in
Lemma~\ref{lem:step-015-traces}, gives (9.2). The factor \(4\) absorbs the
factor \(2\) in the high branch and the stronger low-branch exponent:
since \(d\ge1\),
\[
\frac{\gamma_{\mathrm o}}{30}
=\frac{\alpha}{480}
\ge\frac{\alpha}{3600d^2}.
\tag{9.4}
\]

It remains to discharge the fixed point and the ceiling. The unrounded
quantity in (2.6) exceeds one, so (3.10), now at \(k_{\mathrm o}\), gives
\[
\frac{en_{\mathrm o}}d
\le\frac{2eC_{\mathrm o}k_{\mathrm o}d
a_{\mathrm o}\log Q_{\mathrm o}}{\alpha}
\le2C_{\mathrm o}Q_{\mathrm o}\log Q_{\mathrm o}.
\tag{9.5}
\]
Because \(Q_{\mathrm o}\ge e\),
\[
\log(en_{\mathrm o}/d)
\le C_5\log Q_{\mathrm o},\qquad
C_5:=2+\log(2C_{\mathrm o}).
\tag{9.6}
\]
Also \(m_{\mathrm o}\alpha/d^2
\ge C_{\mathrm o}a_{\mathrm o}\log Q_{\mathrm o}\), while
\[
\begin{aligned}
\log(4k_{\mathrm o})+\log(1/\beta_{\mathrm{tr,o}})
&=\log(16k_{\mathrm o}/\beta)\\
&\le2a_{\mathrm o}\log Q_{\mathrm o},\\
d\log(en_{\mathrm o}/d)
&\le C_5a_{\mathrm o}\log Q_{\mathrm o}.
\end{aligned}
\tag{9.7}
\]
The fixed choice \(C_{\mathrm o}=2^{20}\) satisfies
\[
\frac{C_{\mathrm o}}{3600}>C_5+2.
\tag{9.8}
\]
Combining (9.6)-(9.8) yields
\[
\frac{m_{\mathrm o}\alpha}{3600d^2}
-\log(4k_{\mathrm o})-d\log(en_{\mathrm o}/d)
\ge\log(1/\beta_{\mathrm{tr,o}}).
\tag{9.9}
\]
Insert (9.9) into (9.2), then integrate over the master sample. This proves
(9.3) without using a VC-arm fixed point or event. \(\square\)


\end{proof}

\begin{lemma}[Exact old source endpoint, half-scale map, and inclusion]\label{lem:step-015-source-map}
Under Propositions~\ref{prop:step-015-interfaces} and
\ref{prop:step-015-good}, define the source-indexed
restrictions and scales, including the endpoint omitted from the source's
pre-algorithm display, by
\[
H_{\mathrm o,i,\mathrm{src}}^s
:=\{h\in\bar C:
e_{\bar S_{\mathrm o,i}}(h)\le\rho^s\gamma_{\mathrm o}\},
\quad
p_{\mathrm o,s,\mathrm{src}}:=2^sn_{\mathrm o}d,
\quad 1\le s\le d+1.
\tag{10.1}
\]
Then the current construction (5.1) obeys the exact map
\[
H_{\mathrm o,i}^r=H_{\mathrm o,i,\mathrm{src}}^{r+1},
\qquad
p_{\mathrm o,r}=\tfrac12p_{\mathrm o,r+1,\mathrm{src}},
\qquad 0\le r\le d.
\tag{10.2}
\]
On \(E_{\mathrm{good,o}}\), for every \(0\le r<d\), every \(i_*\), and
every \(i\),
\[
H_{\mathrm o,i_*}^{r+1}\subseteq H_{\mathrm o,i}^r.
\tag{10.3}
\]
The two constants needed for this inclusion are
\[
\rho^d\ge\frac12,\qquad
\frac{1+\xi_d}{1-\xi_d}\rho\le1.
\tag{10.4}
\]


\end{lemma}

\begin{proof}
Equation (10.2) follows by substituting \(s=r+1\) into (10.1) and comparing
with (5.1) and (4.5). The uniform factor \(1/2\) preserves the exact
doubling relation
\(p_{\mathrm o,r+1}=2p_{\mathrm o,r}\) used by (C2).

Bernoulli's inequality gives
\[
\rho^d=(1-1/(2d))^d\ge1-d/(2d)=1/2.
\]
For the second part of (10.4), direct multiplication reduces the desired
inequality to
\[
\frac{2}{5d}\le\frac1{2d}\left(1+\frac1{5d}\right),
\]
which holds for \(d\ge1\).

Fix \(h\in H_{\mathrm o,i_*}^{r+1}\), and write
\(\mu=e_{\bar S_{\mathrm o}}(h)\). If
\(\mu\le\gamma_{\mathrm o}/3\), (9.1) gives
\[
e_{\bar S_{\mathrm o,i}}(h)\le\gamma_{\mathrm o}/2
\le\rho^d\gamma_{\mathrm o}
\le\rho^{r+1}\gamma_{\mathrm o},
\]
where \(r+1\le d\) and \(0<\rho<1\). If
\(\mu>\gamma_{\mathrm o}/3\), then (9.1), membership in the left side,
and (10.4) give
\[
\begin{aligned}
e_{\bar S_{\mathrm o,i}}(h)
&\le(1+\xi_d)\mu\\
&\le\frac{1+\xi_d}{1-\xi_d}
e_{\bar S_{\mathrm o,i_*}}(h)\\
&\le\frac{1+\xi_d}{1-\xi_d}
\rho^{r+2}\gamma_{\mathrm o}
\le\rho^{r+1}\gamma_{\mathrm o}.
\end{aligned}
\]
Thus \(h\in H_{\mathrm o,i}^r\), proving (10.3). Both concentration
branches, rather than an unstated relative estimate at zero, are needed.
\(\square\)


\end{proof}

\begin{lemma}[Exact old essential lists and irreducible leaf witnesses]\label{lem:step-015-lists}
Under Assumption~\ref{assump:finite-littlestone},
Proposition~\ref{prop:step-015-teacher},
Lemma~\ref{lem:step-015-source-map}, and
Lyu~\citep[Claim~4.1, Definition~4.3, Lemmas~4.1 and~4.3, and
Corollary~4.1]{lyu2025}, every nonempty \(H_{\mathrm o,i}^r\) has the
exact finite set \(\mathcal L_{\mathrm o,i}^r\) of
\((p_{\mathrm o,r},d)\)-essential hypotheses, with
\[
|\mathcal L_{\mathrm o,i}^r|
\le p_{\mathrm o,r}^d2^{d^2}\le L_{\mathrm o}.
\tag{11.1}
\]
If \(f\in\mathcal L_{\mathrm o,i}^r\) and
\(\mathcal G\) is a maximum-dimensional leaf of dimension \(t\) witnessing
\(f=\operatorname{SOA}_{\mathcal G}\), then
\[
\mathcal G\subseteq H_{\mathrm o,i}^r,\qquad
\mathcal G\text{ is }
p_{\mathrm o,r}2^{d-t}\text{-irreducible},
\tag{11.2}
\]
and
\[
p_{\mathrm o,r}2^{d-t}\ge p_{\mathrm o,0}
=n_{\mathrm o}d\ge\max\{n_{\mathrm o},d+1\}.
\tag{11.3}
\]
Empty restrictions retain the fixed empty-list convention and are not
passed to a nonempty-class source assertion.


\end{lemma}

\begin{proof}
Every restriction is a subclass of \(\bar C\) and therefore has
Littlestone dimension at most \(d\).
Lyu~\citep[Claim~4.1 and Definition~4.3]{lyu2025} gives a valid optimal
decomposition for each nonempty restriction, and
Lyu~\citep[Lemma~4.1 and Corollary~4.1]{lyu2025} gives (11.1). By the
definition of essentiality, \(f\) occurs as the SOA of a
maximum-dimensional leaf in every optimal decomposition; validity of that
leaf gives (11.2). Since \(0\le t\le d\), \(r\ge0\), and
\(p_{\mathrm o,r}=2^rp_{\mathrm o,0}\), its irreducibility scale is at
least \(p_{\mathrm o,0}\), and Proposition
\ref{prop:step-015-teacher} gives (11.3).

If a class is \(q\)-irreducible, it is \(q'\)-irreducible for
\(1\le q'\le q\): extend any \(q'\)-tuple to length \(q\) by repeating a
fixed quotient point, apply \(q\)-irreducibility, and use monotonicity of
Littlestone dimension under restriction. Thus every leaf in (11.2) is in
particular both \(n_{\mathrm o}\)- and \((d+1)\)-irreducible. \(\square\)


\end{proof}

\begin{proposition}[Finite old DDim descent and common-score stage]\label{prop:step-015-descent}
Under Assumptions~\ref{assump:finite-littlestone} and
\ref{assump:realizable-iid},
Proposition~\ref{prop:step-015-interfaces}, on \(E_{\mathrm{good,o}}\), and under
Lemmas~\ref{lem:step-015-source-map} and
\ref{lem:step-015-lists}, define
\[
q_{\mathrm o,r}:=
\max_{h\in H_C}|\{i:h\in\mathcal L_{\mathrm o,i}^r\}|,
\qquad
M_{\mathrm o,r}:=
\max_i\operatorname{DDim}_{p_{\mathrm o,r},d}
(H_{\mathrm o,i}^r).
\tag{12.1}
\]
Then, for \(0\le r<d\),
\[
q_{\mathrm o,r}<k_{\mathrm o}
\quad\Longrightarrow\quad
M_{\mathrm o,r+1}\le M_{\mathrm o,r}-1,
\tag{12.2}
\]
and
\[
\max_{0\le r\le d}q_{\mathrm o,r}=k_{\mathrm o}.
\tag{12.3}
\]


\end{proposition}

\begin{proof}
Realizability puts \(\bar c\) in every restriction (5.1), because its
block empirical error is zero. Thus all classes used in the descent are
nonempty.

Suppose \(q_{\mathrm o,r}<k_{\mathrm o}\), and choose \(i_*\) attaining
\(M_{\mathrm o,r+1}\). Lemma~\ref{lem:step-015-source-map} gives
\[
H_{\mathrm o,i_*}^{r+1}\subseteq H_{\mathrm o,i}^{r},
\qquad
p_{\mathrm o,r+1}=2p_{\mathrm o,r}
\]
for every \(i\). Applying the decomposition bound (C2) yields
\[
M_{\mathrm o,r+1}
\le\operatorname{DDim}_{p_{\mathrm o,r},d}
(H_{\mathrm o,i}^r)
\le M_{\mathrm o,r}.
\tag{12.4}
\]
If equality held in (12.4), fix any optimal next-stage decomposition and
one of its maximum-dimensional leaves, and write \(f\) for that leaf's
SOA. The equality clause of (C2), applied separately to every \(i\) and
to an arbitrary optimal current decomposition, says that \(f\) occurs at a
maximum-dimensional leaf of every such current decomposition. Hence \(f\)
is \((p_{\mathrm o,r},d)\)-essential to every
\(H_{\mathrm o,i}^r\), so \(f\in\mathcal L_{\mathrm o,i}^r\) for all
\(i\). This gives \(q_{\mathrm o,r}=k_{\mathrm o}\), a contradiction.
DDim is integer-valued, proving the one-unit drop (12.2).

If \(M_{\mathrm o,r}=0\), every current restriction has DDim zero.
Lyu~\citep[Corollary~4.1, item~4]{lyu2025} says its essential list is
the entire finite restriction. The common target \(\bar c\) therefore lies
in every list, and \(q_{\mathrm o,r}=k_{\mathrm o}\).
Finally \(0\le M_{\mathrm o,0}\le d\). If every
\(q_{\mathrm o,r}\) were smaller than \(k_{\mathrm o}\), applying (12.2)
for \(r=0,\ldots,d-1\) would give \(M_{\mathrm o,d}=0\), which forces
\(q_{\mathrm o,d}=k_{\mathrm o}\), again a contradiction. This proves
(12.3) by a finite potential argument with no imported success event.
\(\square\)


\end{proof}

\begin{lemma}[Old AboveThreshold detects a sufficient stage]\label{lem:step-015-at}
Under Assumption~\ref{assump:approximate-dp-regime},
Propositions~\ref{prop:step-015-teacher} and
\ref{prop:step-015-descent}, and conditional on
\(E_{\mathrm{good,o}}\), let \(E_{\mathrm{AT,o}}\) be the event that all
\(d+1\) Laplace noises in the stopped old AboveThreshold transcript have
magnitude at most \(\tau_{\mathrm{AT,o}}\). Then
\[
\Pr(E_{\mathrm{AT,o}}^c)\le\beta_{\mathrm{AT,o}},
\tag{13.1}
\]
and on \(E_{\mathrm{AT,o}}\) the process selects a legal stage
\(r_{\mathrm o,*}\in\{0,\ldots,d\}\) such that
\[
q_{\mathrm o,r_{\mathrm o,*}}
\ge\frac{k_{\mathrm o}}2-\tau_{\mathrm{AT,o}}
\ge B_{\mathrm o}+\tau_{\mathrm{SS,o}}+2.
\tag{13.2}
\]
Every earlier reported Below additionally satisfies
\[
q_{\mathrm o,r}<\frac{k_{\mathrm o}}2+\tau_{\mathrm{AT,o}}
<k_{\mathrm o}.
\tag{13.3}
\]


\end{lemma}

\begin{proof}
For \(Z\sim\operatorname{Lap}(1/\eta_{\mathrm o})\),
\(\Pr(|Z|>u)=e^{-\eta_{\mathrm o}u}\). The definition (2.9) and a union
bound over at most \(d+1\) noises prove (13.1); independence of those
noises is more than is needed for this union bound.

On \(E_{\mathrm{AT,o}}\), a Below report means
\(q_{\mathrm o,r}+Z_r<k_{\mathrm o}/2\), giving the first inequality in
(13.3). The margin (4.6) implies
\(\tau_{\mathrm{AT,o}}<k_{\mathrm o}/2\), which gives the second.
Proposition~\ref{prop:step-015-descent} supplies a stage with score exactly
\(k_{\mathrm o}\). At that stage, even noise
\(-\tau_{\mathrm{AT,o}}\) leaves the noisy score above
\(k_{\mathrm o}/2\), so the stopped transcript reports Above no later than
that stage. At the selected stage, the Above inequality gives
\[
q_{\mathrm o,r_{\mathrm o,*}}
\ge k_{\mathrm o}/2-\tau_{\mathrm{AT,o}},
\]
and (4.6) gives (13.2). Thus no no-success or illegal-stage fallback is
reached on the stated event. \(\square\)


\end{proof}

\begin{proposition}[Actual old Sparse Sample output]\label{prop:step-015-sparse}
Under Lemmas~\ref{lem:step-015-lists} and
\ref{lem:step-015-at}, conditional on
\(E_{\mathrm{good,o}}\cap E_{\mathrm{AT,o}}\), let
\(E_{\mathrm{SS,o}}\) be the event that Sparse Sample at
\(r_{\mathrm o,*}\) returns an actual member of
\(\bigcup_i\mathcal L_{\mathrm o,i}^{r_{\mathrm o,*}}\), rather than
\(\perp\). Then
\[
\Pr(E_{\mathrm{SS,o}}^c
\mid\text{the selected old lists})
\le
\frac{\beta_{\mathrm{SS,o}}}{k_{\mathrm o}L_{\mathrm o}+1}
\le\beta_{\mathrm{SS,o}},
\tag{14.1}
\]
and, on \(E_{\mathrm{SS,o}}\),
\[
\bar H_{\mathrm o}\in
\mathcal L_{\mathrm o,i}^{r_{\mathrm o,*}}
\quad\text{for at least one }i.
\tag{14.2}
\]


\end{proposition}

\begin{proof}
The effective union has size at most \(k_{\mathrm o}L_{\mathrm o}\).
By (13.2), an item \(h_*\) has score
\(q_{\mathrm o,r_{\mathrm o,*}}\). In the exact Sparse Sample law, the
failure symbol has weight
\(e^{\varepsilon_{\mathrm{SS,o}}B_{\mathrm o}}\), while \(h_*\) has
weight
\(e^{\varepsilon_{\mathrm{SS,o}}q_{\mathrm o,r_{\mathrm o,*}}}\).
Discarding every other positive denominator term,
\[
\begin{aligned}
\Pr(\perp\mid\text{selected lists})
&\le
\exp[-\varepsilon_{\mathrm{SS,o}}
(q_{\mathrm o,r_{\mathrm o,*}}-B_{\mathrm o})]\\
&\le
\exp(-\varepsilon_{\mathrm{SS,o}}\tau_{\mathrm{SS,o}})\\
&=\frac{\beta_{\mathrm{SS,o}}}
{k_{\mathrm o}L_{\mathrm o}+1}.
\end{aligned}
\tag{14.3}
\]
Here the unused \(+2\) in (13.2) is harmless positive slack. Every
nonfailure point in the exact effective domain is an actual union item,
giving (14.2). This proves utility directly from the displayed weights and
does not rely on any ambiguity about the failure symbol.
\(\square\)


\end{proof}

\begin{lemma}[Old list locality and query sensitivity on all inputs]\label{lem:step-015-locality}
Under Proposition~\ref{prop:step-015-total}, fix any old partition and any
two quotient datasets differing in at most one record. At every stage, the
two totalized list tuples differ in at most one block coordinate; all lists
at every other coordinate are identical. Consequently every maximum-
occurrence query (5.2) has global sensitivity at most one, and, conditional
on every transcript-selected legal stage, the two Sparse Sample inputs
have one-list replacement adjacency. These statements hold without
realizability or \(E_{\mathrm{good,o}}\).


\end{lemma}

\begin{proof}
If the quotient datasets differ, the changed record lies in one fixed
partition block \(i_0\). Every restriction, validity check, fixed
decomposition choice, essential list ordering, sanitizer, and failure token
for a block is a deterministic function of that block's local labeled
state and public parameters. Thus all local states and lists for
\(i\ne i_0\) are identical at all stages; the \(i_0\)-list may be
replaced arbitrarily, including by or from the empty list.

Extend each list-occurrence count to all \(h\in H_C\), assigning zero when
an item is absent from every list. Replacement of one list changes every
such count by at most one. If \(s,s'\) are the two count functions, then
\[
\max_hs(h)\le\max_hs'(h)+1
\quad\text{and}\quad
\max_hs'(h)\le\max_hs(h)+1,
\]
including the all-empty convention. Hence the maximum query has
sensitivity one. An adaptively selected stage merely chooses one of these
already local tuples; for each fixed first-transcript value, its pair of
second-mechanism inputs still differs in at most list \(i_0\). No event or
stopping-stage equality across neighboring executions is assumed.
\(\square\)


\end{proof}

\begin{proposition}[All-input quotient and raw replacement privacy]\label{prop:step-015-dp}
Under Assumptions~\ref{assump:countable-evaluation-quotient} and
\ref{assump:approximate-dp-regime}, and under the conclusions of
Lemma~\ref{lem:step-015-locality} and
Propositions~\ref{prop:step-002-record-map},
\ref{prop:step-015-interfaces}, and
\ref{prop:step-015-kernel}, the released old quotient kernel is
\((\varepsilon/2,\delta)\)-DP. Its raw pullback
\(A_{\mathrm o,N_{\mathrm o}}\) is therefore
\((\varepsilon,\delta)\)-DP for every pair of raw replace-one labeled
inputs.


\end{proposition}

\begin{proof}
For the stopped AboveThreshold process, (C5), (2.4), and
\(\log(1/\delta_{\mathrm{AT,o}})=\log(2/\delta)\le g_\delta\) give
\[
c_{\mathrm{AT}}\eta_{\mathrm o}
[\sqrt{\log(1/\delta_{\mathrm{AT,o}})}
+\log(1/\delta_{\mathrm{AT,o}})]
\le\varepsilon/4.
\tag{16.1}
\]
Lemma~\ref{lem:step-015-locality} discharges sensitivity one, so this first
transcript is \((\varepsilon/4,\delta/2)\)-DP. At every transcript-selected
stage, the list cap (11.1), threshold (2.8), and one-list adjacency
from Lemma~\ref{lem:step-015-locality} discharge every premise of (C4).
Thus Sparse Sample is
\[
(2\varepsilon_{\mathrm{SS,o}},\delta_{\mathrm{SS,o}})
=(\varepsilon/4,\delta/2)\text{-DP}.
\tag{16.2}
\]
The finite-transcript composition and postprocessing interface proved in
Proposition~\ref{prop:step-015-interfaces} therefore gives
\[
(\varepsilon/4,\delta/2)+(\varepsilon/4,\delta/2)
=(\varepsilon/2,\delta)\preceq(\varepsilon,\delta).
\tag{16.3}
\]
Mixing over the same data-independent partition preserves (16.3).

This argument applies to the total transcript on every labeled input.
In particular, empty restrictions, sanitized invalid lists, different
stopping stages, \(\perp\), no-success, exhaustion, and fallback are
already covered by deterministic postprocessing of the private transcript;
neither realizability nor a utility event was used.

Finally, Proposition~\ref{prop:step-002-record-map} maps raw
neighbors to equal or quotient replace-one inputs. Applying (16.3) to
\(K_{\mathrm o}(T_{N_{\mathrm o}}(\cdot),E)\) proves the claimed raw
\((\varepsilon,\delta)\)-DP inequality for every
\(E\in\mathcal H_C\). \(\square\)


\end{proof}

\begin{proposition}[Exact old SOA identity, fixed family, and empirical certificate]\label{prop:step-015-soa}
Under Assumptions~\ref{assump:finite-littlestone} and
\ref{assump:realizable-iid} and
Proposition~\ref{prop:step-015-interfaces}, on
\[
\mathsf{Core}_{\mathrm o}:=
E_{\mathrm{good,o}}\cap E_{\mathrm{AT,o}}\cap E_{\mathrm{SS,o}},
\tag{17.1}
\]
and under Lemma~\ref{lem:step-015-lists} and
Proposition~\ref{prop:step-015-sparse}, the actual old output has a
maximum-dimensional leaf witness \(\mathcal G\subseteq\bar C\) satisfying
\[
\bar H_{\mathrm o}=\operatorname{SOA}_{\mathcal G},\qquad
\mathcal G\text{ is both }n_{\mathrm o}\text{- and }
(d+1)\text{-irreducible},
\tag{17.2}
\]
\[
\bar H_{\mathrm o}\in\widehat{\bar C}_{d+1}\subseteq H_C,\qquad
\operatorname{LD}(\widehat{\bar C}_{d+1})\le d,
\tag{17.3}
\]
and
\[
e_{\bar S_{\mathrm o}}(\bar H_{\mathrm o})
\le2\gamma_{\mathrm o}=\alpha/8.
\tag{17.4}
\]


\end{proposition}

\begin{proof}
By Proposition~\ref{prop:step-015-sparse}, the actual output belongs to
some selected list. Lemma~\ref{lem:step-015-lists} and the definition of
essentiality produce a maximum-dimensional leaf \(\mathcal G\) with the
exact identity and irreducibility in (17.2). The fixed-family
result (C3) then gives (17.3). This is an improper family in general; no
claim that \(\bar H_{\mathrm o}\in\bar C\) is made.

It remains to prove the empirical certificate. Every
\(g\in\mathcal G\subseteq H_{\mathrm o,i}^{r_{\mathrm o,*}}\) has producer
block error at most
\(\rho^{r_{\mathrm o,*}+1}\gamma_{\mathrm o}\le\gamma_{\mathrm o}\).
Write \(\mu=e_{\bar S_{\mathrm o}}(g)\). If
\(\mu\le\gamma_{\mathrm o}/3\), then already \(\mu<2\gamma_{\mathrm o}\).
If \(\mu>\gamma_{\mathrm o}/3\), the lower relative clause of
\(E_{\mathrm{good,o}}\) gives
\[
\mu\le\frac{e_{\bar S_{\mathrm o,i}}(g)}{1-\xi_d}
\le\frac{\gamma_{\mathrm o}}{1-\xi_d}
\le\frac54\gamma_{\mathrm o}<2\gamma_{\mathrm o}.
\tag{17.5}
\]
Thus every \(g\in\mathcal G\) has master error at most
\(2\gamma_{\mathrm o}\).

Suppose instead that
\(e_{\bar S_{\mathrm o}}(\operatorname{SOA}_{\mathcal G})
>2\gamma_{\mathrm o}\). Any \(g\in\mathcal G\) agreeing with this SOA on
all indexed master points would have exactly the same prediction vector
and hence the same master error, contradicting (17.5). Therefore
\[
\mathcal G|_{(q_1,\operatorname{SOA}_{\mathcal G}(q_1)),\ldots,
(q_{n_{\mathrm o}},
\operatorname{SOA}_{\mathcal G}(q_{n_{\mathrm o}})))}
=\varnothing.
\tag{17.6}
\]
But \(n_{\mathrm o}\)-irreducibility in (17.2) says this restriction
preserves \(\operatorname{LD}(\mathcal G)\), whereas the empty class has
strictly smaller dimension. This contradiction proves (17.4). Equation
(17.6) is the valid contradiction pattern in Lyu~\citep{lyu2025} Theorem 3's proof; the
malformed displayed output and unsupported probability-one sentence are
not used. \(\square\)


\end{proof}

\begin{proposition}[Old occurrence-mark kernel and exact projection]\label{prop:step-015-mark}
Under Assumption~\ref{assump:countable-evaluation-quotient},
Lemmas~\ref{lem:step-003-countable-promotion} and
\ref{lem:step-004-occurrence}, and
Propositions~\ref{prop:step-015-total},
\ref{prop:step-015-kernel},
\ref{prop:step-015-sparse},
\ref{prop:step-015-interfaces},
\ref{prop:step-004-lift}, and
\ref{prop:step-004-projection}, the old internal-state/output law admits a
Markov-kernel lift
\[
\widetilde K_{\mathrm o}:
Z_Q^{N_{\mathrm o}}\leadsto
H_C\times\{0,1,\ldots,k_{\mathrm o}\}.
\tag{18.1}
\]
For every quotient input \(\bar s\) and \(E\in\mathcal H_C\),
\[
\sum_{i=0}^{k_{\mathrm o}}
\widetilde K_{\mathrm o}(\bar s,E\times\{i\})
=K_{\mathrm o}(\bar s,E).
\tag{18.2}
\]
On actual paths, positive marks are uniform over all distinct producer
blocks whose all-stage list union contains the output; every nonactual
path has mark \(0\).


\end{proposition}

\begin{proof}
Let \(R_{\mathrm o,\bar s}(d\omega,dh)\) be the joint law of the complete
old transcript state and terminal \(H_C\)-output. For each block define
\[
\mathcal G_{\mathrm o,i}(\omega):=
\bigcup_{r=0}^d\mathcal L_{\mathrm o,i}^r(\omega).
\tag{18.3}
\]
On the measurable actual-output status event, set
\[
I_{\mathrm o}(\omega,h):=
\{i\in[k_{\mathrm o}]:h\in\mathcal G_{\mathrm o,i}(\omega)\}.
\tag{18.4}
\]
The exact Sparse Sample support in Proposition
\ref{prop:step-015-sparse} makes this set nonempty on every actual path.
Define
\[
w_i(\omega,h):=
\begin{cases}
\mathbf1\{i\in I_{\mathrm o}(\omega,h)\}/|I_{\mathrm o}(\omega,h)|,
&\text{on an actual path},\\
0,&\text{on a nonactual path},
\end{cases}
\quad i\ge1,
\tag{18.5}
\]
and let \(w_0=0\) on actual paths and \(w_0=1\) on nonactual paths.
Finite-list membership, actual status, the finite sum
\(|I_{\mathrm o}|\), and every \(w_i\) are measurable by the
countable-input coding in Lemma~\ref{lem:step-003-countable-promotion} and
the finite-occurrence construction of Lemma~\ref{lem:step-004-occurrence}
and Proposition~\ref{prop:step-004-lift}.
Repeated occurrences within a list or across stages do not duplicate a
block coordinate in (18.4).

For \(i=0,\ldots,k_{\mathrm o}\), define
\[
\widetilde K_{\mathrm o}(\bar s,E\times\{i\})
:=\int\mathbf1_E(h)w_i(\omega,h)
R_{\mathrm o,\bar s}(d\omega,dh),
\tag{18.6}
\]
and extend by finite additivity over mark sections. The weights are
nonnegative and sum pointwise to one, so (18.6) is a probability law.
The quotient input is countable discrete, and the output product is
standard Borel with a finite mark, so
Lemma~\ref{lem:step-003-countable-promotion} gives the kernel (18.1).
Summing (18.6) over marks and using \(\sum_iw_i=1\) is the exact projection
argument of Proposition~\ref{prop:step-004-projection} and proves (18.2).
Thus the mark is analysis-only and changes
neither the released marginal nor the privacy proved in
Proposition~\ref{prop:step-015-dp}. \(\square\)


\end{proof}

\begin{lemma}[Producer-local old reconstruction and pathwise inclusion]\label{lem:step-015-reconstruct}
Under Assumption~\ref{assump:realizable-iid},
Propositions~\ref{prop:step-015-soa} and
\ref{prop:step-015-mark}, fix the data-independent old partition and a
positive mark \(i\). Define
\[
\mathcal G_{\mathrm o,i}(\bar S_{\mathrm o,i})
:=\bigcup_{r=0}^d\mathcal L_{\mathrm o,i}^r,\qquad
|\mathcal G_{\mathrm o,i}|\le(d+1)L_{\mathrm o}.
\tag{19.1}
\]
This finite list is a measurable function only of the \(m_{\mathrm o}\)
records in producer block \(i\). On
\(\mathsf{Core}_{\mathrm o}\cap\{J_{\mathrm o}=i\}\),
\[
\bar H_{\mathrm o}\in\mathcal G_{\mathrm o,i},
\qquad
\widehat{\operatorname{err}}_{\mathrm o,-i}(\bar H_{\mathrm o})
\le\frac{k_{\mathrm o}\alpha}{8(k_{\mathrm o}-1)}
\le\frac\alpha4.
\tag{19.2}
\]
Consequently the following inclusion holds pathwise, before conditioning
on either the mark or block data:
\[
\begin{aligned}
&\{\mathsf{Core}_{\mathrm o},J_{\mathrm o}=i,
\operatorname{err}_{\bar D}(\bar H_{\mathrm o},\bar c)>\alpha\}\\
&\quad\subseteq
\bigcup_{\substack{h\in\mathcal G_{\mathrm o,i}(\bar S_{\mathrm o,i}):\\
\operatorname{err}_{\bar D}(h,\bar c)>\alpha}}
\{\widehat{\operatorname{err}}_{\mathrm o,-i}(h)\le\alpha/4\}.
\end{aligned}
\tag{19.3}
\]


\end{lemma}

\begin{proof}
Every restriction, decomposition choice, and essential list for coordinate
\(i\) is block-local by the construction in Proposition
\ref{prop:step-015-total}; concatenating the \(d+1\) measurable finite
lists proves the locality and measurability in (19.1). The size bound is
the sum of the \(d+1\) caps (11.1); duplicate functions only reduce it.

On an actual path with \(J_{\mathrm o}=i\), the mark definition (18.4)
puts \(\bar H_{\mathrm o}\) in this block's union. On the core,
Proposition~\ref{prop:step-015-soa} bounds the number of mistakes on the
whole master sample by \(n_{\mathrm o}\alpha/8\). The complement has
\(n_{\mathrm o}-m_{\mathrm o}=(k_{\mathrm o}-1)m_{\mathrm o}\) records,
so its number of mistakes is no larger than the full-sample number and
\[
\widehat{\operatorname{err}}_{\mathrm o,-i}(\bar H_{\mathrm o})
\le\frac{n_{\mathrm o}}{n_{\mathrm o}-m_{\mathrm o}}\frac\alpha8
=\frac{k_{\mathrm o}\alpha}{8(k_{\mathrm o}-1)}
\le\frac\alpha4,
\]
where \(k_{\mathrm o}\ge2\). If the left event of (19.3) occurs, its
realized \(\bar H_{\mathrm o}\) is therefore one of the finitely many
candidates indexing the right union and satisfies both defining
conditions. This proves the inclusion without conditioning on an adaptive
output. \(\square\)


\end{proof}

\begin{lemma}[Fixed-candidate complement lower tail]\label{lem:step-015-lower-tail}
Under Assumption~\ref{assump:realizable-iid}, condition on the fixed
partition, a producer block \(\bar S_{\mathrm o,i}\), and any fixed
\(h\in\mathcal G_{\mathrm o,i}(\bar S_{\mathrm o,i})\). If
\(p=\operatorname{err}_{\bar D}(h,\bar c)>\alpha\), then the independent
complement of \(M=(k_{\mathrm o}-1)m_{\mathrm o}\) records satisfies
\[
\Pr[
\widehat{\operatorname{err}}_{\mathrm o,-i}(h)\le\alpha/4
\mid\bar S_{\mathrm o,i}]
\le\exp(-9\alpha M/32).
\tag{20.1}
\]


\end{lemma}

\begin{proof}
After conditioning on the producer block, \(h\) is fixed because of the
locality in Lemma~\ref{lem:step-015-reconstruct}. The complement records
remain iid from \(P_{\bar D,\bar c}\) and independent of that block.
Their mistake indicators are therefore independent Bernoulli variables of
mean \(p\). Since \(p>\alpha\),
\[
\{\widehat{\operatorname{err}}_{\mathrm o,-i}(h)\le\alpha/4\}
\subseteq
\{\widehat{\operatorname{err}}_{\mathrm o,-i}(h)\le p/4\}.
\]
We now derive (C7). Put \(S=\sum_{j=1}^M X_j\), where the \(X_j\)'s are
these mistake indicators. Here
\(M=(k_{\mathrm o}-1)m_{\mathrm o}\ge1\) and \(p>\alpha>0\). For
completeness, if \(M=0\) or \(p=0\), the event in (C7) has probability one
and its right side equals one. Otherwise fix \(0<\theta<1\), set
\(q=(1-\theta)p\), and take
\(\lambda=-\log(1-\theta)>0\). Exponential Markov applied to
\(e^{-\lambda S}\) gives
\[
\begin{aligned}
\Pr(S\le Mq)
&=\Pr(e^{-\lambda S}\ge e^{-\lambda Mq})\\
&\le e^{\lambda Mq}\,\mathbb E e^{-\lambda S}\\
&=\exp\!\left(M[\lambda q+\log(1-p\theta)]\right)\\
&\le\exp\!\left(M[\lambda(1-\theta)p-p\theta]\right)\\
&=\exp\!\left(-pM[\theta+(1-\theta)\log(1-\theta)]\right),
\end{aligned}
\tag{20.2}
\]
where the penultimate line uses \(\log(1-u)\le-u\) for \(0\le u<1\).
To finish the scalar estimate, define
\[
g(\theta):=\theta+(1-\theta)\log(1-\theta)-\frac{\theta^2}{2}.
\]
Then \(g(0)=0\) and
\[
g'(\theta)=-\log(1-\theta)-\theta\ge0,
\tag{20.3}
\]
because \(u(\theta):=-\log(1-\theta)-\theta\) satisfies
\(u(0)=0\) and \(u'(\theta)=\theta/(1-\theta)\ge0\). Hence
\(g(\theta)\ge0\), and (20.2) proves the locally derived inequality
\[
 \Pr[S\le M(1-\theta)p]
 \le\exp(-\theta^2pM/2).
\tag{C7}
\]
This does not invoke an external lower-tail theorem. Taking
\(\theta=3/4\), then using \(p>\alpha\), yields
\[
\exp[-(3/4)^2pM/2]
=\exp(-9pM/32)
\le\exp(-9\alpha M/32),
\]
which is (20.1). \(\square\)


\end{proof}

\begin{proposition}[Independent old-arm PAC ledger]\label{prop:step-015-pac}
Under Assumptions~\ref{assump:finite-littlestone},
\ref{assump:countable-evaluation-quotient},
\ref{assump:realizable-iid}, and
\ref{assump:approximate-dp-regime}, Propositions
\ref{prop:step-015-good}, \ref{prop:step-015-sparse},
\ref{prop:step-015-mark}, and \ref{prop:step-015-soa}, and
Lemmas~\ref{lem:step-015-at}, \ref{lem:step-015-reconstruct}, and
\ref{lem:step-015-lower-tail}, the released old raw learner satisfies
\[
\sup_D\sup_{c\in C}
\Pr[
\operatorname{err}_D(\operatorname{Dec}_C(\bar H_{\mathrm o}),c)
>\alpha]
\le\beta.
\tag{21.1}
\]


\end{proposition}

\begin{proof}
We first discharge the finite multiplicity with the same old block
constant already used for the trace event. From (9.5),
\[
n_{\mathrm o}\le(2C_{\mathrm o}/e)dQ_{\mathrm o}\log Q_{\mathrm o}.
\tag{21.2}
\]
Hence, with
\(C_p:=5+\log(2C_{\mathrm o})\),
\[
\begin{aligned}
\log p_{\mathrm o,d}
&=d\log2+\log n_{\mathrm o}+\log d\\
&\le C_p d\log Q_{\mathrm o},
\end{aligned}
\tag{21.3}
\]
where \(Q_{\mathrm o}\ge e\), \(\log d\le d\), and
\(\log\log Q_{\mathrm o}\le\log Q_{\mathrm o}\) were used explicitly.
Therefore
\[
\begin{aligned}
\log\frac{k_{\mathrm o}(d+1)L_{\mathrm o}}
{\beta_{\mathrm{gen,o}}}
&=\log\frac{4k_{\mathrm o}(d+1)}\beta
+d\log p_{\mathrm o,d}+d^2\log2\\
&\le C_6d\,a_{\mathrm o}\log Q_{\mathrm o},
\qquad C_6:=C_p+3.
\end{aligned}
\tag{21.4}
\]
Indeed \(a_{\mathrm o}\ge d\) and
\(a_{\mathrm o}\ge\log(4k_{\mathrm o}/\beta)\), while
\(\log(d+1)\le d\).

The unrounded lower bound in (2.6) gives
\[
\frac9{32}\alpha(k_{\mathrm o}-1)m_{\mathrm o}
\ge\frac{9C_{\mathrm o}}{32}
(k_{\mathrm o}-1)d^2a_{\mathrm o}\log Q_{\mathrm o}.
\tag{21.5}
\]
For \(C_{\mathrm o}=2^{20}\), \(k_{\mathrm o}\ge2\), and \(d\ge1\),
the right side of (21.5) is at least the right side of (21.4).
Consequently
\[
k_{\mathrm o}(d+1)L_{\mathrm o}
\exp[-9\alpha(k_{\mathrm o}-1)m_{\mathrm o}/32]
\le\beta_{\mathrm{gen,o}}.
\tag{21.6}
\]

Now work under the marked lift. Sum the pathwise inclusion (19.3) over
\(i=1,\ldots,k_{\mathrm o}\). Only after this finite inclusion is in
place, condition on each producer block and apply
Lemma~\ref{lem:step-015-lower-tail} to each of its at most
\((d+1)L_{\mathrm o}\) candidates. Equation (21.6) gives
\[
\Pr[
\mathsf{Core}_{\mathrm o},\
\operatorname{err}_{\bar D}(\bar H_{\mathrm o},\bar c)>\alpha]
\le\beta_{\mathrm{gen,o}}.
\tag{21.7}
\]
No uncountable supremum over \(H_C\) has been taken.

The other failures are charged independently:
\[
\Pr(E_{\mathrm{good,o}}^c)\le\beta_{\mathrm{tr,o}},
\quad
\Pr(E_{\mathrm{AT,o}}^c)\le\beta_{\mathrm{AT,o}},
\quad
\Pr(E_{\mathrm{good,o}}\cap E_{\mathrm{AT,o}}
\cap E_{\mathrm{SS,o}}^c)\le\beta_{\mathrm{SS,o}}.
\tag{21.8}
\]
The last inequality follows by integrating the conditional bound (14.1)
over the selected lists. Equations (21.7)-(21.8) and the exact confidence
split (2.3) yield
\[
\Pr[
\operatorname{err}_{\bar D}(\bar H_{\mathrm o},\bar c)>\alpha]
\le\beta_{\mathrm{tr,o}}+\beta_{\mathrm{AT,o}}
+\beta_{\mathrm{SS,o}}+\beta_{\mathrm{gen,o}}
=\beta.
\tag{21.9}
\]
Projection (18.2) removes the analysis-only mark without changing the
released law. Proposition~\ref{prop:step-002-iid-pushforward}
identifies the quotient sampling law, and Proposition~\ref{prop:step-002-risk} turns (21.9) into (21.1), uniformly
over arbitrary \(D\) and \(c\). \(\square\)


\end{proof}

\begin{proposition}[Ceiling-aware explicit \(R_{\mathrm{old}}\) rate]\label{prop:step-015-rate}
Under Assumptions~\ref{assump:finite-littlestone} and
\ref{assump:approximate-dp-regime}, and
Proposition~\ref{prop:step-015-teacher}, there is a universal
\(K_{\mathrm O}\ge1\) such that, for \(d\ge1\),
\[
N_{\mathrm{old}}
\le K_{\mathrm O}\Lambda^6
R_{\mathrm{old}}(d,\alpha,\beta,\varepsilon,\delta).
\tag{22.1}
\]
No positive power of \(d\) is hidden in \(K_{\mathrm O}\Lambda^6\).


\end{proposition}

\begin{proof}
The candidate functions in (2.5)-(2.9) are nondecreasing in \(t\).
Thus \(k_{\mathrm o}\le\bar k_{\mathrm o}\), the envelope derivation
(3.8)-(3.10), and (4.4) give universal constants \(C_7,C_8,C_9,C_{10}\)
such that
\[
\begin{gathered}
k_{\mathrm o}\le
C_7\frac{d^2\log(64/(\delta\beta))\Lambda^2}{\varepsilon},
\qquad
a_{\mathrm o}\le C_8d\Lambda^2,\\
\log Q_{\mathrm o}\le C_9\Lambda^2,\qquad
m_{\mathrm o}\le
C_{10}\frac{d^3}{\alpha}\Lambda^4.
\end{gathered}
\tag{22.2}
\]
The last inequality retains the ceiling: its unrounded term exceeds one,
so \(\lceil x\rceil\le2x\), exactly as in (3.10).
Multiplication of the actual counts in (4.5) gives
\[
N_{\mathrm{old}}=k_{\mathrm o}m_{\mathrm o}
\le C_{11}\Lambda^6
\frac{d^5\log(64/(\delta\beta))}
{\varepsilon\alpha}.
\tag{22.3}
\]
Since \(\delta\beta<1/4\),
\[
\log(64/(\delta\beta))
=\log64+\log(1/(\delta\beta))
\le4\log(1/(\delta\beta)).
\tag{22.4}
\]
Substitution into (22.3), followed by the nonnegative second summand in
the setting definition
\[
R_{\mathrm{old}}
=\frac{d^5\log(1/(\delta\beta))}{\varepsilon\alpha}
+\frac{d+\log(1/\beta)}{\alpha},
\]
proves (22.1) with \(K_{\mathrm O}=4C_{11}\).
All appearances of \(d\) outside \(R_{\mathrm{old}}\) in this derivation
are inside logarithms already bounded by \(\Lambda^6\); the exposed
polynomial is exactly \(d^5\). \(\square\)


\end{proof}

\begin{proposition}[Old-arm boundaries and fixed versus scheduled \(\delta\)]\label{prop:step-015-boundaries}
Under Propositions~\ref{prop:step-015-zero},
\ref{prop:step-015-dp}, \ref{prop:step-015-pac}, and
\ref{prop:step-015-rate}, the complete old arm has the following boundary
properties.

1. At \(d=0\), it uses \(N_{\mathrm{old}}=0\), deterministic mark \(0\),
   zero risk, and \((0,0)\)-privacy.
2. At \(d=1\), all positive-branch scales, stages, restrictions, and
   inequalities are legal; in particular \(\rho=1/2\),
   \(\rho^d=1/2\), and \(p_{\mathrm o,0}\ge2=d+1\).
3. The construction applies unchanged to finite and infinite \(C\), and to
   \(v=d\); it never introduces \(|C|\), and it makes no VC-sensitive
   improvement claim.
4. Empty, invalid, failure, and fallback paths remain measurable and
   private on all inputs, while the old core excludes them for utility.
   Actual outputs may be improper elements of \(H_C\).
5. Every integer ceiling and all four confidence shares are retained.
6. For every fixed allowed \(0<\delta<1\), the kernel, privacy, PAC, and
   rate conclusions remain valid. Separately, along an asymptotic sequence,
   \[
   \delta K_{\mathrm O}\Lambda^6R_{\mathrm{old}}\longrightarrow0
   \quad\Longrightarrow\quad
   N_{\mathrm{old}}\delta\longrightarrow0.
   \tag{23.1}
   \]
   The setting's VC-arm schedule is not silently asserted to imply
   (23.1).


\end{proposition}

\begin{proof}
Item 1 is Proposition~\ref{prop:step-015-zero}; the mark-\(0\) convention
is the unique inactive occurrence convention. For item 2, substitute
\(d=1\) in (2.2), (4.6), and (10.4); no denominator or stage disappears.
Items 3 and 4 follow because the only class-size control is the finite
trace restriction (7.2) and finite source list cap (11.1), both valid for
infinite classes, while Propositions~\ref{prop:step-015-total} and
\ref{prop:step-015-dp} cover every totalized path. At \(v=d\), the old
dictionary is literally unchanged and (22.1) is the \(d^5\) source
baseline. Proposition~\ref{prop:step-015-soa} explicitly permits improper
SOA outputs.

Item 5 follows from (2.6), (2.8), (4.1), (4.5), and the exact identity
\(4(\beta/4)=\beta\). None of Propositions
\ref{prop:step-015-kernel}, \ref{prop:step-015-dp},
\ref{prop:step-015-pac}, or \ref{prop:step-015-rate} assumes a limiting
relation involving \(\delta\), proving fixed-\(\delta\) validity. Finally,
multiplying (22.1) by \(\delta\) proves (23.1). When \(v\ll d\),
\(R_{\mathrm{old}}\) can be polynomially larger than \(R_{\mathrm{VC}}\),
so the distinct VC-arm schedule in the setting cannot in general replace
the explicit old-arm premise in (23.1). 
Proposition~\ref{prop:step-015-zero} closes the exact \(d=0\) branch. On
\(d\ge1\), Lemma~\ref{lem:step-015-dictionary},
Lemma~\ref{lem:step-015-envelope}, and
Proposition~\ref{prop:step-015-teacher} define the independent old tuple,
prove a finite witness before minimization, retain all ceilings, and export
the source margin. Proposition~\ref{prop:step-015-total} then defines every
valid and fallback path, and Proposition
\ref{prop:step-015-kernel} turns that pointwise law into the required
quotient and raw Markov kernels with exact decoded risk.

The old utility proof is independent of the VC-sensitive arm:
Lemma~\ref{lem:step-015-traces} uses only the coarse exponent \(d\);
Lemma~\ref{lem:step-015-tails} proves both fixed-trace branches; and
Proposition~\ref{prop:step-015-good} closes the old fixed point and charges
its own trace share. Lemma~\ref{lem:step-015-source-map} supplies the exact
source endpoint and half-scale map. Lemma~\ref{lem:step-015-lists} and
Proposition~\ref{prop:step-015-descent} discharge all decomposition,
essentiality, irreducibility, and finite-potential obligations.
Lemma~\ref{lem:step-015-at} and
Proposition~\ref{prop:step-015-sparse} then produce an actual old list
output with their own two confidence charges.

Privacy is event-free: Lemma~\ref{lem:step-015-locality} proves
sensitivity one and one-list adjacency on arbitrary inputs, and
Proposition~\ref{prop:step-015-dp} composes the exact old allocations and
transfers them to raw neighbors. Utility continues through
Proposition~\ref{prop:step-015-soa}, which proves the actual SOA identity,
fixed-family membership, and empirical error. Proposition
\ref{prop:step-015-mark} constructs the old marked lift and exact
projection. Lemmas~\ref{lem:step-015-reconstruct} and
\ref{lem:step-015-lower-tail} reduce the adaptive population-error event to
a finite producer-local union, and Proposition
\ref{prop:step-015-pac} absorbs its multiplicity and sums exactly the four
old confidence shares.

Finally, Proposition~\ref{prop:step-015-rate} eliminates every old
auxiliary and proves (22.1), with the explicit choice
\(q_{\mathrm{old}}=6\). The boundary calculations at the start of this
proof verify the null, first-stage,
finite/infinite-class, \(v=d\), improper-output, fallback, ceiling,
confidence, and \(\delta\)-schedule boundaries. Combining these results
gives the old-arm interface
\[
\boxed{
(K_C^{\mathrm{old\text{-}Lyu}},
A_{\mathrm o,N_{\mathrm o}},
\widetilde K_{\mathrm o},
N_{\mathrm{old}})
\text{ is measurable, raw }(\varepsilon,\delta)\text{-DP, }
(\alpha,\beta)\text{-PAC, and }
N_{\mathrm{old}}\le K_{\mathrm O}\Lambda^6R_{\mathrm{old}}.}
\]



\end{proof}

\subsection{Finite Arm and Baseline Comparison}\label{app:step_016}


 \begin{proposition}[Shared exact zero-sample learner]\label{prop:step-016-zero}
Under Assumptions~\ref{assump:finite-littlestone} and
\ref{assump:countable-evaluation-quotient}, Lemma~\ref{lem:step-002-ld}, 
Propositions~\ref{prop:step-002-factorization},
\ref{prop:step-002-risk}, and
\ref{prop:step-014-vc-arm}, and Proposition~\ref{prop:step-015-zero}, if \(d=0\), then \(v=0\),
\(C\) and \(\bar C\) are singletons, and the finite, old, and VC arms all
use the same \(N=0\) Dirac quotient law at the unique
\(\bar c_0\in\bar C\). Its raw pullback along \(T_0\) is the same Dirac
kernel, is \((0,0)\)-DP, and has decoded risk zero for every \(D,c\).
Consequently \(m_C(\alpha,\beta;\varepsilon,\delta)=0\).


\end{proposition}

\begin{proof}
Lemma~\ref{lem:step-002-ld} proves that a nonempty
Littlestone-dimension-zero class cannot contain two distinct functions: two
functions differing at one point would shatter a depth-one tree. Its stated
conclusion also gives \(v=0\). Proposition~\ref{prop:step-002-factorization} gives the corresponding
singleton conclusion for \(\bar C\). Thus there is only one possible quotient
concept \(\bar c_0\).

Define the finite-arm law on the one-point empty-input space by
\[
K^{\mathrm{fin}}_{C,0}(\varnothing,E)=\mathbf 1_E(\bar c_0).
\tag{1.1}
\]
This is a Markov kernel, and its raw pullback along the unique \(T_0\) is
identical. Since the kernel is data independent, it is \((0,0)\)-DP.
 Propositions~\ref{prop:step-002-factorization} and
\ref{prop:step-002-risk} show that its decoded output is the unique raw target
and hence has risk zero for every \(D\). 
Propositions~\ref{prop:step-014-vc-arm} and
\ref{prop:step-015-zero} define their null branches by this same unique
quotient output. Therefore the three laws are literally identical, not merely
equivalent in risk. The definition of \(m_C\) then gives \(m_C=0\). No
positive-dimensional denominator, ceiling, or limiting argument is used.
\(\square\)


\end{proof}

\begin{proposition}[Measurable finite quotient exponential-weight kernel]\label{prop:step-016-finite-kernel}
Under Assumptions~\ref{assump:countable-evaluation-quotient} and
\ref{assump:approximate-dp-regime}, 
Propositions~\ref{prop:step-002-factorization},
\ref{prop:step-002-record-map}, and
\ref{prop:step-003-quotient-kernel}, and
\ref{prop:step-003-raw-pullback}, and Lemma~\ref{lem:step-003-countable-promotion}, if
\(M:=|C|<\infty\) and \(N\in\mathbb N_0\), then empirical-mistake
exponential weights on the fixed finite class \(\bar C\) define an
everywhere-total Markov kernel
\(K^{\mathrm{fin}}_{C,N}:Z_Q^N\leadsto H_C\). Its raw pullback
\[
A_N^{\mathrm{fin}}(s,E)
:=K^{\mathrm{fin}}_{C,N}(T_N(s),E)
\tag{2.1}
\]
is a Markov kernel \(Z_X^N\leadsto H_C\). Both laws release in the setting's
output space and use \(\operatorname{Dec}_C\) only after release. The law is
defined on every labeled input, including nonrealizable inputs.


\end{proposition}

\begin{proof}
Proposition~\ref{prop:step-002-factorization} is a
bijection, so \(|\bar C|=M\ge1\). For a quotient sample
\(\bar s=((q_r,y_r))_{r=1}^N\) and \(\bar h\in\bar C\), define the
empirical mistake count
\[
L_{\bar s}(\bar h)
:=\sum_{r=1}^N\mathbf 1\{\bar h(q_r)\ne y_r\}.
\tag{2.2}
\]
This definition uses arbitrary labels and has \(L_\varnothing=0\). Put
\[
w_{\bar s}(\bar h):=
\exp\!\left(-\frac{\varepsilon}{2}L_{\bar s}(\bar h)\right),
\qquad
Z_{\bar s}:=\sum_{\bar g\in\bar C}w_{\bar s}(\bar g),
\tag{2.3}
\]
and, for \(E\in\mathcal H_C\),
\[
K^{\mathrm{fin}}_{C,N}(\bar s,E)
:=\frac{\displaystyle
\sum_{\bar h\in\bar C}\mathbf1_E(\bar h)w_{\bar s}(\bar h)}
{Z_{\bar s}}.
\tag{2.4}
\]
Every weight is strictly positive, and the fixed sum has exactly \(M\)
terms. Thus \(0<Z_{\bar s}<\infty\), so (2.4) is an everywhere-defined
Borel probability measure supported on the fixed set \(\bar C\subseteq H_C\).
No empirical subclass or data-dependent range is used.

The domain \(Z_Q^N\) is countable discrete. Lemma~\ref{lem:step-003-countable-promotion} makes every event-mass coordinate
of the pointwise family (2.4) measurable on its full power-set sigma-field.
Each support singleton is Borel because \(H_C\) is standard Borel. Thus
(2.4) is precisely a separately supplied, everywhere-defined pointwise law
to which Proposition~\ref{prop:step-003-quotient-kernel} applies;
in particular, \(K^{\mathrm{fin}}_{C,N}\) is a quotient kernel produced by
that proposition. Proposition~\ref{prop:step-002-record-map} makes
\(T_N\) measurable, so the exact premise of Proposition~\ref{prop:step-003-raw-pullback} is met and that proposition gives
(2.1). The codomain remains exactly \(H_C\); decoding is the unchanged
setting map. \(\square\)


\end{proof}

\begin{proposition}[Pure all-input privacy of finite private ERM]\label{prop:step-016-finite-dp}
Under Assumptions~\ref{assump:countable-evaluation-quotient} and
\ref{assump:approximate-dp-regime}, Proposition~\ref{prop:step-002-record-map}, and
Proposition~\ref{prop:step-016-finite-kernel}, if \(M<\infty\), then for
every \(N\in\mathbb N_0\), the quotient kernel
\(K^{\mathrm{fin}}_{C,N}\) and its raw pullback
\(A_N^{\mathrm{fin}}\) are pure \(\varepsilon\)-DP on all replace-one
labeled inputs. Hence they are \((\varepsilon,\delta)\)-DP for every allowed
\(0<\delta<1\).


\end{proposition}

\begin{proof}
Let \(\bar s\sim\bar s'\) be quotient replace-one samples. For
every \(\bar h\in\bar C\), only one summand in (2.2) can change, whence
\[
|L_{\bar s}(\bar h)-L_{\bar s'}(\bar h)|\le1.
\tag{3.1}
\]
Thus the score \(-L_{\bar s}(\bar h)\) has replacement sensitivity one, and
(2.3) gives the two pointwise inequalities
\[
e^{-\varepsilon/2}w_{\bar s'}(\bar h)
\le w_{\bar s}(\bar h)
\le e^{\varepsilon/2}w_{\bar s'}(\bar h).
\tag{3.2}
\]
For an arbitrary \(E\in\mathcal H_C\), denote the numerator in (2.4) by
\(W_{\bar s}(E)\). Summing the upper inequality in (3.2) over
\(\bar C\cap E\), and the lower inequality over all of \(\bar C\), yields
\[
W_{\bar s}(E)\le e^{\varepsilon/2}W_{\bar s'}(E),
\qquad
Z_{\bar s}\ge e^{-\varepsilon/2}Z_{\bar s'}.
\tag{3.3}
\]
Consequently
\[
K^{\mathrm{fin}}_{C,N}(\bar s,E)
=\frac{W_{\bar s}(E)}{Z_{\bar s}}
\le e^\varepsilon
\frac{W_{\bar s'}(E)}{Z_{\bar s'}}
=e^\varepsilon K^{\mathrm{fin}}_{C,N}(\bar s',E).
\tag{3.4}
\]
This is the complete exponential-mechanism calculation; no privacy theorem
is cited. It uses neither realizability nor a utility event.

Proposition~\ref{prop:step-002-record-map} maps every raw neighbor
pair to equality or a quotient neighbor pair. Applying (3.4) after that map
proves the same pure inequality for \(A_N^{\mathrm{fin}}\), including
arbitrary nonrealizable labels and same-cell replacements. Finally, a pure
\(\varepsilon\)-DP inequality implies the setting's
\((\varepsilon,\delta)\)-DP inequality by adding any \(\delta>0\) to its
right side. \(\square\)


\end{proof}

\begin{lemma}[Direct unconditional realizable risk tail]\label{lem:step-016-finite-tail}
Under Assumptions~\ref{assump:countable-evaluation-quotient},
\ref{assump:realizable-iid}, and
\ref{assump:approximate-dp-regime}, 
Propositions~\ref{prop:step-002-iid-pushforward} and
\ref{prop:step-002-risk}, and
Proposition~\ref{prop:step-016-finite-kernel}, if \(M<\infty\), then for
every \(N\in\mathbb N_0\), every probability measure \(D\), and every
\(c\in C\),
\[
\Pr_{S\sim P_{D,c}^N,\,
\bar H\sim A_N^{\mathrm{fin}}(S,\cdot)}
\!\left[
\operatorname{err}_D(\operatorname{Dec}_C(\bar H),c)>\alpha
\right]
\le M\exp\!\left(-\frac{N\varepsilon\alpha}{4}\right).
\tag{4.1}
\]
The probability is unconditional over both iid sampling and mechanism
randomness, and the bound is uniform in \(D,c\).


\end{lemma}

\begin{proof}
Work first on the quotient experiment. Fix \(D,c\), write
\(\bar c\in\bar C\) for the quotient target, and for
\(\bar h\in\bar C\) put
\[
p_{\bar h}:=
\operatorname{err}_{\bar D}(\bar h,\bar c).
\tag{4.2}
\]
Let \(\bar S\sim P_{\bar D,\bar c}^N\). Realizability gives
\(L_{\bar S}(\bar c)=0\) for every realized sample. Hence the target has
weight one and
\[
Z_{\bar S}\ge w_{\bar S}(\bar c)=1.
\tag{4.3}
\]
For each fixed \(\bar h\), (2.4) and (4.3) imply the pointwise conditional
bound
\[
K^{\mathrm{fin}}_{C,N}(\bar S,\{\bar h\})
\le \exp\!\left(-\frac\varepsilon2L_{\bar S}(\bar h)\right).
\tag{4.4}
\]
The \(N\) mistake indicators of this fixed \(\bar h\) are iid Bernoulli
with mean \(p_{\bar h}\). Therefore
\[
\begin{aligned}
\mathbb E\exp\!\left(-\frac\varepsilon2L_{\bar S}(\bar h)\right)
&=\left(1-p_{\bar h}+p_{\bar h}e^{-\varepsilon/2}\right)^N\\
&=\left[1-p_{\bar h}(1-e^{-\varepsilon/2})\right]^N\\
&\le
\exp\!\left[-Np_{\bar h}(1-e^{-\varepsilon/2})\right].
\end{aligned}
\tag{4.5}
\]
The last line uses \(1-u\le e^{-u}\), which follows directly from
\(e^t\ge1+t\) at \(t=-u\). To lower-bound the remaining coefficient, set
\(x=\varepsilon/2\in(0,1/2]\). Since \(e^x\ge1+x\),
\[
1-e^{-x}\ge1-\frac1{1+x}=\frac{x}{1+x}
\ge\frac x2=\frac\varepsilon4.
\tag{4.6}
\]

The bad subset
\(\mathcal B:=\{\bar h\in\bar C:p_{\bar h}>\alpha\}\) is finite. Taking
expectations in (4.4), using (4.5)-(4.6), and summing only over this fixed
finite set gives
\[
\begin{aligned}
\Pr[\operatorname{err}_{\bar D}(\bar H,\bar c)>\alpha]
&=\mathbb E_{\bar S}
\sum_{\bar h\in\mathcal B}
K^{\mathrm{fin}}_{C,N}(\bar S,\{\bar h\})\\
&\le\sum_{\bar h\in\mathcal B}
e^{-Np_{\bar h}(1-e^{-\varepsilon/2})}\\
&\le |\mathcal B|e^{-N\varepsilon\alpha/4}
\le Me^{-N\varepsilon\alpha/4}.
\end{aligned}
\tag{4.7}
\]
This calculation does not condition on a generalization event and does not
introduce a separate empirical-to-population sampling term.

Proposition~\ref{prop:step-002-iid-pushforward} identifies the raw
sample's quotient image with the quotient iid law, and Proposition~\ref{prop:step-002-risk} identifies the measurable bad-risk
events pointwise. Thus (4.7) is exactly the raw probability in (4.1).
Nothing in the bound depends on \(D\) or \(c\), so both suprema are valid.
\(\square\)


\end{proof}

\begin{proposition}[Ceiling-aware complete finite-class arm]\label{prop:step-016-finite-arm}
Under Assumptions~\ref{assump:finite-littlestone},
\ref{assump:countable-evaluation-quotient},
\ref{assump:realizable-iid}, and
\ref{assump:approximate-dp-regime}, and
Propositions~\ref{prop:step-016-zero},
\ref{prop:step-016-finite-kernel}, and
\ref{prop:step-016-finite-dp}, and
Lemma~\ref{lem:step-016-finite-tail}, if \(M=|C|<\infty\), then the finite
arm is a measurable quotient-first learner whose raw pullback is pure
\(\varepsilon\)-DP, hence \((\varepsilon,\delta)\)-DP, and whose uniform
realizable decoded-risk failure probability is at most \(\beta\). It uses
\(N_{\mathrm{fin}}=0\) when \(d=0\). When \(d\ge1\), it uses
\[
N_{\mathrm{fin}}
:=\left\lceil
\frac{4(\log M+\log(1/\beta))}{\varepsilon\alpha}
\right\rceil
\tag{5.1}
\]
and satisfies the exact normalized bound
\[
N_{\mathrm{fin}}
\le8R_{\mathrm{fin}}(M,\alpha,\beta,\varepsilon).
\tag{5.2}
\]
Thus one may take \(K_{\mathrm{fin}}=8\) and
\(q_{\mathrm{fin}}=0\), with no hidden positive power or sampling term.

If \(|C|=\infty\), define the finite arm's certified cost to be
\(B_{\mathrm{fin}}:=+\infty\) and construct no finite-arm law or finite
surrogate.


\end{proposition}

\begin{proof}
The \(d=0\) case is exactly
Proposition~\ref{prop:step-016-zero}. Here \(M=1\), \(\log M=0\), and
\(\log^+M=1\), but the shared exact learner uses no samples; in particular,
\(0\le8R_{\mathrm{fin}}\).

Suppose \(d\ge1\) and \(M<\infty\). Then \(C\) is not a singleton, so
\(M\ge2\), although the following normalization needs only \(M\ge1\).
Write \(b:=\log(1/\beta)\). Definition (5.1) gives
\[
\frac{N_{\mathrm{fin}}\varepsilon\alpha}{4}
\ge\log M+b.
\tag{5.3}
\]
Lemma~\ref{lem:step-016-finite-tail} and (5.3) therefore give
\[
\sup_D\sup_{c\in C}
\Pr[\operatorname{err}_D(\operatorname{Dec}_C(\bar H),c)>\alpha]
\le M e^{-\log M-b}=\beta.
\tag{5.4}
\]
Propositions~\ref{prop:step-016-finite-kernel} and
\ref{prop:step-016-finite-dp} supply the complete kernel and all-input
privacy interfaces at this actual integer sample size.

It remains to pay the ceiling and normalize to the setting's exact rate. Put
\[
x:=\frac{4(\log M+b)}{\varepsilon\alpha}.
\tag{5.5}
\]
Because \(b>\log4>1\), \(0<\varepsilon\le1\), and
\(0<\alpha<1/4\), one has \(x>16>1\). Hence
\[
N_{\mathrm{fin}}=\lceil x\rceil
\le x+1\le2x
=\frac{8(\log M+b)}{\varepsilon\alpha}.
\tag{5.6}
\]
For every finite nonempty class,
\(\log M\le\log^+M=\max\{1,\log M\}\). Therefore
\[
\begin{aligned}
N_{\mathrm{fin}}
&\le\frac{8(\log^+M+b)}{\varepsilon\alpha}\\
&\le8\left[
\frac{\log^+M+b}{\varepsilon\alpha}+\frac b\alpha
\right]
=8R_{\mathrm{fin}}.
\end{aligned}
\tag{5.7}
\]
The second summand in the setting's \(R_{\mathrm{fin}}\) is used only as
nonnegative slack in the last inequality. The direct tail (4.7) did not
create a \(\log M/\alpha\) or any other sampling term. All constants in
(5.1)-(5.7) are numerical, and no factor \(\Lambda\) is needed.

For infinite \(C\), the setting already defines \(\log^+|C|=+\infty\).
Assigning cost \(+\infty\) disables precisely this arm and does not define a
data-dependent finite subset of \(C\) or \(\bar C\). \(\square\)


\end{proof}

\begin{proposition}[Minimum of three complete common-interface learners]\label{prop:step-016-minimum}
Under Assumptions~\ref{assump:finite-littlestone},
\ref{assump:countable-evaluation-quotient},
\ref{assump:realizable-iid}, and
\ref{assump:approximate-dp-regime},
Proposition~\ref{prop:step-002-risk},
Propositions~\ref{prop:step-016-zero} and
\ref{prop:step-016-finite-arm}, Proposition~\ref{prop:step-014-vc-arm}, and 
Propositions~\ref{prop:step-015-kernel},
\ref{prop:step-015-dp}, \ref{prop:step-015-pac}, and
\ref{prop:step-015-rate}, there is a universal
\[
K_*:=\max\{8,K_{\mathrm O},K_{\mathrm V}\}\ge1
\tag{6.1}
\]
such that, for every fixed allowed \(0<\delta<1\),
\[
m_C(\alpha,\beta;\varepsilon,\delta)
\le K_*\Lambda^6
\min\{R_{\mathrm{fin}},R_{\mathrm{old}},R_{\mathrm{VC}}\}.
\tag{6.2}
\]
The minimum selects one complete learner before data are observed; it does
not mix learner laws. For infinite \(C\), the extended-real value
\(R_{\mathrm{fin}}=+\infty\) removes only the finite arm.


\end{proposition}

\begin{proof}
The \(d=0\) case is already exact: by
Proposition~\ref{prop:step-016-zero}, \(m_C=0\), so (6.2) follows from the
nonnegativity of the three displayed rates. Assume henceforth \(d\ge1\).

The completed arm certificates are as follows. When \(M<\infty\), the finite
arm uses \(K^{\mathrm{fin}}_{C,N_{\mathrm{fin}}}\) and its raw pullback,
outputs in \(H_C\) followed by \(\operatorname{Dec}_C\), is pure
\(\varepsilon\)-DP (hence \((\varepsilon,\delta)\)-DP), has raw decoded-risk
failure at most \(\beta\), and has certified cost
\(B_{\mathrm{fin}}:=8R_{\mathrm{fin}}\). The old arm uses \(K_{\mathrm o}\)
and its raw pullback, with the same output, privacy, and PAC interfaces, and
has cost \(B_{\mathrm o}:=K_{\mathrm O}\Lambda^6R_{\mathrm{old}}\). The
VC arm has the same interfaces and cost
\(B_{\mathrm V}:=K_{\mathrm V}\Lambda^4R_{\mathrm{VC}}\).

Every row is a complete Markov-kernel learner on all raw labeled inputs. The
PAC statement in each row is unconditional, is uniform over the same
\(D,c\), uses binary population risk, and has confidence \(1-\beta\).
Proposition~\ref{prop:step-002-risk}, recorded as (D1), is the
common zero-residual quotient/raw risk bridge. The laws need not be equal on
the positive-dimensional branch; their mathematical interfaces are equal.

If \(M=\infty\), set \(B_{\mathrm{fin}}=+\infty\) and omit that row from
selection. Both Lyu~\citep{lyu2025} bounds are finite under the fixed parameter regime, so at
least one arm remains. If \(M<\infty\), all three rows are available. Choose
an index \(j_*\) minimizing the certified real thresholds
\(B_{\mathrm{fin}},B_{\mathrm o},B_{\mathrm V}\), with a fixed
data-independent tie rule. This choice is a function only of
\(C,d,v,\alpha,\beta,\varepsilon,\delta\), on which the definition of
\(m_C\) permits the learner to depend. Run the single complete learner in
row \(j_*\), using its own integer sample size \(N_{j_*}\). No random arm
choice, sample splitting across arms, post-selection, or mixture occurs.

Each row's actual sample size is bounded by its own threshold, so
\[
m_C\le N_{j_*}\le B_{j_*}
=\min\{B_{\mathrm{fin}},B_{\mathrm o},B_{\mathrm V}\}.
\tag{6.3}
\]
The setting definition gives \(\Lambda\ge1\). With (6.1),
\[
\begin{aligned}
B_{\mathrm{fin}}&\le K_*\Lambda^6R_{\mathrm{fin}},\\
B_{\mathrm o}&\le K_*\Lambda^6R_{\mathrm{old}},\\
B_{\mathrm V}&\le K_*\Lambda^6R_{\mathrm{VC}}.
\end{aligned}
\tag{6.4}
\]
Taking the minimum in (6.4) and applying (6.3) proves (6.2), including the
extended-real infinite-class convention. None of these finite-parameter
inequalities uses an asymptotic condition on \(\delta\). \(\square\)


\end{proof}

\begin{proposition}[Boundary-preserving fixed-parameter frontier]\label{prop:step-016-frontier}
Under Assumptions~\ref{assump:finite-littlestone},
\ref{assump:countable-evaluation-quotient},
\ref{assump:realizable-iid}, and
\ref{assump:approximate-dp-regime},
Propositions~\ref{prop:step-016-zero},
\ref{prop:step-016-finite-arm}, and
\ref{prop:step-016-minimum}, Proposition~\ref{prop:step-014-vc-arm}, and 
Propositions~\ref{prop:step-015-rate} and
\ref{prop:step-015-boundaries}, the following exact boundary and
specialization statements hold.

1. If \(d=0\), then \(v=0\), \(|C|=1\), and the selected law is the shared
   zero-sample law of Proposition~\ref{prop:step-016-zero}; in particular,
   \(m_C=0\). No positive-dimensional rate is evaluated.
2. If \(d\ge1\), then \(1\le v\le d\). At \(v=1\), the VC arm's exposed
   structural factor is \(d^4\). At \(v=d\),
   \[
   vd^4=d\,d^4=d^5
   \tag{7.1}
   \]
   exactly, with no hidden positive power changing the equality. More
   generally \(v=\Theta(d)\) gives the \(d^5\) polynomial scale.
3. If \(C\) is finite, all three complete arms are available and the finite
   arm exposes \(\log^+|C|\). If \(C\) is infinite, then
   \(R_{\mathrm{fin}}=+\infty\), no finite surrogate is constructed, and the
   old and VC arms remain unchanged and available.
4. For every fixed \(0<\delta<1\), the kernel, privacy, PAC, and full
   finite-parameter minimum (6.2) are valid. No statement
   \(N\delta\to0\) follows from fixed positive \(\delta\).
5. On \(d\ge1\), let \(N_*:=N_{j_*}\) be the actual sample size selected in
   the proof of Proposition~\ref{prop:step-016-minimum}. Define \(q_*:=6\).
   Because the selected threshold is minimal, \(\Lambda\ge1\), and
   \(K_*=\max\{8,K_{\mathrm O},K_{\mathrm V}\}\ge K_{\mathrm V}\),
   the proved VC exponent \(4\) is dominated by \(q_*\), and
   \[
   0\le N_*\delta
   \le\delta\min\{B_{\mathrm{fin}},B_{\mathrm o},B_{\mathrm V}\}
   \le\delta K_{\mathrm V}\Lambda^4R_{\mathrm{VC}}
   \le\delta K_*\Lambda^{q_*}R_{\mathrm{VC}}.
   \tag{7.2}
   \]
   Hence the public schedule
   \(\delta K_*\Lambda^{q_*}R_{\mathrm{VC}}\to0\) implies
   \(N_*\delta\to0\), whichever complete arm wins the deterministic
   comparison. This implication is not asserted under a weaker premise.
   If the standalone old arm's own sample size is being studied rather than
   the minimizing selected learner, Proposition~\ref{prop:step-015-boundaries} supplies only
   \[
   \delta K_{\mathrm O}\Lambda^6R_{\mathrm{old}}\to0
   \quad\Longrightarrow\quad
   N_{\mathrm o}\delta=N_{\mathrm{old}}\delta\to0;
   \tag{7.3}
   \]
   the VC schedule is not silently substituted for (7.3). The finite arm is
   pure DP and needs no \(\delta\)-schedule for privacy; no numerical
   \(N_{\mathrm{fin}}\delta\) limit is claimed unless a corresponding scalar
   domination is supplied.
6. Fix \(\alpha,\beta,\varepsilon\) and restrict the
   asymptotic comparison to the declared scheduled-small-\(\delta\)
   convention. On \(d\ge1\), \(\log(1/\beta)\) is fixed and \(v\ge1\), so
   the three displayed rate expressions have polynomial class-complexity
   profiles
   \[
   R_{\mathrm{fin}}:\ \log^+|C|,
   \qquad
   R_{\mathrm{old}}:\ d^5,
   \qquad
   R_{\mathrm{VC}}:\ vd^4,
   \tag{7.4}
   \]
   with only the setting-permitted displayed logarithmic factors suppressed.
   Thus (6.2) is exactly the formalized shorthand
   \[
   m_C(\alpha,\beta;\varepsilon,\delta)
   =\widetilde O\!\left(
   \min\{\log^+|C|,d^5,vd^4\}
   \right).
   \tag{7.5}
   \]
   Here the first entry is \(+\infty\) for infinite \(C\).
7. The comparison between the two Lyu~\citep{lyu2025} certificates is an exposed structural
   upper-bound comparison:
   \[
   \frac{vd^4}{d^5}=\frac vd.
   \tag{7.6}
   \]
   Consequently \(v=o(d)\) gives \(vd^4=o(d^5)\), while
   \(v=\Theta(d)\) returns to the \(d^5\) scale and \(v=d\) gives the exact
   identity (7.1). This does not assert a lower bound on \(m_C\), and
   arm-specific displayed logarithms remain governed by the setting's tilde
   convention.
8. All conclusions remain conditional on the primitive finite-or-countable
   measurable evaluation quotient. They do not cover uncountably many
   evaluation types, characterize private sample complexity, or prove a
   universal \(\operatorname{poly}(v,\log d)\) or
   \(\operatorname{poly}(v,\log^*d)\) bound.


\end{proposition}

\begin{proof}
Item 1 is Proposition~\ref{prop:step-016-zero}. Conversely, a
singleton binary class shatters no one-point set or depth-one tree, so the
singleton and \(d=0\) branches coincide. If \(d\ge1\), two concepts differ
at some point, and that point is VC-shattered; hence \(v\ge1\). If a set
\(\{x_1,\ldots,x_r\}\) is VC-shattered, label every node at tree depth
\(j<r\) by \(x_{j+1}\). Every root-to-leaf binary labeling is realized by
a concept, so this is a Littlestone-shattered tree of depth \(r\). Therefore
\(v\le d\). Direct substitution now proves Items 2 and 7. Proposition
\ref{prop:step-016-finite-arm} and the extended-real convention in
Proposition~\ref{prop:step-016-minimum} prove Item 3. The final sentence of
that minimum proposition, together with pure privacy from
Proposition~\ref{prop:step-016-finite-arm}, the fixed-parameter statement of
Proposition~\ref{prop:step-014-vc-arm}, and the fixed-parameter
statement of Proposition~\ref{prop:step-015-boundaries}, proves
Item 4.

For Item 5, (6.3) makes \(N_*\) no larger than every available certified
threshold, in particular \(B_{\mathrm V}\); multiplying by \(\delta\) gives
(7.2). The squeeze proves the selected-learner limit only under the exact
setting schedule. Equation (7.3) is the separate old-arm statement,
not a consequence newly inferred here.

For Items 6 and 7, with \(\alpha,\beta,\varepsilon\) fixed, the coefficients
and additive confidence quantities are fixed. Since \(v\ge1\), the first
term of \(R_{\mathrm{VC}}\) has exposed polynomial factor \(vd^4\); Proposition~\ref{prop:step-014-vc-arm} proves that the lower-order displayed
term introduces no larger positive power. Proposition~\ref{prop:step-015-rate} exposes exactly \(d^5\), and
Proposition~\ref{prop:step-016-finite-arm} exposes exactly
\(\log^+|C|\). Equation (6.2) then gives (7.5), with the universal
\(K_*\Lambda^6\) hiding only permitted logarithms. The quotient and
open-problem scope in Item 8 is precisely the unchanged formalized setting.
Proposition~\ref{prop:step-016-zero} first isolates the exact common
\(d=0\) law. On the finite positive-dimensional branch,
Proposition~\ref{prop:step-016-finite-kernel} constructs the total
finite-range quotient law and raw pullback on the same \(H_C,T_N\), and
decoder interface as the other arms. Proposition
\ref{prop:step-016-finite-dp} proves pure all-input privacy by the two
sensitivity-one weight comparisons (3.2)-(3.4). Lemma
\ref{lem:step-016-finite-tail} uses the realizable target's unit weight and
the exact Bernoulli transform to prove the unconditional uniform risk tail
(4.1), without a separate sampling event. Proposition
\ref{prop:step-016-finite-arm} chooses the actual integer sample size, pays
the ceiling, handles \(M=1\) and \(\log^+M\), and proves
\(N_{\mathrm{fin}}\le8R_{\mathrm{fin}}\); it assigns cost \(+\infty\) and
constructs no law when \(C\) is infinite.

Proposition~\ref{prop:step-014-vc-arm} supplies the complete
VC-arm tuple. Propositions~\ref{prop:step-015-kernel},
\ref{prop:step-015-dp}, \ref{prop:step-015-pac}, and
\ref{prop:step-015-rate} supply the independently complete old-arm tuple.
Consequently, the deterministic threshold comparison in
Proposition~\ref{prop:step-016-minimum} selects one entire certified learner
with bound
\[
 K_*\Lambda^6\min\{R_{\mathrm{fin}},R_{\mathrm{old}},R_{\mathrm{VC}}\}.
\]
No probability law is averaged or spliced across arms.

Finally, Items 1--8 and their preceding derivations prove each null,
dimension, cardinality, fixed-\(\delta\), scheduled-\(\delta\), and scaling
specialization. It separates the selected learner's setting schedule from
the standalone old arm's condition and states the exact conditional frontier
without upgrading it to a characterization or unrestricted solution. Together,
these results complete the common-interface comparison.



\end{proof}

\subsection{Proof of the Main Theorem}\label{app:main-proof}

\begin{proof}[Proof of the main theorem]
The quotient factorization, dimension preservation, raw-neighbor transport,
iid pushforward, and exact risk identity are given by
Propositions~\ref{prop:step-002-factorization},
\ref{prop:step-002-record-map},
\ref{prop:step-002-iid-pushforward}, and
\ref{prop:step-002-risk}. The countable-domain kernel and raw pullback are
provided by Propositions~\ref{prop:step-003-quotient-kernel} and
\ref{prop:step-003-raw-pullback}. For the VC arm,
Proposition~\ref{prop:step-014-vc-arm} supplies the measurable kernel,
all-input privacy, PAC guarantee, and
$K_{\mathrm V}\Lambda^4R_{\mathrm{VC}}$ sample bound. For the comparison
arm, Propositions~\ref{prop:step-015-kernel},
\ref{prop:step-015-dp}, \ref{prop:step-015-pac},
\ref{prop:step-015-rate}, and
\ref{prop:step-015-boundaries} supply the old-Lyu kernel, privacy, PAC,
rate, and boundary statements. The finite arm is supplied by
Propositions~\ref{prop:step-016-finite-kernel},
\ref{prop:step-016-finite-dp},
\ref{prop:step-016-finite-arm}, and
\ref{prop:step-016-zero}. Proposition~\ref{prop:step-016-minimum}
selects one complete arm before data are observed, and
Proposition~\ref{prop:step-016-frontier} proves the null, dimension,
cardinality, and scheduled-\(\delta\) reductions. All three arms use the
same output space, decoder, privacy quantifier, and risk identity, so the
deterministic minimum is a single learner. This proves Theorem~\ref{thm:main}.
\end{proof}
